% concatenated from zgqq222.tex
\documentclass[11pt,a4paper]{amsart}
\IfFileExists{glyphtounicode.tex}{\input{glyphtounicode}\pdfgentounicode=1}{}
\usepackage{cmap}
\usepackage[T1]{fontenc}
\usepackage[utf8]{inputenc}
\usepackage{lmodern,microtype}
\usepackage[a4paper,margin=1in]{geometry}
\usepackage{amsmath,amssymb,amsthm,mathtools,mathrsfs,bm}
\usepackage{booktabs,array,enumitem,needspace}
\usepackage{cite}
\usepackage[bookmarks=true,bookmarksnumbered=true,hidelinks]{hyperref}
\hypersetup{pdftitle={Gaussian Operator Algebras: Observability, Reconstruction, and Microscopic Limits},pdfauthor={Guangqian Zhao}}
\emergencystretch=3em
\hyphenpenalty=4000
\exhyphenpenalty=4000
\Urlmuskip=0mu plus 1mu
\allowdisplaybreaks[2]
\raggedbottom
\numberwithin{equation}{section}
\newtheorem{theorem}{Theorem}[section]
\newtheorem{proposition}[theorem]{Proposition}
\newtheorem{lemma}[theorem]{Lemma}
\newtheorem{corollary}[theorem]{Corollary}
\theoremstyle{definition}
\newtheorem{definition}[theorem]{Definition}
\newtheorem{example}[theorem]{Example}
\theoremstyle{remark}
\newtheorem{remark}[theorem]{Remark}
\newcommand{\R}{\mathbb R}
\newcommand{\C}{\mathbb C}
\newcommand{\N}{\mathbb N}
\newcommand{\Z}{\mathbb Z}
\newcommand{\E}{\mathbb E}
\newcommand{\Prob}{\mathbb P}
\newcommand{\HH}{\mathcal H}
\newcommand{\cC}{\mathcal C}
\newcommand{\cE}{\mathcal E}
\newcommand{\Sch}{\mathcal S}
\newcommand{\Coeff}{\mathfrak C}
\newcommand{\Wc}{\mathfrak W}
\newcommand{\Alg}{\mathfrak A}
\newcommand{\Data}{\mathfrak D}
\newcommand{\Gaus}{\mathfrak G}
\newcommand{\Flat}{\mathcal F}
\newcommand{\Prof}{\mathfrak P}
\newcommand{\PP}{\mathcal P}
\newcommand{\LL}{\mathfrak L}
\DeclareMathOperator{\UCP}{UCP}
\DeclareMathOperator{\Id}{Id}
\DeclareMathOperator{\Tr}{Tr}
\DeclareMathOperator{\tr}{tr}
\DeclareMathOperator{\rank}{rank}
\DeclareMathOperator{\ran}{ran}
\DeclareMathOperator{\Span}{span}
\DeclareMathOperator{\diag}{diag}
\DeclareMathOperator{\dist}{dist}
\DeclareMathOperator{\Law}{Law}
\DeclareMathOperator{\Sym}{Sym}
\DeclareMathOperator{\sym}{sym}
\DeclareMathOperator{\ad}{ad}
\DeclareMathOperator*{\esssup}{ess\,sup}
\newcommand{\op}{\mathrm{op}}
\newcommand{\norm}[1]{\left\lVert#1\right\rVert}
\newcommand{\abs}[1]{\left|#1\right|}
\newcommand{\ip}[2]{\left\langle#1,#2\right\rangle}
\newcommand{\Nn}[2]{\mathcal N_{#1,s}(#2)}
\newcommand{\En}[2]{\mathcal E_{#1,s}(#2)}
\newcommand{\Dens}{\mathscr D_s}
\newcommand{\Jdu}{\mathscr J_s}
\newcommand{\arxiv}[1]{\href{https://arxiv.org/abs/#1}{\nolinkurl{arXiv:#1}}}
\newcommand{\doi}[1]{\href{https://doi.org/#1}{\nolinkurl{doi:#1}}}
\title[Gaussian operator algebras]{Gaussian Operator Algebras:\\
Observability, Reconstruction, and Microscopic Limits}
\author{Guangqian Zhao}
\address{School of Mathematical Sciences, University of Science and Technology of China, Hefei, Anhui 230026, China}
\email{zhaoguangqian@mail.ustc.edu.cn}
\subjclass[2020]{60B20, 60G15, 47B10, 46B70, 47A57}
\keywords{Gaussian operator algebras, observability, intrinsic innovations, Schatten classes, stable reconstruction, microscopic limits, spectral quotients}
\date{}
\begin{document}
\begin{abstract}
We reconstruct Gaussian operator algebras from finite source-labelled observations, recovering coefficients, ordinary multiplication, and completion in independently specified Schatten norms. Gaussian duality and ordered contractions give a reconstruction theorem including the trace-class endpoint, while entire Mehler responses describe the resulting topology. In the Hilbert coefficient geometry, symmetric Rademacher systems realize the same algebra as a fluctuation limit: microscopic products converge at inverse site-number rate in the original norm scale, and finite decoding commutes with this limit. Relative Gram residuals determine target-dependent inverse costs. For bounded tuples, positive word observations close into completely positive maps, with contextual Schwarz defects controlling realization and multiplication. Positive mass and source innovations characterize compactness. Spectral invariants, Wick operators on closed manifolds, and Gaussian networks provide applications with explicit noise and tail bounds.
\end{abstract}
\maketitle
\pagestyle{plain}
\raggedbottom

\section{Introduction}\label{sec:intro}
A fluctuation limit can retain a multiplication law even when it forgets much of the underlying microscopic system. For symmetric sums of independent signs, shared sites disappear in pairs and produce the contractions of Gaussian chaos. Recovering this limiting algebra requires more than convergence in distribution: the product must converge in a topology that supports its operations, and finite observations must determine its coefficients with controlled error.

We develop such a reconstruction theory for Gaussian operator algebras. Its analytic core identifies source-labelled observations, independently defined Schatten coefficient spaces, and ordinary random-operator multiplication. The construction includes every finite Schatten exponent, including trace class. In the Hilbert coefficient geometry, a symmetric Rademacher model realizes the same product as a quantitative microscopic limit. A second part of the theory treats represented algebras and their observation boundaries: ordered moments recover relations and multiplication tables, while anchored positive kernels distinguish actual tuples from more general completely positive limits.

The central distinction is between an observable relation, an inverse estimate, and a realization of limiting data. These requirements have separate mathematical content. The responses $gP$ and $-gP$, for a standard Gaussian $g$ and nonzero rank-one projection $P$, have the same marginal law but opposite source correlations. Likewise, $H_m(g)P/\sqrt{m!}$ has eventually zero observations at every fixed degree, while the constant coefficient of its square remains $P$. Source-labelled tests remove the first ambiguity; a norm scale controlling all contraction depths removes the second.

\subsection{Reconstruction in the original coefficient geometry}
Fix a real separable Hilbert space $H$, an isonormal Gaussian process $W$ over $H$, and a complex separable physical Hilbert space $E$. On finite-source, finite-matrix kernels, write
\begin{equation}\label{eq:intro-eval}
 F=I(K)=\sum_{m\ge0}I_m(K_m),\qquad
 \E|I_m(f)|^2=m!\norm f_2^2.
\end{equation}
All products are taken in this one Gaussian realization. Symmetry of the source variables leaves the order of the operator coefficients unchanged. Consequently, ordinary multiplication contracts source indices, multiplies physical coefficients in their original order, and produces both higher and lower chaos degrees.

For a degree-$m$ kernel, a subset $S\subset[m]$ specifies which source legs are placed on the input side of a flattening $\Flat_S$. The coefficient space $\Coeff_{m,s}$ is defined using the quotient sum of these Schatten-cut spaces for $s<2$, their compatible intersection for $s>2$, and their common Hilbert norm for $s=2$. The precise construction is given in Section~\ref{sec:coeff}. Let $\Wc_{m,s}$ be its source-symmetric subspace. Before applying random evaluation, define
\begin{equation}\label{eq:intro-scale}
 \Nn\rho K=\sum_{m\ge0}\rho^m\sqrt{m!}\norm{K_m}_{\Wc_{m,s}},\qquad
 \Alg_s=\{K:\Nn\rho K<\infty\text{ for every }\rho\ge1\}.
\end{equation}
The weighted norm is defined for every $\rho>0$; radii $\rho\ge1$ specify the topology. Integer radii are cofinal, so this is a countable norm scale. The degree-zero space is $\Sch_s(E)$. When $E$ is infinite-dimensional, constants proportional to the identity are treated in the unitization.

Independently, let $\mathcal H_m(\Sch_s)$ be the closed space of $\Sch_s$-valued homogeneous Gaussian chaoses of degree $m$, and set
\begin{equation}\label{eq:intro-random}
 \begin{aligned}
 \Gaus_s&=\left\{F=\sum_{m\ge0}F_m:F_m\in\mathcal H_m(\Sch_s),\
                      \En\rho F<\infty\ \text{for all }\rho\ge1\right\},\\
 \En\rho F&=\sum_{m\ge0}\rho^m\norm{F_m}_{L^2(\Sch_s)}.
 \end{aligned}
\end{equation}
The sum converges in $L^2(\Sch_s)$, and the homogeneous components are unique.

For observations, choose a positive injective Hilbert--Schmidt operator $R$ with $Re_k=r_ke_k$, $r_k>0$. Let $(h_\alpha)_{|\alpha|=m}$ be the orthonormal symmetric tensor basis normalized by $I_m(h_\alpha)=\sqrt{m!}\Psi_\alpha$. A finite budget $a=(M,N,L)$ records
\begin{equation}\label{eq:intro-marks}
 z_{\alpha;k\ell}=\E[\Psi_\alpha(W)(RFR)_{k\ell}],
 \quad |\alpha|\le M,\quad\operatorname{supp}\alpha\subset[L],\quad k,\ell\le N.
\end{equation}
These are joint source--response observations. Their inverse $D_a$ divides each coordinate by $r_kr_\ell\sqrt{|\alpha|!}$ and sets all other coefficients to zero. Let $P_n$ denote the corresponding kernel truncations along an increasing cofinal sequence of budgets $a_n$, let $O_n$ denote the observations, and write $\widehat K_n(y)=D_{a_n}y_{a_n}$. Define $\Data_s$ to consist of compatible arrays satisfying
\begin{equation}\label{eq:intro-admissible}
 \lim_{N\to\infty}\sup_{n\ge N}
 \Nn\rho{\widehat K_n(y)-\widehat K_N(y)}=0
 \qquad(\rho\ge1).
\end{equation}
Thus admissibility is expressed entirely through finite inverses and the prescribed coefficient norms.

\begin{theorem}[Reconstruction of operations and completion]\label{thm:main}
Let $1\le s<\infty$. The ordered Gaussian contraction product on finite kernels extends uniquely to $\Alg_s$. Random evaluation and finite decoding give topological $*$-algebra isomorphisms
\begin{equation}\label{eq:main-diagram}
 \Data_s\ \xleftrightarrow[\ (O_n)_n\ ]{\ \lim_nD_{a_n}\ }\
 \Alg_s\ \xleftrightarrow[\ \text{coefficient recovery}\ ]{\ I\ }\ \Gaus_s.
\end{equation}
The data product is obtained from finite computations:
\[
 (y\diamond z)_n=\lim_{M\to\infty}
 O_n\big(\widehat K_M(y)\star\widehat K_M(z)\big).
\]
For every output radius $\sigma\ge1$, there are explicit $R_s(\sigma)$ and $C_{s,\sigma}<\infty$ such that
\begin{equation}\label{eq:intro-product}
 \Nn\sigma{K\star L}
 \le C_{s,\sigma}\Nn{R_s(\sigma)}K\Nn{R_s(\sigma)}L.
\end{equation}
The data norms and tails recover the original ones exactly:
\begin{align}\label{eq:intro-tail-identity}
 \sup_n\Nn\rho{D_{a_n}O_nK}&=\Nn\rho K,\notag\\
 \sup_{n\ge N}\Nn\rho{D_{a_n}O_nK-D_{a_N}O_NK}
 &=\Nn\rho{K-P_NK}.
\end{align}
If the total squared observation error at degree $m$ is at most $\delta_m^2$, the finite decoded output satisfies
\begin{equation}\label{eq:intro-noise}
 \Nn\rho{\widehat K_a-K}
 \le N^{(1/s-1/2)_+}\kappa_N\sum_{m=0}^M\rho^m\delta_m
       +\Nn\rho{K-P_aK},\qquad
 \kappa_N=(\min_{k\le N}r_k)^{-2}.
\end{equation}
If the true inputs have budgets $B_K,B_L$ and reconstruction errors $\varepsilon_K,\varepsilon_L$ at the required input radius, their product error at radius $\sigma$ is at most
\begin{equation}\label{eq:intro-bilinearerror}
 C_{s,\sigma}\big(B_L\varepsilon_K+B_K\varepsilon_L+
                  \varepsilon_K\varepsilon_L\big).
\end{equation}
The same reconstruction therefore determines every fixed finite $*$-polynomial.
\end{theorem}


The coefficient norms in this theorem are fixed before random evaluation. Directed-cut duality gives the Gaussian inverse, and an exponential predual estimate covers trace class. Ordered contractions then give the product bound, while compatible finite inverses identify the completion. The proof is completed in Section~\ref{sec:data}. At the Hilbert exponent, the exact multi-input region is
\begin{equation}\label{eq:intro-arity}
 \sum_{i=1}^k\frac1{1+\rho_i^2}\le\frac1{1+\sigma^2},\qquad\sigma\ge1,
\end{equation}
with optimal constant one (Theorem~\ref{thm:arity}). A single norm cannot make both ordinary Gaussian multiplication and expectation bounded; the changing input radius is therefore intrinsic to this form of algebraic control.

The recovered topology also has an analytic description. A source-measurable $F$ belongs to $\Gaus_s$ precisely when its Mehler response extends to an entire $L^2(\Sch_s)$-valued curve
\[
 \mathscr M_F(z)=\sum_{m\ge0}z^mF_m.
\]
The compact-open topology of these curves agrees with the coefficient topology (Theorem~\ref{thm:mehler_entire}). Compactness further requires control of physical and source directions. Theorem~\ref{thm:fullcompact} characterizes relatively compact families by boundedness at every radius, uniformly vanishing positive physical mass tails, and uniformly vanishing anchored source innovations at each fixed degree. Finite corner energies give measurable versions of the source certificate.

\subsection{Microscopic realization and finite decoding}
The Gaussian product can be obtained from a discrete multiplication law. For each source direction take $N$ independent signs and define the symmetric response
\[
 Q_{m,N}=m!N^{-m/2}\sum_{|I|=m}\prod_{a\in I}\xi_a.
\]
Multiplication groups terms according to $r=|I\cap J|$. For $k=m+n-2r$, the corresponding coefficient is
\[
 r!\binom mr\binom nr\frac{(N-k)_{\underline r}}{N^r}.
\]
The remaining-site factor tends to one, leaving the Gaussian contraction coefficient. Tensoring the source blocks while retaining the order of the physical operator coefficients defines a product $\star_N$ in the original Hilbert coefficient coordinates.

Let $p_\rho$ denote $\mathcal N_{\rho,2}$ in these normalized coordinates. Theorem~\ref{mic:algebra_limit} proves a bound uniform over $N$ in exactly the radius region \eqref{eq:intro-arity}, and for the binary product gives
\begin{equation}\label{eq:intro_micro_rate}
 p_\sigma(A\star_NB-A\star B)
 \le\frac{q^2}{N(1-q)^3}p_{R_+}(A)p_{S_+}(B),
\end{equation}
where $R_+>\rho$, $S_+>\eta$, $q=\max(\rho/R_+,\eta/S_+)<1$, and $(\rho,\eta,\sigma)$ satisfies the binary radius condition. Thus multiplication converges in the original norm scale, uniformly on stronger-radius bounded sets. A quadratic response attains the inverse site-number order.

Finite normalized site marks determine the microscopic coefficients exactly. For every fixed observation budget, Theorem~\ref{mic:decoder} gives
\[
 D_b^{(N)}O_b^{(N)}\bigl(\Phi_N(A)\Phi_N(B)\bigr)
 =P_b(A\star_NB)\longrightarrow
 P_b(A\star B)=D_bO_b\bigl(\Phi(A)\Phi(B)\bigr).
\]
The full error is the sum of amplified observation noise, a target tail, and the collision bias \eqref{eq:intro_micro_rate}. The inverse cost of observing degrees through $M\le N$ includes the exact factor
\[
 \chi_{M,N}=\left(\frac{N^M}{(N)_{\underline M}}\right)^{1/2}.
\]
It tends to $e^{c^2/4}$ when $M/\sqrt N\to c$, and is unbounded when $M/\sqrt N\to\infty$. This identifies how microscopic scale and structural resolution interact under the specified normalized-mark noise model. The response laws converge jointly, and an explicit auxiliary coupling gives convergence of fixed noncommutative polynomials in every finite probability norm. The microscopic theorem concerns $\Alg_2$; the all-exponent Gaussian reconstruction is established independently.

\subsection{Observable relations, target costs, and genuine limits}
An observation protocol selects a quotient only when the forgotten relations are preserved by the permitted operations. Closing its dual tests under partial adjoints gives the smallest invariant family containing the initial tests, and the corresponding largest compatible invisible subspace. Section~\ref{sec:principle} proves this dual description and the factorization criterion for stable target recovery.

For an increasing represented test family, the new information is its quotient by the preceding observations. In Hilbert coordinates this is the relative Gram residual
\[
 \mathfrak S_{V,W}=W^*(I-P_{\overline{\ran V}})W.
\]
The same distance construction removes target-null nuisance directions. Its optimal inverse constant depends on the target and, in an innovation tower, on the complete triangular regression. If $G_N=U_N^*D_NU_N$, a target $L_N$ is recoverable precisely when it annihilates the represented kernel; its exact cost is then
$\norm{L_NU_N^{-1}D_N^{\dagger/2}}$.
Section~\ref{sec:ordered} applies this mechanism to ordered moment data, PBW coordinates, and multiplication tables. The finite constants distinguish a visible layer from a stable tower.

Fixing an ambient Hilbert space introduces a further realization question. The full anchored word observations of bounded tuples have compact closure
\[
 \overline{\{\mathcal Q(\pi_T):\max_j\norm{T_j}\le M\}}^{\norm\cdot_1}
 =\mathcal Q\bigl(\UCP(\mathcal U_M,B(H))\bigr).
\]
Here $\mathcal U_M$ is the universal algebra of norm-bounded generators. The Schwarz score
\[
 \delta_R(\phi)=\sum_a\Tr\!\left(R[\phi(a^*a)-\phi(a)^*\phi(a)]R\right),
\]
summing over generators and adjoints, vanishes exactly at representations. Theorem~\ref{thm:strongclosure} identifies it with output mass escaping from the anchored realization. An exact two-boundary word identity makes the criterion quantitative: contextual defects control the error of replacing $\phi(w)$ by the corresponding word in its compressed generators (Theorem~\ref{thm:contextual_multiplicativity}). Finite matrix observations provide noisy upper certificates for that error.

Spectral invariants retain a different quotient. For a skew Hilbert--Schmidt operator with positive block singular values $(a_j)$, compatible exterior coordinates give $c_k=e_k(a^2)$. On $a_j\le Cj^{-1/\rho}$, $0<\rho<2$, finite moment data yield a rank- and gap-independent inverse in every strictly subcritical weak ideal. Critical recovery uses a relative-tail class. With only a square-mass bound, the determinant closure consists of
\[
 e^{\tau z}\prod_j(1+za_j^2),\qquad\tau\ge0,
\]
where $\tau$ is mass transferred to the spectral point zero. Section~\ref{sec:pfaffian} proves these inverse estimates and identifies their completions.

\subsection{Proof organization and related work}
Sections~\ref{sec:principle}--\ref{sec:product} establish the quotient principles, Gaussian inverse, and ordinary product. Section~\ref{sec:microscopic} constructs the discrete fluctuation realization. Section~\ref{sec:data} proves finite reconstruction, compactness, and completion. Ordered moments and positive realizations occupy Sections~\ref{sec:ordered}--\ref{sec:realizationboundary}; common-noise evolution, spectral recovery, and geometric Wick systems follow in Sections~\ref{sec:time}--\ref{sec:realizations}. Appendix~\ref{sec:fockcosts} computes sharp Fock-space observation costs, including
\begin{equation}\label{eq:intro_branch}
 \sum_i c_i^{-2}<1
\end{equation}
for the existence of a faithful trace-one anchor with the prescribed creation-operator multiplier costs. Appendices~\ref{sec:obsappendix} and \ref{sec:physicalboundary} treat bounded joint-law protocols and Gaussian network continuations.

The analytic inputs are the noncommutative Khintchine and predual inequalities of Lust-Piquard--Pisier, Buchholz, and Haagerup--Pisier~\cite{LP,Buchholz,HP}, together with Gaussian chaos and decoupling~\cite{Janson,Maas}. The binary boundary in \eqref{eq:intro-arity} agrees with the scalar weighted Gaussian multiplication boundary of Da Pelo--Lanconelli--Stan~\cite{DPLS}; here the positive sum controls each ordered contraction in an $\ell^1$ degree scale. Its scalar all-radius topology is related to the smooth-random-variable spaces of Potthoff--Timpel~\cite{PT}. Krawtchouk--Hermite asymptotics, discrete product formulas, and universality of homogeneous sums provide the background for the microscopic construction~\cite{Min,NPRR,NPR}.

PBW symmetrization, shorted operators, positive-kernel representations, and multiplicative domains supply the structural inputs~\cite{CP,AT,Ando,Stinespring,Choi,CJK}. Their use here retains the original coefficient norms and quantifies the finite inverse, the transport of operations, and the tails that select a genuine limit. The drift-generated commutant in common-noise evolution is related to the decoherence-free algebra~\cite{DFSU}, while spectral moments and continuation experiments connect the construction to classical invariant and graph-algebra methods~\cite{KongValiant,Lovasz}.

Throughout, Hilbert spaces are separable and inner products are linear in the first variable. Matrix entries are $T_{kl}=\ip{Te_l}{e_k}$. Banach adjoints are denoted by a prime and Hilbert adjoints by a star. Coefficient pairings are linear in the first argument and conjugate-linear in the second, with the usual conjugate identification for Banach duals. The coefficient algebra and its compactness criterion are defined for $1\le s<\infty$; the full-data linear retraction and source-convexity inequality use $1<s<\infty$. Chaos degree, word length, spectral moment order, and observation resolution have their stated independent roles.

\section{Observable quotients and stable recovery}\label{sec:principle}\label{sec:targetquotients}
An observation protocol and its allowed operations determine the relations that may be forgotten. Closing the dual tests under partial adjoints produces the operational quotient. Distances from nuisance observation ranges then determine which targets admit a bounded inverse. The constructions in this section keep these algebraic and norm requirements separate.

\subsection{Operational quotients and invariant tests}
Let $(X_s)$ be a family of Banach spaces indexed by types. A basic operation is a prescribed bounded multilinear map
\[
 \mu:X_{s_1}\times\cdots\times X_{s_k}\longrightarrow X_t.
\]
Choose closed subspaces $N_s\subset X_s$, put $V_s=N_s^\perp\subset X_s'$, and let $q_s:X_s\to X_s/N_s$ be the quotient maps. Fixing every argument of $\mu$ except the $j$th gives a bounded linear map $T_{\mu,j}:X_{s_j}\to X_t$.

\begin{proposition}[Operational descent and dual invariance]\label{prop:congruence}
The basic operations descend to the quotient spaces if and only if, for every partial map,
\begin{equation}\label{eq:invariance}
 T_{\mu,j}N_{s_j}\subset N_t.
\end{equation}
Equivalently, $T_{\mu,j}'V_t\subset V_{s_j}$. The norm of each quotient operation is at most the norm of the original operation. For any bounded $T:X_s\to X_t$, the leakage of invisible inputs has the exact dual expression
\begin{equation}\label{eq:leakage-dual}
 \norm{q_tT|_{N_s}}
 =\sup_{\lambda\in V_t,\ \norm\lambda\le1}
       \dist(T'\lambda,V_s).
\end{equation}
\end{proposition}
\begin{proof}
Changing the representative in one argument changes the output by the corresponding partial map applied to an element of $N_{s_j}$. A telescoping replacement of the arguments proves that well-definedness is equivalent to~\eqref{eq:invariance}. Since $N_t=(N_t^\perp)_\perp$, this condition is equivalent to the stated adjoint invariance. Choosing representatives with norms arbitrarily close to the quotient norms proves the norm bound.

The dual of $X_t/N_t$ is isometrically $V_t$, so the left-hand side of~\eqref{eq:leakage-dual} equals
\[
 \sup_{\lambda\in V_t,\ \norm\lambda\le1}
       \norm{(T'\lambda)|_{N_s}}.
\]
The restriction map $X_s'\to N_s'$ is a metric quotient by Hahn--Banach, with kernel $N_s^\perp=V_s$. Its quotient norm is the distance to that kernel, which gives the formula.
\end{proof}

Given initial observations $O_s:X_s\to Y_s$, define the successive test spaces by
\[
 \begin{split}
 V_s^{(0)}&=\overline{\ran O_s'}^{\,w^*},\\
 V_s^{(n+1)}&=\overline{\Span\big(V_s^{(n)}\cup
 \{T_{\mu,j}'\lambda:s_j=s,\ \lambda\in V_t^{(n)}\}\big)}^{\,w^*}.
 \end{split}
\]
For a conjugate-linear involution the additional descent condition is invariance of the corresponding invisible subspace under that involution. It may equivalently be included in the saturation by working over the underlying real spaces. The partial maps range over all allowed choices of the other arguments. Put
$V_s^{(\infty)}=\overline{\bigcup_nV_s^{(n)}}^{\,w^*}$ and
$N_s^{(n)}=(V_s^{(n)})_\perp$. Adjoint maps are weak-$*$ continuous, so $V^{(\infty)}$ is invariant. It is the smallest weak-$*$ closed invariant family containing the initial tests. Dually, $N^{(\infty)}$ is the largest operationally compatible closed family contained in the initial observation kernels: any such family is annihilated by every stage of the construction, by induction.

For linear operations $A_i:X\to X$ of one type, this reads
\begin{equation}\label{eq:word-observability}
 N_n=\bigcap_{|w|\le n}\ker(OA_w),\qquad
 V_n=\overline{\Span\{A_w'O'y':|w|\le n\}}^{\,w^*}=N_n^\perp.
\end{equation}
The invisible spaces decrease as the test spaces increase. For operations on a set, the analogous equivalence relation is equality of $OA_wx$ for every allowed observation context $w$. For multilinear operations, contexts include all choices of the other arguments. When the allowed operations on an algebra include every left multiplication, every right multiplication, and the involution, a linear invisible class is a two-sided $*$-ideal.
The quotient has a useful universal property. If $(M_s)$ is any closed operationally compatible family with $M_s\subset\ker O_s$, saturation gives $M_s\subset N_s^{(\infty)}$. Hence the identity on each original space induces a unique contractive quotient map
\[
 X_s/M_s\longrightarrow X_s/N_s^{(\infty)},\qquad
 [x]_{M_s}\longmapsto[x]_{N_s^{(\infty)}}.
\]
These maps intertwine every descended basic operation and every contextual observation. Thus $X_s/N_s^{(\infty)}$ is the coarsest compatible quotient retaining all permitted readings, not a claim that the initial readings already compute the saturated ones.

The saturated tests may require additional contextual experiments. They describe the information sufficient for the allowed operations; they are not automatically computable from the initial datum $Ox$. The annihilator description separates this protocol requirement from the subsequent norm cost of decoding.

For a complex algebra, the involution is conjugate-linear. Its descent condition is $N^*=N$; equivalently one may apply the same invariance argument on the underlying real spaces. No complex-linear adjoint of the involution is being used.

\begin{proposition}[Duality of successive layers]\label{prop:layerdual}
If $N_n\subset N_{n-1}$ are closed, restriction gives an isometric identification
\begin{equation}\label{eq:layerdual}
 (N_{n-1}/N_n)'\simeq N_n^\perp/N_{n-1}^\perp=V_n/V_{n-1}.
\end{equation}
Suppose a bounded linear map $J:V_n\to\HH$ is used to represent this dual layer in a Hilbert space. Its orthogonal residual realizes the Banach quotient with equivalent norms precisely when it has a lower bound
\begin{equation}\label{eq:layer-angle}
 \dist(Jf,\overline{J(V_{n-1})})
 \ge c\,\dist(f,V_{n-1}),\qquad c>0.
\end{equation}
\end{proposition}
\begin{proof}
Let $r:V_n\to(N_{n-1}/N_n)'$ be restriction to $N_{n-1}$. Every $f\in V_n=N_n^\perp$ annihilates $N_n$, so this restriction defines a functional on the quotient, and $\ker r=N_{n-1}^\perp=V_{n-1}$. Conversely, a functional on $N_{n-1}/N_n$ lifts isometrically to $N_{n-1}$ and extends to $X$ with unchanged norm by Hahn--Banach. The extension still annihilates $N_n$. Hence $r$ is a metric quotient and gives \eqref{eq:layerdual}.

For the represented layer put $P=P_{\overline{J(V_{n-1})}}$. The map
\[
 \widetilde J:V_n/V_{n-1}\longrightarrow\HH,
 \qquad \widetilde J[f]=(I-P)Jf
\]
is well defined. For any $g\in V_{n-1}$, its norm is at most $\norm J\norm{f-g}$, so it is bounded by $\norm J\norm{[f]}$. Its range is dense in
$\overline{J(V_n)}\ominus\overline{J(V_{n-1})}$: approximate a vector in $\overline{J(V_n)}$ by $Jf_k$ and then apply $I-P$. Condition~\eqref{eq:layer-angle} is exactly the lower bound $\norm{\widetilde J[f]}\ge c\norm{[f]}$. It makes the range closed, since the Banach quotient is complete, and therefore equal to the entire orthogonal difference. Conversely, an isomorphism onto that difference has precisely this lower bound.
\end{proof}

\subsection{Stable inverses and linear retractions}
Let $E:X\to Y$ and $L:X\to Z$ be bounded linear maps, with $Z$ complete. The inequality and factorization
\begin{equation}\label{eq:stablefactor}
 \norm{Lx}\le C\norm{Ex}
 \quad\Longleftrightarrow\quad
 L=DE,\qquad D:\overline{E(X)}\to Z,\quad\norm D\le C
\end{equation}
are equivalent. They are also equivalent to the following lifting property: for each $z'\in Z'$ there is $y'\in Y'$ with $E'y'=L'z'$ and $\norm{y'}\le C\norm{z'}$. Indeed, the inequality makes $D(Ex)=Lx$ well-defined and continuous; each $z'D$ extends by Hahn--Banach; and taking the supremum over $z'$ recovers the inequality. This pointwise lifting property is distinct from a bounded linear choice of all the lifts.

\begin{theorem}[Paired reconstruction and retraction]\label{thm:paired}
Suppose the pairing of $X_0$ and $X_0^\sharp$ norms both spaces, while $Z$ and $Z^\sharp$ satisfy $|[z,w]|\le\norm z\norm w$. Let linear encodings $J:X_0\to Z$ and $K:X_0^\sharp\to Z^\sharp$ have norm bounds $\alpha,\beta$, and suppose
\[
 |\ip{x}{\ell}-[Jx,K\ell]|
 \le\eta\norm x\norm\ell,\qquad0\le\eta<1.
\]
Then
\begin{equation}\label{eq:pairedreverse}
 \frac{1-\eta}{\beta}\norm x\le\norm{Jx}\le\alpha\norm x,
 \qquad
 \frac{1-\eta}{\alpha}\norm\ell\le\norm{K\ell}\le\beta\norm\ell.
\end{equation}
In complete target spaces, both maps extend to isomorphisms from the domain completions onto closed images.

In addition, let $X$ be reflexive, and let $E:X\to Y$ and $F:X'\to Y'$ be bounded linear maps with $E'F=I$. If $J_X,J_Y$ are the canonical bidual embeddings, then
\begin{equation}\label{eq:retract}
 R=J_X^{-1}F'J_Y:Y\to X,
 \qquad RE=I,\qquad\norm R\le\norm F.
\end{equation}
Thus $ER$ is a bounded projection onto $E(X)$. Conversely, a bounded retraction $RE=I$ gives $F=R'$; this converse holds without reflexivity.
\end{theorem}
\begin{proof}
For $x\in X_0$ and $\ell\in X_0^\sharp$, the pairing estimate gives
\[
 |\ip{x}{\ell}|
 \le\norm{Jx}\norm{K\ell}+\eta\norm x\norm\ell
 \le\beta\norm{Jx}\norm\ell+\eta\norm x\norm\ell.
\]
Taking the norming supremum over $\norm\ell\le1$ proves the first lower bound in \eqref{eq:pairedreverse}. Taking the supremum over $\norm x\le1$ proves the second. The upper bounds are the assumed bounds for the encodings. Thus the encodings and their inverses on the images are uniformly continuous. They extend to the completions, and the lower bounds make their extended images closed.

For the retraction, $F'J_Y$ maps $Y$ into $X''$. Reflexivity is used exactly here: it allows us to regard this map as an $X$-valued map by applying $J_X^{-1}$. For $x\in X$ and $\ell\in X'$,
\[
 \ip{F'J_YEx}{\ell}=(F\ell)(Ex)=(E'F\ell)(x)=\ell(x).
\]
Consequently $RE=I$, and the isometry of the canonical embeddings gives $\norm R\le\norm F$. The operator $ER$ is bounded, has range $E(X)$, and satisfies $(ER)^2=ER$. Conversely, if $RE=I$, taking adjoints gives $E'R'=I$; this direction uses no reflexivity.
\end{proof}

A retraction gives an actual output for arbitrary noisy data: $\norm{R(Ex+e)-x}\le\norm F\norm e$. A closed-image inverse alone is defined only on compatible data.

For an encoded operation, suppose $R_YJ_Y=I$ and put $\mathfrak e=J_YA-BJ_X$. Then
\begin{equation}\label{eq:operator-defect}
 R_YBJ_X-A=-R_Y\mathfrak e,\qquad
 (R_YBJ_X-A)'=-\mathfrak e'R_Y'.
\end{equation}
The same compatibility defect therefore controls both the reconstructed output and the pulled-back test. If the maps are bounded only on an encoded image, the dual in this identity is the dual of that image.

\subsection{The residual space and its universal inverse}
The Hilbert geometry used below is a distance construction. Its formulation does not require a closed range or a bounded pseudoinverse.

\begin{lemma}[Relative residual and elimination in stages]\label{lem:relative_residual}
Let $V:E_0\to\mathcal K$ and $W:E_1\to\mathcal K$ be bounded linear maps between Hilbert spaces. Put
\[
 P=P_{\overline{\ran V}},\qquad R_{V,W}=(I-P)W,
 \qquad \mathfrak S_{V,W}=R_{V,W}^*R_{V,W}.
\]
Then
\begin{equation}\label{eq:residual_distance}
 \norm{R_{V,W}y}^2=\inf_{x\in E_0}\norm{Vx+Wy}^2.
\end{equation}
The completion of $E_1/\ker R_{V,W}$ in the norm $[y]\mapsto\norm{R_{V,W}y}$ is canonically $\overline{\ran R_{V,W}}$. For a bounded linear map $L:E_1\to Z$ into a Banach space, the following are equivalent, with the same least constant $C$:
\begin{equation}\label{eq:residual_factorization}
 \begin{gathered}
 \norm{Ly}\le C\norm{Vx+Wy}\quad(x\in E_0,\ y\in E_1),\\
 L=D R_{V,W},\qquad
 D:\overline{\ran R_{V,W}}\to Z\ \text{bounded},\quad\norm D\le C.
 \end{gathered}
\end{equation}
The map $D$, when it exists, is unique.

For an additional bounded column map $U:E_2\to\mathcal K$, let $Q$ be the projection onto $\overline{\ran R_{V,W}}$. Then
\begin{equation}\label{eq:residual_staged}
 P_{\overline{\ran[V,W]}}=P+Q,
 \qquad
 (I-P_{\overline{\ran[V,W]}})U=(I-Q)(I-P)U.
\end{equation}
Thus removing the combined range of $V,W$ is the same as first removing $V$ and then the residual range of $W$.
\end{lemma}
\begin{proof}
The infimum in \eqref{eq:residual_distance} is the squared distance of $-Wy$ to $\ran V$, equivalently to its closure. Orthogonal projection computes this distance. The map $[y]\mapsto R_{V,W}y$ is therefore an isometry of the stated normed quotient onto $\ran R_{V,W}$, and extends to the asserted completion.

Taking the infimum in $x$ turns the first inequality in \eqref{eq:residual_factorization} into $\norm{Ly}\le C\norm{R_{V,W}y}$. It makes $D(R_{V,W}y)=Ly$ well defined and bounded. Completeness of $Z$ extends $D$ uniquely to the closure without changing its norm. Conversely, a factorization gives the inequality because an orthogonal residual has norm at most $\norm{Vx+Wy}$. This also identifies the least constant.

Let $\mathcal M=\overline{\ran[V,W]}$. It contains $\ran P$ and each $Wy$, hence also $(I-P)Wy$. Conversely every $Vx+Wy$ is a sum of a vector in $\ran P$ and one in $\overline{\ran R_{V,W}}$. These two closed subspaces are orthogonal, so
$\mathcal M=\ran P\oplus\ran Q$. Consequently $PQ=QP=0$, which gives both identities in \eqref{eq:residual_staged}.
\end{proof}

In Gram coordinates this is the distance form of shorting~\cite{AT,Ando}. It identifies the residual norm, not an independently specified target norm: the latter is recovered only when the factorization bound holds. The projection onto a closed range also matters. For example, on $\ell^2(\N)$ let $Ve_j=j^{-1}e_j$ and $w=(j^{-1})_j$. The map $O(x,t)=Vx+tw$ is injective, because $w\notin\ran V$, but $\overline{\ran V}=\ell^2$. The target $t$ is therefore algebraically determined and has zero residual norm. No bounded inverse for this target exists. The nuisance direction can approximate cancellation without achieving it.

\subsection{Eliminating target-null observations}
Let $X,Y$ be Hilbert spaces, let $O\in B(X,Y)$, and let $N\subset X$ be a closed subspace of directions forgotten by the target. Set
\[
 H=N^\perp,\qquad Y_N=\overline{ON},\qquad
 O_N=(I-P_{Y_N})O|_H:H\longrightarrow Y_N^\perp.
\]
The quotient $X/N$ is identified isometrically with $H$ by $[x]\mapsto P_Hx$. Since nuisance vectors in $N$ are unrestricted,
\begin{equation}\label{eq:target_short_variation}
 \inf_{n\in N}\norm{O(h+n)}^2
 =\dist(Oh,Y_N)^2=\norm{O_Nh}^2\qquad(h\in H).
\end{equation}
Consequently the effective Gram operator is
\begin{equation}\label{eq:target_short_gram}
 S_N=O_N^*O_N=\operatorname{Short}_{H}(O^*O)|_H.
\end{equation}
Indeed, in the decomposition $X=H\oplus N$, the variational formula identifies $S_N\oplus0$ as the largest positive operator below $O^*O$ whose range is contained in $H$. This definition applies for arbitrary $ON$. If
$O^*O=\left(\begin{smallmatrix}A&B\\B^*&C\end{smallmatrix}\right)$
in $H\oplus N$ and $C$ has closed range, then $S_N=A-BC^\dagger B^*$.

A bounded positive operator $W$ on $H$ specifies the target seminorm
\[
 \norm{x}_{N,W}=\norm{W^{1/2}P_Hx}.
\]
After quotienting its zero space, its Hilbert completion is canonically
\begin{equation}\label{eq:target_weight_completion}
 \widehat{\mathcal Q}_{N,W}\simeq\overline{\ran W^{1/2}}\subset H.
\end{equation}
The identification sends the class of $x$ to $W^{1/2}P_Hx$ and equips the target with the weighted norm induced by $W$.

\begin{theorem}[Sharp recovery of a weighted target quotient]\label{thm:target_short}
The following conditions are equivalent:
\begin{enumerate}[label=\textup{(\roman*)}]
\item there is $C<\infty$ such that
$\norm{W^{1/2}P_Hx}\le C\norm{Ox}$ for every $x\in X$;
\item $W\le C^2S_N$ for some $C<\infty$;
\item $\ran W^{1/2}\subset\ran O_N^*$;
\item there is a unique bounded linear map
$D:\overline{\ran O}\to\widehat{\mathcal Q}_{N,W}$ satisfying
$D(Ox)=W^{1/2}P_Hx$ for every $x\in X$.
\end{enumerate}
Define
\begin{equation}\label{eq:target_weight_cost}
 \kappa_{N,W}:=
 \begin{cases}
 \displaystyle\sup_{Ox\ne0}
 \frac{\norm{W^{1/2}P_Hx}}{\norm{Ox}},
 & \ker O\subset\ker(W^{1/2}P_H),\\[2mm]
 +\infty,&\text{otherwise}.
 \end{cases}
\end{equation}
The supremum of an empty set of nonnegative ratios is taken to be zero. When the equivalent conditions hold, the common optimal constant satisfies
\begin{equation}\label{eq:target_weight_cost_equiv}
 \kappa_{N,W}
 =\inf\{C\ge0:W\le C^2S_N\}
 =\norm D.
\end{equation} For every $\delta>0$,
\begin{equation}\label{eq:target_unrestricted_risk}
 \inf_{\widehat D:Y\to\widehat{\mathcal Q}_{N,W}}
 \sup_{x\in X,\ \norm e\le\delta}
 \norm{\widehat D(Ox+e)-W^{1/2}P_Hx}
 =\delta\kappa_{N,W},
\end{equation}
where the infimum is over arbitrary deterministic reconstruction maps. If $W=0$, both sides are zero.

For $W=I_H$ and $H\ne\{0\}$, these conditions are equivalent to
$S_N\ge m^2I_H$ for some $m>0$, and also to $H\subset\ran O^*$. In this case
$\kappa_{N,I}=\norm{S_N^{-1/2}}$.
\end{theorem}
\begin{proof}
Write $x=h+n$ in $H\oplus N$. Apply Lemma~\ref{lem:relative_residual} with nuisance columns $O|_N$ and target columns $O|_H$. It gives
\[
 \inf_{n\in N}\norm{O(h+n)}=\norm{O_Nh},\qquad
 \overline{\ran O}=Y_N\oplus\overline{\ran O_N}.
\]
Thus (i) is equivalent to the reduced bound
\begin{equation}\label{eq:target_reduced_bound}
 \norm{W^{1/2}h}\le C\norm{O_Nh}\qquad(h\in H).
\end{equation}
Squaring proves (i)$\Leftrightarrow$(ii). Douglas' range-inclusion theorem, applied to $W^{1/2}$ and $O_N^*$, gives (ii)$\Leftrightarrow$(iii). This is actual range inclusion, not merely inclusion of the closed ranges.

The same reduced bound defines a unique map
\[
 L:\overline{\ran O_N}\longrightarrow\overline{\ran W^{1/2}},
 \qquad L(O_Nh)=W^{1/2}h.
\]
Its norm is the least constant in \eqref{eq:target_reduced_bound}. On the orthogonal decomposition of $\overline{\ran O}$ above, define $D(y_N+z)=Lz$. This gives (iv), $\norm D=\norm L$, and $D(Ox)=W^{1/2}P_Hx$. Uniqueness follows from the density of $\ran O$. Conversely (iv) immediately gives (i). The least constant in (i) is the ratio in \eqref{eq:target_weight_cost}, so \eqref{eq:target_weight_cost_equiv} follows. This also constructs the decoder on every compatible limiting datum, not only on $\ran O$.

Extend $D$ to all of $Y$ by zero on $(\overline{\ran O})^\perp$. The resulting decoder has the same norm and proves the upper risk bound $\delta\kappa_{N,W}$. For the lower bound choose $x$ with $Ox\ne0$ and rescale it to $\norm{Ox}=\delta$. The states $x,-x$, with errors $-Ox,Ox$, give the same datum zero. Every deterministic estimator must therefore miss one of their targets by at least $\norm{W^{1/2}P_Hx}$. Taking the supremum of the ratios proves equality. If the ratios are unbounded, this lower bound is infinite. If some $x\in\ker O$ has a nonzero target, arbitrary scaling of $x$ gives infinite risk even without noise. The zero target is immediate.

For $W=I_H$, domination is exactly $S_N\ge m^2I_H$. Applying Douglas directly to $P_H$ and $O$ gives the equivalent condition $H\subset\ran O^*$. For nonzero $H$, the optimal lower bound is $m=\norm{S_N^{-1/2}}^{-1}$, proving the last assertion.
\end{proof}

The decoder factors the observation through the target completion and may have a nontrivial kernel, reflecting observation directions discarded by the target. If $O$ is compact, an unweighted Hilbert quotient with bounded inverse is necessarily finite-dimensional: otherwise the compact operator $O_N$ would be bounded below on an infinite-dimensional space, forcing the identity on $H$ to be compact. Weighted targets can remain infinite-dimensional.

\subsection{Finite decoding and target-visible tails}
Let $X,Z$ now be Banach spaces, let $T\in B(X,Z)$ be a \emph{bounded linear} target, and let $M_a:X\to Y_a$, $J_a:Y_a\to X$ be bounded linear finite observation and reconstruction maps. Write $\Pi_a=J_aM_a$. The target estimate below depends on the projected residual $T(I-\Pi_a)$ rather than on full-state convergence of $\Pi_a$.
\begin{proposition}[Target conditioning and transverse tails]\label{prop:target_tail}
For $y=M_ax+e$, $\norm e\le\delta$, the output $TJ_ay$ satisfies
\begin{equation}\label{eq:target_tail_bound}
 \norm{TJ_ay-Tx}\le\norm{TJ_a}\delta+\norm{T(I-\Pi_a)x}.
\end{equation}
Consequently, uniform recovery on a class $\mathcal K$ follows from
$\sup_{x\in\mathcal K}\norm{T(I-\Pi_a)x}\to0$ together with vanishing target-amplified noise. For the quotient target $T:X\to X/N$, the tail is exactly
$\dist((I-\Pi_a)x,N)$.
\end{proposition}
\begin{proof}
Linearity gives $TJ_ay-Tx=TJ_ae-T(I-\Pi_a)x$. Take norms; the quotient-norm formula gives the final assertion.
\end{proof}

For the Gaussian finite inverse of Theorem~\ref{thm:finitedecoder}, this becomes
\[
 \norm{TD_a\widetilde z-TK}
 \le\norm{TD_a}\,\delta+\norm{T(K-P_aK)}.
\]
Here $T$ must be bounded on the chosen coefficient space, with any required input radius specified. The spatial and source defects in Section~\ref{sec:data} imply this condition, while a target that annihilates unresolved directions can admit a strictly smaller tail certificate. For a nonlinear target of the form $\Phi\circ T$, compose \eqref{eq:target_tail_bound} with a uniform modulus of $\Phi$ on the relevant target set. General nonlinear invariants are handled by the fibre criterion below.

\begin{proposition}[Stable quotients without stable full reconstruction]\label{prop:target_strict_separation}
Let $G\ne\{0\}$ be a Hilbert space and, on $X=\ell^2(\mathbb N_0;G)$, fix $0<\varrho<1$ and
$O_\varrho x=(\varrho^jx_j)_{j\ge0}$. On the prior ball $\norm x\le M$, with $M>0$ and deterministic data noise at most $\delta>0$, the optimal full-state risk and the optimal risk for the first $r+1$ coordinates are
\begin{equation}\label{eq:target_strict_risks}
 \mathcal R^{\rm full}(M,\delta)=M,\qquad
 \mathcal R_r^{\rm quot}(M,\delta)=\min\{M,\varrho^{-r}\delta\}.
\end{equation}
Nevertheless $O_\varrho$ is injective.
\end{proposition}
\begin{proof}
All weights are positive, proving injectivity. Inverting the first $r+1$ weights has norm $\varrho^{-r}$, while the zero decoder has error at most $M$. For a unit $v\in G$ use
$x=\min\{M,\varrho^{-r}\delta\}e_r\otimes v$.
The states $\pm x$ give common zero data under errors $\mp O_\varrho x$ of norm at most $\delta$; their half-separation proves the quotient lower bound. For full recovery, choose $j$ with $M\varrho^j\le\delta$ and use $x=Me_j\otimes v$. This gives the lower bound $M$, attained by the zero decoder.
\end{proof}

This is a first-chaos operator model: with $G=\Sch_2(E)$, the series
$F_x=\sum_jg_jx_j$ converges in $L^2(\Sch_2)$, and
$\varrho^j\E[g_jF_x]=\varrho^jx_j$.
Thus the strict separation holds within the same Gaussian observation system. The diagonal observation is compact when $G$ is finite-dimensional, but not when $G$ is infinite-dimensional; compactness is not used in \eqref{eq:target_strict_risks}.

The weighted target $T_\tau x=(\tau^jx_j)$, $0<\tau\le1$, is globally Lipschitz recoverable exactly when $\tau\le\varrho$, since
$\sup_j(\tau/\varrho)^j<\infty$ is precisely the operator domination in Theorem~\ref{thm:target_short}. Thus $\tau=1$ separates algebraic identifiability from stable recovery.

\subsection{Nonlinear observation fibres and uniform completion}
Let $\mathcal Q$ be a target set, possibly nonlinear, with a prescribed separated uniform structure generated by countably many pseudometrics $d_\alpha$. Let $j_i:\mathcal Q\to Z_i$ be observations into metric spaces, and write $J_m=(j_1,\ldots,j_m)$ with the finite maximum metric $q_m$. The complete coordinate family is required to separate the target classes themselves.

\begin{theorem}[Finite inverse criteria and invariant completion]\label{thm:invariant_uniform}
Suppose each $j_i$ is uniformly continuous, the complete family separates $\mathcal Q$, and, for every $\alpha$ and $\varepsilon>0$, there are $m$ and $\eta>0$ such that
\begin{equation}\label{eq:invariant_uniform_inverse}
 q_m(J_mu,J_mv)<\eta\quad\Longrightarrow\quad
 d_\alpha(u,v)<\varepsilon\qquad(u,v\in\mathcal Q).
\end{equation}
Then $J=(j_i)_i$ is a uniform isomorphism of $\mathcal Q$ onto $J(\mathcal Q)$. If $\widehat Z_i$ denotes the metric completion of $Z_i$ and the product is equipped with its product uniformity, then $J$ extends uniquely to a uniform isomorphism
\begin{equation}\label{eq:invariant_completion_iso}
 \widehat{\mathcal Q}\ \simeq\
 \overline{J(\mathcal Q)}^{\,\prod_i\widehat Z_i}.
\end{equation}
A sufficient quantitative hypothesis for \eqref{eq:invariant_uniform_inverse} is, on the specified certificate class,
\begin{equation}\label{eq:invariant_inverse_tail}
 d_\alpha(u,v)\le\omega_\alpha\!\big(q_{m}(J_mu,J_mv)+b_{\alpha,m}\big),
 \qquad b_{\alpha,m}\longrightarrow0,
\end{equation}
where $\omega_\alpha$ is nondecreasing and $\omega_\alpha(t)\to0$ as $t\downarrow0$.

For a nonempty fibre of candidates whose $J_m$-data lie within $\delta$ of one observed datum, every pair of candidates has data discrepancy at most $2\delta$. Hence any pairwise inverse estimate yields a bound for the target diameter of the fibre. If a sequence of fibres contains one common target and its target diameters tend to zero, every candidate selection converges to that target. If the fibres are nested, nonempty, and their diameters tend to zero in every $d_\alpha$, every candidate selection is Cauchy and all selections determine the same point of $\widehat{\mathcal Q}$.
\end{theorem}
\begin{proof}
The product uniformity is the uniformity of finite experiments: each basic entourage imposes tolerances on only finitely many coordinates. Forward uniform continuity of every $j_i$ therefore makes $J$ uniformly continuous. Separation makes $J$ injective. Condition~\eqref{eq:invariant_uniform_inverse} says exactly that each target entourage contains the inverse image of one finite-data entourage, and hence that $J^{-1}$ is uniformly continuous on $J(\mathcal Q)$.

A uniform isomorphism preserves Cauchy filters and their equivalence. The product of the complete spaces $\widehat Z_i$ is complete, and the completion of its uniform subspace $J(\mathcal Q)$ is its closure. Extending $J$ and its inverse to completions proves \eqref{eq:invariant_completion_iso}. To check the sufficient quantitative condition, first choose $m$ so that $b_{\alpha,m}$ is within the small-argument range of $\omega_\alpha$, and then choose the tolerance of that finite experiment. This proves \eqref{eq:invariant_uniform_inverse} without interchanging observation order and noise precision.

For the fibre assertions, two candidates within $\delta$ of the same observed datum have finite-data discrepancy at most $2\delta$. A pairwise inverse bound therefore controls the whole target diameter, independently of the candidate-selection rule. If the true target belongs to every fibre, its distance from any selected candidate is bounded by that diameter. If instead the fibres are nested, two selections made after stage $N$ both belong to the $N$th fibre. Vanishing diameters in every $d_\alpha$ make all such selections Cauchy and make any two selections equivalent. Their common limit belongs to the completed target; membership in the original target requires its closedness.
\end{proof}

\begin{proposition}[Compact certificate classes and optimal recovery]\label{prop:invariant_compact}
Let $(\mathcal K,d)$ be a nonempty compact metric target class and let $j_i$ be continuous real coordinates. Define
\begin{equation}\label{eq:invariant_fibre_modulus}
 \omega_m(t)=\max\{d(u,v):u,v\in\mathcal K,\
                 \norm{J_mu-J_mv}_\infty\le t\}.
\end{equation}
The coordinates separate $\mathcal K$ if and only if, for every $\varepsilon>0$, some $m$ and $t>0$ satisfy $\omega_m(t)<\varepsilon$. In that case $J(\mathcal K)$ is compact and closed. For deterministic noise in $\R^m$ with $\ell^\infty$ norm at most $\delta\ge0$, define
\[
 \mathcal R_m(\delta)=\inf_{D:\R^m\to\mathcal K}
 \sup_{u\in\mathcal K,\ \norm e_\infty\le\delta}
 d\big(D(J_mu+e),u\big).
\]
Then
\begin{equation}\label{eq:invariant_compact_risk}
 \tfrac12\omega_m(2\delta)\le\mathcal R_m(\delta)\le\omega_m(2\delta).
\end{equation}
When the coordinates separate $\mathcal K$, for every $\varepsilon>0$ there are finitely many polynomials $P_i$ in a common $J_m$ and a finite $L_\varepsilon$ such that
\begin{equation}\label{eq:invariant_norming_tests}
 d(u,v)\le4\varepsilon+\max_i|P_i(J_mu)-P_i(J_mv)|
 \le4\varepsilon+L_\varepsilon\norm{J_mu-J_mv}_\infty.
\end{equation}
\end{proposition}
\begin{proof}
If uniform recovery fails despite separation, choose $u_m,v_m$ with $d(u_m,v_m)\ge\varepsilon_0$ and $\norm{J_mu_m-J_mv_m}\le1/m$. A common convergent subsequence has limits agreeing in every fixed coordinate, but at distance at least $\varepsilon_0$, a contradiction. The reverse implication is immediate. Since a continuous injection from a compact space to a Hausdorff space is a homeomorphism onto its image, the inverse is uniformly continuous on the compact image.

For the risk upper bound, if the datum lies within $\delta$ of $J_m(\mathcal K)$, select one feasible candidate; the true and selected targets have $J_m$-distance at most $2\delta$ and hence target distance at most $\omega_m(2\delta)$. On infeasible data the decoder may be defined arbitrarily. For the lower bound, a pair realizing the maximum can produce its coordinate midpoint with noise at most $\delta$. By the triangle inequality any estimate misses one target by at least half their distance. This is the standard two-point optimal-recovery argument~\cite{MR,EF}.

Finally choose a finite $\varepsilon$-net $(z_i)$ and set $f_i(u)=d(u,z_i)$. For some $i$, $|f_i(u)-f_i(v)|\ge d(u,v)-2\varepsilon$. The algebra generated by the real separating coordinates contains constants. Stone--Weierstrass approximates each of the finitely many $f_i$ uniformly within $\varepsilon$ by a polynomial in a common finite $J_m$. This gives the first inequality in \eqref{eq:invariant_norming_tests}. Bound the finitely many polynomial gradients on a rectangle containing $J_m(\mathcal K)$ to obtain the second.
\end{proof}

Proposition~\ref{prop:invariant_compact} is qualitative: it does not specify the observation order or the constant $L_\varepsilon$. Section~\ref{sec:pfaffian} supplies these inverse controls from explicit finite-moment estimates and then derives compactness from the same spectral envelope.

\subsection{Which operations survive the target quotient?}
Recovery of a target is separate from closure of its operations. If $A\in B(X)$ preserves $N$, its quotient action is $\overline A=P_HA|_H$. It extends to \eqref{eq:target_weight_completion} whenever
\begin{equation}\label{eq:target_operation_weight}
 \overline A^*W\overline A\le C_A^2W.
\end{equation}
Indeed this inequality preserves the additional zero seminorm space and bounds the induced map by $C_A$ before completion. The extension is unique by density. For multilinear operations, Proposition~\ref{prop:congruence} supplies algebraic descent and the corresponding target-norm estimate supplies continuous extension. For the nonlinear spectral quotient of Section~\ref{subsec:pfaffian_cross}, addition fails to descend in general. Continuous operations on a recovered quotient are therefore obtained by verifying congruence and the corresponding target-norm estimate.


\section{Gaussian coefficient geometry}\label{sec:coeff}
We now realize the abstract inverse mechanism in coefficient spaces defined before random evaluation. Directed cuts specify their norm and duality, Gaussian pairings give a norming lift, and same-source polarization transfers the inverse to Wick chaoses. These analytic estimates recover the original Banach geometry; a Hilbert Gram representation is not substituted for it. The exponential degree costs are the input needed for the all-radius algebra in Section~\ref{sec:product}.

\subsection{Directed cuts and coefficient spaces}
Let $H_1,\ldots,H_m$ be complexifications of real separable Hilbert spaces, and let $\cC,\cE$ be the physical input and output spaces. On the finite-support, finite-matrix core, write
\[
 K=\sum_Jh_{1,j_1}\otimes\cdots\otimes h_{m,j_m}\otimes K_J,
 \qquad K_J:\cC\to\cE.
\]
For $S\subset[m]$, let $H_S=\bigotimes_{j\in S}H_j$, with empty tensor product $\C$. Placing the indices $(J_S,\cC)$ on the input side gives
\begin{equation}\label{eq:cuts}
 \Flat_S(K):\overline{H_S}\otimes_2\cC\longrightarrow H_{S^c}\otimes_2\cE,
 \qquad \Prof_{m,s}(K)=\max_{S\subset[m]}\norm{\Flat_S(K)}_s.
\end{equation}
At $s=2$, all these norms are the norm of the same Hilbert tensor, under unitary rearrangement.

\begin{definition}[Coefficient geometry]\label{def:cutspace}
The space $\Coeff_{m,s}$ is the completion of the core in the norm
\begin{equation}\label{eq:cutnorm}
 \norm K_{\Coeff_{m,s}}=
 \begin{cases}
 \displaystyle\inf_{K=\sum_SK^S}\sum_{S\subset[m]}\norm{\Flat_S(K^S)}_s,&1\le s<2,\\[2mm]
 \norm K_{\bigotimes_jH_j\otimes_2\overline{\cC}\otimes_2\cE},&s=2,\\[1mm]
 \displaystyle\max_{S\subset[m]}\norm{\Flat_S(K)}_s,&2<s<\infty.
 \end{cases}
\end{equation}
When all source legs are labeled copies of one real Hilbert space, $\Wc_{m,s}$ is the closed image of source symmetrization, with the inherited norm. At degree zero, the space is $\Sch_s(\cC,\cE)$.
\end{definition}

Finite-factor compressions are contractions on every cut. Source permutations permute the cuts, so their average is a contractive symmetrization projection. Physical adjoints exchange the input and output spaces and complement the cuts. The high-exponent completion can contain compatible Schatten arrays whose total Hilbert tensor norm is infinite.

\begin{lemma}[Polarity and reflexivity]\label{lem:polarity}
For $1<s<\infty$ and $s'=s/(s-1)$, the core pairing
\begin{equation}\label{eq:coeffpair}
 \langle K,L\rangle=\sum_J\Tr(K_JL_J^*)
 =\Tr\big(\Flat_S(K)\Flat_S(L)^*\big)
\end{equation}
is independent of $S$. Its continuous extension identifies $\Coeff_{m,s'}$ isometrically with the conjugate dual of $\Coeff_{m,s}$. Both spaces are reflexive, as are their symmetric subspaces. Common finite-factor compressions converge strongly to the identity on both sides.
\end{lemma}
\begin{proof}
We first construct the low-exponent space and identify its dual. For $1<s<2$, pull the Schatten space of each cut back to a tensor space $X_{S,s}$. Each embeds continuously in the common Hilbert tensor space $\mathscr K$, since $\norm A_2\le\norm A_s$. The addition map
\[
 \Sigma:\bigoplus_S^{\ell^1}X_{S,s}\longrightarrow\mathscr K,
 \qquad (K^S)_S\longmapsto\sum_SK^S
\]
has closed kernel. Its quotient is a Banach space with the infimum norm in \eqref{eq:cutnorm}.

This quotient agrees with the completion of the finite core. To verify the point that matters here, take a decomposition of a core tensor and apply one sufficiently large common finite-factor compression. It fixes the total tensor, contracts every summand, and turns the decomposition into a core decomposition with no larger cost. Conversely, for an arbitrary quotient element, compress a near-minimal decomposition in each of its finitely many summands. These compressions converge in their Schatten norms. Thus the core has the stated norm and is dense in the quotient.

The dual of the finite $\ell^1$ sum is the $\ell^\infty$ sum of the conjugate-dual Schatten spaces. A family $(B_S)_S$ defines a quotient functional precisely when it annihilates $\ker\Sigma$. Testing a tensor and its negative in the two positions $S,S_0$ gives
\[
 \Tr(\Flat_S(K)B_S^*)
 =\Tr(\Flat_{S_0}(K)B_{S_0}^*)
 \qquad(K\text{ in the core}).
\]
These are exactly the compatibility conditions between the cuts. Indeed, compress all the $B_S$ by the same finite factors and recover a finite tensor $L_n$ from the $S_0$-cut. Testing matrix units forces its other cuts to be the other compressed $B_S$. Since $s'<\infty$, those compressions converge in every Schatten-$s'$ norm. There are finitely many cuts, so $(L_n)$ converges in their maximum norm. Its limit represents the original functional. Conversely, every compatible limit annihilates $\ker\Sigma$. The quotient-dual norm is the maximum of the cut norms, proving the isometric identification.

For a completed low-exponent kernel, the resulting pairing is interpreted as
\[
 \langle K,L\rangle
 =\sum_S\Tr\big(\Flat_S(K^S)\Flat_S(L)^*\big),
 \qquad K=\sum_SK^S.
\]
Its value is independent of the decomposition by the annihilator condition, and its absolute value is at most
$\sum_S\norm{\Flat_S(K^S)}_s\norm L_{\Coeff_{m,s'}}$.
In particular, this formula does not require a low-exponent sum-space element to belong to every individual cut space.

Finite sums and closed quotients of reflexive Schatten spaces are reflexive. Biduality now gives the reverse identification. At $s=2$ all statements are Hilbert-space identities. Source symmetrization is a contractive projection self-adjoint for the coefficient pairing, so restricting to its range proves the symmetric version. Finally, common compressions converge in each of the finitely many cuts in the high branch and in each summand of a quotient decomposition in the low branch. This proves the asserted strong approximation on both sides.
\end{proof}

For example, the degree-one coefficients $K_j=e_je_1^*$, $1\le j\le N$, have cut norms $\sqrt N$ and $N^{1/s}$. For $s<2$, their quotient-sum norm is $\sqrt N$: one cut gives the upper bound, and the common Hilbert norm gives the lower bound. Thus the maximum-cut norm would give a genuinely different low-exponent geometry.

\subsection{Forward estimates and closed-image recovery}
For independent standard real Gaussian families $g_{\nu,j}$, define on the core
\begin{equation}\label{eq:dec-eval}
 E_mK=T_K=\sum_{j_1,\ldots,j_m}
       \left(\prod_{\nu=1}^mg_{\nu,j_\nu}\right)K_{j_1\cdots j_m}.
\end{equation}
The same norm estimates hold for standard circular complex Gaussians normalized by $\E|g|^2=1$.

\begin{lemma}[Forward bounds]\label{lem:forward}
For $1\le s\le2$,
\begin{equation}\label{eq:lowforward}
 \norm{T_K}_{L^2(\Sch_s)}\le\norm K_{\Coeff_{m,s}}.
\end{equation}
For $p,s\ge2$, an absolute constant $C$ gives
\begin{equation}\label{eq:highforward}
 \norm{T_K}_{L^p(\Sch_s)}
 \le C^m(p+s)^{m/2}\max_S\norm{\Flat_S(K)}_s.
\end{equation}
The constants are independent of the source and physical dimensions.
\end{lemma}
\begin{proof}
Fix a cut. At $s=1$, take a singular-value decomposition
$\Flat_SK=\sum_\ell a_\ell v_\ell u_\ell^*$.
Evaluate the source variables on its two sides. The resulting random vectors $V_\ell,U_\ell$ depend on disjoint Gaussian families, and each has second moment one. Hence
\[
 \norm{V_\ell U_\ell^*}_{L^2(\Sch_1)}^2
 =\E(\norm{V_\ell}^2\norm{U_\ell}^2)=1.
\]
Minkowski gives the trace-norm bound $\sum_\ell a_\ell$. At $s=2$, Gaussian orthogonality gives an isometry. Interpolation between rectangular Schatten spaces and Bochner $L^2$ spaces~\cite{BL} yields the single-cut estimate for $1\le s\le2$. Apply it to each summand of a near-minimal cut decomposition and use the triangle inequality to obtain~\eqref{eq:lowforward}.

For one Gaussian leg, the rectangular noncommutative Khintchine inequality~\cite{LP,Buchholz} is
\[
 \left\|\sum_jg_jA_j\right\|_{L^p(\Sch_s)}
 \le C\sqrt{p+s}\max\left\{
 \left\|(\sum_jA_j^*A_j)^{1/2}\right\|_s,
 \left\|(\sum_jA_jA_j^*)^{1/2}\right\|_s\right\}.
\]
The rectangular form follows by a corner embedding or self-adjoint dilation. For $p\le s$, the $L^s$ bound and monotonicity of probability norms suffice; for $p>s$, Gaussian Banach-valued moment comparison supplies the factor $\sqrt{p/s}$.

Set $q=\max(p,s)$. Condition on all but one source leg and apply the one-leg inequality at probability exponent $q$. The two square functions are precisely the row and column linearizations of the remaining rectangular array. In $L^q$, the maximum of these two branches is bounded by their $q$-sum. Repeating for each leg gives at most one absolute factor per step. The $2^m$ terminal row/column choices are the cuts $\Flat_SK$. Since $q+s\le2(p+s)$, all absolute factors are absorbed into $C^m$, proving~\eqref{eq:highforward}.
\end{proof}

Fix a sufficiently large absolute constant $C_0$, and set
\begin{equation}\label{eq:as}
 a_s=\begin{cases}1,&1\le s\le2,\\ C_0\sqrt{s+2},&s>2.\end{cases}
\end{equation}
Thus $E_m:\Coeff_{m,s}\to L^2(\Sch_s)$ has norm at most $a_s^m$.

\begin{theorem}[Dual recovery and a retraction on the full data space]\label{thm:coeffrecover}
For $1<s<\infty$,
\begin{equation}\label{eq:dec-equivalence}
 a_{s'}^{-m}\norm K_{\Coeff_{m,s}}
 \le\norm{E_mK}_{L^2(\Sch_s)}
 \le a_s^m\norm K_{\Coeff_{m,s}}.
\end{equation}
The image of $E_m$ is closed. There is a bounded linear retraction
$R_{m,s}:L^2(\Sch_s)\to\Coeff_{m,s}$ with $R_{m,s}E_m=I$ and
$\norm{R_{m,s}}\le a_{s'}^m$. Writing $G_J=\prod_{\nu=1}^mg_{\nu,j_\nu}$, its coordinates are
\begin{equation}\label{eq:rawdecoder}
 (R_{m,s}Y)_J=
 \begin{cases}
 \E(G_JY),&\text{real Gaussian sources},\\
 \E(\overline{G_J}Y),&\text{circular complex Gaussian sources}.
 \end{cases}
\end{equation}
Consequently, $E_mR_{m,s}$ is a projection onto the image with norm at most $(a_sa_{s'})^m$.
\end{theorem}
\begin{proof}
For finite kernels, Gaussian orthogonality gives the norming identity
\begin{equation}\label{eq:gaussianpair}
 \langle K,L\rangle=\E\Tr\big((E_mK)(E_mL)^*\big).
\end{equation}
It extends by the forward bounds to $K\in\Coeff_{m,s}$ and $L\in\Coeff_{m,s'}$. Schatten H\"older and probability Cauchy--Schwarz therefore yield
\[
 |\langle K,L\rangle|
 \le a_{s'}^m\norm{E_mK}_{L^2(\Sch_s)}\norm L_{\Coeff_{m,s'}}.
\]
The polarity lemma turns this into the lower estimate in \eqref{eq:dec-equivalence}; the upper estimate is the forward bound. In particular $E_m$ is injective and its image is closed.

To recover arbitrary, possibly incompatible data $Y\in L^2(\Sch_s)$, use the same identity as a definition. The functional
\[
 L\longmapsto\E\Tr\big(Y(E_mL)^*\big)
\]
is bounded on $\Coeff_{m,s'}$ with norm at most $a_{s'}^m\norm Y_{L^2(\Sch_s)}$. Polarity and reflexivity give a unique $R_{m,s}Y\in\Coeff_{m,s}$ representing it. This construction is linear in $Y$, has the claimed norm, and satisfies $R_{m,s}E_mK=K$ by \eqref{eq:gaussianpair}. It is the retraction mechanism of Theorem~\ref{thm:paired}, with the Gaussian evaluation of dual kernels supplying the lifts.

Testing a single source coordinate and a physical matrix unit gives \eqref{eq:rawdecoder}; the adjoint supplies the conjugate in the circular model. The Bochner expectations exist because $G_J\in L^2$. Their assembly into an element of the full coefficient space follows from the bounded functional above, not from coordinatewise integrability. Finally $(E_mR_{m,s})^2=E_mR_{m,s}$, and the product of the two norm bounds gives its projection cost.
\end{proof}

For $s\ge2$, the reverse constant in~\eqref{eq:dec-equivalence} is one. The trace endpoint uses a different predual estimate.

\begin{proposition}[Trace recovery with an exponential degree cost]\label{prop:traceendpoint}
For each $m\ge1$, there is a dimension-independent constant $C_m$ such that
\begin{equation}\label{eq:traceendpoint}
 \norm K_{\Coeff_{m,1}}\le C_m\norm{T_K}_{L^1(\Sch_1)}.
\end{equation}
One may take $C_m=(8\sqrt{2/\pi})^m$ for real sources and $C_m=(8/\sqrt\pi)^m$ for circular sources. Together with~\eqref{eq:lowforward}, this gives closed-image recovery with an exponential, rather than an unspecified, degree cost. Degree zero is the identity map.
\end{proposition}
\begin{proof}
First use independent unit-modulus complex phases. The higher-order predual inequality of Haagerup--Pisier~\cite[Theorem~3.5 and (3.8)--(3.11)]{HP} bounds the sum of all predual cuts by $2^{2m}=4^m$ times the phase $L^1$ norm. For rectangular coefficients, embed $K_J$ in the off-diagonal corner of $B(\cC\oplus\cE)$. This embedding preserves every cut norm. Conversely, compressing an ambient decomposition to that corner contracts every trace cut and gives a rectangular decomposition. Thus the predual sum is exactly the coefficient norm \eqref{eq:cutnorm}, without a dimensional or additional degree factor.

For a standard circular Gaussian write $g=\zeta r$, with independent phase and modulus and $\E r=\sqrt\pi/2$. Conditional expectation over all moduli multiplies every degree-$m$ multilinear coefficient by $(\sqrt\pi/2)^m$. Conditional Jensen therefore gives
\[
 \E\norm{\sum_JG_JK_J}_1
 \ge(\sqrt\pi/2)^m\E\norm{\sum_J\zeta_JK_J}_1.
\]
The circular coefficient inverse consequently has constant $(8/\sqrt\pi)^m$.

To compare with real sources, write each circular leg as $(g_\nu+ih_\nu)/\sqrt2$ and expand the product. The resulting $2^m$ real multilinear arrays have the same law. The triangle inequality gives
$\norm{T_K^{\rm circ}}_{L^1(\Sch_1)}\le2^{m/2}\norm{T_K^{\rm real}}_{L^1(\Sch_1)}$.
This proves the stated real constant. The forward bound is in $L^2(\Sch_1)$, and $L^1\le L^2$ turns the reverse inequality into a reverse Cauchy estimate there. Density then gives closed-image recovery. No bounded retraction on arbitrary trace-class-valued random data is used.
\end{proof}

\subsection{Same-source Wick evaluation}
If $X$ is a Banach-valued homogeneous Gaussian chaos of degree $m$, positivity of the Mehler kernel gives $\norm{M_tX}_B\le M_t(\norm X_B)$. Scalar hypercontractivity~\cite{Janson,Maas} consequently yields
\begin{equation}\label{eq:hyper}
 \norm X_{L^q(B)}\le(q-1)^{m/2}\norm X_{L^2(B)},\qquad q\ge2.
\end{equation}
Apply Paley--Zygmund to $\norm X_B^2$ with threshold one half. Since $\norm X_{L^4}\le3^{m/2}\norm X_{L^2}$, the event
$\norm X_B\ge\norm X_{L^2}/\sqrt2$ has probability at least $1/(4\cdot3^{2m})$. Thus
$\norm X_{L^2}\le4\sqrt2\,3^{2m}\norm X_{L^1}$.
All finite probability norms are therefore equivalent at each fixed degree.

Let $K$ be symmetric, and let $W_1,\ldots,W_m$ be independent isonormal processes. Their normalized average $\overline W=m^{-1/2}\sum_jW_j$ satisfies
\begin{align}
 \E[T_K^{\rm dec}\mid\overline W]
 &=m^{-m/2}I_m(K;\overline W),\label{eq:regression}\\
 T_K^{\rm dec}
 &=\frac{m^{m/2}}{2^mm!}\sum_{\epsilon\in\{-1,1\}^m}
       \Big(\prod_j\epsilon_j\Big)
 I_m\left(K;m^{-1/2}\sum_j\epsilon_jW_j\right).\label{eq:polarization}
\end{align}
The first identity follows by differentiating the conditional Gaussian exponential in $m$ distinct source parameters: the conditional covariance produces exactly the Wick correction. In the second, sign averaging removes every term that does not contain all source labels once; the surviving permutations contribute $m!$. The normalization gives~\eqref{eq:polarization}.

\begin{theorem}[Same-source recovery with exponential degree costs]\label{thm:wickrecover}
For fixed $m$ and $1\le p,s<\infty$, $I_m:\Wc_{m,s}\to L^p(\Sch_s)$ is an isomorphism onto a closed image. For $1\le s<\infty$, there are $u_s,v_s\ge1$ such that
\begin{equation}\label{eq:wick-exponential}
 v_s^{-m}\sqrt{m!}\norm K_{\Wc_{m,s}}
 \le\norm{I_mK}_{L^2(\Sch_s)}
 \le u_s^m\sqrt{m!}\norm K_{\Wc_{m,s}}.
\end{equation}
One may take $u_s=\sqrt e\,a_s$ and $v_s=\sqrt e\,a_{s'}$ for $1<s<\infty$; at $s=2$ the exact choice is $u_2=v_2=1$. At $s=1$, one may take $u_1=\sqrt e$ and $v_1=8\sqrt{2e/\pi}$. For $p,s\ge2$,
\begin{equation}\label{eq:wick-high}
 \norm{I_mK}_{L^p(\Sch_s)}
 \le m^{m/2}C^m(p+s)^{m/2}\max_S\norm{\Flat_SK}_s.
\end{equation}
\end{theorem}
\begin{proof}
First take a finite symmetric kernel. The conditional identity \eqref{eq:regression} and the forward estimate for independent sources give
\[
 \norm{I_mK}_{L^2(\Sch_s)}
 \le m^{m/2}\norm{T_K^{\rm dec}}_{L^2(\Sch_s)}
 \le m^{m/2}a_s^m\norm K_{\Wc_{m,s}}.
\]
For the reverse direction when $1<s<\infty$, every normalized signed source in \eqref{eq:polarization} has the law of the original isonormal process. Applying the triangle inequality to that identity, and then the independent-source inverse, gives
\[
 \norm K_{\Wc_{m,s}}
 \le a_{s'}^m\norm{T_K^{\rm dec}}_{L^2(\Sch_s)}
 \le a_{s'}^m\frac{m^{m/2}}{m!}
                    \norm{I_mK}_{L^2(\Sch_s)}.
\]
Since $m!\ge(m/e)^m$, both degree factors are absorbed by $\sqrt{m!}$ and an exponential factor, yielding \eqref{eq:wick-exponential}. At $s=2$, Gaussian orthogonality gives the exact isometry and removes these comparison losses.

At $s=1$, write $b_1=8\sqrt{2/\pi}$. Proposition~\ref{prop:traceendpoint}, the inclusion $L^1\le L^2$, and the same polarization identity give
\[
 \norm K_{\Wc_{m,1}}
 \le b_1^m\norm{T_K^{\rm dec}}_{L^1(\Sch_1)}
 \le b_1^m\frac{m^{m/2}}{m!}\norm{I_mK}_{L^2(\Sch_1)}.
\]
The preceding forward estimate has $a_1=1$. The factorial inequality therefore gives $u_1=\sqrt e$ and $v_1=\sqrt e\,b_1$ in \eqref{eq:wick-exponential}. Hypercontractivity and the preceding $L^1$--$L^2$ comparison then give all fixed finite probability exponents. The forward bound extends evaluation from the dense core, and the reverse bound turns a Cauchy sequence of evaluated kernels into a Cauchy sequence of original kernels. Its image is therefore closed. Finally, applying \eqref{eq:highforward} before the conditional expectation in \eqref{eq:regression} gives \eqref{eq:wick-high}.
\end{proof}

\begin{proposition}[An intrinsic characterization of the homogeneous image]\label{prop:chaosimage}
Fix source coordinates $g_1,g_2,\ldots$. A source-measurable $T\in L^p(\Sch_s(\cC,\cE))$, $1\le p,s<\infty$, belongs to $I_m(\Wc_{m,s})$ if and only if all its finite physical compressions and finite source conditional expectations are homogeneous degree-$m$ Wick polynomials. With $E_L=\E[\,\cdot\mid g_1,\ldots,g_L]$ and source projection $S_L$, one has
\begin{equation}\label{eq:cutintertwine}
 I_m\big((S_L^{\otimes m}\otimes\overline{P_N}\otimes Q_N)K\big)
 =Q_NE_L(I_mK)P_N.
\end{equation}
Along cofinal finite compressions, both sides converge to the original object.
\end{proposition}
\begin{proof}
For a finite kernel, source conditioning removes precisely the Hermite coordinates involving an index greater than $L$, and physical compression removes precisely the corresponding matrix entries. This proves \eqref{eq:cutintertwine} on the core. Both kernel compressions are contractive and converge strongly, so continuity of evaluation extends the identity to the completion. Each finite matrix-valued homogeneous polynomial space is finite-dimensional and closed, which proves necessity.

For sufficiency, choose increasing finite source and physical projections with
\[
 T_n=Q_{N_n}E_{L_n}(T)P_{N_n}\longrightarrow T\quad\text{in }L^p(\Sch_s).
\] Source measurability gives convergence of conditional expectations; density of finite-rank operators and dominated approximation give convergence of physical compressions. By assumption, each $T_n$ is a finite matrix-valued homogeneous polynomial and therefore equals $I_mK_n$ for a finite symmetric kernel. The inverse estimate of Theorem~\ref{thm:wickrecover}, applied to $T_n-T_k$, makes $(K_n)$ Cauchy in the original coefficient norm. Its limit $K$ satisfies $I_mK=T$ by continuity. The same inverse estimate proves uniqueness. This argument also proves the stated convergence along every cofinal family of compressions.
\end{proof}

Finite sums of degrees are handled by one finite inversion. For $F=\sum_{m=0}^MI_mK_m$, choose $M+1$ distinct Mehler parameters. The Vandermonde matrix $(t_j^m)$ is invertible, so each homogeneous component is a fixed linear combination of the contractions $M_{t_j}F$. Consequently,
\begin{equation}\label{eq:finitechaos}
 \norm{\sum_{m=0}^MI_mK_m}_{L^p(\Sch_s)}
 \asymp_{M,p,s}\sum_{m=0}^M\norm{K_m}_{\Wc_{m,s}}.
\end{equation}
At the Hilbert exponent, all degrees satisfy the exact Fock identity
\begin{equation}\label{eq:hilbertfock}
 \norm{\sum_{m\ge0}I_mK_m}_{L^2(\Sch_2)}^2
 =\sum_{m\ge0}m!\norm{K_m}_2^2.
\end{equation}
The space described by this identity will be distinguished from the stronger multiplication scale in the next section.

\subsection{Sharp growth and indispensable cuts}
\begin{proposition}[Two independent sharp growth parameters]\label{prop:sharpfinite}
For fixed $m\ge1$, the order $(p+s)^{m/2}$ in~\eqref{eq:highforward} cannot be reduced uniformly. There is also no dimension-free Gaussian operator-norm estimate depending only on the maximum operator norm of the cuts.
\end{proposition}
\begin{proof}
A single scalar rank-one coefficient has every cut norm equal to one and evaluates to $\prod_{\nu=1}^mg_\nu$, whose $L^p$ norm is $\norm g_p^m\asymp_m p^{m/2}$. For the Schatten direction, take
\[
 K_N=\sum_{j=1}^Nh_{1,j}\otimes\cdots\otimes h_{m,j}\otimes e_je_j^*.
\]
Every cut has $N$ singular values equal to one. The evaluated operator is diagonal with independent Gaussian products on its diagonal. The law of large numbers, Jensen, and Fatou give
\[
 \lim_{N\to\infty}N^{-1/s}\norm{T_{K_N}}_{L^2(\Sch_s)}
 =\norm g_s^m\asymp_m s^{m/2}.
\]
For $s\ge2$, uniform integrability follows by Jensen from a slightly higher moment of one diagonal entry. The two lower bounds together have order $(p+s)^{m/2}$.

For the same diagonal example, the largest diagonal entry has size $(\log N)^{m/2}$. A lower bound follows by requiring all $m$ factors in one coordinate to exceed a sufficiently small multiple of $\sqrt{\log N}$, and using independence across coordinates. The $p$ lower bound follows from one coordinate. Choosing $s\asymp\log(eN)$ in~\eqref{eq:highforward} proves the matching upper bound
\[
 \norm{T_{K_N}}_{L^p(B(E))}\asymp_m(p+\log(eN))^{m/2},
\]
while every operator-norm cut equals one.
\end{proof}

The same-source comparison also has a sharp degree cost. The leading term of $H_m(x)$ is $x^m$, and its lower-degree remainder has $L^p$ norm $O_m(p^{(m-2)/2})$. The Gaussian moment formula and Stirling's formula yield
\[
 \lim_{p\to\infty}\frac{\norm{H_m(g)}_p}
 {\norm{\prod_{j=1}^mg_j}_p}=m^{m/2}.
\]
Thus the constant $m^{m/2}$ in the universal increasing-convex comparison supplied by~\eqref{eq:regression} is optimal.

For decoupled circular complex sources there is an exact fourth-moment identity:
\begin{equation}\label{eq:fourthmoment}
 \E\norm{T_K}_4^4=\sum_{S\subset[m]}\norm{\Flat_SK}_4^4.
\end{equation}
For each source leg,
$\E\overline g_ag_b\overline g_cg_d
 =\delta_{ab}\delta_{cd}+\delta_{ad}\delta_{cb}$.
Recording by $S$ the legs using the second pairing, the expanded trace is precisely
$\Tr[(\Flat_S(K)^*\Flat_S(K))^2]$.
Apply the same identity to differences of finite kernels to extend it to the cut completion. This formula is specific to the circular decoupled model; real sources have a third pairing.

\begin{proposition}[Every directed cut is necessary]\label{prop:cut_sharpness}
Fix $m\ge1$ and $s>2$. Removing any prescribed cut from the maximum destroys dimension-free control of the original $\Coeff_{m,s}$ norm. In the circular model, both constants in
\begin{equation}\label{eq:fourth_sharp_constants}
 \max_S\norm{\Flat_SK}_4\le\norm{T_K}_{L^4(\Sch_4)}
 \le2^{m/4}\max_S\norm{\Flat_SK}_4
\end{equation}
are optimal.
\end{proposition}
\begin{proof}
Fix $S_0\subset[m]$. Give every source leg dimension $D$, and take physical input and output spaces $\bigotimes_{j\in S_0}\C^D$ and $\bigotimes_{j\notin S_0}\C^D$. Define
\[
 K_D=\sum_{i_1,\ldots,i_m=1}^D
 h_{1,i_1}\otimes\cdots\otimes h_{m,i_m}
 \otimes\overline{c_{i_{S_0}}}\otimes e_{i_{S_0^c}}.
\]
Each source leg is paired with its corresponding physical factor. In the $S$-cut, a pairing with $j\notin S\triangle S_0$ stays on one side and has one singular value $\sqrt D$; a pairing with $j\in S\triangle S_0$ crosses the cut and gives an identity factor with $D$ singular values one. Thus
\begin{equation}\label{eq:hamming_cut_profile}
 \norm{\Flat_SK_D}_s
 =D^{m/2}D^{(1/s-1/2)|S\triangle S_0|}.
\end{equation}
For $s>2$, $S_0$ is the unique maximizing cut, and deleting it loses at least $D^{1/2-1/s}$. At $s=4$,~\eqref{eq:fourthmoment} gives
\[
 \frac{\E\norm{T_{K_D}}_4^4}{\max_S\norm{\Flat_SK_D}_4^4}
 =(1+D^{-1})^m\longrightarrow1,
\]
showing optimality of the lower constant. A scalar coefficient has every cut norm one and circular Gaussian product fourth moment $2^m$, attaining the upper constant.
\end{proof}


\section{Ordered multiplication and analytic regularity}\label{sec:product}
Coefficient observability becomes algebra reconstruction only after the prescribed multiplication is shown to be continuous in the original scale. We use one real source space $H$ and one physical Hilbert space $E$. Ordinary same-source multiplication contracts source indices at every depth while multiplying physical coefficients in their given order. The Hilbert estimate gives the exact radius region; the other Schatten exponents use Gaussian recovery to return each output to its original coefficient norm.

\subsection{The finite product and the need for a scale}
\begin{proposition}[A single-norm obstruction]\label{prop:nobanalg}
Let $g$ be standard real Gaussian. There is no norm on $\C[g]$ and no finite constants $C,M$ such that
\begin{equation}\label{eq:banachnogo}
 \norm{fh}\le C\norm f\norm h,\qquad |\E f|\le M\norm f
\end{equation}
for all polynomials. Completeness of the normed space is not required.
\end{proposition}
\begin{proof}
Iterating the product bound gives
\[
 (2n-1)!!=\E g^{2n}\le MC^{2n-1}\norm g^{2n}.
\]
The $2n$th root of the right-hand side remains bounded. The $2n$th root of the left-hand side tends to infinity, as follows already by retaining the last $\lfloor n/2\rfloor$ factors in the product. This is a contradiction.
\end{proof}
A rank-one projection embeds this polynomial algebra into operator-valued polynomials by $f\mapsto fP$. We therefore use a scale of norms for ordinary multiplication.

For $0\le r\le m\wedge n$, put
\[
 b_{m,n,r}=r!\binom mr\binom nr,\qquad k=m+n-2r.
\]
On finite symmetric kernels define
\begin{equation}\label{eq:Gamma}
 \Gamma_r(f\otimes A,g\otimes B)
 =(f\widetilde\otimes_rg)\otimes AB.
\end{equation}
The contraction uses the complex bilinear extension of the real inner product, followed by symmetrization of the uncontracted source legs. The physical product is $AB$. The ordinary Gaussian product formula~\cite{Nourdin} gives
\begin{equation}\label{eq:starcoeff}
 (K\star L)_k=\sum_{m+n-2r=k}b_{m,n,r}\Gamma_r(K_m,L_n),
 \qquad I(K\star L)=I(K)I(L).
\end{equation}
On the core, injectivity of $I$ proves associativity and
$(K\star L)^*=L^*\star K^*$. Already one source variable detects noncommutativity: for $A=E_{12}$ and $B=E_{21}$,
\[
 [gA,gB]=(H_2(g)+1)[A,B].
\]
Wick multiplication is the $r=0$ part of~\eqref{eq:starcoeff}; the remaining contractions are essential for ordinary multiplication.

\subsection{A positive contraction estimate at the Hilbert exponent}
Write
$\mathcal K_m=H_\C^{\odot m}\otimes_2\Sch_2(E)$ and
$\mathcal N_{\rho,2}(K)=\sum_m\rho^m\sqrt{m!}\norm{K_m}_{\mathcal K_m}$.

\begin{lemma}[Ordered contractions and a positive kernel]\label{lem:positivekernel}
Each $\Gamma_r$ extends continuously with
$\norm{\Gamma_r(K,L)}\le\norm K\norm L$.
If
\[
 c_{m,n,r}=b_{m,n,r}\sqrt{\frac{k!}{m!n!}},
 \qquad c_{m,n,r}^2=\binom mr\binom nr\binom{k}{m-r},
\]
then, for $\sigma,t>0$,
\begin{equation}\label{eq:positivekernel}
 \sum_{r=0}^{m\wedge n}c_{m,n,r}\sigma^{m+n-2r}
 \le\left(\sigma^2+\frac{1+\sigma^2}{t}\right)^{m/2}
     \left(\sigma^2+(1+\sigma^2)t\right)^{n/2}.
\end{equation}
\end{lemma}
\begin{proof}
Denote the uncontracted indices by $a,b$ and the contracted indices by $j$. Before source symmetrization, the coefficients are $\sum_jK_{a,j}L_{j,b}$. The ideal inequality and Cauchy--Schwarz give
\[
 \begin{split}
 \sum_{a,b}\left\|\sum_jK_{a,j}L_{j,b}\right\|_2^2
 &\le\sum_{a,b}\left(\sum_j\norm{K_{a,j}}_2^2\right)
                     \left(\sum_j\norm{L_{j,b}}_2^2\right)\\
 &=\norm K^2\norm L^2.
 \end{split}
\]
Symmetrization is an orthogonal projection on the source Hilbert tensor, proving the contraction estimate.

For the numerical bound, set $a=m-r$, $b=n-r$. The binomial expansion implies
$\binom{a+b}a\le(1+t)^{a+b}t^{-a}$.
With $U=\sigma^2(1+t)/t$ and $V=\sigma^2(1+t)$, the $r$th summand is at most
\[
 \left[\binom mrU^{m-r}t^{-r}\right]^{1/2}
 \left[\binom nrV^{n-r}t^r\right]^{1/2}.
\]
Apply Cauchy--Schwarz in $r$, and enlarge the two nonnegative sums to the full ranges $0\le r\le m$ and $0\le r\le n$. The result is
$(U+t^{-1})^{m/2}(V+t)^{n/2}$, which is the right-hand side of~\eqref{eq:positivekernel}.
\end{proof}

\begin{theorem}[The sharp multi-input radius region]\label{thm:arity}
Assume $H\ne0$, $E\ne0$, $k\ge2$, $\sigma\ge1$, and $\rho_i>0$. Let $\mathcal C_{\rho,2}$ be the kernel space with finite $\mathcal N_{\rho,2}$ norm. Ordered multiplication on the finite core is bounded from $\prod_{i=1}^k\mathcal C_{\rho_i,2}$ into $\mathcal C_{\sigma,2}$ if and only if
\begin{equation}\label{eq:arityregion}
 \sum_{i=1}^k\frac1{1+\rho_i^2}\le\frac1{1+\sigma^2}.
\end{equation}
In that region, its best multilinear constant is one:
\begin{equation}\label{eq:aritybound}
 \mathcal N_{\sigma,2}(K_1\star\cdots\star K_k)
 \le\prod_{i=1}^k\mathcal N_{\rho_i,2}(K_i).
\end{equation}
For equal input radii, the exact condition is $\rho^2\ge k(1+\sigma^2)-1$.
\end{theorem}
\begin{proof}
\emph{Binary sufficiency and absolute convergence.}
For two inputs,~\eqref{eq:arityregion} is equivalent to
\begin{equation}\label{eq:binaryregion}
 \rho^2>\sigma^2,\quad\eta^2>\sigma^2,\quad
 (\rho^2-\sigma^2)(\eta^2-\sigma^2)\ge(1+\sigma^2)^2.
\end{equation}
Indeed, clearing the positive denominators gives the last inequality after adding $\sigma^4$; either of the first two failures makes the reciprocal sum too large. Choose
\[
 \frac{1+\sigma^2}{\rho^2-\sigma^2}\le t
 \le\frac{\eta^2-\sigma^2}{1+\sigma^2}.
\]
Both factors on the right of~\eqref{eq:positivekernel} are then bounded by $\rho^m$ and $\eta^n$. The contraction estimate yields the stronger absolute bound
\begin{equation}\label{eq:absolutecontraction}
 \begin{split}
 &\sum_{m,n\ge0}\sum_{r\le m\wedge n}
 \sigma^{m+n-2r}\sqrt{(m+n-2r)!}\,b_{m,n,r}
       \norm{\Gamma_r(K_m,L_n)}\\
 &\hspace{25mm}\le\mathcal N_{\rho,2}(K)\mathcal N_{\eta,2}(L).
 \end{split}
\end{equation}
All contractions can therefore be grouped by output degree. Finite kernels are dense in the input spaces, so the bound supplies the unique continuous extension.

\emph{Ordered multi-input multiplication.}
Set $c_i=(1+\rho_i^2)^{-1}$. Choose an arbitrary binary bracketing preserving input order. For a node containing the consecutive input set $A$, assign the radius
\[
 r_A=\left(\left(\sum_{i\in A}c_i\right)^{-1}-1\right)^{1/2}.
\]
The root radius is at least $\sigma$, and every intermediate radius is at least the root radius. For each binary node, the reciprocal weights of its two children add to that of their parent, so the binary estimate holds at equality in its radius condition. Iterating with constant one proves~\eqref{eq:aritybound}; monotonicity in the output radius gives the prescribed $\sigma$. Associativity on the finite core and continuity make the result independent of bracketing. Taking every input to be the same constant rank-one projection of Hilbert--Schmidt norm one proves that the constant cannot be smaller than one.

\emph{Necessity.}
Fix a unit source vector $h$ and a rank-one projection $P$. For $z\ge0$, consider the coherent kernel
\[
 \kappa(z)_m=\frac{z^m}{m!}h^{\otimes m}\otimes P,
 \qquad\mathcal A(u)=\sum_{m\ge0}\frac{u^m}{\sqrt{m!}}.
\]
Direct contraction gives
\[
 \kappa(z_1)\star\cdots\star\kappa(z_k)
 =\exp\left(\sum_{i<j}z_iz_j\right)\kappa\left(\sum_i z_i\right),
 \qquad\mathcal N_{\rho,2}(\kappa(z))=\mathcal A(\rho z).
\]
These tests are justified even when boundedness is initially assumed only on the finite core: truncate each input, and pass to the limit in the nonnegative coefficient sums.

We have $\mathcal A(u)^2\ge e^{u^2}$. For $0<q<1$, Cauchy--Schwarz gives
\[
 \mathcal A(u)\le(1-q^2)^{-1/2}\exp\big(u^2/(2q^2)\big).
\]
Letting $u\to\infty$ and then $q\uparrow1$ shows
$\log\mathcal A(u)/u^2\to1/2$.
If multiplication has any finite bound, substitute $z_i=ta_i$, take logarithms, divide by $t^2$, and let $t\to\infty$. One obtains
\[
 \sum_i(1+\rho_i^2)a_i^2-(1+\sigma^2)\left(\sum_i a_i\right)^2\ge0
 \qquad(a_i\ge0).
\]
Taking $a_i=(1+\rho_i^2)^{-1}$ gives~\eqref{eq:arityregion}.
\end{proof}

Writing $b(\rho)=(1+\rho^2)^{-1}$, the sharp region says that the input budgets add: $\sum_i b(\rho_i)\le b(\sigma)$. Equivalently,
\[
 \diag(1+\rho_1^2,\ldots,1+\rho_k^2)
 -(1+\sigma^2)\mathbf1\mathbf1^*\ge0.
\]
Indeed conjugation by the inverse square root of the diagonal matrix reduces positivity to a rank-one eigenvalue test. This explains why any bracketing preserving input order has the same optimal total radius cost: each internal node simply adds its children's reciprocal budgets.

The binary boundary agrees with the scalar weighted Gaussian $L^2$ condition of Da Pelo--Lanconelli--Stan~\cite[Corollary~5.5]{DPLS}: take their product parameter to be the identity and the three probability exponents to be two. The norm here is the $\ell^1$ Fock coefficient norm, and~\eqref{eq:absolutecontraction} retains each contraction depth and operator order. At the level of all-radius scalar spaces, the comparison with the usual $\ell^2$ scale is
\begin{equation}\label{eq:l1l2}
 \left(\sum_m\rho^{2m}m!\norm{K_m}^2\right)^{1/2}
 \le\mathcal N_{\rho,2}(K)
 \le\frac{\left(\sum_mR^{2m}m!\norm{K_m}^2\right)^{1/2}}
 {\sqrt{1-(\rho/R)^2}},\qquad R>\rho.
\end{equation}
This relates the scalar topology to classical smooth-random-variable spaces~\cite{PT}.

\subsection{Ordinary multiplication in the original Schatten norms}
For $s>2$, a coefficient need not be a Hilbert tensor. To control its contractions, we first isolate the output degrees on the random-variable side and then recover them in the original coefficient norm.

\begin{lemma}[Mehler extraction with an explicit degree cost]\label{lem:chebyshev}
Let $B$ be a Banach space, and let $F=\sum_{j=0}^dF_j$ be a finite Gaussian polynomial with its homogeneous decomposition. For $\lambda\ge1$, put
$\chi_\lambda=\lambda+\sqrt{1+\lambda^2}$ and
$C_\lambda=2\chi_\lambda/(\chi_\lambda-1)$. Then
\begin{equation}\label{eq:chebyshev}
 \sum_{j=0}^d\lambda^j\norm{F_j}_{L^2(B)}
 \le C_\lambda\chi_\lambda^d\norm F_{L^2(B)}.
\end{equation}
\end{lemma}
\begin{proof}
For $t\in[-1,1]$, the correlated Mehler operator is a contraction on $L^2(B)$ and satisfies $M_tF=\sum_jt^jF_j$; a negative parameter adds a source reflection. Expand this Banach-valued polynomial in Chebyshev polynomials:
$M_tF=\sum_{a=0}^dA_aT_a(t)$.
The cosine integral formula gives $\norm{A_0}\le\norm F$ and $\norm{A_a}\le2\norm F$ for $a\ge1$.

Let $L_a(\lambda)$ be the sum of the absolute values of the monomial coefficients of $T_a$, weighted by $\lambda^j$. The recurrence $T_{a+1}=2tT_a-T_{a-1}$ gives
\[
 L_{a+1}\le2\lambda L_a+L_{a-1},\qquad L_0=1,\quad L_1=\lambda.
\]
Since $\chi_\lambda^2=2\lambda\chi_\lambda+1$, induction yields $L_a\le\chi_\lambda^a$. Taking norms of each monomial coefficient in the Chebyshev expansion and summing yields
\[
 \sum_j\lambda^j\norm{F_j}
 \le\left(1+2\sum_{a=1}^d\chi_\lambda^a\right)\norm F
 \le C_\lambda\chi_\lambda^d\norm F.
\]
\end{proof}

\begin{theorem}[The Schatten coefficient algebra]\label{thm:nativeproduct}
Let $1\le s<\infty$ and $\sigma\ge1$, and define
\begin{equation}\label{eq:nativeradius}
 R_s(\sigma)=u_s\sqrt3\,\chi_{v_s\sigma}.
\end{equation}
Ordered finite contractions extend uniquely and satisfy
\begin{equation}\label{eq:nativeproduct}
 \Nn\sigma{K\star L}
 \le C_{v_s\sigma}\Nn{R_s(\sigma)}K\Nn{R_s(\sigma)}L.
\end{equation}
Under the right-hand budget, all $(m,n,r)$ outputs are absolutely summable. The space $\Alg_s$ is a complete Fr\'echet $*$-algebra, and $I$ is a topological $*$-algebra isomorphism from $\Alg_s$ onto $\Gaus_s$, realizing ordinary same-source multiplication.
\end{theorem}
\begin{proof}
The proof recovers the product in the original coefficient norm after separating its homogeneous outputs. For finite homogeneous inputs put
$F_m=I_mK_m$, $G_n=I_nL_n$, $q_m=\norm{F_m}_{L^2(\Sch_s)}$, and $q_n'=\norm{G_n}_{L^2(\Sch_s)}$. Schatten multiplication and probability H\"older give
\[
 \norm{F_mG_n}_{L^2(\Sch_s)}
 \le\norm{F_m}_{L^4(\Sch_s)}\norm{G_n}_{L^4(\Sch_s)}
 \le3^{(m+n)/2}q_mq_n'.
\]
Here $\norm{AB}_s\le\norm A_s\norm B_s$ follows from the ideal inequality and $\norm B_{\op}\le\norm B_s$.

For these fixed input degrees, the output degree $j=m+n-2r$ determines the contraction depth $r$ uniquely. The Mehler estimate with $\lambda=v_s\sigma$, followed by coefficient recovery at each output degree, therefore gives
\begin{equation}\label{eq:nativeabsolute}
 \begin{split}
 &\sum_{r\le m\wedge n}\sigma^j\sqrt{j!}\,b_{m,n,r}
       \norm{\Gamma_r(K_m,L_n)}_{\Wc_{j,s}}\\
 &\quad\le C_{v_s\sigma}(\sqrt3\chi_{v_s\sigma})^{m+n}q_mq_n'\\
 &\quad\le C_{v_s\sigma}R_s(\sigma)^{m+n}\sqrt{m!n!}
       \norm{K_m}_{\Wc_{m,s}}\norm{L_n}_{\Wc_{n,s}}.
 \end{split}
\end{equation}
In particular, each fixed contraction extends by continuity to its two original coefficient spaces. This construction is valid also when a high-exponent kernel has infinite Hilbert tensor norm: the extension uses the cut norm, not a formal Hilbert contraction of an undefined tensor.

Summing the nonnegative bound over $m,n$ proves absolute summability of every $(m,n,r)$ output and proves \eqref{eq:nativeproduct}. It remains to identify the sum with ordinary multiplication. The input evaluations converge in $L^4(\Sch_s)$, because $R_s(\sigma)\ge u_s\sqrt3$. Their finite products consequently converge in $L^2(\Sch_s)$ to the actual random product. Applied before coefficient inversion, the same Mehler estimate makes the corresponding random homogeneous outputs absolutely summable in $L^2$, with weights $(v_s\sigma)^j$. For every fixed $j$, continuity of the homogeneous inverse identifies its kernel with the sum of the degree-$j$ contractions. Summing these identified random outputs gives exactly the ordinary product. The argument applies also when the output radius is smaller than $u_s$.

The intersection over integer radii of the weighted $\ell^1$ coefficient spaces is complete. A Cauchy sequence has a limit at each radius, and continuity of the degree projections makes these limits have the same coefficients. Finite degree, source, and physical compressions are dense. The product bound supplies a finite larger input radius for every output radius, hence joint continuity on the intersection. Adjoints are isometries in each coefficient norm. Associativity and the reversed-product identity pass from the finite core by this continuity.

Finally, fixed-degree evaluation and inversion give
\begin{equation}\label{eq:scaletopology}
 \Nn{\rho/v_s}K\le\En\rho{I(K)}\le\Nn{u_s\rho}K.
\end{equation}
Thus evaluation is continuous and has a continuous inverse in the all-radius topologies. Given $F\in\Gaus_s$, Proposition~\ref{prop:chaosimage} recovers the unique kernel of each homogeneous component. The left inequality at arbitrarily large radii places the resulting sequence in $\Alg_s$, and its evaluation is $F$. This proves surjectivity and the asserted topological $*$-algebra isomorphism.
\end{proof}

At $s=2$, Theorem~\ref{thm:arity} gives the sharper optimal region. At $s=1$, the same proof uses the exponential predual constant from Proposition~\ref{prop:traceendpoint}, rather than reflexivity or a bounded retraction on the full random data space. The entire trace-class coefficient algebra and its finite inverses are therefore available. Section~\ref{sec:data} supplies positive physical masses and anchored source innovations also at $s=1$, giving explicit original-norm tail certificates for this completion. Only the alternative source-convexity inequality retains the strict range $1<s<\infty$.

\subsection{An intrinsic analytic description of the all-radius topology}
The radius scale also has a description that does not choose source coordinates. It is the compact-open topology of an entire Mehler response. This identifies the regularity encoded by the coefficient weights; it does not replace the independently defined cut norms or the proof of their equivalence with Gaussian evaluation.

Write $\mathscr B_s=L^2(\sigma(W);\Sch_s)$ and let $M_t$, $0\le t\le1$, be the source Mehler contraction on this Banach space. Thus $M_tI_mK_m=t^mI_mK_m$. The parameter is the source correlation, rather than the Ornstein--Uhlenbeck time $-\log t$.

\begin{theorem}[Entire Mehler characterization]\label{thm:mehler_entire}
Let $1\le s<\infty$ and let $F\in\mathscr B_s$. Then $F\in\Gaus_s$ if and only if $t\mapsto M_tF$ extends to an entire $\mathscr B_s$-valued function $\mathscr M_F$ on $\C$. The extension is unique and satisfies
\begin{equation}\label{eq:mehler_entire_series}
 \mathscr M_F(z)=\sum_{m\ge0}z^mF_m,
 \qquad M_t\mathscr M_F(z)=\mathscr M_F(tz)
 \quad(0\le t\le1,\ z\in\C).
\end{equation}
Put $\mathfrak m_r(F)=\sup_{|z|\le r}\norm{\mathscr M_F(z)}_{\mathscr B_s}$. For $0<r<R$,
\begin{equation}\label{eq:mehler_compact_open_norms}
 \mathfrak m_r(F)\le\En rF
 \le\frac{\mathfrak m_R(F)}{1-r/R}.
\end{equation}
Consequently the compact-open topology of these curves is the original recovered topology. More explicitly, for $F=I(K)$ and $R>v_s\rho$,
\begin{align}
 \mathfrak m_\rho(F)&\le\Nn{u_s\rho}K,&
 \Nn\rho K&\le\frac{\mathfrak m_R(F)}{1-v_s\rho/R},\label{eq:mehler_native_norms}\\
 \Nn\rho{K^{>M}}&\le
 \frac{(v_s\rho/R)^{M+1}}{1-v_s\rho/R}\,\mathfrak m_R(F).
 \label{eq:mehler_native_tail}
\end{align}
The curves in \eqref{eq:mehler_entire_series} form the closed linear subspace of entire $\mathscr B_s$-valued curves satisfying the displayed Mehler covariance. Evaluation at $z=1$ is the inverse parametrization of this subspace.
\end{theorem}
\begin{proof}
If $F\in\Gaus_s$, its homogeneous series converges absolutely and uniformly on each disk in $\mathscr B_s$, since every radius sum is finite. It therefore defines an entire function. For real $t\in[0,1]$, continuity of the Mehler contraction identifies its value with $M_tF$. Applying the same contraction termwise gives the covariance identity, and the triangle inequality gives the first bound in \eqref{eq:mehler_compact_open_norms}.

Conversely, suppose an entire extension $H(z)$ exists and write its Banach-valued Taylor series as $H(z)=\sum_{m\ge0}z^mA_m$. We identify $A_m$ before using any homogeneous projection on the full Banach space. For a finite Hermite index $\alpha$ and a physical matrix entry, the scalar functional
\[
 Y\longmapsto\E[\Psi_\alpha Y_{k\ell}]
\]
is continuous on $\mathscr B_s$. Scalar Mehler symmetry gives
\[
 \E[\Psi_\alpha(M_tF)_{k\ell}]
   =t^{|\alpha|}\E[\Psi_\alpha F_{k\ell}].
\]
Both sides extend to entire functions of $t$, so the identity theorem and comparison of Taylor coefficients show that $A_m$ has no Hermite coefficient outside degree $m$. Condition on finitely many source coordinates and compress to finitely many physical coordinates. Completeness of scalar Hermite polynomials in each resulting $L^2$ space makes this compression a homogeneous degree-$m$ matrix polynomial. Proposition~\ref{prop:chaosimage} therefore gives $A_m\in\mathcal H_m(\Sch_s)$, including when $s=1$.

Banach-valued Cauchy estimates give $\norm{A_m}_{\mathscr B_s}\le R^{-m}\sup_{|z|\le R}\norm{H(z)}_{\mathscr B_s}$. Summing at any $r<R$ proves the second bound in \eqref{eq:mehler_compact_open_norms}. In particular the homogeneous series is absolutely summable at every radius, and $H(1)=F$ identifies $A_m=F_m$. This proves the characterization and uniqueness, without a bounded homogeneous projection on arbitrary trace-class-valued data.

For $F=I(K)$, use the degreewise bounds
\[
 \sqrt{m!}\norm{K_m}_{\Wc_{m,s}}
 \le v_s^m\norm{F_m}_{\mathscr B_s},
 \qquad
 \norm{F_m}_{\mathscr B_s}
 \le u_s^m\sqrt{m!}\norm{K_m}_{\Wc_{m,s}}.
\]
They give \eqref{eq:mehler_native_norms}; summing the Cauchy estimate only over $m>M$ gives \eqref{eq:mehler_native_tail}.

Finally, the covariance equations are closed under compact-open convergence because each $M_t$ is bounded. Any entire curve satisfying them obeys $H(t)=M_tH(1)$, and hence has the form just characterized. Conversely every $\mathscr M_F$ satisfies the equations. Uniform Cauchy convergence on disks and the Cauchy formula make the entire-curve space complete. This supplies the stated closed realization of the recovered topology.
\end{proof}

This is the exponential Ornstein--Uhlenbeck regularity underlying the scalar smooth-random-variable scale of \cite{PT}; the original Schatten-cut geometry is identified with it by the proved coefficient estimates. A bound on $\mathfrak m_R$ is an analytic input to the finite inverse estimate. The complex parameter is defined by holomorphic continuation.

The multiplication on these curves remains the transported ordinary Gaussian product: it sends $(\mathscr M_F,\mathscr M_G)$ to $\mathscr M_{FG}$. It is not pointwise multiplication at the complex parameter. Already for a scalar Gaussian,
\[
 \mathscr M_{g^2}(z)=z^2H_2(g)+1,
 \qquad \mathscr M_g(z)^2=z^2H_2(g)+z^2.
\]
The difference $1-z^2$ is the contraction term. Thus the analytic description clarifies the topology while leaving the ordered contraction law intact.

\subsection{Truncation, ordered words, and dual weights}
For an admissible Hilbert pair of input radii, degree truncation satisfies
\begin{align}\label{eq:chaostailproduct}
 \mathcal N_{\sigma,2}(K\star L-K^{\le M}\star L^{\le N})
 &\le\mathcal N_{\rho,2}(K^{>M})\mathcal N_{\eta,2}(L)
      +\mathcal N_{\rho,2}(K)\mathcal N_{\eta,2}(L^{>N}),\notag\\
 \mathcal N_{\rho,2}(K^{>M})
 &\le(\rho/R)^{M+1}\mathcal N_{R,2}(K),\qquad R>\rho.
\end{align}
The first bound includes every output degree, including complete contraction to degree zero. For comparison, $F_M=H_M(g)P/\sqrt{M!}$ has eventually zero observations at each fixed low degree but $\E F_M^2=P$. Its coefficient norm at radius $\rho$ is $\rho^M$, which identifies the missing high-degree control.

Chaos degree and operator word length can be recorded separately. First take scalar source kernels $K_{m,w}\in H_\C^{\odot m}$ indexed by formal $*$-words, and set
\[
 \mathcal D_{\rho,R}(K)=\sum_{m,w}\rho^m\sqrt{m!}\,R^{|w|}\norm{K_{m,w}}_2,
 \qquad R>0.
\]
Source indices contract as in~\eqref{eq:starcoeff}, while words concatenate in order. Since $R^{|uv|}=R^{|u|}R^{|v|}$, the Hilbert radius estimate remains valid after summation over words. At fixed $R$, the intersection over source radii is an ordered Fr\'echet $*$-algebra.

For a bounded tuple $A$ with $\norm{A_i}\le R$, define
\[
 \operatorname{ev}_A K=\sum_{m,w}I_m(K_{m,w})\,w(A).
\]
The sum converges absolutely in every finite $L^p(B(E))$: the scalar chaos estimate bounds its norm by $\mathcal D_{\max(1,\sqrt{p-1}),R}(K)$ for $p\ge2$, and $L^p\le L^2$ handles $p<2$. On the finite core the scalar coefficients commute with the evaluated letters, so evaluation preserves multiplication and the involution. The all-radius product estimate and $L^{2p}$ convergence pass these identities to the completion.

One may also retain operator-valued $K_{m,w}$ and multiply them in their own order while concatenating words; the same norm estimate defines the formal algebra. A multiplicative evaluation into one operator space then requires the coefficient algebra to commute with the substituted letters, or a representation on separate tensor factors. Without that requirement, the formal commuting relation between a coefficient and a letter need not survive substitution. For instance, a formal letter commutes with a constant coefficient, whereas $E_{12}$ and $E_{21}$ do not. Ordinary same-space noncommutative products are instead formed in the original algebra by \eqref{eq:starcoeff}.

A finite-degree PBW change of coordinates can always be equipped with the norm pulled back from words. For a separately prescribed coordinate norm, let $B_n,D_n$ be the norms of the degree-$n$ transform and its inverse. A bounded change from word radius $R'$ to $R$ requires
\[
 \sup_n(R/R')^nB_n<\infty,
\]
and the analogous condition with $D_n$ for the inverse. Descent to a represented algebra additionally requires invariance of the representation kernel under the relevant PBW projections.

Finally, the dual weights are reciprocal. For Banach spaces $X_n$ and positive weights $w_n$,
\begin{equation}\label{eq:weighteddual}
 \left(\bigoplus_n^{\ell^1}w_nX_n\right)'
 \simeq\left\{(f_n):\sup_nw_n^{-1}\norm{f_n}_{X_n'}<\infty\right\},
 \qquad\langle x,f\rangle=\sum_nf_n(x_n).
\end{equation}
H\"older proves boundedness, single-coordinate tests prove the norm equality, and finite-support density identifies every functional. Thus the forward coefficient sums and the reverse test bounds retain their respective weights under completion.


\section{Microscopic fluctuation algebras}\label{sec:microscopic}
The Hilbert coefficient algebra admits a discrete realization in which ordinary multiplication is performed before the Gaussian limit. Independent signs produce symmetric site responses. Intersections of site sets generate the contraction terms, and the number of remaining sites gives their finite-scale correction. This section establishes convergence of the product in the original coefficient topology and identifies the finite observations that commute with this limit.

\subsection{Symmetric site responses and their coefficient geometry}
Fix an orthonormal basis $(e_j)$ of the real source space. A finite-dimensional source uses the same construction with finitely many indices. In the normalized Hermite coordinates write
\begin{equation}\label{mic:native}
 A=(A_\alpha),\qquad
 a_m(A)=\left(\sum_{|\alpha|=m}\norm{A_\alpha}_2^2\right)^{1/2},\qquad
 p_\rho(A)=\sum_{m\ge0}\rho^m a_m(A),
\end{equation}
where $\alpha$ ranges over finitely supported multi-indices and $A_\alpha\in\Sch_2(E)$. These are coordinates of the original space $\Alg_2$:
\begin{equation}\label{mic:normalization}
 K_m=\frac1{\sqrt{m!}}\sum_{|\alpha|=m}h_\alpha\otimes A_\alpha,
 \qquad p_\rho(A)=\mathcal N_{\rho,2}(K),\qquad
 \Phi(A)=\sum_\alpha A_\alpha\Psi_\alpha=I(K).
\end{equation}
In particular each homogeneous coefficient space keeps its Hilbert norm. The physical space $E$ may be infinite-dimensional.

Let $(\xi_{j,a})_{j,a\ge1}$ be independent symmetric signs. The source index $j$ and the site index $a$ define
\begin{equation}\label{mic:field}
 S_{j,N}=N^{-1/2}\sum_{a=1}^N\xi_{j,a},\qquad
 W_N(h)=\sum_jh_jS_{j,N},\qquad
 \E[W_N(h)W_N(k)]=\ip hk.
\end{equation}
The last series converges in $L^2$ for $h\in H$. For one source block set
\begin{equation}\label{mic:qdef}
 Q_{k,N}^{(j)}=k!N^{-k/2}
   \sum_{\substack{I\subset\{1,\ldots,N\}\\|I|=k}}\prod_{a\in I}\xi_{j,a},
 \qquad 0\le k\le N.
\end{equation}
Set $Q_{0,N}=1$ and $Q_{k,N}=0$ for $k>N$, and define
\begin{equation}\label{mic:coordinates}
 \Psi_{\alpha,N}=\prod_j\frac{Q_{\alpha_j,N}^{(j)}}{\sqrt{\alpha_j!}},\qquad
 d_{\alpha,N}=\prod_j\frac{(N)_{\underline{\alpha_j}}}{N^{\alpha_j}}.
\end{equation}
Here $(t)_{\underline r}=t(t-1)\cdots(t-r+1)$, with value one at $r=0$ and zero for integers $0\le t<r$. Let
$\mathcal I_N=\{\alpha:\max_j\alpha_j\le N\}$, set $d_{\alpha,N}=0$ outside this set, and let $\Pi_N$ restrict coefficients to $\mathcal I_N$. The space $\Alg_{2,N}=\Pi_N\Alg_2$ retains arbitrarily many source directions even at fixed $N$.

\begin{lemma}[Site Gram and uniform evaluation]\label{mic:gram}
The site coordinates satisfy
\begin{equation}\label{mic:orthogonality}
 \E[\Psi_{\alpha,N}\Psi_{\beta,N}]
   =\mathbf1_{\alpha=\beta}d_{\alpha,N},\qquad
 e^{-|\alpha|}\le d_{\alpha,N}\le1\quad(\alpha\in\mathcal I_N).
\end{equation}
The evaluation $\Phi_N(A)=\sum_\alpha A_\alpha\Psi_{\alpha,N}$ is injective on $\Alg_{2,N}$. Its homogeneous Walsh components obey
\begin{align}
 \norm{\Phi_{N,m}(A)}_{L^2(\Sch_2)}^2
 &=\sum_{|\alpha|=m}d_{\alpha,N}\norm{A_\alpha}_2^2,\label{mic:isometry}\\
 \sum_m\rho^m\norm{\Phi_{N,m}(A)}_{L^2(\Sch_2)}
 &\le p_\rho(A)
 \le\sum_m(\sqrt e\rho)^m\norm{\Phi_{N,m}(A)}_{L^2(\Sch_2)}
 \quad(A\in\Alg_{2,N}).\label{mic:radius_equivalence}
\end{align}
All constants are independent of the site number and both source and physical dimensions.
\end{lemma}
\begin{proof}
Products of signs over different subsets are orthogonal Walsh characters. Thus
$\E[Q_{k,N}Q_{l,N}]=\mathbf1_{k=l}k!(N)_{\underline k}/N^k$.
Independence between blocks gives the Gram formula, and Hilbert-valued Parseval gives \eqref{mic:isometry}. For $1\le k\le N$,
\[
 \frac{(N)_{\underline k}}{N^k}
 =\prod_{a=0}^{k-1}(1-a/N)\ge\frac{k!}{k^k}\ge e^{-k}.
\]
The last inequality follows by integrating $\log t$ from $1$ to $k$. Multiplication over the source blocks and summation over degrees prove \eqref{mic:radius_equivalence}. Finite kernels are dense, so the formulas extend by $L^2$ completeness.
\end{proof}

\begin{lemma}[Uniform moment bounds]\label{mic:moments}
Put $u_p=1$ for $1\le p\le2$ and $u_p=\sqrt{2\lceil p/2\rceil-1}$ for $p>2$. For every $A\in\Alg_2$,
\begin{equation}\label{mic:moment_bound}
 \norm{\Phi_N(A)}_{L^p(\Sch_2)}\le p_{u_p}(A),\qquad
 \norm{\Phi(A)}_{L^p(\Sch_2)}\le p_{u_p}(A).
\end{equation}
For a homogeneous input the right-hand side is $u_p^m a_m(A)$.
\end{lemma}
\begin{proof}
For a finite real scalar Walsh polynomial $f=\sum_{|I|=m}c_I\xi_I$, expand its even moment of order $2q$. Taking absolute values of the coefficients and replacing the signs by independent standard Gaussians can only increase the sum: every surviving sign moment is one, whereas the corresponding even Gaussian moment is at least one. The comparison polynomial has distinct indices in every monomial and belongs to Gaussian chaos $m$. Hence \eqref{eq:hyper} gives
\[
 \E|f|^{2q}\le
 \E\left|\sum_{|I|=m}|c_I|\prod_{i\in I}g_i\right|^{2q}
 \le(2q-1)^{mq}\left(\sum_Ic_I^2\right)^q.
\]
Realify the coefficient Hilbert space, apply this inequality to each real orthonormal coordinate, and use Minkowski in $L^q$ on the squared coordinate norms. This gives the same Hilbert-valued bound. Choose $q=\lceil p/2\rceil$, use \eqref{mic:isometry}, and sum over degrees. Density proves the infinite-dimensional assertion. The Gaussian estimate follows directly from hypercontractivity.
\end{proof}

Each $Q_{k,N}$ depends only on the number of negative signs in its block and is a normalized Krawtchouk polynomial. The classical Krawtchouk--Hermite asymptotics and Rademacher product formulas are treated in \cite{Min,NPRR}; Gaussian universality for homogeneous sums is developed in \cite{NPR}. Here the symmetric block model gives an exact finite-scale coefficient algebra.

\subsection{Collisions and the finite-scale product}
\begin{proposition}[Exact site collision formula]\label{mic:collision}
For $0\le m,n\le N$,
\begin{equation}\label{mic:one_product}
 Q_{m,N}Q_{n,N}
 =\sum_{\substack{0\le r\le m\wedge n\\k=m+n-2r\le N}}
 r!\binom mr\binom nr\frac{(N-k)_{\underline r}}{N^r}Q_{k,N}.
\end{equation}
\end{proposition}
\begin{proof}
Let $I,J$ have sizes $m,n$ and put $r=|I\cap J|$. Their sign products multiply to the character on $I\triangle J$, since $\xi_a^2=1$. Fix this remaining set, of size $k=m+n-2r$. Choose its $m-r$ members belonging to $I\setminus J$, and then choose the $r$ shared sites from its complement. This gives
$\binom{k}{m-r}\binom{N-k}{r}$ pairs. Multiplying by $m!n!N^{-(m+n)/2}$ and dividing by $k!N^{-k/2}$ yields \eqref{mic:one_product}.
\end{proof}
In particular,
\begin{equation}\label{mic:recurrence}
 S_{j,N}Q_{k,N}^{(j)}
 =Q_{k+1,N}^{(j)}+k\left(1-\frac{k-1}{N}\right)Q_{k-1,N}^{(j)}.
\end{equation}
Thus $Q_{1,N}=S_N$, $Q_{2,N}=S_N^2-1$, and
$Q_{3,N}=S_N^3-(3-2/N)S_N$.

For multi-indices $r\le\alpha\wedge\beta$, put $\gamma=\alpha+\beta-2r$ and
\begin{equation}\label{mic:collision_weights}
 c(\alpha,\beta,r)=r!\binom\alpha r\binom\beta r
       \sqrt{\frac{\gamma!}{\alpha!\beta!}},\qquad
 \vartheta_N(\alpha,\beta,r)
       =\prod_j\frac{(N-\gamma_j)_{\underline{r_j}}}{N^{r_j}}.
\end{equation}
Take $\vartheta_N=0$ if any of $\alpha,\beta,\gamma$ lies outside $\mathcal I_N$. Define on finite coefficients
\begin{equation}\label{mic:product}
 (A\star_NB)_\gamma
 =\sum_{\alpha+\beta-2r=\gamma}
 c(\alpha,\beta,r)\vartheta_N(\alpha,\beta,r)A_\alpha B_\beta.
\end{equation}
The physical product retains its order. The Gaussian product $\star$ in these coordinates replaces every $\vartheta_N$ by one.

\begin{corollary}[The represented site algebra]\label{mic:represented}
Finite coefficients satisfy
\begin{equation}\label{mic:evaluation_product}
 \Phi_N(A\star_NB)=\Phi_N(A)\Phi_N(B),\qquad
 (A\star_NB)^*=B^*\star_NA^*.
\end{equation}
On coefficients supported in $\mathcal I_N$, this product is associative and represents faithfully the polynomial responses invariant under permutations of the sites in each source block.
\end{corollary}
\begin{proof}
Apply \eqref{mic:one_product} independently in each source block. The Gram in \eqref{mic:orthogonality} is strictly positive on its support, so evaluation is injective there and transfers associativity from ordinary multiplication. The $N+1$ orthogonal polynomials $Q_{0,N},\ldots,Q_{N,N}$ span all functions of the $N+1$ block-sum values. Tensoring over finitely many source directions identifies the represented polynomial algebra.
\end{proof}

\begin{lemma}[Uniform depletion error]\label{mic:depletion}
For $d=|\alpha|+|\beta|$ and $N\ge1$,
\begin{equation}\label{mic:depletion_bound}
 0\le1-\vartheta_N(\alpha,\beta,r)
 \le\min\left\{1,\frac{d(d-1)}{2N}\right\}.
\end{equation}
If $N\ge d$ and $\zeta=\sum_j[r_j\gamma_j+r_j(r_j-1)/2]$, then
\begin{equation}\label{mic:first_correction}
 \vartheta_N=1-\zeta/N+\mathcal R_N,\qquad
 0\le\mathcal R_N\le\zeta^2/(2N^2).
\end{equation}
\end{lemma}
\begin{proof}
When $N\ge d$, the factors are $1-(\gamma_j+a)/N$, with $0\le a<r_j$. Their losses lie in $[0,1]$, sum to $\zeta/N$, and $\zeta\le\binom d2$. The elementary product bounds
\[
 0\le\sum_i x_i-\left(1-\prod_i(1-x_i)\right)
 \le\sum_{i<j}x_ix_j\le\tfrac12\left(\sum_i x_i\right)^2
\]
prove both conclusions. If $d>N$, then $\binom d2/N\ge1$, so \eqref{mic:depletion_bound} follows from $0\le\vartheta_N\le1$. Degrees zero and one have no loss.
\end{proof}
The finite correction retains the occupancy of the chosen source blocks. Its limit is the basis-independent Gaussian contraction law.

\subsection{Convergence of multiplication in the original topology}
\begin{theorem}[Microscopic product limit]\label{mic:algebra_limit}
Let $\sigma\ge1$ and let $\rho,\eta>0$ satisfy
\begin{equation}\label{mic:radius}
 \frac1{1+\rho^2}+\frac1{1+\eta^2}\le\frac1{1+\sigma^2}.
\end{equation}
For all $N\ge1$, the products extend continuously and satisfy
\begin{equation}\label{mic:uniform_product}
 p_\sigma(A\star_NB)\le p_\rho(A)p_\eta(B),\qquad
 p_\sigma(A\star B)\le p_\rho(A)p_\eta(B).
\end{equation}
Each $\Alg_{2,N}$ is a complete Fr\'echet $*$-algebra, and $\Phi_N$ realizes ordinary multiplication. If $R_+>\rho$, $S_+>\eta$ and
$q=\max(\rho/R_+,\eta/S_+)<1$, then
\begin{equation}\label{mic:product_rate}
 p_\sigma(A\star_NB-A\star B)
 \le\frac{q^2}{N(1-q)^3}p_{R_+}(A)p_{S_+}(B).
\end{equation}
The radius region \eqref{mic:radius} is necessary for a bound uniform over all $N$, and its optimal uniform constant is one.
\end{theorem}
\begin{proof}
For finite coefficients set $a_\alpha=\norm{A_\alpha}_2$ and $b_\beta=\norm{B_\beta}_2$. Positivity of $c$, the bound $0\le\vartheta_N\le1$, and the Hilbert--Schmidt ideal inequality give
\[
 \norm{(A\star_NB)_\gamma}_2\le(a\star b)_\gamma.
\]
The right side is the Gaussian product of nonnegative scalar coefficient arrays. At each degree its input Hilbert norms are exactly $a_m(A)$ and $a_m(B)$. The positive contraction estimate of Lemma~\ref{lem:positivekernel}, summed as in Theorem~\ref{thm:arity}, therefore proves \eqref{mic:uniform_product} without a source-dimension factor.

For fixed input degrees $m,n$, Lemma~\ref{mic:depletion} multiplies this positive bound by at most $\binom{m+n}{2}/N$. Hence
\[
 p_\sigma(A\star_NB-A\star B)
 \le\frac1N\sum_{m,n}\binom{m+n}{2}\rho^m\eta^n a_m(A)a_n(B).
\]
Use $\rho^m\eta^n\le R_+^mS_+^n q^{m+n}$ and
$\sup_{d\ge0}\binom d2q^d\le q^2/(1-q)^3$ to obtain \eqref{mic:product_rate}. This estimate includes coefficients removed by $\Pi_N$.

Finite degree, source, and physical cutoffs are dense at each radius. The estimates give continuous products with absolutely summable contraction outputs. By \eqref{mic:moment_bound}, approximating the inputs at sufficiently large radii gives $L^4$ convergence, so their finite products converge in $L^2$ and identify \eqref{mic:evaluation_product} on the completion. The range of the continuous projection $\Pi_N$ is closed; associativity and the involution identity extend by continuity.

A uniform bound on finite coefficients can be passed to $N\to\infty$ termwise, producing the Gaussian bound at the same radii. Its necessity follows from Theorem~\ref{thm:arity}. Equal constant rank-one projections show that the uniform constant is at least one.
\end{proof}
For any fixed number of ordered inputs, the same binary-tree allocation gives the region \eqref{eq:arityregion}. Telescoping at its vertices gives an $O(N^{-1})$ coefficient error for every fixed noncommutative polynomial, with constants depending on the polynomial and the radius margins.

The order in \eqref{mic:product_rate} is attained by a quadratic response. For $N\ge4$,
\begin{equation}\label{mic:quadratic}
 Q_{2,N}^2=Q_{4,N}+4(1-2/N)Q_{2,N}+2(1-1/N).
\end{equation}
Taking $A_{2e_1}=\sqrt2P$ for a rank-one projection and all other coefficients zero gives
\begin{equation}\label{mic:sharp_rate}
 p_\sigma(A\star_NA-A\star A)=\frac{2+8\sqrt2\sigma^2}{N}.
\end{equation}
With two noncommuting physical coefficients, the same collision polynomial multiplies their ordered product; its antisymmetrization multiplies their commutator.

\subsection{Joint fluctuation limits}
\begin{theorem}[Joint realization of the coefficient algebra]\label{mic:clt}
For every finite family $A^1,\ldots,A^k\in\Alg_2$,
\begin{equation}\label{mic:law_limit}
 (\Phi_N(A^1),\ldots,\Phi_N(A^k))
 \Longrightarrow(\Phi(A^1),\ldots,\Phi(A^k))
 \quad\text{in }\Sch_2(E)^k.
\end{equation}
There is a common auxiliary coupling of the source block sums and these responses under which convergence holds in every finite $L^p(\Sch_2)$. Every fixed noncommutative $*$-polynomial with zero constant term converges in the same sense. These conclusions also hold for $A_N^j\to A^j$ in all coefficient radii.
\end{theorem}
\begin{proof}
The characteristic function of one block sum is
$\cos(t/\sqrt N)^N\to e^{-t^2/2}$. Independence gives the joint Gaussian limit for finitely many blocks. For general $h\in H$, the $L^2$ error of a finite source truncation of \eqref{mic:field} is exactly the omitted Hilbert norm, uniformly in $N$.

Let $F_N$ and $F_G$ be the distribution functions of a block sum and a standard Gaussian. With independent uniforms $U_j$, couple
$\widetilde S_{j,N}=F_N^{-1}(U_j)$ and $g_j=F_G^{-1}(U_j)$.
Continuity and strict increase of $F_G$ give almost-sure convergence. The recurrence \eqref{mic:recurrence} shows inductively that each fixed $Q_{k,N}$ is a polynomial in $S_N$ whose coefficients tend to those of $H_k$. Thus any finite family of site coordinates converges almost surely to the corresponding Hermite coordinates. Conditional on each block sum, uniform choice of a sign arrangement realizes the correct microscopic marginal law as well.

Finite coefficient responses now converge almost surely. The moment estimates at a larger exponent give uniform integrability and hence convergence in each finite $L^p$. For general coefficients, approximate at radius $u_p$ by a finite coefficient array. Equation~\eqref{mic:moment_bound} controls both approximation errors uniformly, proving the conclusion. For varying coefficients use the same bound on $A_N^j-A^j$. A word difference is a telescoping sum; Schatten multiplication and probability H\"older, with exponent $p$ times its word length for each factor, give the polynomial conclusion.
\end{proof}
The coefficient rate \eqref{mic:product_rate} is deterministic and is separate from the coupling above. On the original nested sign sequence,
$\E|S_N-S_{2N}|^2=2-\sqrt2$, so those block sums are not $L^2$-Cauchy. Higher-order limits are Gaussian chaoses, such as $Q_{2,N}\Longrightarrow g^2-1$.

\subsection{Finite decoding and the observation-depth threshold}
Fix a normalized faithful anchor $Rf_l=r_lf_l$ on the physical space. A budget $b=(M,d,L)$ retains degrees through $M$, the first $d$ source directions, and a physical corner of size $L$; let $P_b$ be its coefficient projection. For $N\ge M$ define normalized tests and observations by
\begin{equation}\label{mic:marks}
 \psi^{\rm on}_{\alpha,N}=d_{\alpha,N}^{-1/2}\Psi_{\alpha,N},\qquad
 (O_b^{(N)}F)_{\alpha;kl}
 =\E[\psi^{\rm on}_{\alpha,N}(RFR)_{kl}].
\end{equation}
The test and the response use the same sample. Orthogonality gives
\begin{equation}\label{mic:inverse}
 (O_b^{(N)}\Phi_N(C))_{\alpha;kl}
   =r_kr_l\sqrt{d_{\alpha,N}}(C_\alpha)_{kl}.
\end{equation}
Let $D_b^{(N)}$ divide by these weights and set unobserved coefficients to zero. At the Gaussian limit, the tests are $\Psi_\alpha$ and $d_{\alpha,N}$ is replaced by one; denote these maps by $O_b,D_b$. In the normalization \eqref{mic:normalization}, these are the finite inverses of Theorem~\ref{thm:finitedecoder}.

\begin{theorem}[Decoding commutes with the product limit]\label{mic:decoder}
For $A,B\in\Alg_2$, fixed $b$, and $N\ge M$,
\begin{equation}\label{mic:commuting}
 \begin{aligned}
 D_b^{(N)}O_b^{(N)}(\Phi_N(A)\Phi_N(B))&=P_b(A\star_NB),\\
 P_b(A\star_NB)&\longrightarrow P_b(A\star B)
   =D_bO_b(\Phi(A)\Phi(B))
 \end{aligned}
\end{equation}
in every coefficient radius. Suppose the retained degree-$m$ product marks have total Euclidean error at most $\delta_m$. Set
\[
 \kappa_L=(\min_{k\le L}r_k)^{-2},\qquad
 \chi_{m,N}=\left(\frac{N^m}{(N)_{\underline m}}\right)^{1/2},\quad\chi_{0,N}=1.
\]
For $\widehat C=D_b^{(N)}\widehat y$ and radii as in Theorem~\ref{mic:algebra_limit},
\begin{equation}\label{mic:total_error}
 \begin{split}
 p_\sigma(\widehat C-A\star B)\le{}&
 \kappa_L\sum_{m=0}^M\sigma^m\chi_{m,N}\delta_m
   +p_\sigma((I-P_b)(A\star B))\\
 &+\frac{q^2}{N(1-q)^3}p_{R_+}(A)p_{S_+}(B).
 \end{split}
\end{equation}
Along cofinal budgets $b_N$ with $M_N\le N$, vanishing of the noise term at every radius implies $\widehat C_N\to A\star B$ in $\Alg_2$.
\end{theorem}
\begin{proof}
Equation~\eqref{mic:inverse} gives $D_b^{(N)}O_b^{(N)}\Phi_N=P_b\Pi_N$. Apply it to the exact product \eqref{mic:evaluation_product}; then \eqref{mic:product_rate} proves \eqref{mic:commuting}.

For $|\alpha|=m\le N$,
\begin{equation}\label{mic:worst_gram}
 d_{\alpha,N}\ge(N)_{\underline m}/N^m.
\end{equation}
Indeed the right side is the probability that $m$ independent uniform site choices are all distinct. The left side imposes distinctness only within each group carrying the same source label. Equality holds when all degree is in one source direction. The degree-$m$ inverse therefore costs at most $\kappa_L\chi_{m,N}$ in the coefficient Hilbert norm. Sum the noise estimates with weights $\sigma^m$ and use
\[
 \widehat C-A\star B
 =[\widehat C-P_b(A\star_NB)]
  +P_b(A\star_NB-A\star B)-(I-P_b)(A\star B).
\]
This proves \eqref{mic:total_error}. Cofinal finite cutoffs remove the fixed target tail, and the other two terms vanish as stated.
\end{proof}

\begin{proposition}[Exact cost of increasing the observed degree]\label{mic:threshold}
For at least one retained source direction and $M\le N$, the inverse norm of the factors $\sqrt{d_{\alpha,N}}$ on degrees through $M$ is exactly $\chi_{M,N}$. Moreover
\begin{equation}\label{mic:threshold_bounds}
 \frac{M(M-1)}{4N}\le\log\chi_{M,N}
 \le\frac{M(M-1)}{4(N-M+1)}.
\end{equation}
As $N\to\infty$, $M/\sqrt N\to c<\infty$ implies
$\chi_{M,N}\to e^{c^2/4}$, while $M/\sqrt N\to\infty$ implies divergence. For sequences with $M\le N$, these inverse norms are asymptotically bounded exactly when $M^2/N$ is bounded.
\end{proposition}
\begin{proof}
The inverse is diagonal in orthogonal coordinates. Equality in \eqref{mic:worst_gram} for $\alpha=Me_1$ gives its exact norm. Use
\[
 \log\chi_{M,N}=\frac12\sum_{a=0}^{M-1}-\log(1-a/N),\qquad
 x\le-\log(1-x)\le\frac{x}{1-x}.
\]
Summation proves \eqref{mic:threshold_bounds}. If $M/\sqrt N$ is bounded, then $M/N\to0$, and the two bounds have the same limit. The lower bound gives the remaining assertions.
\end{proof}
This threshold refers to variance-one source marks with Euclidean noise. At $M>N$, a single-source degree-$N+1$ coordinate is zero and the corresponding target is invisible.

\begin{proposition}[Sampling the product marks]\label{mic:sampling}
Suppose each sample supplies the source block sums and the finite anchored matrix corner of $F_N=\Phi_N(A)\Phi_N(B)$. Put
$B_4=p_{\sqrt7}(A)p_{\sqrt7}(B)$ and
$v_m(d)=\binom{d+m-1}{m}$ for $m\ge1$, with $v_0(d)=1$.
For the degree-$m$ empirical mark vector from $J$ independent samples,
\begin{equation}\label{mic:sampling_bound}
 \begin{aligned}
 \E\norm{\overline y_{m,J}-y_m}_{\ell^2}^2
 &\le v_m(d)3^mB_4^2/J,\\
 \Prob\{\norm{\overline y_{m,J}-y_m}_{\ell^2}>\delta_m\}
 &\le v_m(d)3^mB_4^2/(J\delta_m^2),\qquad\delta_m>0.
 \end{aligned}
\end{equation}
\end{proposition}
\begin{proof}
The moment bound and H\"older give $\norm{F_N}_{L^4(\Sch_2)}\le B_4$. Each normalized test has $L^2$ norm one and pure Walsh degree $m$, so its squared $L^4$ norm is at most $3^m$. With $Z=P_LRF_NRP_L$, sum the matrix entries first to get
\[
 \E\sum_{|\alpha|=m,\,\operatorname{supp}\alpha\subset[d]}
   |\psi^{\rm on}_{\alpha,N}|^2\norm Z_2^2
 \le v_m(d)3^mB_4^2.
\]
Independent centered sample errors have zero cross terms, so averaging divides this bound by $J$. Markov's inequality gives the probability estimate.
\end{proof}
There are $L^2\binom{d+M}{M}$ complex marks in the finite budget. A union bound combines the degreewise guarantees. The target tails in \eqref{mic:total_error} are supplied by the compactness criteria of Section~\ref{sec:data} or by a concrete coefficient estimate.

The microscopic construction is in the Hilbert coefficient geometry. On a common physical corner of dimension $L$, it also yields
$\mathcal N_{\sigma,s}(K_C)\le L^{(1/s-1/2)_+}p_\sigma(C)$:
for $s\ge2$ use the Hilbert norm of each cut, and for $s<2$ use an extreme cut of rank at most $L$. Thus the same finite experiments give the corresponding local Schatten estimates. The general Gaussian algebra of Theorem~\ref{thm:main} has its independent all-exponent construction.

The represented microscopic target consists of symmetric site responses with their source labels and ordered physical coefficients. For example, signs and standard Gaussians share the same first-order central limit but have fourth moments one and three; their microscopic distributions are distinguished by information beyond that limit. In the site model the finite correction $\E S_N^4=3-2/N$ records such information. Source-labelled observations in \eqref{mic:marks} retain the coefficient target throughout the product limit.

\section{Finite reconstruction, tightness, and completion}\label{sec:data}
The finite source and physical observations now determine both reconstruction errors and the topology of their limits. We begin with the completion selected by compatible finite inverses. Hermite marks realize these inverses in the original coefficient norms, and positive mass together with anchored source innovations supplies the corresponding compactness criterion. The same construction works at the trace-class endpoint.

\subsection{The completion selected by finite inverses}
Let $A_0$ carry increasing norms $p_1\le p_2\le\cdots$. Suppose it has finite-rank projections $P_n$ satisfying
\begin{equation}\label{eq:abstractcuts}
 P_nP_m=P_{\min(n,m)},\qquad p_j(P_nx)\le p_j(x),\qquad
 A_0=\bigcup_n A_n,\quad A_n=\ran P_n.
\end{equation}
Let $A$ be its Hausdorff completion. Each $P_n$ extends continuously to $A$, and $P_nx\to x$ in every $p_j$: this holds eventually on $A_0$, and the uniform estimate $p_j((I-P_n)x)\le2p_j(x)$ extends the conclusion by density.

For each $n$, prescribe a finite-dimensional normed data space $Y_n$ and a linear isomorphism $W_n:A_n\to Y_n$. Set
\[
 O_n=W_nP_n,\qquad D_n=W_n^{-1},\qquad
 r_{nm}=W_nP_nD_m\quad(m\ge n).
\]
An array $y=(y_n)$ is compatible if $r_{nm}y_m=y_n$. For such an array define
\begin{equation}\label{eq:finitecauchy}
 a_n(y)=D_ny_n,\qquad
 q_j(y)=\sup_n p_j(a_n(y)),\qquad
 e_{j,N}(y)=\sup_{n\ge N}p_j(a_n(y)-a_N(y)).
\end{equation}
The suprema are allowed to be infinite. Admissibility is a condition on the full array; a finite experiment supplies it through a separate tail bound.

\begin{theorem}[Exact completion of Cauchy data]\label{thm:completion}
Let $Y_{\rm adm}$ consist of compatible arrays for which $e_{j,N}(y)\to0$ for every $j$. Then
\begin{equation}\label{eq:completionmap}
 O:A\longrightarrow Y_{\rm adm},\qquad O(x)=(O_nx)_n,
 \qquad Ry=\lim_nD_ny_n
\end{equation}
is a linear bijection with inverse $R$, and
\begin{equation}\label{eq:completionnorm}
 q_j(Ox)=p_j(x),\qquad e_{j,N}(Ox)=p_j(x-P_Nx).
\end{equation}
In particular, the data norms give a complete space and recover the uniform structure specified by the original norms.
\end{theorem}
\begin{proof}
Compatibility is the identity that allows all finite reconstructions to belong to one object. Namely, for $m\ge n$,
\[
 P_n a_m(y)=P_nD_my_m=D_nr_{nm}y_m=D_ny_n=a_n(y).
\]
If $y$ is admissible, then for $m,n\ge N$,
$p_j(a_m-a_n)\le2e_{j,N}(y)$. The sequence $(a_n)$ is therefore Cauchy in every defining norm. Let $x$ be its limit in $A$. Continuity of $P_N$ in each norm and the compatibility identity imply $P_Nx=a_N(y)$, so $O_Nx=y_N$ for every $N$.

Conversely, for $x\in A$, its observations are compatible and satisfy $D_nO_nx=P_nx$. Since $P_nx\to x$ in every $p_j$, the decoded sequence is admissible and reconstructs $x$. If two elements have the same observations, all their finite projections agree; passage to the limit makes the elements equal. This proves the claimed inverse formulas.

It remains to identify the uniform structure rather than merely the underlying set. Contractivity and convergence give
$\sup_n p_j(P_nx)=p_j(x)$. For $n\ge N$, nesting gives
\[
 P_nx-P_Nx=P_n(I-P_N)x,
\]
so the corresponding norm is at most $p_j((I-P_N)x)$. Its limit as $n\to\infty$ is that same number. This proves both identities in \eqref{eq:completionnorm}. Applied to differences, they identify every defining data norm with the original norm. Completeness of the data space follows from this isometry and completeness of $A$.
\end{proof}

Coordinatewise convergence generally gives a weaker topology. The Hermite example in the introduction has vanishing finite coordinates but a nonvanishing constant product response. Equally, a homeomorphism need not preserve a completion: $x\mapsto\arctan x$ is a homeomorphism onto its image, whereas $(\arctan n)$ is Cauchy and $(n)$ is not. Formula~\eqref{eq:completionnorm} records the stronger uniform statement needed here.

\begin{proposition}[Common-tail classes and uniform inversion]\label{prop:certclass}
Let $B_j<\infty$ and $b_{j,N}\downarrow0$. The compatible arrays satisfying $q_j(y)\le B_j$ and $e_{j,N}(y)\le b_{j,N}$ form a compact set $Y(B,b)$ in the coordinatewise topology, contained in $Y_{\rm adm}$. On this set, the coordinatewise and reconstructed topologies coincide, and
\begin{equation}\label{eq:unifinverse}
 p_j(Ry-Rz)\le\inf_N\left\{
 \kappa_{j,N}\norm{y_N-z_N}_{Y_N}+2b_{j,N}\right\},
 \qquad\kappa_{j,N}=\norm{D_N:Y_N\to(A_N,p_j)}.
\end{equation}
More generally, $\mathcal K\subset A$ is relatively compact if and only if every $p_j$ is bounded on $\mathcal K$ and
\begin{equation}\label{eq:tailcompact}
 \lim_N\sup_{x\in\mathcal K}p_j(x-P_Nx)=0\qquad\text{for every }j.
\end{equation}
\end{proposition}
\begin{proof}
For fixed $n$, the possible $y_n$ lie in the image under $W_n$ of the finite-dimensional closed ball $\{a\in A_n:p_1(a)\le B_1\}$. Each compatibility identity and each finite inequality defining $q_j$ and $e_{j,N}$ is closed. Thus $Y(B,b)$ is a closed subset of a countable product of compact sets. The tail assumptions imply admissibility. Splitting $Ry-Rz$ into its two tails and its finite decoded difference proves~\eqref{eq:unifinverse}. Choose $N$ to make the tails small, and then use convergence of the $N$th coordinate; this proves continuity of the inverse on the compact class.

Under~\eqref{eq:tailcompact}, any finite collection of norms can be approximated uniformly by one finite-dimensional projected family. Since the norms increase, it is enough to use their largest index. The projected family is bounded, hence totally bounded, in that norm. This gives total boundedness for the metrizable locally convex topology, and completeness yields relative compactness. Conversely, strong convergence of $P_n$ is uniform on compact sets because $I-P_n$ is uniformly bounded in each norm. A finite-net argument proves~\eqref{eq:tailcompact}.
\end{proof}

\subsection{Operations through completion}
Assume in addition that $A_0$ is a $*$-algebra, that $P_n(x^*)=(P_nx)^*$, and that for each $j$ there are indices $k(j),h(j)$ and finite constants $C_j,L_j$ with
\begin{equation}\label{eq:abstractproduct}
 p_j(xy)\le C_jp_{k(j)}(x)p_{k(j)}(y),\qquad
 p_j(x^*)\le L_jp_{h(j)}(x).
\end{equation}
Here multiplication on the finite core is already specified, or has been reconstructed from finite observations.

\begin{theorem}[Compatibility of operations and completion]\label{thm:operationcompletion}
The completion $A$ has a unique Fr\'echet $*$-algebra structure extending the core operations. The formulas
\begin{equation}\label{eq:dataproduct}
 (y\diamond z)_n=\lim_{M\to\infty}O_n(a_M(y)a_M(z)),
 \qquad (y^\dagger)_n=W_n(a_n(y)^*)
\end{equation}
define the corresponding operations on $Y_{\rm adm}$, and $O,R$ are inverse topological $*$-algebra isomorphisms. Put $t_{j,N}(x)=p_j(x-P_Nx)$. There is an increasing $\nu(N)\ge N$ with $A_NA_N\subset A_{\nu(N)}$, and, writing $k=k(j)$,
\begin{align}
 p_j\big(P_N(xy)-P_N((P_Nx)(P_Ny))\big)
 &\le C_j\{t_{k,N}(x)p_k(y)+p_k(x)t_{k,N}(y)\},\label{eq:truncdefect}\\
 t_{j,\nu(N)}(xy)
 &\le2C_j\{t_{k,N}(x)p_k(y)+p_k(x)t_{k,N}(y)\}.\label{eq:tailproduct}
\end{align}
Thus products of bounded common-tail families again have common tails, with the displayed change of resolution and norm.
\end{theorem}
\begin{proof}
Fix an output norm $p_j$ and put $k=k(j)$. If $(x_n)$ and $(y_n)$ are Cauchy in all the defining norms, they are bounded in $p_k$, and
\[
 p_j(x_ny_n-x_my_m)
 \le C_j\{p_k(x_n-x_m)p_k(y_n)+p_k(x_m)p_k(y_n-y_m)\}.
\]
Thus $(x_ny_n)$ is Cauchy in $p_j$. Applying the same estimate to two different approximating pairs shows that the product of their limits is independent of the approximations. The involution extends by its bound in \eqref{eq:abstractproduct}. These extensions are continuous, and all algebra identities pass from the dense core by continuity. Any continuous extension must have these values, proving uniqueness.

For admissible $y,z$, Theorem~\ref{thm:completion} gives $a_M(y)\to Ry$ and $a_M(z)\to Rz$. Continuity of multiplication and of $O_n$ shows that the first limit in \eqref{eq:dataproduct} equals $O_n((Ry)(Rz))$. It is consequently an admissible array. Since $P_n$ commutes with the involution, its data formula is $O_n((Ry)^*)=W_n(a_n(y)^*)$. Hence the inverse maps $O,R$ preserve both operations, with the norm identities already proved.

Finally, the finite approximation error is the actual difference
\[
 xy-(P_Nx)(P_Ny)=(x-P_Nx)y+(P_Nx)(y-P_Ny).
\]
The product bound and contractivity of $P_N$ give \eqref{eq:truncdefect}. Choose a basis of $A_N$. Its finitely many pairwise products lie in some common $A_{\nu(N)}$ because $A_0=\bigcup_mA_m$. Increasing $\nu$ if necessary makes it monotone and at least $N$. Applying $I-P_{\nu(N)}$ to the preceding difference removes $(P_Nx)(P_Ny)$ and costs at most two in $p_j$. This proves \eqref{eq:tailproduct}, including the stated uniform-tail conclusion.
\end{proof}

Finite projected multiplication itself need not be associative. With the Hermite projection through degree two and $g=H_1$,
\[
 \mu_2(\mu_2(g,g),H_2)-\mu_2(g,\mu_2(g,H_2))=3H_2,
 \qquad\mu_2(f,h)=P_2(fh).
\]
Indeed, $g^2=H_2+1$, $gH_2=H_3+2g$, and $H_2^2=H_4+4H_2+2$. The associative operation is recovered by the controlled completion, with the finite projection defect retained in~\eqref{eq:truncdefect}.

For arbitrary finite noisy data $\widetilde y_N=O_Nx+\eta_N$, the actual output $\widehat x_N=D_N\widetilde y_N$ lies in $A_N$ and satisfies
\begin{equation}\label{eq:generalnoise}
 p_j(\widehat x_N-x)\le\kappa_{j,N}\norm{\eta_N}_{Y_N}+t_{j,N}(x).
\end{equation}
If true inputs have budgets $B_x,B_y$ and errors $\varepsilon_x,\varepsilon_y$ in $p_{k(j)}$, expanding the product of the perturbed inputs gives
\begin{equation}\label{eq:abstractnoiseproduct}
 p_j(\widehat x\widehat y-xy)
 \le C_j(B_y\varepsilon_x+B_x\varepsilon_y+\varepsilon_x\varepsilon_y).
\end{equation}
Infinite compatibility of the noisy array is unnecessary for this finite estimate: the decoder always produces a core object, while its error and tail certificates determine its limiting behavior.

\subsection{Hermite observations and their finite inverse}
For orthonormal source coordinates $g_1,g_2,\ldots$, let
\[
 \Psi_\alpha(g)=\prod_j\frac{H_{\alpha_j}(g_j)}{\sqrt{\alpha_j!}},
 \qquad I_m(h_\alpha)=\sqrt{m!}\Psi_\alpha\quad(|\alpha|=m).
\]
In the rectangular setting, choose positive injective Hilbert--Schmidt anchors of Hilbert--Schmidt norm one,
$R_\cC e_\ell=a_\ell e_\ell$ and $R_\cE f_k=b_kf_k$.
For $a=(M,N_C,N_E,L)$ define
\begin{equation}\label{eq:finiteprojection}
 (P_aK)_m=\mathbf1_{m\le M}
 Q_{N_E}(S_L^{\otimes m}K_m)P_{N_C}.
\end{equation}
These projections are compatible contractions in every original radius norm and converge strongly along cofinal budgets. When input and output spaces coincide,
\[
 \ran P_{M,N,N,L}\star\ran P_{M,N,N,L}
 \subset\ran P_{2M,N,N,L}.
\]
Orthogonality yields the exact observations
\begin{equation}\label{eq:finiteHermite}
 (O_aK)_{\alpha;k\ell}
 =\E[\Psi_\alpha(g)(R_\cE I(K)R_\cC)_{k\ell}]
 =b_ka_\ell\sqrt{m!}(K_\alpha)_{k\ell}.
\end{equation}
The identity extends from finite kernels by $L^2$ continuity. The number of complex coordinates is
$N_CN_E\binom{L+M}{M}$, including degree zero; this is a coordinate count, separate from the number of independent samples.

\begin{theorem}[Finite reconstruction in the original norm]\label{thm:finitedecoder}
For $1\le s<\infty$, let $D_a$ divide each coordinate in~\eqref{eq:finiteHermite} by its nonzero weight and set all other coordinates to zero. Then $D_aO_a=P_a$. Put
\[
 \kappa_a=\left(\min_{\ell\le N_C}a_\ell\min_{k\le N_E}b_k\right)^{-1},
 \quad r_a=\min(N_C,N_E),\quad\gamma_s=(1/s-1/2)_+.
\]
If the total squared error of the retained degree-$m$ observations is at most $\delta_m^2$, then $\widehat K_a=D_a\widetilde z$ satisfies
\begin{equation}\label{eq:finiteerror}
 \Nn\rho{\widehat K_a-K}
 \le r_a^{\gamma_s}\kappa_a\sum_{m=0}^M\rho^m\delta_m
       +\Nn\rho{K-P_aK}\qquad(\rho\ge1).
\end{equation}
Every such output is a finite kernel in the original model.
\end{theorem}
\begin{proof}
The weights in \eqref{eq:finiteHermite} are nonzero, so coordinatewise division gives $D_aO_a=P_a$. Write $E_m$ for the degree-$m$ kernel produced by applying this inverse to the observation error alone. In the orthonormal symmetric tensor basis,
\[
 \sqrt{m!}\norm{E_m}_2\le\kappa_a\delta_m.
\]
This is a bound for the entire retained Hilbert tensor, including all source and matrix entries.

For $s\ge2$, every Schatten-$s$ flattening has norm at most its Hilbert--Schmidt norm. For $s<2$, use the one extreme cut whose smaller side is the physical space of dimension $r_a$. This cut has rank at most $r_a$, independently of the source dimension, and
$\norm A_s\le r_a^{1/s-1/2}\norm A_2$. Its cost is an admissible single-summand decomposition in the quotient-sum norm. In both cases,
\[
 \sqrt{m!}\norm{E_m}_{\Wc_{m,s}}
 \le r_a^{\gamma_s}\kappa_a\delta_m.
\]
Summing with the weights $\rho^m$ bounds the decoded noise. Adding $P_aK-K$ gives \eqref{eq:finiteerror}. The decoder uses finitely many source and matrix coordinates, so every output is an actual finite kernel, even when the noisy observations are not mutually compatible across resolutions.
\end{proof}

At $s=2$, the finite de-anchoring norm on normalized coefficients is exactly $\kappa_a$. The low-exponent rank power is also necessary: one source monomial and a matrix with equal nonzero singular values attain the finite-dimensional norm-conversion factor. For unequal anchor weights, their product in~\eqref{eq:finiteerror} is a general upper bound. Given certified errors and tails on a finite grid of budgets, one may choose the budget minimizing this bound, with a fixed rule for ties.

\subsection{A sampling protocol for the observations}\label{subsec:empiricalmarks}
Assume independent repetitions of the original Gaussian sample are available. In each sample, the required source coordinates are observed, $F$ can be applied to prescribed input vectors, and its output can be combined coherently with known reference vectors before measuring the squared norm. These access assumptions produce the population observations in~\eqref{eq:finiteHermite}.

Indeed, for $v=FR_\cC e_\ell$ and $u=R_\cE f_k$, polarization with our first-variable-linear convention gives
\begin{equation}\label{eq:energy_polarization}
 \begin{split}
 (R_\cE FR_\cC)_{k\ell}=\frac14\big\{&\norm{v+u}^2-\norm{v-u}^2\\
 &+i(\norm{v+iu}^2-\norm{v-iu}^2)\big\}.
 \end{split}
\end{equation}
Four coherent energy readings thus give one complex matrix entry, which is multiplied by $\Psi_\alpha$ from that same sample.

\begin{proposition}[A finite-sample certificate]\label{prop:empiricalmarks}
Let $F$ belong to the degree-$m$ $\Sch_s$-valued Gaussian chaos, $1\le s<\infty$, and suppose $\norm F_{L^2(\Sch_s)}\le B$. Retain a finite physical corner, the first $L$ source directions, and all Hermite observations of degree $m$. Put
\[
 v_m(L)=\binom{L+m-1}{m}\quad(m\ge1),\qquad v_0(L)=1.
\]
For the empirical mean $\overline z_J$ of the full retained observation vector over $J$ independent original samples,
\begin{align}
 \E_{\rm obs}\norm{\overline z_J-O_aK}_{\ell^2}^2
 &\le\frac{v_m(L)9^mB^2}{J},\label{eq:markvariance}\\
 \Prob_{\rm obs}\{\norm{\overline z_J-O_aK}_{\ell^2}>\delta\}
 &\le\frac{v_m(L)9^mB^2}{J\delta^2}\quad(\delta>0).\label{eq:markconfidence}
\end{align}
An additional instrument error vector of norm at most $\delta_{\rm inst}$ increases the certified total error to $\delta+\delta_{\rm inst}$.
\end{proposition}
\begin{proof}
For one sample, let
$Z=Q_{N_E}R_\cE FR_\cC P_{N_C}$.
The normalized anchors give $\norm Z_2\le\norm F_s$, and $Z$ remains a degree-$m$ chaos. Form the finite Hilbert vector
$Y=(\Psi_\alpha Z_{k\ell})_{\alpha,k,\ell}$. Summing the physical entries first gives
\[
 \E\norm Y_{\ell^2}^2
 =\sum_{|\alpha|=m,\ \operatorname{supp}\alpha\subset[L]}
     \E[|\Psi_\alpha|^2\norm Z_2^2].
\]
For each term, probability Cauchy--Schwarz and hypercontractivity give
$\norm{\Psi_\alpha}_{L^4}^2\le3^m$ and
$\norm Z_{L^4(\Sch_2)}^2\le3^mB^2$.
The sum is at most $v_m(L)9^mB^2$. Physical dimension is already included in the Hilbert--Schmidt sum.

For independent copies $Y^{(r)}$, all centered cross terms have zero expectation. Therefore
\[
 \E_{\rm obs}\left\|J^{-1}\sum_{r=1}^J(Y^{(r)}-\E Y)\right\|^2
 =J^{-1}\E\norm{Y-\E Y}^2\le J^{-1}\E\norm Y^2.
\]
This proves~\eqref{eq:markvariance}; Markov's inequality proves~\eqref{eq:markconfidence}. The instrument error is added by the triangle inequality.
\end{proof}

For finitely many responses or degrees, the same original samples may be used. A union bound combines their certificates, without independence between responses. For example, a failure budget $\varepsilon_j$ and tolerance $\delta_j$ for response $j$ are ensured by
$J\ge v_{m_j}(L)9^{m_j}B_j^2/(\varepsilon_j\delta_j^2)$.
This is a second-moment sampling bound. Once the data are fixed, the decoded coefficients define a polynomial on the original canonical Gaussian space; the reconstruction error is its population $L^2$ norm, rather than an empirical fitting error.

\subsection{Positive norm mass, including trace class}
Fix $1\le s<\infty$, let $s'=s/(s-1)$ when $s>1$ and $s'=\infty$ when $s=1$, and put $q=\max(2,s)$. For $T=U|T|\in\Sch_s(\cC,\cE)$ define, with value zero at $T=0$,
\begin{align}
 \Jdu(T)&=\norm T_s^{2-s}U|T|^{s-1},\label{eq:dualitymap}\\
 \Dens(T)&=\tfrac12\big(\Jdu(T)^*T\oplus T\Jdu(T)^*\big)
 =\tfrac12\norm T_s^{2-s}\big(|T|^s\oplus|T^*|^s\big).\label{eq:positivemass}
\end{align}
At $s=1$ the first formula means $\mathscr J_1(T)=\norm T_1U$, using the polar partial isometry. It is a choice of norming functional, not a differentiability assertion for the trace norm. The positive mass is unambiguous:
\[
 \mathscr D_1(T)=\tfrac12\norm T_1\big(|T|\oplus|T^*|\big).
\]
The two blocks act on the input and output spaces. Singular-value calculus gives
\begin{equation}\label{eq:massidentity}
 \norm{\Jdu(T)}_{s'}=\norm T_s,
 \qquad\Tr(\Jdu(T)^*T)=\Tr\Dens(T)=\norm T_s^2.
\end{equation}

We will use one elementary positive-compression estimate both here and at the realization boundary. For $A\ge0$ in trace class and an orthogonal projection $P$, put $a=\Tr((I-P)A)$. Then
\begin{equation}\label{eq:positivecompression}
 \norm{A-PAP}_1\le a+2\sqrt{(\Tr A)a}.
\end{equation}
Indeed, the omitted diagonal block has trace $a$, while
$PA(I-P)=(PA^{1/2})(A^{1/2}(I-P))$ has trace norm at most
$\sqrt{\Tr(PA)\,a}$ by Hilbert--Schmidt H\"older. The other off-diagonal block is its adjoint. Adding the three blocks proves the estimate.

The mass map $\Dens:\Sch_s\to\Sch_1(\cC\oplus\cE)$ is continuous for the full range $1\le s<\infty$. For $s>1$ this follows from continuity of the noncommutative Mazur maps~\cite{Ricard} and Schatten H\"older. Here is a direct endpoint argument. If $T_j\to T$ in $\Sch_1$, then $T_j^*T_j\to T^*T$ in operator norm. Polynomial approximation to the square-root function on a common bounded spectral interval gives $|T_j|\to|T|$ in operator norm. Also
$\Tr|T_j|=\norm{T_j}_1\to\norm T_1=\Tr|T|$.
For a fixed finite-rank $P$, the compressed trace norms converge and
\[
 \Tr((I-P)|T_j|)\longrightarrow\Tr((I-P)|T|).
\]
Choose $P$ with small limiting tail and apply \eqref{eq:positivecompression} to both positive operators. Their finite corners converge in trace norm, so $|T_j|\to|T|$ in trace norm. The same argument applies to $|T_j^*|$, proving continuity of $\mathscr D_1$. No continuity of the polar partial isometry is required.

If $T_j\to T$ in $L^2(\Sch_s)$, the squared norms are uniformly integrable and
\[
 \norm{\Dens(T_j)-\Dens(T)}_1\le\norm{T_j}_s^2+\norm T_s^2.
\]
Continuity gives convergence in probability of the masses; uniform integrability then gives convergence in $L^1(\Sch_1)$. Thus the averaged positive mass remains a continuous observation also at trace class.

For physical projections $P_N,Q_N$, set
\[
 \begin{gathered}
 C_NT=Q_NTP_N,\qquad
 e_N(T)=\Tr[(I-(P_N\oplus Q_N))\Dens(T)],\\
 A_N(T)=\E e_N(T),\qquad B(T)=\norm T_{L^2(\Sch_s)}.
 \end{gathered}
\]

\begin{proposition}[Spatial tails in the original norm]\label{prop:massbound}
For $1\le s<\infty$, one has
\begin{align}
 \norm{T-C_NT}_s&\le2\norm T_s^{1-2/q}e_N(T)^{1/q},\label{eq:spatialpoint}\\
 \norm{T-C_NT}_{L^2(\Sch_s)}&\le2B(T)^{1-2/q}A_N(T)^{1/q}.\label{eq:spatialrandom}
\end{align}
When $q=2$, the factor with exponent zero is taken to be one; the zero operator is treated directly.
\end{proposition}
\begin{proof}
Let $t_o=\Tr((I-Q_N)|T^*|^s)$ and $t_i=\Tr((I-P_N)|T|^s)$.
For $s\ge2$, take an orthonormal eigenbasis of the positive compression $(I-Q_N)TT^*(I-Q_N)$. Scalar Jensen for $x^{s/2}$ on each unit eigenvector, followed by summation, gives
$\norm{(I-Q_N)T}_s^s\le t_o$.
For $1\le s<2$, factor
\[
 (I-Q_N)T=(I-Q_N)|T^*|^{s/2}|T^*|^{1-s/2}U
\]
and use Schatten H\"older with exponents $2$ and $2s/(2-s)$ to obtain
$\norm{(I-Q_N)T}_s\le t_o^{1/2}\norm T_s^{1-s/2}$.
The input tail is treated identically. Now
\[
 T-C_NT=(I-Q_N)T+Q_NT(I-P_N),\qquad
 t_o+t_i=2\norm T_s^{s-2}e_N(T).
\]
The inequality $x^{1/q}+y^{1/q}\le2^{1-1/q}(x+y)^{1/q}$ gives~\eqref{eq:spatialpoint}, with the stated constant two in either exponent range. Squaring this inequality and applying probability H\"older with conjugate exponents $q/(q-2)$ and $q/2$ proves~\eqref{eq:spatialrandom}; at $q=2$ it is immediate.
\end{proof}

\begin{theorem}[A closed norm encoding for coupled random operators]\label{thm:massclosed}
Let $1\le s<\infty$. Put $Z(T)=R_\cE TR_\cC$. The map
\begin{equation}\label{eq:masscode}
 T\longmapsto\big(Z(T),\E\Dens(T)\big)
\end{equation}
is a topological embedding of $L^2(\Sch_s)$ into
$L^2(\Sch_1)\oplus\Sch_1$, with closed image. The same conclusions hold when the first data space is weakened to the Hilbert space $L^2(\Sch_2)$. The operators may be arbitrarily coupled random variables, without Gaussian or deterministic samplewise boundedness assumptions.
\end{theorem}
\begin{proof}
The estimate $\norm{R_\cE TR_\cC}_1\le\norm T_s$ gives continuity of the first coordinate in $L^2$, and the previously established $L^1$ continuity of $\Dens$ gives continuity of the second. If the first coordinates agree, their matrix entries agree after multiplication by nonzero anchor weights; hence $T=S$ almost surely. Thus the encoding is injective.

For the inverse, first recover only a finite physical corner. Its two anchor inverses are bounded, so Proposition~\ref{prop:massbound} and a head--tail decomposition give
\begin{equation}\label{eq:massinverse}
 \begin{split}
 \norm{T-S}_{L^2(\Sch_s)}\le{}&
 \kappa_N\norm{Z(T)-Z(S)}_{L^2(\Sch_1)}\\
 &+2B(T)^{1-2/q}A_N(T)^{1/q}
  +2B(S)^{1-2/q}A_N(S)^{1/q}.
 \end{split}
\end{equation}
At $q=2$ the zero exponents have the convention already stated in Proposition~\ref{prop:massbound}.

For the Hilbert first coordinate, the same finite-corner argument replaces the first term in \eqref{eq:massinverse} by
\[
 r_N^{(1/s-1/2)_+}\kappa_N\norm{Z(T)-Z(S)}_{L^2(\Sch_2)},
 \qquad r_N=\min(\rank P_N,\rank Q_N).
\]
This uses only the finite-dimensional Schatten norm comparison; it does not assume a bounded inverse of the global Hilbert observation.

Suppose both encoded coordinates of $(T_n)$ are Cauchy, using either stated first-coordinate norm. The traces of the positive masses bound $B(T_n)$ uniformly. A trace-norm Cauchy family of positive trace-class operators has uniformly vanishing physical tails: compare every sufficiently late term with one fixed term, and then truncate that term and the finitely many exceptions. The last two terms in \eqref{eq:massinverse} are consequently uniformly small for large $N$. Fix this $N$ before using the Cauchy property of the first coordinates. It follows that $(T_n)$ is Cauchy in $L^2(\Sch_s)$. Its limit exists, and continuity identifies both prescribed data limits with its encoding. This proves closedness. Applying the same argument to an encoded sequence converging to the encoding of a fixed $T$ proves continuity of the inverse on the image.
\end{proof}
The mass coordinate is additional information. In finite observations it must be supplied either by independent mass measurements or by a separate tail certificate.

\subsection{Unresolved-source convexity}
In this subsection let $1<s<\infty$ and $q=\max(2,s)$. Let $E_LT=\E[T\mid\mathcal F_L]$, where $\mathcal F_L=\sigma(g_1,\ldots,g_L)$. For $T\in L^q(\Sch_s)$ define
\begin{equation}\label{eq:sourcegap}
 b_L(T)=\E\norm T_s^q-\E\norm{E_LT}_s^q\ge0.
\end{equation}
For a source-measurable $T$, these quantities decrease to zero by conditional Jensen and convergence of conditional expectations.

\begin{proposition}[Source defects control source error]\label{prop:sourcebound}
Let $1<s<\infty$ and $q=\max(2,s)$. Set
\[
 c_s=\begin{cases}(s-1)^{-1/2},&1<s<2,\\1,&s=2,\\2^{1-1/s},&s>2.\end{cases}
\]
Then
\begin{equation}\label{eq:sourceerror}
 \norm{T-E_LT}_{L^2(\Sch_s)}\le c_s b_L(T)^{1/q}.
\end{equation}
With $S=E_LT$ and $\mathscr J_{s,q}(S)=\norm S_s^{q-2}\Jdu(S)$, the defect also has the Bregman form
\begin{equation}\label{eq:bregman}
 b_L(T)=\E\left[\norm T_s^q-\norm S_s^q
 -q\operatorname{Re}\Tr(\mathscr J_{s,q}(S)^*(T-S))\right].
\end{equation}
For $1<s\le2$, the displayed $c_s$ is optimal over all such random operators and source conditional expectations, already on a two-dimensional physical space.
\end{proposition}
\begin{proof}
Put $S=E_LT$ and $D=T-S$. We first express the defect as the expected Bregman remainder of $F(A)=\norm A_s^q$. Its derivative is
\[
 DF(S)[D]=q\operatorname{Re}\Tr(\mathscr J_{s,q}(S)^*D).
\]
The dual element is $\mathcal F_L$-measurable and has $\Sch_{s'}$ norm $\norm S_s^{q-1}$. H\"older makes its pairing with $D$ integrable, and conditioning makes the expectation of that pairing zero. This proves \eqref{eq:bregman}; convexity also makes its integrand nonnegative.

For $1<s<2$, the Ball--Carlen--Lieb inequality~\cite{BCL,RX} implies that, for fixed $S,D$, the real function
\[
 t\longmapsto\norm{S+tD}_s^2-(s-1)t^2\norm D_s^2
\]
is midpoint convex and continuous, hence convex. Its supporting line at zero gives
\[
 \norm{S+D}_s^2-\norm S_s^2
 -2\operatorname{Re}\Tr(\Jdu(S)^*D)
 \ge(s-1)\norm D_s^2.
\]
The derivative exists also at $S=0$, with value zero. Taking expectations yields
$b_L(T)\ge(s-1)\norm{T-E_LT}_{L^2(\Sch_s)}^2$,
which proves the stated constant. At $s=2$, conditional expectation is an orthogonal projection and the same statement is an equality.

For $s>2$, Clarkson--McCarthy~\cite{FK} gives
\[
 \norm{(T+S)/2}_s^s+\norm{(T-S)/2}_s^s
 \le\tfrac12(\norm T_s^s+\norm S_s^s).
\]
Conditional Jensen and $E_L((T+S)/2)=S$ bound the expectation of the first term below by $\E\norm S_s^s$. Rearrangement yields
$\E\norm{T-S}_s^s\le2^{s-1}b_L(T)$,
and $L^2\le L^s$ proves the stated bound.

For optimality when $1<s<2$, take a sign $\epsilon=\pm1$ independent of $\mathcal F_L$ and
$T_t=I_2+t\epsilon\diag(1,-1)$, $0<t<1$. Such a sign can be obtained from an unresolved Gaussian coordinate. Then $E_LT_t=I_2$,
\[
 \norm{T_t-E_LT_t}_{L^2(\Sch_s)}^2=2^{2/s}t^2,
 \qquad
 b_L(T_t)=2^{2/s}(s-1)t^2+O_s(t^4).
\]
Their ratio forces $c_s\ge(s-1)^{-1/2}$. At $s=2$, any nonzero centered example forces $c_2\ge1$. These are consequences of the classical convexity constant, now expressed in the source-defect normalization.
\end{proof}

Physical compression and source conditioning commute. Splitting
$T-C_NE_LT=(T-C_NT)+C_N(T-E_LT)$ gives
\begin{equation}\label{eq:doubletail}
 \norm{T-C_NE_LT}_{L^2(\Sch_s)}
 \le2B(T)^{1-2/q}A_N(T)^{1/q}+c_sb_L(T)^{1/q}.
\end{equation}
The two positive quantities measure different unresolved directions.

\subsection{Source innovations without uniform convexity}\label{subsec:anchored_source}
The preceding convexity argument uses $s>1$. At trace class, the positive norm mass still controls spatial escape, but its conditional norm loss can vanish without source recovery. The source is instead tested after a faithful Hilbert-valued observation. This weaker observation returns to the original norm only after physical truncation, with its finite inverse cost retained.

Let $1\le s<\infty$ and let $T\in L^2(\sigma(W);\Sch_s(\cC,\cE))$. Retain the normalized anchors and their spectral projections from the finite decoder. Set $Z(T)=R_\cE TR_\cC$, $Z_N(T)=Q_NZ(T)P_N$, and
\begin{align}
 \iota_L(T)&=\norm{Z(T)-E_LZ(T)}_{L^2(\Sch_2)}^2
 =\norm{Z(T)}_{L^2(\Sch_2)}^2-\norm{E_LZ(T)}_{L^2(\Sch_2)}^2,
 \label{eq:anchored_source_defect}\\
 \iota_{N,L}(T)&=\norm{Z_N(T)-E_LZ_N(T)}_{L^2(\Sch_2)}^2
 \le\iota_L(T).\label{eq:localized_source_defect}
\end{align}
These are orthogonal-projection residuals in $L^2(\Sch_2)$, not convexity defects of $\Sch_1$. For possibly different finite input and output dimensions, write
\[
 r_N=\min(\rank P_N,\rank Q_N),\quad
 \kappa_N=\left(\min_{e_\ell\in P_N\cC}a_\ell
                   \min_{f_k\in Q_N\cE}b_k\right)^{-1},\quad
 \zeta_{N,s}=r_N^{(1/s-1/2)_+}\kappa_N.
\]
We take nonzero finite corners, so these quantities are finite.

\begin{proposition}[Anchored innovations and original-norm tails]\label{prop:anchored_source}
For every $1\le s<\infty$, the source residuals decrease to zero as $L\to\infty$, and
$\iota_L(T)=0$ if and only if $T=E_LT$ almost surely. With $q=\max(2,s)$,
\begin{equation}\label{eq:anchored_doubletail}
 \norm{T-C_NE_LT}_{L^2(\Sch_s)}
 \le2B(T)^{1-2/q}A_N(T)^{1/q}
       +\zeta_{N,s}\iota_{N,L}(T)^{1/2}.
\end{equation}
In particular the last term can be replaced by
$\zeta_{N,s}\iota_L(T)^{1/2}$. At $s=1$ the first term is $2A_N(T)^{1/2}$.
For a family with bounded $B(T)$, uniformly vanishing spatial masses and global anchored source residuals therefore give uniform finite physical--source approximation in the original norm. The order of choice is first $N$ and then $L$. A finite number of source directions does not itself impose a finite chaos degree.
\end{proposition}
\begin{proof}
The ideal inequality gives $\norm{Z(T)}_2\le\norm T_s$ pointwise, so the Hilbert residuals are well defined. Conditional expectation is an orthogonal projection in $L^2(\sigma(W);\Sch_2)$. Nested source sigma-fields and their exhaustion of the source prove the energy identity, monotonicity, and convergence to zero. The deterministic anchors commute with conditioning. If $\iota_L(T)=0$, then $R_\cE(T-E_LT)R_\cC=0$ almost surely. Countably many matrix entries and the strictly positive anchor eigenvalues give $T=E_LT$ on one common event. The converse is immediate.

Decompose the original error before changing its norm:
\[
 T-C_NE_LT=(T-C_NT)+C_N(T-E_LT).
\]
Proposition~\ref{prop:massbound} bounds the first term. On the finite corner, left and right anchor inversion cost $\kappa_N$ in Hilbert--Schmidt norm. The rank bound then gives
\[
 \norm{C_N(T-E_LT)}_{L^2(\Sch_s)}
 \le r_N^{(1/s-1/2)_+}\kappa_N
          \norm{Z_N(T)-E_LZ_N(T)}_{L^2(\Sch_2)}.
\]
This proves \eqref{eq:anchored_doubletail}. Once $N$ has made the spatial term small, its inverse cost is a fixed finite number, so one can choose $L$ to control the source term. There is no global de-anchoring in this argument.
\end{proof}

The distinction is already visible on two physical coordinates. Let $\epsilon=\operatorname{sign}(g_{L+1})$ and
\[
 T=\diag\big((1+\epsilon)/2,(1-\epsilon)/2\big).
\]
Then $E_LT=I_2/2$ and both $\norm T_1$ and $\norm{E_LT}_1$ equal one, whereas
$\norm{T-E_LT}_{L^2(\Sch_1)}=1$. Thus the endpoint squared-norm loss is zero but the source is unresolved. If the two anchors coincide and have entries $r_1,r_2>0$ on this block, then
\begin{equation}\label{eq:trace_source_flat_face}
 \iota_L(T)=(r_1^4+r_2^4)/4>0.
\end{equation}
This example concerns general source-measurable operators; it rules out simply extending the preceding conditional-expectation inequality to $s=1$. Conversely anchored residuals alone do not control the original norm: for $T_j=g_je_je_j^*$ and one diagonal anchor $R$, $\norm{T_j}_1=|g_j|$ and, when $j>L$, $\iota_L(T_j)=r_j^4\to0$. The spatial mass escapes. Both parts of \eqref{eq:anchored_doubletail} are needed.

\begin{proposition}[A finite corner-energy certificate]\label{prop:finite_source_energy}
Suppose $T=I_mK$ is homogeneous of a fixed degree $m\ge0$, and put
$\mathscr E_N=\E\norm{Z_N(T)}_2^2$. Let $z$ be the finite vector of degree-$m$ Hermite marks in the chosen physical corner and first $L$ source coordinates. Then
\begin{equation}\label{eq:finite_source_energy_identity}
 \iota_{N,L}(T)=\mathscr E_N-\norm z_{\ell^2}^2.
\end{equation}
Suppose $\norm{\widehat z-z}_{\ell^2}\le\delta$ and $\mathscr E_N\le\mathscr E_N^+$. Define
\begin{equation}\label{eq:finite_source_energy_budget}
 U_{N,L}=\mathscr E_N^+-\norm{\widehat z}_{\ell^2}^2
                      +2\delta\norm{\widehat z}_{\ell^2}.
\end{equation}
The actual finite Hermite decoder $\widehat T=I_m\widehat K$ then satisfies, whenever the stated data are consistent,
\begin{align}
 U_{N,L}&\ge0,\qquad
 \norm{\widehat T-T}_{L^2(\Sch_s)}
 \le2B(T)^{1-2/q}A_N(T)^{1/q}+\zeta_{N,s}\sqrt{U_{N,L}},
 \label{eq:finite_source_response_certificate}\\
 \rho^m\sqrt{m!}\norm{\widehat K-K}_{\Wc_{m,s}}
 &\le(\rho v_s)^m
 \left(2B(T)^{1-2/q}A_N(T)^{1/q}+\zeta_{N,s}\sqrt{U_{N,L}}\right).
 \label{eq:finite_source_native_certificate}
\end{align}
A negative $U_{N,L}$ certifies inconsistency of the marks, their noise bound, and the corner-energy bound.
\end{proposition}
\begin{proof}
At fixed physical resolution the source conditional expectation is a finite homogeneous Hermite polynomial, with coefficient vector exactly $z$ from \eqref{eq:finiteHermite}. Parseval gives \eqref{eq:finite_source_energy_identity}. Write $\widehat Z=\sum_\alpha\Psi_\alpha\widehat z_\alpha$, interpreting each coefficient as its finite matrix. Since the unresolved Hermite part is orthogonal to every retained coordinate,
\begin{align*}
 \norm{Z_N(T)-\widehat Z}_{L^2(\Sch_2)}^2
 &=\mathscr E_N-\norm z^2+\norm{\widehat z-z}^2\\
 &=\mathscr E_N-\norm{\widehat z}^2
       +2\operatorname{Re}\ip{\widehat z}{\widehat z-z}
 \le U_{N,L}.
\end{align*}
This is the Hilbert posterior-energy identity, used before finite de-anchoring. The inverse on the corner costs $\zeta_{N,s}$, while the unresolved physical part is bounded by Proposition~\ref{prop:massbound}. This proves the response bound. The difference $\widehat T-T$ is homogeneous of degree $m$, so the original coefficient inverse from Theorem~\ref{thm:wickrecover} gives \eqref{eq:finite_source_native_certificate}. The nonnegativity assertion follows from the squared-error identity, not by replacing a negative bound by zero.
\end{proof}

The extra quantity $\mathscr E_N$ is the expectation of a finite physical energy, not an unobserved infinite trace. With $J$ independent repetitions and $B(T)\le B$, its empirical mean $\widehat{\mathscr E}_{N,J}$ satisfies
\begin{equation}\label{eq:finite_source_energy_sampling}
 \Prob_{\rm obs}\{|\widehat{\mathscr E}_{N,J}-\mathscr E_N|>\varepsilon\}
 \le\frac{9^mB^4}{J\varepsilon^2}\qquad(\varepsilon>0).
\end{equation}
Indeed, hypercontractivity gives
$\E\norm{Z_N(T)}_2^4\le9^mB^4$, and independence bounds the variance of the average. On the certified event one may take
$\mathscr E_N^+=\widehat{\mathscr E}_{N,J}+\varepsilon$.
The same samples may supply the marks: a union bound combines the two failure probabilities, without independence between these two readouts. As before, a realized decoder is evaluated in the population norm after fixing its measured coefficients. The physical tail $A_N(T)$ and the original size bound remain separate model or measurement certificates; the corner energy does not recover them for free.

\subsection{Compactness from physical mass and source innovations}
\begin{theorem}[Compactness at fixed degree and at all radii]\label{thm:fullcompact}
Let $1\le s<\infty$. At a fixed degree $m$, suppose $F_m=I_mK_m$ is a uniformly bounded family in $L^2(\Sch_s)$. Its coefficient family is relatively compact in $\Wc_{m,s}$ if and only if $A_N(F_m)$ and $\iota_L(F_m)$ tend uniformly to zero.

For $\mathcal K\subset\Alg_s$, relative compactness is equivalent to
\begin{equation}\label{eq:allradii-bounded}
 \sup_{K\in\mathcal K}\Nn\rho K<\infty\qquad\text{for every }\rho\ge1
\end{equation}
and, for every fixed $m$,
\begin{equation}\label{eq:allradii-defects}
 \lim_N\sup_{K\in\mathcal K}A_N(I_mK_m)=0,
 \qquad\lim_L\sup_{K\in\mathcal K}\iota_L(I_mK_m)=0.
\end{equation}
For $1<s<\infty$, the same criteria hold with $\iota_L$ replaced by the original convexity defect $b_L$. This replacement is an equivalence of the criteria including the spatial condition, not of the two source quantities in isolation.
Under these conditions, convergence of every fixed Hermite observation of a sequence implies convergence to a unique $K\in\Alg_s$ in every radius norm. Every fixed finite $*$-polynomial of the reconstructed objects converges to the corresponding genuine operation.
\end{theorem}
\begin{proof}
At fixed degree, first choose $N$ to make the spatial term in \eqref{eq:anchored_doubletail} uniformly small. Its de-anchoring cost $\zeta_{N,s}$ is now fixed, so the condition on $\iota_L$ allows a choice of $L$ making the remaining term small. The compressed, conditioned responses $C_NE_LF_m$ form a bounded subset of a finite-dimensional space of homogeneous matrix polynomials. Its dimension is at most
\[
 (\rank P_N)(\rank Q_N)\binom{L+m-1}{m},
\]
with the binomial interpreted as one for $m=0$. Thus the responses are totally bounded in $L^2(\Sch_s)$. The original coefficient inverse transfers this property to $\Wc_{m,s}$, including at $s=1$.

Conversely, coefficient compactness gives response compactness in $L^2(\Sch_s)$. Continuity of the positive mass, proved also at the endpoint, gives a relatively compact family of averaged masses in trace class. A finite-net argument then makes their spatial tails uniformly small. The bounded map $T\mapsto Z(T)$ takes the response family to a relatively compact subset of $L^2(\Sch_2)$. The contractions $E_L$ converge strongly to the identity on that space, hence uniformly on its compact subsets. This proves the condition on $\iota_L$ without uniform convexity of the original space.

For $1<s<\infty$, uniformly small $A_N$ and $b_L$ give the same finite-dimensional approximation by \eqref{eq:doubletail}. For necessity, hypercontractivity applied to homogeneous differences transfers response compactness to $L^q(\Sch_s)$. Conditioning converges uniformly there, and
\[
 0\le b_L(T)\le q\norm T_{L^q(\Sch_s)}^{q-1}
                         \norm{T-E_LT}_{L^q(\Sch_s)}
\]
proves uniform decay of the original source defect. This retains the interior criterion without a finite de-anchoring factor.

For the full scale, whenever $R>\rho$,
\begin{equation}\label{eq:strongradiustail}
 \sum_{m>M}\rho^m\sqrt{m!}\norm{K_m}_{\Wc_{m,s}}
 \le(\rho/R)^{M+1}\Nn R K.
\end{equation}
The stronger-radius bound makes the remaining degree tail uniform. Apply the fixed-degree result to the finitely many degrees up to $M$, and take common physical and source cutoffs. Proposition~\ref{prop:certclass} gives relative compactness in the full scale. Necessity follows from continuity of radius norms and degree projections, together with the fixed-degree criterion.

A sequence under these hypotheses is relatively compact. Its limiting Hermite marks force all cluster points to have the same source and physical coefficients, because the anchor weights are nonzero. Consequently the full sequence has one limit in $\Alg_s$. Continuity of multiplication and adjoints proves the polynomial conclusion.
\end{proof}

The result also describes compact families of entire Mehler responses. By Theorem~\ref{thm:mehler_entire}, boundedness of every coefficient radius is equivalent to boundedness on every complex disk. In this Banach-valued setting disk bounds alone are not a compactness theorem: physical and source directions can still escape. The two positive tail conditions in \eqref{eq:allradii-defects} provide precisely the missing finite-dimensional approximation at each degree. Their equivalence with original compactness also makes the qualitative criterion independent of the chosen faithful anchors; its quantitative costs depend on them.

The degree condition cannot be dropped. For $P=e_1e_1^*$, the projection onto the first anchor eigenvector, the responses
$F_j=H_j(g_1)P/\sqrt{j!}$ have $L^2(\Sch_s)$ norm one and use only one source coordinate and one physical direction. Their spatial and anchored source tails vanish after those coordinates, yet Gaussian orthogonality gives $\norm{F_j-F_k}_{L^2(\Sch_s)}=\sqrt2$ for $j\ne k$. Their original radius norms are $\rho^j$, so the stronger-radius condition detects exactly this remaining escape.

For an explicit tail budget let $B_m=\norm{I_mK_m}_{L^2(\Sch_s)}$ and define
\begin{equation}\label{eq:hybrid_source_budget}
 \Theta_{N,L}^{(s)}(T)=
 \begin{cases}
  \zeta_{N,1}\iota_{N,L}(T)^{1/2},&s=1,\\
  \min\{\zeta_{N,s}\iota_{N,L}(T)^{1/2},\ c_sb_L(T)^{1/q}\},&1<s<\infty.
 \end{cases}
\end{equation}
The two valid source estimates and the stronger-radius tail give, in the square physical setting,
\begin{equation}\label{eq:certifiedfulltail}
 \begin{split}
 \Nn\rho{K-P_{M,N,N,L}K}\le{}&
 (\rho/R)^{M+1}\Nn R K\\
 &+\sum_{m=0}^M(\rho v_s)^m
 \left[2B_m^{1-2/q}A_N(I_mK_m)^{1/q}
              +\Theta_{N,L}^{(s)}(I_mK_m)\right].
 \end{split}
\end{equation}
At $q=2$ the zero-exponent convention is as in Proposition~\ref{prop:massbound}; the degree-zero source term vanishes. The bound preserves the original convexity certificate when it is available, while the anchored alternative includes trace class. Its local source residual can be obtained from the finite corner-energy identity \eqref{eq:finite_source_energy_identity}. For noisy finite outputs, Proposition~\ref{prop:finite_source_energy} gives the corresponding posterior certificate. In either implementation the physical masses and the stronger-radius bound remain separate inputs.

\begin{proof}[Proof of Theorem~\ref{thm:main}]
Theorem~\ref{thm:nativeproduct} gives the complete coefficient $*$-algebra and its isomorphism with $\Gaus_s$. The finite projections \eqref{eq:finiteprojection} are nested, contractive in every radius norm, commute with adjoints, and exhaust the finite core along any cofinal sequence. Theorem~\ref{thm:finitedecoder} gives their exact inverses. Applying Theorems~\ref{thm:completion} and \ref{thm:operationcompletion} therefore identifies $\Data_s$ with the same coefficient completion, including its norms, tails, multiplication, and involution.

The finite-noise estimate is \eqref{eq:finiteerror} in the square physical setting, including the trace-class rank cost $N^{1/2}$ when $s=1$. For every $1\le s<\infty$, the certificate \eqref{eq:certifiedfulltail} supplies its unresolved part from positive spatial mass, anchored source innovations, and a stronger-radius bound. For $s>1$ one may instead use the original convexity source estimate whenever its budget is smaller. At trace class the source estimate is applied after the finite physical cutoff. Expanding the product of two perturbed inputs and applying \eqref{eq:nativeproduct} at the required input radius gives \eqref{eq:intro-bilinearerror}. Iterating the same estimate proves convergence and finite error control for every fixed $*$-polynomial. These steps construct the data algebra from the independently specified coefficient geometry and recover the original ordinary random product.
\end{proof}

\subsection{Energy certificates and exact recovery risks}
Let $\mathscr H$ be a Hilbert coefficient space, $P$ a finite-rank orthogonal projection, and $O=WP$, where $W$ is positive diagonal on $P\mathscr H$ with entries $w_\lambda$. Let $D=W^{-1}$. For same-source Wick kernels, use the normalized coefficients $\sqrt{m!}K_\alpha$ or the Hilbert Fock space~\eqref{eq:hilbertfock}.

If the total energy $\mathcal E=\norm x^2$ is an additional exact observation, then
\begin{equation}\label{eq:energytail}
 \norm{x-Px}^2=\mathcal E-\norm{DOx}^2.
\end{equation}
If only $\norm x^2\le E_+$ is known and $y=Ox+e$, $\norm e\le\delta$, the output $\widehat x=Dy$ satisfies
\begin{equation}\label{eq:posterior}
 0\le\norm{x-\widehat x}^2
 \le E_+-\norm{\widehat x}^2+2\delta\norm{D^*\widehat x}.
\end{equation}
Indeed, with $z=De$, orthogonality gives
$\norm{x-\widehat x}^2=\norm x^2-\norm{\widehat x}^2+
2\operatorname{Re}\ip{\widehat x}{z}$.
A negative right-hand side certifies incompatibility of the observations with the supplied bounds. A right-hand side at most $\varepsilon^2$ gives a finite stopping certificate.

For noiseless data and exact energy, the fibre is a sphere in $\ker P$ translated by $Dy$; with an energy upper bound, it is the corresponding translated ball. Provided the fibre is nonempty, its optimal enclosing radius is the square root of the energy remaining after subtracting $\norm{Dy}^2$ (with the evident singleton case when the invisible space is zero). Formula~\eqref{eq:posterior} is a valid noisy certificate, independently of whether $Dy$ is the optimal centre for a particular noisy fibre.

\begin{proposition}[Exact risk for three resources]\label{prop:threeresource}
Let $S,U$ be commuting orthogonal projections with $P=SU$, and let $O=WP$ be finite and diagonally weighted as above. Define
\[
 \mathcal K(B,\tau,\nu)=\{x:\norm x\le B,
      \ \norm{(I-S)x}\le\tau,\ \norm{(I-U)x}\le\nu\}.
\]
Assume $P\mathscr H$, $(I-S)U\mathscr H$, and $S(I-U)\mathscr H$ are nonzero. For deterministic noise $\norm e\le\delta$,
\begin{equation}\label{eq:threeresource}
 \inf_{\widehat x}\sup_{x\in\mathcal K,\ \norm e\le\delta}
 \norm{\widehat x(Ox+e)-x}
 =\min\{B,(\kappa^2\delta^2+\tau^2+\nu^2)^{1/2}\},
 \quad\kappa=\max_\lambda w_\lambda^{-1}.
\end{equation}
The infimum ranges over all reconstruction maps.
\end{proposition}
\begin{proof}
The four orthogonal sectors of $S,U$ give
\[
 \norm{(I-SU)x}^2
 =\norm{(I-S)x}^2+\norm{(I-U)x}^2
      -\norm{(I-S)(I-U)x}^2\le\tau^2+\nu^2.
\]
The finite inverse error lies in the visible sector and is orthogonal to this tail, proving the square-root upper bound. The zero output gives the bound $B$.

Choose a unit visible coordinate $v_0$ of minimal weight and unit vectors
$v_s\in(I-S)U\mathscr H$, $v_u\in S(I-U)\mathscr H$. Set
$h_0=\kappa\delta v_0+\tau v_s+\nu v_u$ and
$h=\min(1,B/\norm{h_0})h_0$, taking $h=0$ when $h_0=0$.
Both $h$ and $-h$ lie in the class, and noises $-Oh$ and $Oh$ of norm at most $\delta$ make their observations identical. Every reconstruction errs by at least $\norm h$ on one of the two models, proving equality.
\end{proof}

For a real scalar target $f\in X'$ and bounded observation $M:X\to Y$, $\delta>0$, the corresponding dual formula is
\begin{equation}\label{eq:scalar-minimax}
 \begin{split}
 \inf_a\sup_{\norm x\le1,\ \norm e\le\delta}|a(Mx+e)-f(x)|
 &=\sup_{\norm x\le1,\ \norm{Mx}\le\delta}|f(x)|\\
 &=\min_{y'\in Y'}\{\norm{f-M'y'}+\delta\norm{y'}\}.
 \end{split}
\end{equation}
The indistinguishable pair $\pm x$ gives the lower bound. On $X\oplus Y$ with norm $\max(\norm x,\norm y/\delta)$, extend the functional $(x,Mx)\mapsto f(x)$ from the graph by Hahn--Banach. The dual norm is the sum $\norm{x'}+\delta\norm{y'}$, and on the graph $x'=f-M'y'$. This proves attainment of the minimum and gives a linear optimal estimator. This is the classical optimal-recovery mechanism~\cite{MR,EF}.

\subsection{A matching rate in the original Schatten geometry}
Fix $1\le s<\infty$, $u>0$, and $\beta>1/2$. On $\ell^2$, use the trace-normalized anchor
\[
 Re_j=c_\beta j^{-\beta}e_j,
 \qquad c_\beta=\left(\sum_{j\ge1}j^{-2\beta}\right)^{-1/2}.
\]
For one fixed normalized colored Gaussian monomial $G$, put
$T_x=G\diag(x_j)$. Both its original coefficient norm and its response $L^2(\Sch_s)$ norm equal $\norm x_{\ell^s}$. Consider
\[
 \mathcal X_{s,u}(A)=\{x:\norm x_s\le A,\quad
       \norm{x_{>N}}_s\le AN^{-u}\text{ for every }N\ge1\},\quad A>0,
\]
with observations $y_j=c_\beta^2j^{-2\beta}x_j+e_j$ and $\norm e_{\ell^2}\le\delta$.
The weighted true data belong to $\ell^2$ also when $s>2$, by H\"older.

\begin{proposition}[Matching upper and lower exponents]\label{prop:matching_diagonal_rate}
Let $\zeta_s=2\beta+(1/s-1/2)_+$. The minimax $\ell^s$ risk from all weighted coordinates satisfies, for sufficiently small $\delta/A$,
\begin{equation}\label{eq:matching_diagonal_rate}
 \mathcal R(\delta)\asymp_{s,u,\beta}
 A^{\zeta_s/(u+\zeta_s)}\delta^{u/(u+\zeta_s)}.
\end{equation}
The upper bound uses only the first
$N\asymp(A/\delta)^{1/(u+\zeta_s)}$ coordinates.
\end{proposition}
\begin{proof}
Invert the first $N$ weights and set the rest to zero. The exact finite-dimensional inclusion $\ell^2_N\hookrightarrow\ell^s_N$ gives
\[
 \norm{x-\widehat x}_s\le AN^{-u}+c_\beta^{-2}\delta N^{\zeta_s}.
\]
Balancing the two terms and choosing a neighboring integer proves the upper bound.

For $s\ge2$, take $x=An^{-u}e_n$. It satisfies every tail constraint and has data norm $c_\beta^2An^{-(u+2\beta)}$. Choose an integer $n$ making this at most $\delta$. The pair $\pm x$ with opposite noises has common zero data and gives the matching lower bound.

For $s<2$, use the block
\[
 x_j=A(2N)^{-u}N^{-1/s}\mathbf1_{\{N<j\le2N\}}.
\]
Its $\ell^s$ norm is $A(2N)^{-u}$. For a cutoff below $N$, its total norm bounds the tail; between $N$ and $2N$, use $(2N)^{-u}\le n^{-u}$; above $2N$, the tail is zero. Thus all constraints hold. The weighted data norm is at most
$c_\beta^2A2^{-u}N^{-(u+\zeta_s)}$.
Choose $N$ so this is at most $\delta$, and again use $\pm x$ with common zero data. Their half-distance gives~\eqref{eq:matching_diagonal_rate}.
\end{proof}
The block witness shows why the low-exponent rank cost enters the rate. This model fixes the source, so its only unresolved tail is physical.

\subsection{Independent escape mechanisms}
Physical escape is exhibited by $T_n=g_1e_ne_n^*$: its $L^2(\Sch_s)$ norm is one, its source is resolved, and its compactly anchored finite observations vanish, while its physical tails do not vanish uniformly. Source escape is exhibited by $S_n=g_ne_1e_1^*$: its physical mass stays fixed and all marginal laws agree, but $E_LS_n=0$ for $n>L$. Both examples have degree one and are bounded at every coefficient radius.

Without fixed-degree or stronger-radius control, amplitude can also escape. On $(0,1)$, the random operators $X_n=\sqrt n\,\mathbf1_{(0,1/n)}P$ converge in probability to zero with constant $L^2$ norm. By contrast, for degree-$m$ responses of $L^2$ norm at most $B$, hypercontractivity gives
\[
 \E[\norm T_s^2\mathbf1_{\{\norm T_s>\Lambda\}}]
 \le3^{2m}B^4/\Lambda^2.
\]
Thus fixed-degree moment bounds already control amplitude tails; the all-radius condition supplies the corresponding control across degrees.

\subsection{Compatible changes of finite observations}
On the same completed space $(A,(p_j))$ of Theorem~\ref{thm:completion}, suppose there are two finite reconstruction systems
$(P_n^{(r)},W_n^{(r)},O_n^{(r)},D_n^{(r)})$, $r=1,2$.
Within each system, the projections are nested, contractive in every $p_j$, and converge strongly to the identity. No commutation between the two systems is assumed. Suppose the finite cross matrices $O_n^{(2)}D_m^{(1)}$ and $O_n^{(1)}D_m^{(2)}$ are known, and denote the two admissible data spaces by $Y_{\rm adm}^{(r)}$.

\begin{proposition}[Protocol change]\label{prop:protocolchange}
For $y\in Y_{\rm adm}^{(1)}$, the finite-data limit
\begin{equation}\label{eq:protocolchange}
 (\mathcal T_{21}y)_n=\lim_{m\to\infty}O_n^{(2)}D_m^{(1)}y_m
\end{equation}
exists and belongs to $Y_{\rm adm}^{(2)}$. The analogous map $\mathcal T_{12}$ is its inverse, and
$q_j^{(2)}(\mathcal T_{21}y)=q_j^{(1)}(y)$.
These maps preserve every continuous original operation represented by both finite systems. With
$\lambda_{n,j}=\norm{O_n^{(2)}:(A,p_j)\to Y_n^{(2)}}$,
\begin{equation}\label{eq:protocolchange_error}
 \norm{(\mathcal T_{21}y)_n-O_n^{(2)}D_m^{(1)}y_m}_{Y_n^{(2)}}
 \le\lambda_{n,j}e_{j,m}^{(1)}(y).
\end{equation}
\end{proposition}
\begin{proof}
The first finite inverse sequence converges to $x=R_1y$. Continuity of the fixed finite observation $O_n^{(2)}$ gives \eqref{eq:protocolchange}, and its value is $O_n^{(2)}x$. The whole array is therefore the admissible observation of the same object under the second protocol. Applying the reverse construction recovers the same $x$, proving the inverse relation. Both data norms equal $p_j(x)$ by the completion theorem, which also proves compatibility with continuous operations.

The exact tail identity gives
$p_j(x-D_m^{(1)}y_m)=e_{j,m}^{(1)}(y)$.
Apply $O_n^{(2)}$ to this difference to obtain \eqref{eq:protocolchange_error}. Its norm $\lambda_{n,j}$ is finite because $O_n^{(2)}=W_n^{(2)}P_n^{(2)}$, the projection is contractive, and $W_n^{(2)}$ acts on a finite-dimensional space.
\end{proof}
Thus changes of basis, cofinal resolution, or faithful anchor preserve the original reconstruction whenever the finite cross maps are available. Their finite implementation retains the displayed tail and conditioning costs.

The entire Mehler realization of Theorem~\ref{thm:mehler_entire} describes the same all-radius uniform structure, not a weaker observation completion. The compactness criterion of Theorem~\ref{thm:fullcompact} adds the physical and source tails needed beyond disk boundedness. In particular, the Hilbert residual in \eqref{eq:anchored_source_defect} is one certificate inside this original reconstruction, not a replacement target topology. This contrasts with the genuinely different targets below.

\subsection{Monotone defects and genuine limits}
The Cauchy criterion above starts with the intended completion. In a weaker observation topology, the closure can be larger. A useful way to distinguish the genuine image is to express its defect as a decreasing limit of continuous nonnegative readouts. The same elementary limit argument will apply to anchored output mass and to spectral mass, after each model has identified its own defect.

\begin{lemma}[Monotone defects and escape]\label{lem:monotone_defect}
Let $\mathcal K$ be a metric space and let $r_N:\mathcal K\to[0,\infty)$ be continuous, with $r_{N+1}\le r_N$. Put
\[
 d(z)=\inf_N r_N(z),\qquad \mathcal X=\{z:d(z)=0\}.
\]
Then $d$ is upper semicontinuous and $\mathcal X$ is a $G_\delta$ subset of $\mathcal K$. If $z_l\in\mathcal X$ and $z_l\to z$ in $\mathcal K$, then
\begin{equation}\label{eq:monotone_escape}
 d(z)=\lim_N\lim_l r_N(z_l)=\lim_N\sup_l r_N(z_l).
\end{equation}
In particular the limit belongs to $\mathcal X$ exactly when the readout tails vanish uniformly over the sequence.

Suppose also that $r_1\le B<\infty$ and that Borel probability measures $\nu_l$, concentrated on $\mathcal X$, converge weakly to $\nu$ on $\mathcal K$. Then
\begin{equation}\label{eq:monotone_law_escape}
 \int d\,d\nu
 =\lim_N\lim_l\int r_N\,d\nu_l
 =\lim_N\sup_l\int r_N\,d\nu_l.
\end{equation}
This quantity vanishes exactly when $\nu(\mathcal X)=1$.
\end{lemma}
\begin{proof}
For every $a\in\R$, the set $\{d<a\}=\bigcup_N\{r_N<a\}$ is open. Thus $d$ is upper semicontinuous, and
\[
 \mathcal X=\bigcap_{m\ge1}\bigcup_N\{r_N<1/m\}
\]
is $G_\delta$. Continuity at each fixed $N$ gives the first equality in \eqref{eq:monotone_escape} and the bound $\sup_l r_N(z_l)\ge r_N(z)\ge d(z)$.

For the reverse bound fix $\varepsilon>0$. Choose $N_0$ with $r_{N_0}(z)<d(z)+\varepsilon$, and then $l_0$ such that $r_{N_0}(z_l)<d(z)+2\varepsilon$ for $l\ge l_0$. Monotonicity preserves this bound for $N\ge N_0$. Each of the finitely many earlier $z_l$ belongs to $\mathcal X$, so its $r_N$ tends to zero. A further increase of $N$ controls them simultaneously. This proves the second equality.

For measures, each $r_N$ is bounded and continuous, so its integral converges under weak convergence. Bounded monotone convergence gives $\int r_N\,d\nu\downarrow\int d\,d\nu$ and $\int r_N\,d\nu_l\downarrow0$ for each fixed $l$. Apply the same finite-exception argument to these scalar integrals to obtain \eqref{eq:monotone_law_escape}. Nonnegativity of $d$ proves the last assertion.
\end{proof}

This lemma does not identify $\mathcal X$ with a particular operator or spectral model; that identification is a separate realization theorem. Nor does a positive defect necessarily admit a continuous detector. If $\mathcal X$ is dense and is not all of $\mathcal K$, every continuous function vanishing on $\mathcal X$ vanishes on its closure. The decreasing continuous readouts can certify small upper bounds on a defect without providing a continuous positive lower certificate at every boundary point.


\section{Ordered moments and represented algebras}\label{sec:ordered}\label{sec:minimality}
Ordered moments recover represented relations and legal multiplication. Positivity builds the cyclic space; uniform localizing bounds extend the generators; a faithful state identifies the generated algebra. PBW coordinates and Hilbert innovations organize this information, while state, trace, and replica experiments determine the distinctions retained by each protocol.

\subsection{Truncated moments and partial multiplication}
Let $\PP=\C\langle x_1,\ldots,x_d\rangle$, with $x_i=x_i^*$. The involution conjugates coefficients and reverses words. A depth-$N$ moment functional $\ell_N:\PP_{\le2N}\to\C$ satisfies
\begin{align}
 \ell_N(1)=1,\qquad \ell_N(p^*)=\overline{\ell_N(p)},\qquad
 \ell_N(p^*p)&\ge0 \quad(\deg p\le N),\label{eq:momentpositive}\\
 \ell_N(p^*x_i^2p)&\le M_i^2\ell_N(p^*p)
 \quad(\deg p\le N-1).\label{eq:momentlocalizing}
\end{align}
Write $\Gamma_N(p,q)=\ell_N(q^*p)$, linear in its first argument, and let $\mathcal N_N$ be its null space. The localizing inequality is exactly the bound needed to define multiplication on the quotient.

\begin{theorem}[Finite moments and partial multiplication]\label{thm:finitejet}
Functionals satisfying \eqref{eq:momentpositive}--\eqref{eq:momentlocalizing} correspond bijectively to labelled unitary equivalence classes of the following objects: a Hilbert space $H_N$ spanned by vectors $v_w$, $|w|\le N$, with $\norm{v_1}=1$, and operators
\[
 S_i^{(N)}:\Span\{v_w:|w|\le N-1\}\longrightarrow H_N,
 \qquad S_i^{(N)}v_w=v_{x_iw},\qquad \norm{S_i^{(N)}}\le M_i,
\]
which are symmetric on their domains. The reconstruction is
\begin{equation}\label{eq:jetinverse}
 H_N=\PP_{\le N}/\mathcal N_N,\qquad v_w=[w],\qquad
 S_i^{(N)}[p]=[x_ip].
\end{equation}
\end{theorem}
\begin{proof}
Positivity makes $\Gamma_N$ a positive semidefinite form. Cauchy--Schwarz shows that its zero-norm vectors form its radical $\mathcal N_N$, so the quotient $H_N=\PP_{\le N}/\mathcal N_N$ is a finite-dimensional Hilbert space. If $p\in\PP_{\le N-1}$ has zero class, the localizing inequality gives
$\norm{[x_ip]}^2\le M_i^2\norm{[p]}^2=0$.
Thus the partial multiplication $S_i^{(N)}[p]=[x_ip]$ is well defined and bounded by $M_i$. For $p,q$ in its domain,
\[
 \Gamma_N(x_ip,q)=\ell_N(q^*x_ip)=\Gamma_N(p,x_iq),
\]
which proves symmetry. The labelled word vectors span by construction.

Conversely, start with the stated vectors and partial operators. For a word of length at most $2N$, choose a split $u^*v$ with $|u|,|v|\le N$ and prescribe
\[
 \ell_N(u^*v)=\ip{v_v}{v_u}.
\]
We check that this prescription is independent of the split. Two adjacent permissible splits of the word $a^*x_ib$ give $(x_ia)^*b$ and $a^*(x_ib)$. Their permissibility implies $|a|,|b|\le N-1$, so symmetry on the given domain gives
\[
 \ip{v_b}{S_i^{(N)}v_a}=\ip{S_i^{(N)}v_b}{v_a}.
\]
This is exactly equality of the two proposed values. Moving the split one letter at a time connects all permissible splits; at length $2N$ there is only the central split. Hence the functional is well defined by linear extension, and reversing a split proves its Hermitian symmetry.

For $p=\sum_wp_ww$ the construction gives, in the respective degree ranges,
\[
 \ell_N(p^*p)=\Big\|\sum_wp_wv_w\Big\|^2,
 \qquad
 \ell_N(p^*x_i^2p)
 =\Big\|S_i^{(N)}\sum_wp_wv_w\Big\|^2.
\]
Positivity, normalization, and the localizing bounds follow. Finally $[p]\mapsto\sum_wp_wv_w$ is an isometry with spanning range, hence a unitary. It preserves every label and partial multiplication, so it gives exactly the claimed labelled equivalence and is unique with those properties.
\end{proof}

The last layer records vectors of degree $N$; multiplication is specified only on the preceding layers. Positivity alone would not suffice. For example, the one-variable moments $\ell(1)=\ell(x^4)=1$ and $\ell(x)=\ell(x^2)=\ell(x^3)=0$ give the positive Gram matrix $\diag(1,0,1)$ on $(1,x,x^2)$, but $[x]=0$ and $[x^2]\ne0$. The localizing bound excludes exactly this failure.

\begin{theorem}[Compatible moment towers]\label{thm:gnstower}
Suppose $\ell_{N+1}|_{\PP_{\le2N}}=\ell_N$ and the bounds $M_i$ are independent of $N$. The maps $[p]_N\mapsto[p]_{N+1}$ are isometric and intertwine partial multiplication. Completing their union gives bounded self-adjoint operators $T_i$, $\norm{T_i}\le M_i$, and a cyclic vector $\xi=[1]$ such that
\[
 \ell_N(w)=\ip{w(T)\xi}{\xi}\qquad(|w|\le2N).
\]
The cyclic realization is unique up to the unitary preserving all word vectors. Coordinatewise convergence of compatible moment arrays corresponds to weak-$*$ convergence of states on the universal unital $C^*$-algebra with these generator bounds.
\end{theorem}
\begin{proof}
For $p,q\in\PP_{\le N}$, compatibility preserves $\ell_N(q^*p)$. In particular
$\mathcal N_{N+1}\cap\PP_{\le N}=\mathcal N_N$,
so $[p]_N\mapsto[p]_{N+1}$ is a well-defined isometry. It intertwines every partial multiplication already defined at level $N$.

Form the algebraic union of these Hilbert spaces under the isometric inclusions, and define $T_i[p]=[x_ip]$ there. A vector belongs to a level on which the relevant partial multiplication is defined. The uniform bound $M_i$ therefore gives $\norm{T_i[p]}\le M_i\norm{[p]}$ on the entire union. Each $T_i$ extends to the Hilbert completion. Its symmetry extends from the dense union, so this bounded extension is self-adjoint. The vector $\xi=[1]$ is cyclic, since $w(T)\xi=[w]$ and word vectors span the union. This proves the moment identities.

If another cyclic realization has the same moments, sending $[w]$ to its word vector preserves all inner products. It extends to the unique unitary preserving the labelled cyclic vectors and intertwining the generators. The bounds $M_i$ allow the representation to extend to the universal unital $C^*$-algebra, where the cyclic vector defines a state. Conversely every state on that algebra supplies such a tower. States have norm one, and polynomials are norm dense. Convergence on every word thus extends uniformly by polynomial approximation to convergence on each algebra element, which is precisely weak-$*$ convergence of the states.
\end{proof}

For the full GNS construction the Gram radical $\mathcal N$ is a left ideal. The kernel of its representation is the two-sided $*$-ideal
\begin{equation}\label{eq:representationkernel}
 \ker\pi=\{a\in\PP:ab\in\mathcal N\text{ for every }b\in\PP\}.
\end{equation}
These two kernels serve different purposes: the first defines the cyclic space, while the second records identities of the represented algebra.

\subsection{Recovering multiplication tables}
Let $T$ be a bounded self-adjoint tuple on $H$, and let $R>0$ be injective, Hilbert--Schmidt, and normalized by $\norm R_2=1$. Set
\[
 \varphi_R(a)=\Tr(RaR),\qquad J_Rp=p(T)R.
\]
The density of $\ran R$ gives $\ker J_R=\{p:p(T)=0\}$. Choose a basis $e_a=q_a(T)R$ of $J_R(\PP_{\le N})$, with its first $r_{N-1}$ vectors spanning the previous layer. Define
\begin{equation}\label{eq:multiplicationmatrices}
 A_N(a,c)=\varphi_R(q_a^*q_c),\qquad
 H_i^{(N)}(a,b)=\varphi_R(q_a^*x_iq_b),\qquad
 M_i^{(N)}=A_N^{-1}H_i^{(N)}.
\end{equation}
Only columns $b\le r_{N-1}$ occur in $H_i^{(N)}$, so all required moments have degree at most $2N$.

\begin{proposition}[Multiplication and stabilization]\label{prop:multiplicationtable}
The $b$th column of $M_i^{(N)}$ is the coordinate vector of $x_iq_b(T)R$ in the basis $(e_a)$. Moments of degree at most $2N$ therefore determine the represented relations of degree at most $N$ and all legal compositions of partial multiplication within this degree range. If
$\dim J_R(\PP_{\le N})=\dim J_R(\PP_{\le N-1})$,
then the generated algebra is finite-dimensional and all subsequent layers stabilize.
\end{proposition}
\begin{proof}
For $b\le r_{N-1}$, the vector $x_iq_b(T)R$ lies in $J_R(\PP_{\le N})$. Write it as $\sum_ce_cm_{cb}$. Pairing with each $e_a$ gives $A_Nm_b=H_{i,b}^{(N)}$, hence $m_b=A_N^{-1}H_{i,b}^{(N)}$. The full Gram matrix also gives the null space of $J_R$ on $\PP_{\le N}$. Since $R$ has dense range, $p(T)R=0$ is equivalent to $p(T)=0$. Thus these null vectors are exactly the represented polynomial relations in the stated degree range. Composing partial multiplication is legitimate only while its intermediate vectors remain in the specified domains, giving the claimed legal products.

For stabilization let $\mathcal A_N(T)=\{p(T):\deg p\le N\}$. Right multiplication by $R$ is injective, so equality of the two consecutive dimensions gives $\mathcal A_N(T)=\mathcal A_{N-1}(T)$. Left multiplication by every generator maps $\mathcal A_{N-1}(T)$ into $\mathcal A_N(T)$, and hence into itself. This finite-dimensional space contains the identity. Induction on word length shows that it contains every word, so it is the entire generated algebra and every later layer agrees with it.
\end{proof}

The finite inverse retains its conditioning. If $A_N\ge\gamma I$, $\norm{\widehat A-A_N}\le\eta_A<\gamma$, and $\norm{\widehat H-H}\le\eta_H$, then
\begin{equation}\label{eq:multiplication-noise}
 \norm{\widehat A^{-1}\widehat H-A_N^{-1}H}
 \le\frac{\eta_H+\eta_A\norm{A_N^{-1}H}}{\gamma-\eta_A}.
\end{equation}
This follows by multiplying the difference by $\widehat A$. Stable detection of a relation space additionally requires a gap separating the nonzero Gram eigenvalues from zero.

The state $\varphi_R$ is faithful on $C^*(I,T)$. Indeed, $\varphi_R(a^*a)=\norm{aR}_2^2=0$ implies $a=0$. Its GNS representation is therefore faithful, and the full moments determine the labelled abstract $C^*$-algebra, including the norm formula
\begin{equation}\label{eq:faithfulnorm}
 \norm{p(T)}^2=
 \sup_{\varphi_R(q^*q)>0}
 \frac{\varphi_R(q^*p^*pq)}{\varphi_R(q^*q)}.
\end{equation}
This is the norm of left multiplication on the dense cyclic subspace. Recovering an ambient Hilbert-space frame uses the operator-valued anchors of Section~\ref{sec:realizationboundary}.

\subsection{PBW coordinates and represented innovations}
Let $V=\Span\{x_1,\ldots,x_d\}$ and let $\LL$ be the free Lie algebra on $V$. The PBW symmetrization $\sym:S(\LL)\to T(V)$ is a linear and coalgebra isomorphism. With the primitive coproduct $\Delta x=x\otimes1+1\otimes x$, write $\ast_c$ for convolution and set
\[
 e=\log^{\ast_c}(\Id),\qquad e^{(k)}=e^{\ast_c k}/k!,\qquad
 P_{n,\delta}=e^{(n-\delta)}|_{V^{\otimes n}},\quad0\le\delta<n.
\]
Here $n\ge1$ and $k\ge1$; each fixed degree involves finitely many terms. The maps $e^{(k)}$ project onto symmetrized products of exactly $k$ Lie primitives. To see the projection property explicitly, for a primitive $h$ evaluate the convolution series on the group-like element $\exp(th)$:
\[
 e(\exp(th))=th,\qquad
 e^{(k)}(\exp(th))=\frac{t^kh^k}{k!}.
\]
Comparison of powers of $t$ shows that $e^{(k)}(h^j)$ equals $h^j$ for $j=k$ and zero otherwise. Polarization in primitive elements gives the same statement on symmetrized products. These products span the enveloping algebra by PBW, proving that the projections are mutually annihilating and sum to the identity on positive degrees. On degree zero one adds $e^{(0)}=u\epsilon$, the unit followed by the counit projection. Thus $P_{n,0}$ is ordinary symmetrization in the original letters; other components can contain products of several nontrivial Lie primitives. These are the classical Eulerian projections of \cite{CP}.

The transported multiplication is
\begin{equation}\label{eq:pbwtransport}
 a\star_{\rm PBW}b=\sym^{-1}(\sym(a)\sym(b)).
\end{equation}
For a finite change-of-basis matrix $C$, the word Gram matrix becomes $C^*GC$, and multiplication and localizing forms transform at the same time. The dual coordinates use $P_{n,\delta}'$. These projections are generally nonorthogonal: in degree four, the Eulerian coefficient of a permutation $\sigma$ is
$(-1)^{d(\sigma)}/[4\binom3{d(\sigma)}]$.
The inverse permutations $2413$ and $3142$ have respectively one and two descents, giving transposed coefficients $-1/12$ and $1/12$.

For represented information, let $F_n$ be an increasing family of test spaces and let $J$ be the specified Hilbert-space observation. Define
\begin{equation}\label{eq:innovation}
 \mathcal O_n=\overline{J(F_n)},\qquad
 \mathcal J_n=\mathcal O_n\ominus\mathcal O_{n-1},\qquad
 Q_n=P_{\mathcal O_n}-P_{\mathcal O_{n-1}}.
\end{equation}
The quotient $\mathcal O_n/\mathcal O_{n-1}$ is canonically unitarily isomorphic to $\mathcal J_n$. If columns $W$ supply new tests, their innovation Gram matrix is $W^*Q_nW$. Replacing $W$ by $WS+VB$, where $V$ has range in the preceding layer and $S$ is invertible, changes this Gram matrix by congruence:
\[
 W^*Q_nW\longmapsto S^*(W^*Q_nW)S.
\]
Consequently the innovation space and its finite dimension are intrinsic to the filtration, while numerical eigenvalues depend on the chosen coordinates.

A formal PBW projection descends to a represented quotient exactly when it preserves $\ker\pi$. For a coordinate transformation $B$, preservation of the represented cumulative filtration means
\begin{equation}\label{eq:represented_filtration_invariance}
 \overline{J(BF_n)}=\overline{J(F_n)}\qquad\text{for every }n.
\end{equation}
Preserving each formal $F_n$ is sufficient, but is not necessary when $J$ has a kernel. For example, take $J:\C^2\to\C$, $J(a,b)=a$, $F_0=\Span\{e_1\}$, and $Be_1=e_1+e_2$, $Be_2=e_2$. The formal first layer changes, but its represented image does not. Under \eqref{eq:represented_filtration_invariance} the orthogonal innovations in \eqref{eq:innovation} are unchanged. This is the invariance relevant to observed information.

The identity $xy=(xy+yx)/2+[x,y]/2$ mixes two PBW components while preserving total word degree. For Pauli matrices $X,Z$, the represented dimensions through total degrees $0,1,2,3$ are $1,3,4,4$. At degree two the symmetric part adds no direction, whereas $[X,Z]=2XZ$ supplies the remaining one.
A vanishing innovation is therefore interpreted in the represented quotient: the relevant tests add no vector beyond the preceding layer. It need not indicate missing observations, and it cannot by itself certify failure of algebra reconstruction. Nondegeneracy and conditioning are asked only for the independent target directions that the specified representation and protocol retain.

\subsection{Shorting and noisy innovations}
The target-null shorting of Section~\ref{sec:targetquotients} also describes removal of tests already present in a filtration. Here the additional issue is perturbation of the Gram matrix itself, rather than noise in a fixed linear observation. For bounded column operators $V:E_0\to\HH$ and $W:E_1\to\HH$, put
\[
 G=\begin{pmatrix}A&B\\B^*&C\end{pmatrix}
 =\begin{pmatrix}V^*V&V^*W\\W^*V&W^*W\end{pmatrix},\qquad
 P=P_{\overline{\ran V}},\qquad \mathfrak S=W^*(I-P)W.
\]
Polar decomposition gives $V(A+\lambda I)^{-1}V^*\to P$ strongly, and hence
\begin{equation}\label{eq:shorting}
 \ip{\mathfrak Sy}y=\inf_x\norm{Vx+Wy}^2,\qquad
 \mathfrak S=C-\operatorname*{s-lim}_{\lambda\downarrow0}
 B^*(A+\lambda I)^{-1}B.
\end{equation}
The distance identity is Lemma~\ref{lem:relative_residual}. Taking this infimum in the quadratic-form order also proves that $0\oplus\mathfrak S$ is the largest positive operator below $G$ whose range lies in the second summand. This is the shorted operator \cite{AT,Ando}. If $\ran V$ is closed, the expression uses the bounded Moore--Penrose inverse.

In finite dimensions positivity gives $\ran B\subset\ran A$. Indeed, if $x\in\ker A$, positivity of the quadratic form at $(tx,y)$ for all real and complex $t$ forces $B^*x=0$. Thus, with $L=A^\dagger B$,
\begin{equation}\label{eq:innovation_triangular_gram}
 G=
 \begin{pmatrix}I&L\\0&I\end{pmatrix}^{\!*}
 \begin{pmatrix}A&0\\0&\mathfrak S\end{pmatrix}
 \begin{pmatrix}I&L\\0&I\end{pmatrix}.
\end{equation}
This identity retains the cross information: $B=AL$ and $C=\mathfrak S+L^*AL$. Storing only the innovation would discard the regression. Lemma~\ref{lem:relative_residual} shows that the underlying orthogonal elimination can be performed in stages, independently of a closed-range assumption.

\begin{proposition}[Relations and the inverse cost of an innovation tower]\label{prop:innovation_decoder}
Let $C_N:E_N\to\mathcal K$ be a column map from a finite-dimensional Hilbert space, with a fixed filtration-adapted orthonormal coordinate system. No independence of its columns is assumed. Iterating \eqref{eq:innovation_triangular_gram} gives
\[
 G_N=C_N^*C_N=U_N^*D_NU_N,\qquad
 D_N=\diag(\mathfrak S_0,\ldots,\mathfrak S_N)\ge0,
\]
where $U_N$ is block upper triangular with identity diagonal blocks. Write $D_N^{\dagger/2}=(D_N^\dagger)^{1/2}$, zero on $\ker D_N$. For a linear target $L_N:E_N\to Z$ into a Banach space, recovery from $C_Nx$ is possible exactly when
\begin{equation}\label{eq:innovation_relation_test}
 \ker C_N\subset\ker L_N
 \quad\Longleftrightarrow\quad
 L_NU_N^{-1}|_{\ker D_N}=0.
\end{equation}
Under this condition, the least norm of a bounded linear decoder is
\begin{equation}\label{eq:innovation_exact_cost}
 \kappa_N(L_N)=\norm{L_NU_N^{-1}D_N^{\dagger/2}}
              =\norm{L_NG_N^{\dagger/2}}.
\end{equation}
If the condition fails, set $\kappa_N(L_N)=\infty$. With unrestricted $x\in E_N$ and deterministic data noise of norm at most $\delta>0$ in $\mathcal K$, the optimal worst-case target error, over arbitrary decoders, is $\delta\kappa_N(L_N)$. In the incompatible case even the noiseless worst-case error is infinite.

Suppose $E_N$ increase with dense union in a Hilbert space $E$, and $C_N=J|_{E_N}$ and $L_N=L|_{E_N}$ for bounded linear maps $J:E\to\mathcal K$ and $L:E\to Z$. Then
\begin{equation}\label{eq:innovation_tower_criterion}
 L\text{ is stably recoverable from }J
 \quad\Longleftrightarrow\quad
 \sup_N\kappa_N(L_N)<\infty.
\end{equation}
The constants are nondecreasing; their supremum is the least global decoder norm. Multiplying it by $\delta>0$ gives the unrestricted deterministic noise risk.
\end{proposition}
\begin{proof}
Positivity of each leading Gram block gives the range inclusion used in \eqref{eq:innovation_triangular_gram}. Hence the elimination can be iterated even when a leading block is singular, using its Moore--Penrose inverse. Each elementary triangular factor is invertible, with identity diagonal. Induction gives the displayed factorization and $D_N\ge0$. In particular,
\[
 \norm{C_Nx}^2=\norm{D_N^{1/2}U_Nx}^2,
 \qquad \ker C_N=U_N^{-1}\ker D_N.
\]
Equality of observations determines a linear target exactly when the target annihilates this kernel, proving \eqref{eq:innovation_relation_test}.

Assume that condition and put $P_N^D=P_{(\ker D_N)^\perp}$. The map
\[
 Q_N=C_NU_N^{-1}D_N^{\dagger/2}
\]
satisfies $Q_N^*Q_N=P_N^D$ and is a partial isometry onto $\ran C_N$. Since $C_NU_N^{-1}$ annihilates $\ker D_N$,
\[
 C_N=Q_ND_N^{1/2}U_N.
\]
Write $K_N=L_NU_N^{-1}D_N^{\dagger/2}$. The relation condition implies $L_NU_N^{-1}=L_NU_N^{-1}P_N^D$. Consequently the decoder
\begin{equation}\label{eq:innovation_canonical_decoder}
 \mathscr D_N=K_NQ_N^*
\end{equation}
satisfies $\mathscr D_NC_N=L_N$. It vanishes on $(\ran C_N)^\perp$ and has norm $\norm{K_N}$: $K_N=K_NP_N^D$ and $Q_N$ is isometric on that support. Every decoder must agree with this one on $\ran C_N$, so this is the least possible norm. Applying the same argument to the polar factorization $C_N=\widetilde Q_NG_N^{1/2}$ proves the second expression in \eqref{eq:innovation_exact_cost}. In either form the constant is the intrinsic ratio
\[
 \sup_{C_Nx\ne0}\frac{\norm{L_Nx}}{\norm{C_Nx}},
\]
with value zero if both maps vanish.

The linear decoder gives the upper noise bound. For the lower bound, the states $\pm x$ with $\norm{C_Nx}\le\delta$ can both produce datum zero, using errors $\mp C_Nx$. Every reconstruction then misses one target by at least $\norm{L_Nx}$. Taking the supremum proves the exact risk. If $C_Nx=0$ but $L_Nx\ne0$, arbitrary scalar multiples of this pair give infinite error already without noise.

In the tower, the finite constants are the best constants in $\norm{Lx}\le C\norm{Jx}$ on the increasing subspaces $E_N$, and therefore are nondecreasing. A common finite bound passes by density and continuity to all of $E$. The rule $D(Jx)=Lx$ then extends uniquely to a bounded map on $\overline{\ran J}$, and orthogonal projection extends it to all data. Conversely a global decoder bounds each finite restriction. The same two-point argument identifies the global risk, whether the supremum is finite or infinite.
\end{proof}

The relation test is an exact-data condition. Small absolute Gram error alone does not determine the null space: for $C_\varepsilon=\diag(1,\varepsilon)$ and $L(x_1,x_2)=x_2$, the cost is $\varepsilon^{-1}$ when $\varepsilon>0$, while at $\varepsilon=0$ the target is not determined at all. The Grams nevertheless converge in norm. A noisy relation decision therefore needs a rank or gap certificate; regularizing the residual does not make this algebraic issue disappear.

The norm in \eqref{eq:innovation_exact_cost} is unchanged by a simultaneous invertible coordinate change $(C_N,L_N)\mapsto(C_NB,L_NB)$: the ratio ranges over the same pairs of observation and target values. Individual innovation eigenvalues and triangular factors need not be invariant. The formulas concern a specified column observation into $\mathcal K$; scalar Gram data alone determine the cost, not an unobserved orientation of these columns in an ambient space.

The kernel test cannot be replaced by applying a pseudoinverse. For $C(x_0,x_1)=x_0+x_1$,
\[
 G=\begin{pmatrix}1&1\\1&1\end{pmatrix}
 =\begin{pmatrix}1&1\\0&1\end{pmatrix}^{\!*}
   \begin{pmatrix}1&0\\0&0\end{pmatrix}
   \begin{pmatrix}1&1\\0&1\end{pmatrix}.
\]
The target $x_0+x_1$ respects the relation and has recovery cost one. The target $x_1$ does not: $C(t,-t)=0$ for every $t$. For this second target the expression $LU^{-1}D^{\dagger/2}$ is nevertheless zero. It is not a recovery cost because \eqref{eq:innovation_relation_test} fails. A zero innovation is harmless for a target that factors through the represented relation, and fatal for a target that does not.

\begin{corollary}[Finite decoders and uniform approximation]\label{cor:decoder_convergence}
In the tower of Proposition~\ref{prop:innovation_decoder}, suppose $\sup_N\kappa_N(L_N)<\infty$. Let $D:\mathcal M\to Z$ be the global inverse, where $\mathcal M=\overline{\ran J}$, and let $\Pi_N$, $\Pi$ be the orthogonal projections onto $J(E_N)$ and $\mathcal M$. The zero-extended finite decoders in \eqref{eq:innovation_canonical_decoder} satisfy
\begin{equation}\label{eq:decoder_projection_identity}
 \mathscr D_N=D\Pi_N,\qquad
 \mathscr D_Ny\longrightarrow D\Pi y\quad(y\in\mathcal K),\qquad
 \norm{\mathscr D_N}\uparrow\norm D.
\end{equation}
Convergence is uniform on each compact subset of $\mathcal K$. Moreover
\begin{equation}\label{eq:decoder_compact_criterion}
 \norm{\mathscr D_N-D\Pi}_{B(\mathcal K,Z)}\longrightarrow0
 \quad\Longleftrightarrow\quad D:\mathcal M\to Z\text{ is compact}.
\end{equation}
\end{corollary}
\begin{proof}
The finite decoder equals $D$ on $J(E_N)$, and both sides of the first identity vanish on its orthogonal complement. The spaces $J(E_N)$ increase and have dense union in $\mathcal M$, since $J$ is continuous and $\bigcup_NE_N$ is dense in $E$. Therefore $\Pi_N\to\Pi$ strongly. This proves pointwise convergence; the uniformly bounded decoder norms and a finite-net argument give uniform convergence on compact sets. The norm limit follows also from the increasing optimal constants in Proposition~\ref{prop:innovation_decoder}.

If the operator-norm convergence in \eqref{eq:decoder_compact_criterion} holds, $D\Pi$ is a norm limit of finite-rank operators and is compact, so its restriction $D$ is compact. Conversely, suppose $D$ is compact but the norm convergence fails. There are unit vectors $y_N$ along a subsequence with $\norm{D(\Pi-\Pi_N)y_N}\ge\varepsilon>0$. Set $z_N=(\Pi-\Pi_N)y_N\in\mathcal M$. These vectors are bounded and converge weakly to zero, because for each fixed $v\in\mathcal M$,
\[
 |\ip{z_N}v|\le\norm{(\Pi-\Pi_N)v}\longrightarrow0.
\]
A compact operator maps a bounded weakly null sequence to a norm-null sequence. Indeed every norm-convergent subsequence of its images has weak limit zero, so compactness allows no other cluster point. This contradicts the lower bound and proves the assertion.
\end{proof}

Even $J=L=I$ on an infinite-dimensional Hilbert space separates these conclusions: the global inverse is stable, but $\mathscr D_N=\Pi_N$ and $\norm{I-\Pi_N}=1$ at every finite stage. Stable inversion alone is not uniform finite-dimensional recovery on the entire data unit ball. The common-tail classes of Proposition~\ref{prop:certclass} supply a different, class-dependent way to obtain uniform recovery.

\begin{example}[Uniformly positive layers without a uniform inverse]\label{ex:innovation_conditioning}
On $\ell^2(\mathbb N_0)$ let $Se_j=e_{j+1}$ and $J=I-S$. The map $J$ is bounded and injective. For $F_N=\Span\{e_0,\ldots,e_N\}$ its columns are $v_j=e_j-e_{j+1}$, and their Gram matrix has diagonal entries two and adjacent off-diagonal entries minus one. If $d_N$ denotes its determinant, expansion at an end gives
$d_N=2d_{N-1}-d_{N-2}$ with $d_{-1}=1$, $d_0=2$, hence $d_N=N+2$. Its scalar innovation at stage $N$ is therefore
\begin{equation}\label{eq:innovation_shift_pivots}
 \mathfrak S_N=\frac{d_N}{d_{N-1}}=\frac{N+2}{N+1}\ge1.
\end{equation}
Nevertheless the sine vectors diagonalize this tridiagonal matrix: its eigenvalues are $2-2\cos(k\pi/(N+2))$, $1\le k\le N+1$. Thus
\begin{equation}\label{eq:innovation_shift_inverse}
 \norm{(J|_{F_N})^{-1}}
 =\frac1{2\sin(\pi/(2(N+2)))}\sim\frac{N+2}{\pi}.
\end{equation}
For completeness, substitution of the vector with entries
$\sin((j+1)k\pi/(N+2))$, $0\le j\le N$, gives the displayed eigenvalue by the addition formula for sine, with zero boundary entries at $j=-1,N+1$. Injectivity of $J$ follows from $x_0=0$ and then $x_j=x_{j-1}$ whenever $Jx=0$.

The example is already a first-chaos observation model. For a rank-one projection $P$ and $x\in\ell^2$, let $F_x=\sum_{j\ge0}x_jg_jP$. With $g_{-1}=x_{-1}=0$, the source marks
$\E[(g_k-g_{k-1})F_x]=(x_k-x_{k-1})P$
are exactly $Jx$. Every new column contributes a nondegenerate orthogonal residual, but the original coefficient inverse is unbounded. In \eqref{eq:innovation_exact_cost} the loss is carried by the triangular regression factor, not by a vanishing diagonal innovation. The target $x\mapsto x_0$ behaves differently: its finite optimal cost is $\sqrt{(N+1)/(N+2)}$, because the first diagonal entry of $G_N^{-1}$ is the ratio of the determinants of the length-$N$ and length-$(N+1)$ tridiagonal matrices. Its full-space decoder is simply the first observed coordinate, with norm one. Thus the same tower separates stable target recovery from unstable full recovery.
\end{example}

Averaging and shorting are also different operations: for a random scalar $T=\pm1$, the sample Gram matrices of $(1,T)$ have zero innovation, whereas their average is the identity and has innovation one.

\begin{theorem}[Noise-calibrated shorting]\label{thm:noiseshort}
Let $G,\widehat G\ge0$, $\norm{\widehat G-G}\le\eta$, and $\eta>0$. Set
$\mathfrak S_\lambda(G)=C-B^*(A+\lambda I)^{-1}B$ for $\lambda>0$, with $\mathfrak S_0(G)=\mathfrak S$. For $\lambda\ge\eta$,
\begin{equation}\label{eq:shortsandwich}
 \mathfrak S_{\lambda-\eta}(G)-\eta I
 \le\mathfrak S_\lambda(\widehat G)
 \le\mathfrak S_{\lambda+\eta}(G)+\eta I.
\end{equation}
In particular $\widehat{\mathfrak S}=\mathfrak S_\eta(\widehat G)$ obeys
\begin{equation}\label{eq:shorterror}
 \norm{\widehat{\mathfrak S}-\mathfrak S}
 \le\eta+\beta(2\eta),\qquad
 \beta(t)=\norm{(PW)^*t(VV^*+tI)^{-1}PW}.
\end{equation}
If $PW=(VV^*)^{\theta/2}Z$, $\norm Z\le R_0$, and $0<\theta\le1$, then
$\beta(t)\le c_\theta R_0^2t^\theta$, where
$c_\theta=\theta^\theta(1-\theta)^{1-\theta}$ for $\theta<1$ and $c_1=1$.
The exponent $\theta$ is optimal in this deterministic perturbation model.
\end{theorem}
\begin{proof}
Completing the square gives
\[
 \ip{\mathfrak S_\lambda(G)y}y
 =\inf_x\{\norm{Vx+Wy}^2+\lambda\norm x^2\}.
\]
The two quadratic forms differ by at most $\eta(\norm x^2+\norm y^2)$. Add the penalty and take the infimum in $x$ to obtain \eqref{eq:shortsandwich}, including $\lambda=\eta$ through the variational formula at zero. The identity
\[
 V(A+tI)^{-1}V^*=VV^*(VV^*+tI)^{-1}
\]
shows that $\mathfrak S_t(G)-\mathfrak S(G)$ is the positive operator defining $\beta(t)$. Substitution of the source condition gives
\[
 \beta(t)\le\norm Z^2\sup_{u\ge0}\frac{tu^\theta}{u+t}
 =c_\theta R_0^2t^\theta.
\]

For optimality take, for $0<t\le1$,
\[
 G_t^{(1)}=\begin{pmatrix}t&t^{(1+\theta)/2}\\t^{(1+\theta)/2}&1\end{pmatrix},\qquad
 G_t^{(2)}=\begin{pmatrix}2t&t^{(1+\theta)/2}\\t^{(1+\theta)/2}&1\end{pmatrix}.
\]
Both are positive, their distance is $t$, and their shorted operators are $1-t^\theta$ and $1-t^\theta/2$. Realize them with
\[
 V_j1=\sqrt{a_j}e_1,\qquad
 W_j1=(b/\sqrt{a_j})e_1+\sqrt{1-b^2/a_j}\,e_2,
 \quad a_1=t,\ a_2=2t,\ b=t^{(1+\theta)/2}.
\]
The source-condition factors have norms $1$ and $2^{-(1+\theta)/2}$. Their common midpoint is admissible observed data at noise level $\eta=t/2$, and every estimator has error at least $t^\theta/4$ on one of these two models. A uniform $o(\eta^\theta)$ estimate is therefore impossible.
\end{proof}

The norm bias $\beta(t)$ need not vanish without an additional condition. For $Ve_j=j^{-1}e_j$ and $W=I$ on $\ell^2(\N)$, the range of $V$ is dense, so $\mathfrak S=0$, but
$\norm{t(VV^*+tI)^{-1}}=1$ for every $t>0$. The source condition in the theorem is one quantitative way to obtain norm convergence; strong convergence of the resolvent alone is insufficient.

The regularized decomposition retains its penalty:
\[
 \mathfrak S_\lambda=(W-VL_\lambda)^*(W-VL_\lambda)
                    +\lambda L_\lambda^*L_\lambda,
 \qquad L_\lambda=(A+\lambda I)^{-1}B.
\]
For raw finite data within $\eta$ of a positive Gram matrix, taking the self-adjoint part and then its positive part gives a positive estimate within $2\eta$ of the true matrix. Additional moment or anchor constraints remain separate finite-data conditions.

\subsection{Minimality of finite state information}
For self-adjoint free generators, let $\PP_{\le N}^{\rm sa}$ be the real vector space of self-adjoint polynomials of degree at most $N$.
\begin{theorem}[Finite state information]\label{thm:positive_minimality}
Let $S$ be a real linear subspace of $\PP_{\le N}^{\rm sa}$ containing $1$, and let $p\in\PP_{\le N}^{\rm sa}$. Over all self-adjoint contraction tuples and states, the value $\varphi(p(A))$ is determined by the values $\varphi(s(A))$, $s\in S$, if and only if $p\in S$. If $p\notin S$, the two indistinguishable models can use the same positive contraction tuple on a space of dimension
\[
 D=1+d+\cdots+d^N
\]
and two density matrices $\rho_\pm\ge I/(2D)$.
\end{theorem}
\begin{proof}
Membership in $S$ gives a linear readout. For the converse, on the truncated full Fock space
$\HH_N=\bigoplus_{k=0}^N(\C^d)^{\otimes k}$
let $L_ie_w=e_{iw}$ when $|w|<N$ and let $L_i$ vanish on the last layer. Put
\begin{equation}\label{eq:positive_fock_witness}
 A_i=(2I+L_i+L_i^*)/4.
\end{equation}
Since $\norm{L_i}\le1$, these are positive contractions. If $q\ne0$ has highest homogeneous part $q_n$, then
\begin{equation}\label{eq:fock_top_injective}
 P_nq(A)\Omega=4^{-n}\sum_{|w|=n}(q_n)_we_w\ne0.
\end{equation}
Only a product of $n$ creation terms can reach the $n$th layer from the vacuum. Consequently evaluation is injective on $\PP_{\le N}$.

In the real Hilbert space of Hermitian matrices with inner product $\Tr BC$, let
$R_p=p(A)-\operatorname{Proj}_{S(A)}p(A)$.
Then $R_p\ne0$ and $\Tr R_p=0$. For $0<\epsilon\le(2D\norm{R_p})^{-1}$ set
$\rho_\pm=I/D\pm\epsilon R_p$.
Both have trace one and are bounded below by $I/(2D)$. Their retained observations agree, whereas their target values differ by
\[
 \Tr((\rho_+-\rho_-)p(A))=2\epsilon\Tr R_p^2>0.
\]
\end{proof}

The finite witness embeds into any sufficiently large, in particular infinite-dimensional, Hilbert space. Strict positivity is asserted on its supporting block. It also yields explicit dual tests. If $q_0=1,q_1,\ldots,q_{D-1}$ is a real basis of $\PP_{\le N}^{\rm sa}$, put $B_j=q_j(A)$ and $G_{jk}=\Tr B_jB_k$. Then
\begin{equation}\label{eq:finite_dual_witness}
 H_j=\sum_k(G^{-1})_{jk}B_k,\qquad \Tr(H_jB_k)=\delta_{jk}.
\end{equation}
For $j\ne0$, $\Tr H_j=0$. Taking
$\epsilon_j=[2D\max(1,\Tr H_j^2)]^{-1}$
in $I/D\pm\epsilon_jH_j$ changes only the $j$th basis observation. Thus even strict positivity leaves each nonconstant coordinate potentially necessary within a fixed complete Hermitian PBW basis and fixed degree budget.

\subsection{Trace observations and the cyclic quotient}
Define the linear cyclic subspace and quotient
\begin{equation}\label{eq:cyclic_quotient}
 \mathscr C_{\le N}=\Span_\C\{uv-vu:|u|+|v|\le N\},\qquad
 \mathscr Q_N=\PP_{\le N}/\mathscr C_{\le N}.
\end{equation}
This quotient identifies cyclic rotations of words. Its degree-$n$ dimension is the necklace count
\begin{equation}\label{eq:necklace_count}
 c_n(d)=\frac1n\sum_{a\mid n}\phi(a)d^{n/a},\qquad n\ge1,
\end{equation}
by Burnside's lemma; here $\phi$ is Euler's totient function.

\begin{lemma}[The finite-degree trace kernel]\label{lem:cyclic_trace_kernel}
Let $M\ge\max(1,N)$. For $p\in\PP_{\le N}$,
\begin{equation}\label{eq:trace_kernel_iff}
 \tr_Mp(A)=0\text{ for all positive contraction tuples }A
 \quad\Longleftrightarrow\quad p\in\mathscr C_{\le N},
 \qquad \tr_M=M^{-1}\Tr.
\end{equation}
\end{lemma}
\begin{proof}
Cyclicity of trace proves one direction. Conversely, strictly positive tuples strictly below the identity form a nonempty open subset of the real space of Hermitian tuples. The trace polynomial therefore vanishes on every Hermitian tuple. Scalar dilation separates homogeneous degrees, including the constant term.

Suppose the coefficient sum of a cyclic class $[w]$ in a homogeneous part $p_n$ is nonzero, where $w=i_1\cdots i_n$. Use an oriented $n$-cycle, with independent complex edge parameters $z_1,\ldots,z_n$, and set
\[
 X_a(z)=\sum_{j:i_j=a}(z_jE_{j,j+1}+\bar z_jE_{j+1,j}),
\]
with cyclic indices, embedded in $M$ dimensions. The coefficient of $z_1\cdots z_n$ in $\tr_Mp_n(X(z))$ is
\[
 \frac1M\sum_{j=0}^{n-1}(p_n)_{\operatorname{rot}^jw}
 =\frac{n}{M|[w]|}\sum_{v\in[w]}(p_n)_v\ne0.
\]
A closed walk using each positively oriented edge exactly once must follow one of these rotations. For $n=1,2$, loops or parallel matrix entries in this testing construction are still distinguished by their separate parameters. Since $(z,\bar z)$ is an invertible complex-linear change of coordinates from the real and imaginary parts, the resulting polynomial cannot vanish on all parameter values. The contradiction shows that every cyclic coefficient sum is zero, which is exactly membership in $\mathscr C_{\le N}$.
\end{proof}
The relation between trace identities and cyclic equivalence is classical \cite{BCKP,KS}. The dimension hypothesis above excludes additional low-dimensional trace identities in the stated degree range.

\begin{theorem}[Minimal random trace information]\label{thm:trace_minimality}
Let $M\ge N\ge1$, let $S$ be a real linear subspace of $\PP_{\le N}^{\rm sa}$ containing $1$, and let $p\in\PP_{\le N}^{\rm sa}$. Over arbitrary laws of positive contraction tuples in dimension $M$,
\begin{equation}\label{eq:trace_minimality_iff}
 \E\tr_Mp(A)\text{ is determined by }\{\E\tr_Ms(A):s\in S\}
 \quad\Longleftrightarrow\quad[p]\in[S]\subset\mathscr Q_N.
\end{equation}
When determination fails, two witness laws can be supported on the same set of at most $\dim_\R[S]+1$ tuples, with strictly positive weights at every support point.
\end{theorem}
\begin{proof}
On the compact space of positive contraction tuples write $f_p(A)=\tr_Mp(A)$. Lemma~\ref{lem:cyclic_trace_kernel} identifies the kernel of $p\mapsto f_p$, so the condition on the right gives linear determination. Otherwise choose a real basis $f_1=1,f_2,\ldots,f_k$ of $f_S$. Linear independence of $f_1,\ldots,f_k,f_p$ provides $k+1$ points $A^{(j)}$ with invertible full evaluation matrix. Indeed, if their evaluation columns never spanned the full space, a nonzero linear combination of the functions would vanish everywhere.

Choose a nonzero vector $c$ in the null space of the first $k$ evaluation rows. Then
\[
 \sum_jc_jf_i(A^{(j)})=0,\qquad
 \sum_jc_jf_p(A^{(j)})\ne0,\qquad \sum_jc_j=0.
\]
For sufficiently small $\epsilon>0$, the weights
$\mu_\pm(A^{(j)})=(k+1)^{-1}\pm\epsilon c_j$
are strictly positive. They define the required two probability laws.
\end{proof}

For example,
\[
 \Tr[A,B]=0,\qquad
 \Tr\big(B[A,[A,B]]\big)=-\Tr([A,B]^2).
\]
The second identity follows from $\Tr(X[Y,Z])=\Tr([X,Y]Z)$. In the cyclic quotient, both expressions are represented by PBW terms of factor deficit two. Before taking traces, $B[A,[A,B]]$ also has the primitive term $\tfrac12[B,[A,[A,B]]]$; its trace is zero. Counting independent trace information therefore requires first passing to the cyclic quotient.

\subsection{Matrix-valued sources and ordered traces}
For Hermitian matrices $G_1,\ldots,G_d\in M_M(\C)$ define the analytic germ
\begin{equation}\label{eq:matrix_source_germ}
 \mathcal M_G^{(a)}(X)=\det\left(I-\sum_jX_j\otimes G_j\right)^{-1},
 \qquad X_j\in M_a(\C).
\end{equation}
For Hermitian sources with $\norm{\sum_jX_j\otimes G_j}<1$, this is a circular Gaussian quadratic-form exponential moment. Complex sources below are interpreted through its analytic continuation near zero.

\begin{proposition}[Cycle extraction]\label{prop:matrix_word_source}
For $w=i_1\cdots i_\ell$, take $a=\ell$ and
$X_j(t)=\sum_{b:i_b=j}t_bE_{b,b+1}$, where $E_{\ell,\ell+1}=E_{\ell,1}$. Then
\begin{align}
 \det\left(I-\sum_jX_j(t)\otimes G_j\right)
 &=\det(I-t_1\cdots t_\ell G_{i_1}\cdots G_{i_\ell}),\label{eq:cycle_determinant}\\
 [t_1\cdots t_\ell]\log\mathcal M_G^{(\ell)}(X(t))
 &=\Tr(G_{i_1}\cdots G_{i_\ell}).\label{eq:cycle_extract}
\end{align}
For arbitrary square matrices $C_b$, independent standard circular Gaussian vectors also give
\begin{equation}\label{eq:gaussian_ring}
 \E_\xi\prod_{b=1}^{\ell}(\xi_b^*C_b\xi_{b+1})
 =\Tr(C_1\cdots C_\ell),\qquad \xi_{\ell+1}=\xi_1.
\end{equation}
\end{proposition}
\begin{proof}
The block source follows a single oriented cycle. Only powers $h\ell$ have nonzero traces, equal to
$\ell(t_1\cdots t_\ell)^h\Tr[(G_{i_1}\cdots G_{i_\ell})^h]$.
Insert this in $-\log\det(I-A)=\sum_{k\ge1}\Tr(A^k)/k$. Both identities follow near zero; the first then holds identically as a polynomial. In \eqref{eq:gaussian_ring}, each Gaussian vector appears once in a conjugated and once in an unconjugated coordinate. Its covariance identifies the adjacent indices, leaving exactly the ordered trace.
\end{proof}
Matrix sources resolve order up to the intrinsic cyclic equivalence of trace. In particular they distinguish the transposed Pauli example in Appendix~\ref{sec:obsappendix}, whose scalar determinant sources agree but whose oriented cubic traces are $16+2i$ and $16-2i$.

\begin{theorem}[Simultaneous finite-dimensional orbit recovery]\label{thm:finite_orbit}
Two self-adjoint tuples $A,B$ of the same known dimension $M$ have equal traces of all words if and only if $B_i=UA_iU^*$ for one unitary $U$. Words of degree at most $2M^2-1$ suffice.

For rectangular families $T_a,S_a$ with the same respective finite input and output dimensions, equality of matrix-source germs for the complete cross-Gram families is equivalent to
\begin{equation}\label{eq:biunitary_orbit}
 S_a=VT_aU^*\qquad\text{for every }a,
\end{equation}
with common input and output unitaries. If the identity response $T_0=S_0=I$ is retained, then $V=U$.
\end{theorem}
\begin{proof}
Equality of word traces gives equality of $\Tr p(A)^*p(A)$, and hence equality of the two polynomial kernels. The map $p(A)\mapsto p(B)$ is a trace-preserving isomorphism of finite-dimensional $*$-algebras. Decomposing these algebras into full matrix blocks, the ambient trace of a minimal projection records the multiplicity of each represented block. The isomorphism preserves these multiplicities and is consequently implemented by an ambient unitary. This is the Specht-type mechanism of \cite{FHS}.

For the finite degree bound, let $V_n$ be the span of words of degree at most $n$ in the block tuple $A_i\oplus B_i$. Its dimension is at most $2M^2$, and $\dim V_0=1$. Equality $V_n=V_{n+1}$ makes $V_n$ invariant under multiplication by every generator, and then it contains every word. If no earlier equality occurs, dimension growth forces stabilization by $2M^2-1$. The trace-difference functional vanishing on that layer thus vanishes on the full generated algebra.

Replace the cross-Gram family by its Hermitian coordinates: the diagonal blocks and the real and imaginary parts of the off-diagonal blocks. Cycle extraction and the first assertion give a common $U$ such that
$S_a^*S_b=UT_a^*T_bU^*$.
The map
$\sum_aT_ah_a\mapsto\sum_aS_aUh_a$
is a well-defined isometry on the response span, because all cross inner products agree. Equal finite output dimensions allow its extension to a unitary $V$, proving \eqref{eq:biunitary_orbit}. The converse is immediate. See also \cite{Jing,FHS}.
\end{proof}

Cross terms are essential: $T_1=T_2=I_2$ and $S_1=I_2$, $S_2=\diag(1,-1)$ have identical individual squares but different cross traces. In infinite dimensions, Hilbert--Schmidt cross-Grams admit Fredholm and cycle formulas by finite-rank approximation; ambient orbit recovery additionally depends on kernel dimensions and representation multiplicities.

\subsection{Random orbits and replica depth}
For fixed $M,d$, the unitary orbit space of self-adjoint contraction tuples is compact. The real and imaginary parts of normalized word traces are bounded continuous separating coordinates; Theorem~\ref{thm:finite_orbit} allows a finite degree bound. Products of these coordinates from the same random tuple determine its orbit law by Stone--Weierstrass.

Lie pulses provide equivalent coordinates. For a homogeneous real Lie polynomial $h$, put $H_h(A)=i^{|h|-1}h(A)$, which is self-adjoint, and consider
\[
 z_t(A)=\tr_M\exp\left(i\sum_ht_hH_h(A)\right)
\]
for finite rational parameter lists. Continuity recovers real parameters; mixed derivatives give symmetrized products of Lie factors, and PBW recovers word traces. Their same-sample products determine the same orbit law.

\begin{theorem}[Word degree and replica depth are independent]\label{thm:replica_obstruction}
For every $r\ge1$ there are two laws of positive projections in dimension $M=r+1$ such that, for polynomials of arbitrary degrees,
\begin{equation}\label{eq:replica_equal}
 \E_+\prod_{a=1}^b\tr_Mp_a(P)
 =\E_-\prod_{a=1}^b\tr_Mp_a(P)\qquad(b\le r),
\end{equation}
but some $(r+1)$-fold observation differs. The same statement holds for trace pulses when every letter is the same projection.
\end{theorem}
\begin{proof}
Set $P_k=\diag(I_k,0)$, $0\le k\le M$, and give its rank the laws
\begin{equation}\label{eq:replica_rank_laws}
 \mu_+(k)=2^{-r}\binom Mk\mathbf1_{\{k\ \text{even}\}},\qquad
 \mu_-(k)=2^{-r}\binom Mk\mathbf1_{\{k\ \text{odd}\}}.
\end{equation}
Both are probability laws. Since $P_k^a=P_k$ for $a\ge1$, every polynomial trace is affine in $k/M$. The finite-difference identities
\[
 \sum_{k=0}^M(-1)^k\binom Mk k^a=0\quad(a<M),\qquad
 \sum_{k=0}^M(-1)^k\binom Mk k^M=(-1)^MM!
\]
prove the equality of all products of at most $r$ traces, and also give
\begin{equation}\label{eq:replica_gap}
 \E_+(\tr_MP)^M-\E_-(\tr_MP)^M
 =\frac{(-1)^MM!}{2^rM^M}\ne0.
\end{equation}
Finally $\tr_Me^{itP_k}=1+(k/M)(e^{it}-1)$ is affine as well.
\end{proof}
For $r=2$, the even-rank law assigns masses $1/4,3/4$ to $0,2$, and the odd-rank law assigns $3/4,1/4$ to $1,3$. The first two normalized rank moments are $1/2,1/3$ under both laws; the third moments are $2/9$ and $5/18$. The witness dimension grows with $r$. Thus the result concerns a uniform replica cutoff across dimensions.


\section{Positive limits and quantitative realization}\label{sec:realizationboundary}
We next fix the ambient Hilbert space. A faithful anchor identifies it inside a positive Gram realization. The residual outside that space determines whether the observation is a genuine tuple. Full-word data close into completely positive maps, and a two-boundary identity controls their multiplication defects through the contexts used by each target.

\subsection{An anchored flat extension}
Let $R$ be a bounded positive injective operator on $H$. Its range is dense. Suppose a positive block matrix $G$ has $G_{00}=R^2$, and choose a minimal Gram realization $G_{ab}=V_a^*V_b$. There is an isometry $U:H\to\mathcal K$ such that $V_0=UR$. Shorting against the unit block gives
\begin{equation}\label{eq:anchorflat}
 \Delta_{ab}=G_{ab}-\operatorname*{s-lim}_{\varepsilon\downarrow0}
 G_{a0}(R^2+\varepsilon I)^{-1}G_{0b}
 =V_a^*(I-UU^*)V_b.
\end{equation}
The identity follows from $R(R^2+\varepsilon I)^{-1}R\to I$ strongly.

\begin{theorem}[A flat border determines all words]\label{thm:flatborder}
The conditions
\[
 G_{0i}=G_{0i}^*,\qquad G_{ii}\le M_i^2R^2,\qquad \Delta_{ii}=0
 \quad(i\ge1)
\]
are necessary and sufficient for the existence of unique self-adjoint operators $T_i$, $\norm{T_i}\le M_i$, such that
$G_{ab}=RT_a^*T_bR$, with $T_0=I$.
The border has a unique positive full-word Hankel extension,
\begin{equation}\label{eq:flatfull}
 K(u,v)=Ru(T)^*v(T)R,\qquad K(iu,v)=K(u,iv).
\end{equation}
Uniqueness holds among all positive Hankel extensions of the given border.
\end{theorem}
\begin{proof}
The unit column identifies the space in which a genuine realization must live. The map $Rh\mapsto V_0h$ is isometric and extends, by density of $\ran R$, to an isometry $U:H\to\mathcal K$. Put $P=UU^*$. The shorting formula \eqref{eq:anchorflat} says $\Delta_{ii}=V_i^*(I-P)V_i$, so flatness is equivalent to $V_i=PV_i$.

Under flatness define $T_i$ first on $\ran R$ by $T_iRh=U^*V_ih$. Injectivity of $R$ gives well-definedness, and the diagonal bound gives
\[
 \norm{T_iRh}^2=\norm{V_ih}^2
 =\ip{G_{ii}h}h\le M_i^2\norm{Rh}^2.
\]
Hence $T_i$ has a unique bounded extension. The identity $G_{0i}=RT_iR$ and self-adjointness of $G_{0i}$ make this extension self-adjoint. Since $V_i=UT_iR$, every cross block is $RT_a^*T_bR$. Conversely these formulas verify the hypotheses. The first row uniquely determines each $T_i$ on the dense anchor range, and thus on $H$.

We must also exclude a different positive Hankel extension that adds hidden higher-word vectors. Take a minimal realization $K(u,v)=W_u^*W_v$ of an extension, writing $W_e=UR$ for its unit-column isometry and $P=UU^*$. On the algebraic span of formal word columns define the shifts by $S_iW_wh=W_{iw}h$. Hankel symmetry makes them symmetric for the positive form. If a formal vector $z$ has zero norm, then Cauchy--Schwarz and symmetry give
\[
 \norm{S_iz}^2=\ip z{S_i^2z}=0.
\]
Thus the shifts descend to a well-defined common dense domain $D$ in the positive quotient. No boundedness of these shifts has yet been assumed.

The determined border gives $W_e=UR$ and $W_i=UT_iR$. Testing symmetry against $URh$ for every $h$ shows that, for $z\in D$,
\[
 PS_iz=UT_iU^*Pz.
\]
If in addition $z\in\ran U$, then for any $y\in D$,
\[
 \ip{S_iz}y=\ip z{S_iy}
 =\ip z{UT_iU^*Py}=\ip{UT_iU^*z}y.
\]
Density of $D$ implies $S_iz=UT_iU^*z\in\ran U$. Starting with the unit-word columns, induction on word length now gives $W_w=Uw(T)R$ for every word. Minimality leaves no further subspace. Therefore the full kernel is necessarily \eqref{eq:flatfull}, which is positive and Hankel by direct evaluation. This proves both existence and uniqueness of the extension.
\end{proof}

Relative to the anchor, $\Delta$ measures a realization defect rather than a missing polynomial direction. It must vanish even when the abstract word filtration has nonzero higher innovations: those innovations may all be represented inside the same anchored Hilbert space. This is an anchored form of flat extension; see \cite{MS} for the general positive-functional setting. An operator-valued border contains all its scalar matrix entries. The scalar positive matrix $\diag(1,1)$ with unit anchor is a simple nonflat example: its first row prescribes $T=0$, while its square entry equals one.

The quantitative extension of flatness is most transparent before taking norms. A word can leave the anchored subspace and return to it; the following identity isolates the two boundary crossings while retaining the entire ordered operator between them.

\begin{lemma}[Two-boundary identity for compressed words]\label{lem:compressed_word_identity}
Let $U:H\to\mathcal K$ be an isometry and put $P=UU^*$. For bounded operators $A_1,\ldots,A_n$ on $\mathcal K$, define
\[
 T_j=U^*A_jU,\qquad
 B_j=(I-P)A_jU,\qquad C_j=(I-P)A_j^*U.
\]
Empty products below are identities. Then
\begin{equation}\label{eq:two_boundary_identity}
 \begin{split}
 U^*A_1\cdots A_nU-T_1\cdots T_n
 =\sum_{1\le i<j\le n}
  T_1\cdots T_{i-1}C_i^*
  A_{i+1}\cdots A_{j-1}B_jT_{j+1}\cdots T_n.
 \end{split}
\end{equation}
For every positive Hilbert--Schmidt anchor $R$ on $H$,
\begin{equation}\label{eq:two_boundary_anchored}
 \begin{split}
 &\norm{R(U^*A_1\cdots A_nU-T_1\cdots T_n)R}_1\\
 &\quad\le\sum_{i<j}
 \norm{C_i(T_1\cdots T_{i-1})^*R}_2
 \left(\prod_{i<k<j}\norm{A_k}\right)
 \norm{B_jT_{j+1}\cdots T_nR}_2.
 \end{split}
\end{equation}
\end{lemma}
\begin{proof}
First telescope the difference before compression:
\[
 A_1\cdots A_nU-UT_1\cdots T_n
 =\sum_{j=1}^n A_1\cdots A_{j-1}B_jT_{j+1}\cdots T_n.
\]
For each prefix, telescope from the other end:
\[
 U^*A_1\cdots A_{j-1}-T_1\cdots T_{j-1}U^*
 =\sum_{i<j}T_1\cdots T_{i-1}C_i^*A_{i+1}\cdots A_{j-1}.
\]
Multiply this identity by $B_j$, use $U^*B_j=0$, and substitute into the first one after left compression by $U^*$. This proves \eqref{eq:two_boundary_identity}. The middle product is not replaced by its compression: it can contain further departures and returns.

After sandwiching a summand by $R$, it is $X^*AY$, with
\[
 X=C_i(T_1\cdots T_{i-1})^*R,
 \quad A=A_{i+1}\cdots A_{j-1},
 \quad Y=B_jT_{j+1}\cdots T_nR.
\]
Both $X,Y$ are Hilbert--Schmidt. Schatten H\"older gives
$\norm{X^*AY}_1\le\norm X_2\norm A\norm Y_2$, proving the second estimate.
\end{proof}

For a self-adjoint tuple $A_i$, put $T_i=U^*A_iU$, $B_i=(I-UU^*)A_iU$, $\eta_i=\norm{B_i}$, and assume $\norm{A_i}\le M_i$. Taking operator norms in the same identity gives, for $w=i_1\cdots i_n$,
\begin{equation}\label{eq:twoleakage}
 \norm{U^*w(A)U-w(T)}
 \le\sum_{a<b}\eta_{i_a}\eta_{i_b}\prod_{k\ne a,b}M_{i_k}.
\end{equation}
If $M_i\le M$ and $\eta_i^2\le\delta$, the bound is $\binom n2M^{n-2}\delta$ for $n\ge2$. Its coefficient is sharp: take
$A_\varepsilon=(1+\varepsilon)^{-1}\left(\begin{smallmatrix}1&\varepsilon\\\varepsilon&1\end{smallmatrix}\right)$ and $U1=e_1$, divide the error of its $n$th power by $\varepsilon^2/(1+\varepsilon)^2$, and let $\varepsilon\downarrow0$.
The anchored estimate \eqref{eq:two_boundary_anchored} has a different hypothesis: it records the defect after the prefix and suffix have acted on the anchor. Section~\ref{subsec:contextual_multiplicativity} identifies when ordinary generator scores suffice to control these contextual defects.

\subsection{\texorpdfstring{Strong-$*$}{Strong-*} reconstruction and output tightness}
Let $H$ be infinite-dimensional and separable, and let $R\in\Sch_2(H)$ be positive, injective, and normalized by $\norm R_2=1$. For tuples satisfying $\max_j\norm{T_j}\le M$, include adjoint labels:
\[
 A_0=I,\qquad A_j=T_j,\qquad A_{d+j}=T_j^*,\qquad b=2d.
\]
Define
\begin{equation}\label{eq:QH}
 B(T)=(A_aR)_{a=0}^b,\qquad Q(T)=B(T)^*B(T),\qquad
 H(T)=\sum_{a=1}^bA_aR^2A_a^*.
\end{equation}
Here $B(T):H^{b+1}\to H$ is a Hilbert--Schmidt row operator. For $Re_k=r_ke_k$ put $D_N=\sum_{k\le N}r_k^{-1}e_ke_k^*$ and let $P_N$ project onto the first $N$ basis vectors. Set
\begin{align}
 d_R(T,S)^2&=\sum_{a=1}^b\norm{(A_a(T)-A_a(S))R}_2^2,\label{eq:strongmetric}\\
 h_N(T)&=\Tr((I-P_N)H(T)),\qquad
 s_N(G)=\sum_{a=1}^b\left(\Tr G_{aa}-\norm{D_NG_{0a}}_2^2\right).
 \label{eq:flatmass}
\end{align}
On the bounded tuple ball, $d_R$ is a complete separable metric for the strong-$*$ topology. Indeed, finite-rank approximation of $R$ proves that strong-$*$ convergence implies $d_R$ convergence. Conversely, every $r_k>0$ recovers convergence on each basis vector; the uniform operator bound extends it to all vectors. The same argument applied to a Cauchy sequence constructs its bounded limit and its adjoint.

\begin{theorem}[Flatness and escaping output mass]\label{thm:strongclosure}
Suppose $Q(T_n)\to G$ in trace norm for a sequence in this bounded ball. Then
\begin{equation}\label{eq:escapemass}
 s_\infty(G):=\lim_Ns_N(G)=\lim_N\sup_nh_N(T_n).
\end{equation}
The limit equals zero if and only if $G=Q(T)$ for a genuine tuple, equivalently if and only if $T_n\to T$ strongly-$*$. The map $T\mapsto Q(T)$ is a topological embedding onto its image, and the enhanced map
\begin{equation}\label{eq:closedstrongcode}
 T\longmapsto(Q(T),H(T))
\end{equation}
has closed image in the product of the corresponding trace-class spaces.
\end{theorem}
\begin{proof}
Let $G_{ab}=V_a^*V_b$ be a minimal Gram realization of the limiting border, with $V_0=UR$. Since $D_NR=P_N$, the shorting residual at physical resolution $N$ is
\[
 s_N(G)=\sum_{a=1}^b\norm{(I-UP_NU^*)V_a}_2^2,
 \qquad
 s_\infty(G)=\sum_{a=1}^b\norm{(I-UU^*)V_a}_2^2.
\]
In particular the finite residuals are nonnegative and decrease to the anchored defect. For each fixed $N$, every term defining $s_N$ is continuous in the trace-class border: $D_N$ is bounded. Thus
$h_N(T_n)=s_N(Q(T_n))\to s_N(G)$.

On the closure of these borders, the functions $s_N$ are continuous, nonnegative, and decreasing. At a genuine border their limit is zero, since $s_N(Q(T))=h_N(T)$ is an output trace tail. Lemma~\ref{lem:monotone_defect}, applied with $r_N=s_N$ and $z_n=Q(T_n)$, therefore gives \eqref{eq:escapemass}. Put $\delta=s_\infty(G)$ for the rest of the proof.

If $\delta=0$, every Gram column lies in $UH$. The border inequalities, which pass to the trace-norm limit, define bounded operators $A_a$ by $A_aRh=U^*V_ah$. Their first-row relations give $A_{d+j}=A_j^*$ on the dense anchor range and therefore on $H$. Hence $G=Q(T)$ for a genuine tuple. For each $a$ and $h$, first-row convergence gives weak convergence of $A_a(T_n)Rh$ against the dense set $\ran R$; the common operator bound extends it to all test vectors. Diagonal-block convergence also gives
\[
 \norm{A_a(T_n)Rh}^2\longrightarrow\norm{A_a(T)Rh}^2.
\]
Weak convergence together with convergence of norms is strong convergence. Density of $\ran R$ and the adjoint labels prove strong-$*$ convergence on $H$.

Conversely, strong-$*$ convergence with a common operator bound gives $A_a(T_n)R\to A_a(T)R$ in Hilbert--Schmidt norm. To see this, truncate $R$ first; the finite head converges strongly and the remainder is bounded by the uniform operator bound times its Hilbert--Schmidt tail. The Gram inequality
\[
 \norm{B^*B-C^*C}_1\le(\norm B_2+\norm C_2)\norm{B-C}_2
\]
proves continuity of $Q$; its analogous form for $BB^*$ proves continuity of $H$. This also proves inverse continuity at every genuine datum.

Finally suppose $(Q(T_n),H(T_n))$ converges in the product of trace-class spaces. Convergence of the positive $H(T_n)$ gives uniformly vanishing output tails. The escape identity forces the limit of $Q(T_n)$ to be genuine, and the argument above gives a strong-$*$ limit $T$. Continuity then identifies both prescribed limits with $(Q(T),H(T))$, proving that the enhanced image is closed.
\end{proof}

For $\eta=\norm{Q(T)-Q(S)}_1$, finite left inversion of the first block row gives the explicit modulus
\begin{equation}\label{eq:strongmodulus}
 d_R(T,S)^2\le r_{*,N}^{-2}\eta^2+2h_N(T)+2h_N(S),
 \qquad r_{*,N}=\min_{k\le N}r_k.
\end{equation}
The Hilbert--Schmidt norm of that entire row is bounded by the trace norm of the full block matrix. The remaining high-output part is bounded by the two tails. A subfamily of the bounded tuple ball is strongly-$*$ relatively compact if and only if its $h_N$ vanish uniformly: the input tails are already bounded by $bM^2\norm{R(I-P_K)}_2^2$, so output tightness leaves uniformly finite matrix approximations. Necessity follows from continuity of $H$ on compact sets.

For $T_n=e_ne_1^*$, linear anchor blocks vanish while a quadratic block retains $r_1^2e_1e_1^*$. Its positive Gram limit is not a true observation. Moreover $T_n$ and $-T_n$ have the same $H$, and their $Q$ distance tends to zero although their $d_R$ distance does not. Uniform inverse moduli thus belong to common-tail classes rather than to the entire bounded ball.

\subsection{The completely positive closure}
Let $\mathcal U_M$ be the universal unital $C^*$-algebra generated by $d$ elements of norm at most $M>0$, with no self-adjointness requirement. Choose a homogeneous $*$-polynomial dictionary $(p_i)_{i\ge0}$ spanning all words, with $p_0=1$. For $i\ge1$, let $n_i\ge1$ be its degree and $a_i$ its sum of absolute word coefficients. Put
\[
 c_0=1,\qquad c_i=\frac{2^{-i}}{n_i!\max(1,a_i)},\qquad
 L=1+e^{2M}/3.
\]
For a unital completely positive map $\phi:\mathcal U_M\to B(H)$, define
\begin{equation}\label{eq:ucpcode}
 \mathcal Q(\phi)_{ij}=c_ic_jR\phi(p_i^*p_j)R.
\end{equation}
The products inside $\phi$ retain the full correlation data.

\begin{theorem}[Positive closure and multiplicativity]\label{thm:ucp}
The map $\mathcal Q$ is a continuous injection from the pointwise weak-operator topology on $\UCP(\mathcal U_M,B(H))$ into the positive trace-class operators, with compact image. On infinite-dimensional separable $H$,
\begin{equation}\label{eq:ucpclosure}
 \overline{\{\mathcal Q(\pi_T):\max_j\norm{T_j}\le M\}}^{\norm\cdot_1}
 =\mathcal Q(\UCP(\mathcal U_M,B(H))).
\end{equation}
Each boundary point is approximated simultaneously in all words by a single sequence of tuples, with error at most $4L\sqrt{\tau_N}$, where $\tau_N=\norm{R(I-P_N)}_2^2$.

For generators and their adjoints $a$, set
\begin{equation}\label{eq:schwarzscore}
 \delta_R(\phi)=\sum_a\Tr\left(R[\phi(a^*a)-\phi(a)^*\phi(a)]R\right).
\end{equation}
Then $\delta_R(\phi)=0$ if and only if $\phi$ is a unital $*$-representation. If true tuple observations converge to $\mathcal Q(\phi)$, they have a genuine strong-$*$ limit realizing all these data if and only if this score vanishes.
\end{theorem}
\begin{proof}
For a unital completely positive map $\phi$, Stinespring's construction~\cite{Stinespring} gives $\phi(a)=V^*\pi(a)V$, where $V:H\to\mathcal K$ is isometric. In the present separable setting the minimal dilation is separable: the vectors $\pi(a)Vh$ from countable dense sets of $a$ and $h$ span a dense subspace. One can construct it directly by completing $\mathcal U_M\odot H$ in the positive form
\[
 \ip{a\otimes h}{b\otimes k}=\ip{\phi(b^*a)h}k.
\]
Complete positivity makes the form positive. Applying it to the positive matrices obtained from $\norm a^2I-a^*a$ bounds left multiplication by $a$ by $\norm a$, so multiplication descends to the quotient and extends to the completion. The column $h\mapsto1\otimes h$ is the isometry $V$.

The encoding is the Gram operator of the Hilbert--Schmidt row $(c_i\pi(p_i)VR)_i$. Since $c_i\norm{p_i}\le2^{-i}e^M$ for $i\ge1$, its squared norm is at most $L$. The first row gives $c_iR\phi(p_i)R$. The anchor is injective with dense range and all $c_i$ are nonzero, so it determines $\phi$ on the polynomial span and hence on all of $\mathcal U_M$.

We next identify the topology of this encoding. The UCP space is compact and metrizable in the pointwise weak-operator topology: diagonal extraction on countable dense lists preserves linearity, the unit, and all finite matrix positivity conditions, and contractivity extends the limit to the whole algebra. The encoding is continuous in trace norm. Indeed, its omitted dictionary diagonal mass is at most $e^{2M}4^{-J}/3$, and its omitted physical diagonal mass is at most $L\tau_N$. The positive-compression estimate \eqref{eq:positivecompression} controls both omissions in trace norm. Truncate the dictionary and physical coordinates to make these bounds small, then use entrywise convergence on the remaining finite matrix. Continuity, injectivity, and compactness give the asserted compact image.

Every point of this image is approximated by genuine tuples on the original infinite-dimensional $H$. Choose a unitary $U_N:\mathcal K\to H$ such that $U_NVh=h$ for $h\in P_NH$. The finite isometry extends because both orthogonal complements are countably infinite-dimensional. Set $T_j^{(N)}=U_N\pi(x_j)U_N^*$. This one tuple has $\norm{T_j^{(N)}}\le M$ and satisfies $p_i(T^{(N)})=U_N\pi(p_i)U_N^*$ for every $i$. Moreover
\[
 \norm{U_N^*R-VR}_2\le2\norm{R(I-P_N)}_2=2\sqrt{\tau_N}.
\]
After applying $U_N^*$ to the true observation row, its distance from the dilation row is at most $2\sqrt{L\tau_N}$. Both row norms are at most $\sqrt L$, so the Gram inequality gives $4L\sqrt{\tau_N}$. This proves density simultaneously in all word coordinates. Compactness of the UCP image gives the reverse inclusion in \eqref{eq:ucpclosure}.

Finally the multiplicativity residual has the dilation factorization
\[
 \phi(p^*q)-\phi(p)^*\phi(q)
 =V^*\pi(p)^*(I-VV^*)\pi(q)V.
\]
Thus the score is the sum of $\norm{(I-VV^*)\pi(a)VR}_2^2$ over generators and adjoints. If it is zero, density of $\ran R$ gives $(I-VV^*)\pi(a)V=0$ for every one of those labels. The subspace $VH$ is reducing, not merely invariant, and compression of $\pi$ to it is a $*$-representation. Conversely a representation has zero residual. This is the multiplicative-domain mechanism~\cite{Choi,CJK}. For a convergent sequence of true encodings, the corresponding finite border converges and is flat exactly when this score vanishes. Theorem~\ref{thm:strongclosure} then gives a genuine strong-$*$ limit. Conversely such a limit determines every word, and the uniform dictionary tails identify the full limiting encoding.
\end{proof}

Averages $\phi(p)=\E p(T)$ are another source of UCP maps. Their Schwarz defects measure random fluctuations. For deterministic tuple sequences, the same dilation formula describes loss of multiplicativity through the limiting observation.

\subsection{Finite-border closure and limiting laws}
Let $\mathfrak K_{M,R}$ consist of positive trace-class borders satisfying
\[
 Q_{00}=R^2,\qquad Q_{aa}\le M^2R^2,\qquad Q_{0,d+j}=Q_{0j}^*.
\]
Their tails under $I_{b+1}\otimes P_N$ are bounded by $(1+bM^2)\tau_N$, so \eqref{eq:positivecompression} makes this set trace-norm compact. Every such border has a unique decomposition
\begin{equation}\label{eq:borderresidual}
 Q=Q(T(Q))+(0\oplus D(Q)),\qquad D(Q)\ge0.
\end{equation}
To obtain it, use the unit-column isometry $U$ in a Gram realization and define $A_aR=U^*V_a$. The first row determines the tuple and its adjoints; the residual is $V_+^*(I-UU^*)V_+$.

Each border in $\mathfrak K_{M,R}$ is a trace-norm limit of true borders. A common factorization gives
$D_{ac}=RF_a^*F_cR$ and $F_a^*F_a\le M^2I-A_a^*A_a$.
For each $T_j$, let $D_f=(M^2I-T_j^*T_j)^{1/2}$ and $D_b=(M^2I-T_jT_j^*)^{1/2}$. Douglas factorization gives contractions $\Gamma_j,\Lambda_j$ with
$F_j=\Gamma_jD_f$ and $F_{d+j}=\Lambda_jD_b$.
The Julia completion \cite{Robinson}
\begin{equation}\label{eq:julia}
 \widetilde T_j=
 \begin{pmatrix}T_j&D_b\Lambda_j^*\\\Gamma_jD_f&-\Gamma_jT_j^*\Lambda_j^*\end{pmatrix}
 =\begin{pmatrix}I&0\\0&\Gamma_j\end{pmatrix}
 \begin{pmatrix}T_j&D_b\\D_f&-T_j^*\end{pmatrix}
 \begin{pmatrix}I&0\\0&\Lambda_j^*\end{pmatrix}
\end{equation}
has norm at most $M$: the middle matrix divided by $M$ is unitary. Its first column and the first column of its adjoint realize the two required labels. The unitary compression construction above then gives true borders on $H$ with error at most $4(1+bM^2)\sqrt{\tau_N}$.

A nonzero residual leaves higher-order data undetermined. On $\Span\{e_1,e_n\}$ take
$T_n^{(t)}=\left(\begin{smallmatrix}0&a\\a&t\end{smallmatrix}\right)$,
with $a=1/4$ and $t=0$ or $1/2$, and zero elsewhere. The linear and quadratic anchored limits agree, but
$\ip{(T_n^{(t)})^3e_1}{e_1}=a^2t$.

\subsection{Finite matrix certificates for the boundary defect}
The quantity $s_N(G)$ has a finite output resolution, but its traces still sum over all input coordinates. To obtain a genuinely finite readout, both sides must be truncated. The known anchor then controls the omitted input mass uniformly over the whole positive-border class, including nongenuine borders.

For $G\in\mathfrak K_{M,R}$, set $d_R^{\rm esc}(G)=s_\infty(G)$ and, for $N,L\ge1$, define
\begin{equation}\label{eq:finite_border_residual}
 s_{N,L}(G)=\sum_{a=1}^b
 \left\{\Tr(P_LG_{aa}P_L)-\norm{D_NG_{0a}P_L}_2^2\right\}.
\end{equation}
This uses only the diagonal blocks through input resolution $L$ and the first-row entries with output resolution $N$ and input resolution $L$.

\begin{proposition}[A finite noisy upper certificate for escape]\label{prop:finite_escape_certificate}
For every $G\in\mathfrak K_{M,R}$,
\begin{equation}\label{eq:finite_border_tail}
 0\le s_N(G)-s_{N,L}(G)\le bM^2\tau_L,\qquad
 0\le d_R^{\rm esc}(G)\le s_{N,L}(G)+bM^2\tau_L,
 \quad \tau_L=\norm{R(I-P_L)}_2^2.
\end{equation}
Let $K\ge\max(N,L)$, $\boldsymbol P_K=I_{b+1}\otimes P_K$, and let a self-adjoint finite matrix $\widehat G$ satisfy
$\norm{\widehat G-\boldsymbol P_KG\boldsymbol P_K}_1\le\eta$.
Put
\[
 \widehat Z=(D_N\widehat G_{0a}P_L)_{a=1}^b,\qquad
 h=r_{*,N}^{-1}\eta,\qquad
 \varepsilon_{N,L}=\eta+2h\norm{\widehat Z}_2+h^2.
\]
Here $\widehat Z$ is a row operator and its Hilbert--Schmidt norm includes all its blocks. Then
\begin{equation}\label{eq:finite_noisy_escape}
 |s_{N,L}(\widehat G)-s_{N,L}(G)|\le\varepsilon_{N,L},\qquad
 0\le d_R^{\rm esc}(G)
 \le s_{N,L}(\widehat G)+\varepsilon_{N,L}+bM^2\tau_L.
\end{equation}
A negative final upper bound certifies incompatibility with the stated error and model bounds. At each genuine border, arbitrarily small upper bounds can be obtained by choosing first $N$, then $L$, and finally the data precision.
\end{proposition}
\begin{proof}
In a minimal Gram realization write $G_{ac}=V_a^*V_c$ and $V_0=UR$. Set $E_N=I-UP_NU^*$, an orthogonal projection. Directly from the finite inverse of the anchor,
\[
 s_{N,L}(G)=\sum_{a=1}^b\norm{E_NV_aP_L}_2^2,\qquad
 s_N(G)-s_{N,L}(G)=\sum_{a=1}^b\norm{E_NV_a(I-P_L)}_2^2.
\]
The first expression is nonnegative. Since $E_N$ is contractive and $V_a^*V_a=G_{aa}\le M^2R^2$, the second is at most $bM^2\tau_L$. Also $d_R^{\rm esc}(G)\le s_N(G)$, proving \eqref{eq:finite_border_tail}.

Let $Z=(D_NG_{0a}P_L)_{a=1}^b$. Compression to the relevant row followed by $D_N$ gives
\[
 \norm{\widehat Z-Z}_2\le\norm{D_N}\eta=h.
\]
The total diagonal trace in \eqref{eq:finite_border_residual} is one compression trace of the finite self-adjoint error and thus changes by at most $\eta$, with no factor $b$. Moreover
\[
 \big|\norm{Z}_2^2-\norm{\widehat Z}_2^2\big|
 \le2\norm{\widehat Z}_2h+h^2.
\]
Combining these estimates proves the first noisy bound; \eqref{eq:finite_border_tail} gives the second. Positivity of $\widehat G$ is not needed for these error estimates.

For a genuine border $s_N(G)\downarrow0$. Choose $N$ with this residual small, then $L$ with $bM^2\tau_L$ small. Since $s_{N,L}(G)\le s_N(G)$, the exact upper certificate is small. At these fixed resolutions its finite readout is continuous, and $\varepsilon_{N,L}\to0$ as $\eta\to0$. This proves the final assertion.
\end{proof}

The one-sided nature of this certificate is essential. Genuine borders are dense in $\mathfrak K_{M,R}$, while nonflat borders exist when $M>0$. Therefore no continuous statistic on the full closure can vanish on every genuine border and be positive on every nongenuine one. The score is instead an upper semicontinuous defect, obtained from decreasing continuous readouts by Lemma~\ref{lem:monotone_defect}. A finite experiment can bound its possible size; equality to zero and convergence to a genuine limit are certified by the corresponding sequence of bounds.

\subsection{Limiting laws and uniform tails}

For probability laws $\mu_n$ on the bounded strong-$*$ tuple ball, suppose $Q_\#\mu_n\Rightarrow\nu$ in $\mathfrak K_{M,R}$. Then
\begin{equation}\label{eq:lawescape}
 \int\Tr D(Q)\,d\nu(Q)
 =\lim_N\lim_n\int h_N(T)\,d\mu_n(T)
 =\lim_N\sup_n\int h_N(T)\,d\mu_n(T).
\end{equation}
The functions $s_N$ are continuous, bounded by $bM^2$, and decrease to $\Tr D$. Lemma~\ref{lem:monotone_defect} applied to the pushed-forward laws proves the formula, including its uniform-tail expression. Its vanishing is equivalent to concentration of $\nu$ on the true image, and hence to a genuine weak limit law in the strong-$*$ ball. To justify the last implication, let $\mu$ be the pullback of $\nu$ by the inverse on the true image. This image is Borel, since it is the zero set of the nonnegative Borel residual score. For an open set $U$ in the tuple space, the topological embedding gives an ambient open $V$ whose intersection with the true image is $Q(U)$. Hence $\liminf_n\mu_n(U)=\liminf_nQ_\#\mu_n(V)\ge\nu(V)=\mu(U)$. Portmanteau proves weak convergence, without requiring the genuine image to be closed.

A family $(\mu_n)$ is uniformly tight in that ball if and only if
$\lim_N\sup_n\int h_N\,d\mu_n=0$.
Necessity follows by restricting to a common high-probability compact set, where output tails are uniform. For sufficiency, choose $N_k$ with expected tail at most $\epsilon2^{-2k}$. Markov's inequality gives mass at least $1-\epsilon$ to
$\bigcap_k\{h_{N_k}\le2^{-k}\}$,
which is closed and strongly-$*$ compact by the output-tail criterion.

\subsection{Recovering a known multiplier}
Let $R>0$ be injective and Hilbert--Schmidt, and put $Q=R^2$ and $q_R(X)=\norm{XR}_2$. For a known $B\in B(H)$ define
\[
 \vartheta_R(B)=\inf\{c\ge0:BQB^*\le c^2Q\}.
\]
\begin{theorem}[Right multiplication in the observation norm]\label{thm:multiplier}
For $c\ge0$, the following conditions are equivalent:
$q_R(XB)\le cq_R(X)$ for every finite-rank $X$;
$BQB^*\le c^2Q$;
and $BR=RC_B$ for a bounded $C_B$ with $\norm{C_B}\le c$.
When they hold, $C_B$ is unique and
\begin{equation}\label{eq:multiplierconstant}
 \sup_{X\ne0\ \text{finite rank}}\frac{q_R(XB)}{q_R(X)}
 =\vartheta_R(B)=\norm{C_B}.
\end{equation}
For observations $Y=XR+E$, $\norm E_2\le\delta$, the estimator $YC_B$ recovers $XBR$ with optimal worst-case error $\delta\vartheta_R(B)$ over all finite-rank $X$, when $\delta>0$. If $\vartheta_R(B)=\infty$, this worst-case error is infinite.
\end{theorem}
\begin{proof}
For $X=uv^*$ the two observation norms are $\norm u\norm{RB^*v}$ and $\norm u\norm{Rv}$. Rank-one testing therefore gives the quadratic-form inequality. It makes $D(Rv)=RB^*v$ bounded by $c$ on the dense range of $R$. Extend $D$ and take adjoints to get $BR=RD^*$. Conversely this factorization gives the quadratic inequality and
$XBR=(XR)C_B$.
Rank-one testing also gives equality of the optimal constants; injectivity of $R$ gives uniqueness of $C_B$. This is Douglas factorization in the specified observation norm \cite{Douglas}.

The ideal inequality gives the stated upper error. For any nonzero finite-rank $X$, rescale so that $q_R(X)=\delta$. The unknowns $X$ and $-X$, with opposite observation errors, give the same zero datum. Their target separation is $2q_R(XB)$, so every estimator has error at least $q_R(XB)$ on one of them. Take the supremum of the ratios.
\end{proof}

For paired adjoint observations write
\[
 d_R^{\pm}(X)=\max(\norm{XR}_2,\norm{X^*R}_2),\qquad
 \mathfrak M_R=\{B:\vartheta_R(B),\vartheta_R(B^*)<\infty\},
\]
and give this space the norm
$\Lambda_R(B)=\max(\norm B,\vartheta_R(B),\vartheta_R(B^*))$.
It is a unital Banach $*$-algebra. Multiplication follows from $ABR=RC_AC_B$; take norm limits of $B_n,C_{B_n},C_{B_n^*}$ and preserve both factorizations. The map $\Theta_R(B)=C_B$ preserves products and brackets, although generally not adjoints. For $R=\diag(r_1,r_2)$ and $B=E_{12}$, the factors for $B$ and $B^*$ have coefficients $r_2/r_1$ and $r_1/r_2$.

The two ideal estimates imply
\begin{equation}\label{eq:twosidedmodule}
 d_R^{\pm}(AXB)\le\Lambda_R(A)\Lambda_R(B)d_R^{\pm}(X).
\end{equation}
If two generator tuples have $\Lambda_R$-bounds $M$ and $d_R^{\pm}$-distance at most $\varepsilon$, their length-$k$ $*$-words differ by at most $kM^{k-1}\varepsilon$. This protocol observes $XR$ itself; alternative anchored observations require their own inversion costs.

\subsection{Completing a family of observable contexts}
Let $B_1,\ldots,B_d$ be known operators that can be applied as observation contexts. Choose an injective positive trace-class $Q_0=R_0^2$ and numbers $t_i>0$ such that
$\alpha=\sum_it_i^2\norm{B_i}^2<1$.
For a word $w=i_1\cdots i_n$, put $t_w=\prod_at_{i_a}$ and $B_w=B_{i_1}\cdots B_{i_n}$, with empty-word values one and $I$.
\begin{proposition}[A trace-class observation Gramian]\label{prop:context_completion}
The positive series
\begin{equation}\label{eq:context_gramian}
 \widehat Q=\sum_wt_w^2B_wQ_0B_w^*,\qquad \widehat R=\widehat Q^{1/2}
\end{equation}
converges in trace norm and is injective. It satisfies
\begin{align}
 \widehat Q&=Q_0+\sum_it_i^2B_i\widehat QB_i^*,&
 \Tr\widehat Q&\le\frac{\Tr Q_0}{1-\alpha},\label{eq:context_stein}\\
 \vartheta_{\widehat R}(B_i)&\le t_i^{-1},&
 \norm{\widehat Q-\widehat Q_{\le N}}_1
 &\le\frac{\Tr Q_0\,\alpha^{N+1}}{1-\alpha}.\label{eq:context_cost}
\end{align}
It is the least bounded positive supersolution of this Stein equation and its unique bounded solution. Moreover
\begin{equation}\label{eq:context_readings}
 \norm{X\widehat R}_2^2=\sum_wt_w^2\norm{XB_wR_0}_2^2.
\end{equation}
\end{proposition}
\begin{proof}
The sum of traces at length $n$ is bounded by $\alpha^n\Tr Q_0$. This proves convergence and both trace estimates. Since $\widehat Q\ge Q_0$, it is injective. Grouping words by their initial letter gives the Stein equation and
$t_i^2B_i\widehat QB_i^*\le\widehat Q$;
Theorem~\ref{thm:multiplier} yields the multiplier costs. The positive map
$\Psi(H)=\sum_it_i^2B_iHB_i^*$
has norm at most $\alpha$ on both $B(H)$ and $\Sch_1$. Iterating any positive supersolution shows that it dominates every partial sum. Two bounded solutions differ by $D=\Psi(D)$, hence agree. Taking the trace of the positive series for $X\widehat QX^*$ proves \eqref{eq:context_readings}.
\end{proof}
This is the noncommutative observability Gramian mechanism \cite{BBF} for the present observation norm. It uses the expanded family of readings $XB_wR_0$. If $\norm X\le M$, discarding contexts of length greater than $N$ has error at most
\[
 M\norm{R_0}_2\frac{\alpha^{(N+1)/2}}{\sqrt{1-\alpha}}
\]
in the $\ell^2(\Sch_2)$ data norm. An adjoint-closed context alphabet places each $B_i$ in $\mathfrak M_{\widehat R}$.

\subsection{Quantitative multiplicativity under observable contexts}\label{subsec:contextual_multiplicativity}
Exact zero Schwarz defect implies multiplicativity, but a small anchored defect need not give a uniform word estimate. The two-boundary identity specifies the missing data: the generator defects must be tested on the input vectors produced by the target's contexts. This is a quantitative use of the operational saturation from Section~\ref{sec:principle}.

Let $\phi:\mathcal U_M\to B(H)$ be unital and completely positive. Write $\mathcal A$ for its generator and adjoint labels, put $T_a=\phi(a)$, and for any bounded context $C$ define
\begin{equation}\label{eq:contextual_schwarz_reading}
 \begin{split}
 S_\phi(a)&=\phi(a^*a)-T_a^*T_a\ge0,\\
 e_R(a;C)^2&=\Tr\big(RC^*S_\phi(a)CR\big),
 \qquad \delta_R(\phi)=\sum_{a\in\mathcal A}e_R(a;I)^2.
 \end{split}
\end{equation}
Thus the last quantity is the score of \eqref{eq:schwarzscore}. An adjoint label is included even when its evaluated operator happens to be self-adjoint. The anchor is positive, injective, and Hilbert--Schmidt; normalization is not needed for the estimates below. A fixed polynomial target uses only the finitely many prefix and suffix contexts appearing in its words.

\begin{theorem}[Contextual control of all ordered products]\label{thm:contextual_multiplicativity}
For a word $w=a_1\cdots a_n$ of length $n\ge2$, put $T_w=T_{a_1}\cdots T_{a_n}$. Then
\begin{equation}\label{eq:contextual_word_bound}
 \begin{split}
 \norm{R[\phi(w)-T_w]R}_1
 \le\sum_{i<j}
 e_R\big(a_i^*;(T_{a_1}\cdots T_{a_{i-1}})^*\big)
 \left(\prod_{i<k<j}\norm{a_k}_{\mathcal U_M}\right)
 e_R\big(a_j;T_{a_{j+1}}\cdots T_{a_n}\big).
 \end{split}
\end{equation}
If each compressed generator belongs to $\mathfrak M_R$, then the two contextual factors are at most
\[
 e_R(a_i^*;I)\prod_{k<i}\vartheta_R(T_{a_k}^*),
 \qquad
 e_R(a_j;I)\prod_{k>j}\vartheta_R(T_{a_k}),
\]
respectively. In particular, if $K\ge\max(1,M)$ and $\Lambda_R(T_a)\le K$ for every label, then
\begin{equation}\label{eq:contextual_uniform_word}
 \norm{R[\phi(w)-T_w]R}_1
 \le\binom n2K^{n-2}\delta_R(\phi).
\end{equation}
For the homogeneous dictionary and weights of \eqref{eq:ucpcode}, this also gives the full trace-class estimate
\begin{equation}\label{eq:contextual_full_word_code}
 \norm{\mathcal Q(\phi)-\mathcal Q(\pi_T)}_1
 \le\big(e^K+2e^{2K}\big)\delta_R(\phi),
\end{equation}
where $\pi_T$ is the representation of $\mathcal U_M$ defined by the norm-bounded tuple $T$. 
\end{theorem}
\begin{proof}
Take a Stinespring realization $\phi(b)=U^*\pi(b)U$ and put $P=UU^*$. The positive defect factorization gives
\[
 S_\phi(a)=B_a^*B_a,
 \qquad B_a=(I-P)\pi(a)U,
 \qquad e_R(a;C)=\norm{B_aCR}_2.
\]
Apply Lemma~\ref{lem:compressed_word_identity} to $A_j=\pi(a_j)$. Its left boundary factor is $B_{a_i^*}$ and its right boundary factor is $B_{a_j}$. Contractivity of $\pi$ bounds the interior product by the displayed product of algebra norms. This proves \eqref{eq:contextual_word_bound}; it depends only on $\phi$ and the specified contexts, not on a choice of dilation.

If $C\in\mathfrak M_R$, Theorem~\ref{thm:multiplier} gives $CR=RC_C$ with $\norm{C_C}=\vartheta_R(C)$. Therefore
\[
 e_R(a;C)=\norm{B_aRC_C}_2
 \le e_R(a;I)\vartheta_R(C).
\]
Multiplicativity of the context factors gives the two product bounds. Each squared generator factor is at most $\delta_R(\phi)$, so the uniform hypothesis bounds every summand by $K^{n-2}\delta_R(\phi)$. There are $\binom n2$ summands, proving \eqref{eq:contextual_uniform_word}.

For the full dictionary let $b_0=1$, $n_0=0$, and $b_i=c_ia_i\le2^{-i}/n_i!$ for $i\ge1$, using the absolute word coefficient sums $a_i$. A polynomial $p_i^*p_j$ has degree $n_i+n_j$ and absolute coefficient sum at most $a_ia_j$. Apply \eqref{eq:contextual_uniform_word} to each of its words; terms of degree zero or one have zero defect. The sum of the trace norms of all resulting blocks is bounded by
\[
 \delta_R(\phi)\sum_{i,j\ge0}b_ib_j
        \binom{n_i+n_j}{2}K^{n_i+n_j-2}.
\]
Set $f(x)=\sum_{i\ge0}b_ix^{n_i}$ for $x\ge0$. The coefficient bounds give
$f(K)\le1+e^K$, $f'(K)\le e^K$, and $f''(K)\le e^K$.
The double sum is $f'(K)^2+f(K)f''(K)$, at most $e^K+2e^{2K}$. Absolute summation of trace-class blocks bounds the trace norm of the whole block operator, proving \eqref{eq:contextual_full_word_code}.
\end{proof}

With the normalized anchor of Proposition~\ref{prop:finite_escape_certificate}, the border $G$ of $\phi$ has escape score $\delta_R(\phi)$. If Proposition~\ref{prop:finite_escape_certificate} supplies the finite noisy upper bound
\[
 U_{N,L}=s_{N,L}(\widehat G)+\varepsilon_{N,L}+bM^2\tau_L,
\]
then, under the same certified context bound $K$,
\begin{equation}\label{eq:finite_multiplicativity_certificate}
 \norm{\mathcal Q(\phi)-\mathcal Q(\pi_T)}_1
 \le(e^K+2e^{2K})U_{N,L}.
\end{equation}
This certifies the error of replacing the positive map by the word evaluation of its compressed tuple. Finite recovery of that tuple itself retains the separate anchor-inverse and tail costs.

The zero-defect criterion is the classical multiplicative-domain mechanism~\cite{Choi,CJK}. The estimates above specify its quantitative form in this anchored observation norm. They compare a positive map with the word evaluation of its compressed tuple in the free norm-bounded algebra. The comparison is in the anchored trace norm on the free norm-bounded algebra.

\begin{example}[A small generator score can hide a fixed word defect]\label{ex:anchored_defect_instability}
Let $H=\ell^2(\N)$ and fix $Re_j=r_je_j$, with $r_j>0$, $r_j\to0$, and $\sum_jr_j^2=1$. On $\mathcal K=H\oplus\C f$, for $n\ge2$ take the self-adjoint contractions
\[
 A_n=fe_n^*+e_nf^*,\qquad
 B_n=e_1e_n^*+e_ne_1^*,
\]
extended by zero on the other directions. The embedding $U:H\to\mathcal K$ and the representation $\pi_n(a)=A_n$, $\pi_n(b)=B_n$ define a UCP map $\phi_n=U^*\pi_n(\,\cdot\,)U$ on $\mathcal U_1$. Its compressed tuple is $T_a=0$, $T_b=B_n|_H$. The only nonzero generator defects are
$S_{\phi_n}(a)=S_{\phi_n}(a^*)=e_ne_n^*$; the $b$ and $b^*$ defects vanish. Thus
\begin{equation}\label{eq:anchored_small_score_counterexample}
 \delta_R(\phi_n)=2r_n^2\longrightarrow0,
 \qquad
 \norm{R[\phi_n(ba^2b)-T_bT_a^2T_b]R}_1=r_1^2.
\end{equation}
Indeed $U^*B_nA_n^2B_nU=e_1e_1^*$. The missing reading is already the single context $T_b$:
$e_R(a;T_b)^2=r_1^2$. Correspondingly,
\[
 \vartheta_R(T_b)=\max\{r_1/r_n,r_n/r_1\}\longrightarrow\infty.
\]
There is therefore no uniform modulus tending to zero that bounds even this fixed fourth-order target by the generator score alone on the full UCP class. The contextual hypothesis in Theorem~\ref{thm:contextual_multiplicativity} rules out this uncontrolled transfer into weakly anchored input directions.
\end{example}

\begin{corollary}[Saturating a defect by its observable contexts]\label{cor:saturated_defect}
For the same UCP map, choose a faithful trace-class $Q_0=R_0^2$ and weights $t_a>0$ on the generator--adjoint alphabet with
$\alpha=\sum_at_a^2\norm{T_a}^2<1$.
Apply Proposition~\ref{prop:context_completion} to the compressed tuple:
\[
 \widehat Q=\sum_wt_w^2T_wQ_0T_w^*,\qquad
 \widehat R=\widehat Q^{1/2}.
\]
Then $T_a\in\mathfrak M_{\widehat R}$ for every label, and
\begin{equation}\label{eq:saturated_defect_identity}
 \delta_{\widehat R}(\phi)
 =\sum_wt_w^2\sum_{a\in\mathcal A}e_{R_0}(a;T_w)^2.
\end{equation}
Writing $C_{\mathcal A}=\sum_a\norm{a}_{\mathcal U_M}^2$, the omitted contexts of length greater than $N$ satisfy
\begin{equation}\label{eq:saturated_defect_tail}
 0\le\delta_{\widehat R}(\phi)
 -\sum_{|w|\le N}t_w^2\sum_a e_{R_0}(a;T_w)^2
 \le C_{\mathcal A}\Tr Q_0\,
          \frac{\alpha^{N+1}}{1-\alpha}.
\end{equation}
Thus the expanded readings provide a finite-context certificate for the score entering Theorem~\ref{thm:contextual_multiplicativity}.
\end{corollary}
\begin{proof}
The context Gramian gives $\vartheta_{\widehat R}(T_a)\le t_a^{-1}$. The alphabet contains the adjoint of each label, so it also bounds $\vartheta_{\widehat R}(T_a^*)$, proving membership in $\mathfrak M_{\widehat R}$. In a Stinespring realization,
\[
 \delta_{\widehat R}(\phi)=\sum_a\Tr(B_a\widehat QB_a^*).
\]
Insert the positive trace-norm convergent Gramian series and take the trace term by term. The result is \eqref{eq:saturated_defect_identity}. Since $\norm{B_a}\le\norm{a}_{\mathcal U_M}$, the total contribution at word length $k$ is at most $C_{\mathcal A}\Tr Q_0\alpha^k$. Summing the geometric tail proves \eqref{eq:saturated_defect_tail}.
\end{proof}

The compressed operators must be available as contexts, or obtained by a separate certified reconstruction. This construction enlarges the observation protocol: it uses the contextual readings in \eqref{eq:saturated_defect_identity}, not just the original score $\delta_{R_0}$. The finite word cutoff does not by itself truncate the physical traces; those require a separate physical-resolution and noise budget, as in Proposition~\ref{prop:finite_escape_certificate}. If a trace-one anchor is desired, normalize $\widehat Q$ by its trace; all squared scores and sandwiched trace norms scale by that same factor, while the multiplier costs are unchanged.


\section{Observable operations and common-noise evolution}\label{sec:transport}\label{sec:time}
Recovered multilinear operations carry a common error calculus: primal errors propagate through partial evaluations and dual tests through their adjoints. Common-noise evolution exposes new contexts through drift commutators and transports ordered products through replica semigroups. All drivers in the evolution results are deterministic bounded operators.

\subsection{Primal errors and dual tests}
Consider a finite rooted ordered tree. An internal vertex $v$ applies a bounded multilinear operation $\mu_v$ to its ordered children. Let
\[
 x_v=\mu_v(x_{v1},\ldots,x_{vk_v}),\qquad
 a_v=\mu_v(a_{v1},\ldots,a_{vk_v})+r_v,
 \qquad e_v=a_v-x_v.
\]
The residual $r_v$ includes local computational error. Define
\[
 L_{vj}h=\mu_v(a_{v1},\ldots,a_{v,j-1},h,x_{v,j+1},\ldots,x_{vk_v}).
\]
\begin{proposition}[Exact error and adjoint identities]\label{prop:cascade}
At every internal vertex,
\begin{equation}\label{eq:cascade}
 e_v=r_v+\sum_jL_{vj}e_{vj}.
\end{equation}
Starting with a root functional $\lambda_{\rm root}$ and recursively setting $\lambda_{vj}=L_{vj}'\lambda_v$ gives
\begin{equation}\label{eq:adjointcascade}
 \lambda_{\rm root}(e_{\rm root})
 =\sum_{v\ \text{internal}}\lambda_v(r_v)
  +\sum_{v\ \text{leaf}}\lambda_v(e_v).
\end{equation}
\end{proposition}
\begin{proof}
Replace the true inputs one at a time by their computed values. Multilinearity gives the telescoping identity \eqref{eq:cascade} exactly. Substitute this identity down the tree and use the adjoint pairing along each edge. Every branch ends in a local residual or leaf error, giving \eqref{eq:adjointcascade}.
\end{proof}

If $C_v$ bounds the node operation, $B_{vj}$ bounds its true inputs in the required input spaces, $E_{vj}$ bounds their errors, and $\norm{r_v}\le\delta_v$, then
\begin{equation}\label{eq:cascadebudget}
 B_v=C_v\prod_jB_{vj},\qquad
 E_v=\delta_v+C_v\sum_jE_{vj}
       \prod_{i<j}(B_{vi}+E_{vi})\prod_{i>j}B_{vi}.
\end{equation}
Choose input radii recursively from the desired output radius; at the Hilbert exponent use \eqref{eq:arityregion}. Shared subexpressions can be treated on a directed acyclic graph by adding all incoming dual functionals. The exact adjoint identity contains true states; the budget recurrence is computable from certified bounds alone.

\subsection{Finite-chaos generators and countable dictionaries}
Let $T_j,\widetilde T_j$ belong to the same degree-$m_j$ chaos over the same source. Set
\[
 \Delta=\max_j\norm{T_j-\widetilde T_j}_{L^2(\Sch_s)},\quad
 B\ge\max\left(1,\max_j\norm{T_j}_{L^2(\Sch_s)},
                    \max_j\norm{\widetilde T_j}_{L^2(\Sch_s)}\right),
 \quad m_* =\max_jm_j.
\]
\begin{proposition}[Word and dictionary bounds]\label{prop:weightedwords}
Every length-$k$ $*$-word satisfies
\begin{equation}\label{eq:momentword}
 \norm{w(T)-w(\widetilde T)}_{L^2(\Sch_s)}
 \le kH_kB^{k-1}\Delta,\qquad H_k=(2k-1)^{m_*k/2}.
\end{equation}
For a countable dictionary of homogeneous $*$-polynomials $p_j$ of positive word degrees $k_j$, with absolute coefficient sums $a_j$, let
\[
 c_j=\frac{2^{-j}}{\max(1,a_j)k_j!H_{k_j}}.
\]
Then
\begin{equation}\label{eq:weightedwordbound}
 \sum_{j\ge1}c_j\norm{p_j(T)-p_j(\widetilde T)}_{L^2(\Sch_s)}\le e^B\Delta,
 \qquad
 \sum_{j>J}c_j\norm{p_j(T)}_{L^2(\Sch_s)}\le2^{-J}e^B.
\end{equation}
\end{proposition}
\begin{proof}
The word difference is a sum of $k$ telescoping terms. In each, bound the differing factor in $\Sch_s$ and all others in operator norm, then use $\norm A\le\norm A_s$. Probability H\"older requires $L^{2k}$ for each factor. Hypercontractivity charges at most $(2k-1)^{m_*/2}$ per factor, including the homogeneous difference. This proves \eqref{eq:momentword}; a single word has norm at most $H_kB^k$. After multiplication by $c_j$, the corresponding bounds are
$2^{-j}B^{k_j-1}\Delta/(k_j-1)!$ and $2^{-j}B^{k_j}/k_j!$.
Their sums give \eqref{eq:weightedwordbound}.
\end{proof}

Include $p_0=I,c_0=1$ and a normalized faithful anchor $R$. The row
$\mathsf B(T)=(c_jp_j(T)R)_{j\ge0}$
has $L^2(\Sch_2)$ norm at most $1+e^B$. Consequently
\begin{equation}\label{eq:allgramerror}
 \norm{\mathsf B(T)^*\mathsf B(T)-\mathsf B(\widetilde T)^*\mathsf B(\widetilde T)}_{L^1(\Sch_1)}
 \le2(1+e^B)e^B\Delta.
\end{equation}
This includes all cross-word and cross-PBW entries. The moment factor $H_k$ reflects a genuine integrability cost: for $T=g_1g_2g_3P$ and every $c>0$, $\E e^{c|g_1g_2g_3|}=\infty$. Integrating over $[n,n+1]^3$ compares the growth $cn^3$ with Gaussian decay of order $3(n+1)^2/2$ and proves divergence. Thus polynomial closure and entire functional calculus have different integrability requirements.

\subsection{A bounded random-unitary semigroup}
Let $H_0,H_1,\ldots,H_q$ be bounded self-adjoint operators, and define
\begin{equation}\label{eq:timegenerator}
 \delta_j(A)=i[H_j,A],\qquad
 \mathcal L=\delta_0+\tfrac12\sum_{j=1}^q\delta_j^2,\qquad
 \Phi_t=e^{t\mathcal L}.
\end{equation}
The auxiliary transport noise is independent of the Gaussian source used to reconstruct coefficients. For each diffusion direction,
\[
 e^{t\delta_j^2/2}A
 =\E\big[e^{i\sqrt t\,gH_j}Ae^{-i\sqrt t\,gH_j}\big].
\]
The bounded-generator product formula, together with $e^{t\delta_0}$, shows that $\Phi_t$ is a norm-continuous unital completely positive semigroup on $B(H)$ and a compatible contraction semigroup on every $\Sch_s$, $1\le s<\infty$. The product formula here also follows directly from $F(h)=I+h\mathcal L+O(h^2)$ and a telescoping estimate. This random-unitary setting is part of the standard theory of quantum dynamical semigroups \cite{Lindblad,Aniello}.

\begin{proposition}[Covariance and multiplication defects]\label{prop:timecov}
Define
\[
 \Gamma(A,B)=\mathcal L(A^*B)-\mathcal L(A)^*B-A^*\mathcal L(B),\qquad
 C_t(A,B)=\Phi_t(A^*B)-\Phi_t(A)^*\Phi_t(B).
\]
Then
\begin{align}
 \Gamma(A,B)&=\sum_{j=1}^q\delta_j(A)^*\delta_j(B),\label{eq:timegamma}\\
 C_t(A,B)&=\int_0^t\Phi_{t-u}\Gamma(\Phi_uA,\Phi_uB)\,du.\label{eq:timecov}
\end{align}
Every finite block matrix formed from $C_t$ is positive, and
\begin{align}\label{eq:timebracket}
 \Phi_t([A,B])-[\Phi_tA,\Phi_tB]
 &=C_t(A^*,B)-C_t(B^*,A)\notag\\
 &=\int_0^t\Phi_{t-u}\sum_j[\delta_j(\Phi_uA),\delta_j(\Phi_uB)]\,du.
\end{align}
\end{proposition}
\begin{proof}
Each $\delta_j$ is a $*$-derivation. Applying the Leibniz rule twice gives \eqref{eq:timegamma}. The norm derivative of
$\Phi_{t-u}((\Phi_uA)^*\Phi_uB)$
is $-\Phi_{t-u}\Gamma(\Phi_uA,\Phi_uB)$; integration proves \eqref{eq:timecov}. The gradient blocks are Gram matrices, and complete positivity and integration preserve their positivity. Subtracting the two ordered product defects yields \eqref{eq:timebracket}.
\end{proof}

If $\sigma_H^2=\sum_j\norm{H_j}^2$, the bracket defect is at most
$8\sigma_H^2t\norm A\norm B$.
Positivity of the covariance blocks also bounds it by
\[
 \norm{C_t(A^*,A^*)}^{1/2}\norm{C_t(B,B)}^{1/2}
 +\norm{C_t(B^*,B^*)}^{1/2}\norm{C_t(A,A)}^{1/2}.
\]
For self-adjoint generators $T$, the positive word kernel
$K_t(u,v)=R\Phi_t(u(T)^*v(T))R$
has unit-anchor shorting defect
\begin{equation}\label{eq:timeflat}
 \Delta_{ab}^t=RC_t(T_a,T_b)R.
\end{equation}
The regularized shorting limit holds in trace norm after Hilbert--Schmidt sandwiching, by finite-rank approximation. Vanishing of its generator diagonal defects makes $\Phi_t$ multiplicative on $C^*(I,T)$ by the faithful-anchor multiplicative-domain criterion.

\subsection{Common-noise replicas}
On the Banach projective tensor product
$\mathcal X_k=B(H)^{\widehat\otimes_\pi k}$,
ordered multiplication $m_k(A_1\otimes\cdots\otimes A_k)=A_1\cdots A_k$ is contractive. Define
\begin{equation}\label{eq:replicagenerator}
 D_j^{(k)}=\sum_{r=1}^k\delta_j^{[r]},\qquad
 \mathcal L^{(k)}=D_0^{(k)}+\tfrac12\sum_j(D_j^{(k)})^2
 =\sum_r\mathcal L^{[r]}+\sum_{r<s}\sum_j\delta_j^{[r]}\delta_j^{[s]}.
\end{equation}
\begin{theorem}[Replica intertwining and coefficient transport]\label{thm:replicas}
The operators $\Phi_t^{(k)}=e^{t\mathcal L^{(k)}}$ form a contraction semigroup on $\mathcal X_k$, and
\begin{equation}\label{eq:replicaintertwine}
 m_k\Phi_t^{(k)}=\Phi_tm_k.
\end{equation}
Adjacent-slot multiplication, multilinear Lie words, slot permutations, and unit insertions intertwine the corresponding replica semigroups. Applying $\Phi_t$ to each physical coefficient defines $\widetilde\Phi_t$ on $\Alg_s$, with
\begin{equation}\label{eq:coeftransport}
 I\widetilde\Phi_t=\Phi_tI,\qquad
 \Nn\rho{\widetilde\Phi_tK}\le\Nn\rho K
 \quad(t\ge0,\ \rho\ge1).
\end{equation}
\end{theorem}
\begin{proof}
For every real $a$, the simultaneous conjugation
$e^{aD_j^{(k)}}=(e^{a\delta_j})^{\otimes k}$
is an isometry in the projective tensor norm. Averaging one common Gaussian parameter gives the diffusion step $e^{t(D_j^{(k)})^2/2}$. The bounded-generator product formula, together with the drift step, therefore makes $\Phi_t^{(k)}$ a contraction semigroup.

The Leibniz identity $m_kD_j^{(k)}=\delta_jm_k$ holds on elementary tensors and extends by boundedness. It then holds for the squares, the generators, and their exponentials, proving \eqref{eq:replicaintertwine}. The same identity within a merged pair of slots proves the assertion for adjacent multiplication. Permutations commute with simultaneous slot actions, unit insertions use $\delta_j(I)=0$, and multilinear Lie words are linear combinations of the resulting ordered products.

For coefficients, define $\widetilde\delta_j$ by applying $i[H_j,\cdot]$ to each physical coefficient. Left and right physical multiplication have their usual operator-norm bounds in every directed cut, so $\widetilde\delta_j$ is bounded by $2\norm{H_j}$ in every coefficient norm, including the quotient-sum branch. Consequently
$\widetilde{\mathcal L}=\widetilde\delta_0+\tfrac12\sum_j\widetilde\delta_j^2$
is bounded on each radius space. Its conjugation steps are isometries: physical unitary factors act on the outer legs of each cut. Gaussian averages and the bounded product formula are therefore contractions in those same norms. Their limit is $e^{t\widetilde{\mathcal L}}$, which proves the inequality in \eqref{eq:coeftransport}. On the finite core, applying the physical maps commutes with Gaussian evaluation. Continuity and density extend this identity to $\Alg_s$. The exponential construction, or contraction followed by finite-core approximation, gives strong continuity at every radius.
\end{proof}
The cross terms in \eqref{eq:replicagenerator} distinguish common-noise transport from independent-slot transport. The evolved target is $\Phi_t(A_1\cdots A_k)$; it is generally different from $\Phi_t(A_1)\cdots\Phi_t(A_k)$.

For the homogeneous generators of Proposition~\ref{prop:weightedwords}, contraction gives
\[
 \sup_{t\ge0}\norm{\Phi_t(w(T))-\Phi_t(w(\widetilde T))}_{L^2(\Sch_s)}
 \le kH_kB^{k-1}\Delta.
\]
The weighted response row $(c_j\Phi_t(p_j(T))R)_j$ inherits the bounds $1+e^B$, $e^B\Delta$, and $2(1+e^B)e^B\Delta$ for its size, difference, and Gram difference. These are uniform-in-time Bochner estimates. The actual correlation kernel $R\Phi_t(p_i(T)^*p_j(T))R$ differs from the Gram of the averaged responses by the covariance in \eqref{eq:timecov}.

\subsection{The drift-generated visibility filtration}
Set
\begin{equation}\label{eq:time_ancestors}
 J_{j,k}=(\ad H_0)^kH_j,\qquad
 \mathcal N_r=\{J_{j,k}:1\le j\le q,\ 0\le k<r\}',\qquad
 \mathcal N_0=B(H),\qquad \mathcal N_\infty=\bigcap_r\mathcal N_r.
\end{equation}
Since $J_{j,k}^*=(-1)^kJ_{j,k}$, each $\mathcal N_r$ is a von Neumann algebra. This filtration depends on the prescribed drift and noise. For an injective positive Hilbert--Schmidt anchor define
\[
 O_H(A)=(\delta_1(A)R,\ldots,\delta_q(A)R)\in\Sch_2(H)^q.
\]
\begin{lemma}[Ancestral commutators and dual tests]\label{lem:ancestry_dual_chain}
For $r\ge1$,
\begin{equation}\label{eq:ancestry_observable_chain}
 \mathcal N_r=\bigcap_{k=0}^{r-1}\ker(O_H\delta_0^k),\qquad
 \mathcal N_r^\perp=
 \overline{\Span\{(O_H\delta_0^k)'\lambda:0\le k<r\}}^{w^*}.
\end{equation}
The annihilator is taken in $B(H)'$, and $\lambda$ ranges over the continuous dual of the observation space.
\end{lemma}
\begin{proof}
Conjugation separates the drift from the noise test:
\begin{equation}\label{eq:ancestry_conjugation}
 \delta_j(e^{u\delta_0}A)
 =i e^{u\delta_0}[e^{-u\delta_0}H_j,A].
\end{equation}
The coefficient of $u^k$ in the inner commutator is
$(-i)^k[J_{j,k},A]/k!$. Multiplication by the outer exponential is a triangular transformation of Taylor coefficients with identity diagonal. Hence the first $r$ coefficients vanish on one side exactly when they vanish on the other. The derivatives on the left are $\delta_j\delta_0^kA$, $0\le k<r$. A bounded operator vanishes on $\ran R$ exactly when it vanishes on $H$, because the anchor range is dense. This proves the common-kernel description of $\mathcal N_r$.

For a family of bounded linear maps, the annihilator of the intersection of their kernels is the weak-$*$ closed span of their adjoint ranges. One way to see the equality is to take preannihilators of both sides: they are the same common kernel, and the double-annihilator theorem identifies the weak-$*$ closure. Applying this observation to the displayed maps gives \eqref{eq:ancestry_observable_chain}.
\end{proof}
Thus the ancestral chain is the operational saturation of the initial noise observation under the drift, in the sense of Section~\ref{sec:principle}.

\begin{theorem}[The first nonzero covariance order]\label{thm:temporal_onset}
Let $A\in B(H)$, and suppose $r\ge0$ is the smallest integer for which some $[J_{j,r},A]\ne0$. In operator norm as $t\downarrow0$,
\begin{equation}\label{eq:temporal_onset}
 C_t(A,A)=\frac{t^{2r+1}}{(r!)^2(2r+1)}
     \sum_j[J_{j,r},A]^*[J_{j,r},A]+O(t^{2r+2}).
\end{equation}
For $\norm R_2=1$, the positive scalar observation $v_R(t,A)=\Tr(RC_t(A,A)R)$ satisfies
\begin{equation}\label{eq:temporal_readout}
 v_R(t,A)=\frac{t^{2r+1}}{(r!)^2(2r+1)}\kappa_r(A,R)+O(t^{2r+2}),
 \quad \kappa_r(A,R)=\sum_j\norm{[J_{j,r},A]R}_2^2>0.
\end{equation}
For every $r\ge0$, including targets with no finite first visible order,
\begin{equation}\label{eq:temporal_depth_iff}
 A\in\mathcal N_r\quad\Longleftrightarrow\quad v_R(t,A)=O(t^{2r+1}).
\end{equation}
\end{theorem}
\begin{proof}
The onset order comes from a gradient, its positive square, and one time integration. By the minimality of $r$ and \eqref{eq:ancestry_conjugation},
\[
 \delta_j\delta_0^kA=0\quad(k<r),\qquad
 \delta_j\delta_0^rA=i(-i)^r[J_{j,r},A].
\]
Induction gives $\mathcal L^kA=\delta_0^kA$ for $k\le r$: at each earlier step every diffusion term $\delta_j^2\delta_0^kA$ vanishes. The analytic expansion of the bounded semigroup therefore gives
\[
 \delta_j\Phi_uA=u^rL_j+O(u^{r+1}),\qquad
 L_j=\frac{i(-i)^r}{r!}[J_{j,r},A].
\]
Squaring the gradients gives $\Gamma(\Phi_uA,\Phi_uA)=u^{2r}\sum_jL_j^*L_j+O(u^{2r+1})$. Inserting this into \eqref{eq:timecov}, and using $\Phi_{t-u}=I+O(t-u)$, yields
$C_t(A,A)=t^{2r+1}\sum_jL_j^*L_j/(2r+1)+O(t^{2r+2})$.
This is \eqref{eq:temporal_onset}; positivity prevents cancellation of a nonzero leading commutator.

Here is a uniform bound for the remainder. Fix $t_0>0$, take $c\ge\norm{\mathcal L}$, and put $d_j=2\norm{H_j}$,
\[
 b_j=\frac{d_j\norm A e^{ct_0}c^{r+1}}{(r+1)!},
 \qquad a_j=\norm{L_j},\qquad
 D_r=\sum_j(2a_jb_j+t_0b_j^2).
\]
The exponential series gives $\norm{\delta_j\Phi_uA-u^rL_j}\le b_ju^{r+1}$ for $0\le u\le t_0$. Hence
\[
 \Big\|\Gamma(\Phi_uA,\Phi_uA)-u^{2r}\sum_jL_j^*L_j\Big\|
 \le D_ru^{2r+1}.
\]
Contraction of $\Phi$ and $\norm{\Phi_v-I}\le cv e^{ct_0}$ give, with $G_r=\sum_jL_j^*L_j$,
\[
 \left\|C_t(A,A)-\frac{t^{2r+1}}{2r+1}G_r\right\|
 \le\left(\frac{D_r}{2r+2}
 +\frac{ce^{ct_0}\sum_ja_j^2}{(2r+1)(2r+2)}\right)t^{2r+2}.
\]
The denominators come from the integrals of $u^{2r+1}$ and $(t-u)u^{2r}$, respectively.

Sandwiching by $R$ and taking trace has norm at most $\norm R_2^2=1$. If a leading commutator is nonzero, it cannot vanish on the dense range of $R$, so $\kappa_r(A,R)>0$. This proves \eqref{eq:temporal_readout}. Finally, if $A\in\mathcal N_r$, Taylor expansion through degree $r-1$ gives $\delta_j\Phi_uA=O(u^r)$ and hence the asserted upper order of $v_R$. Otherwise some first nonzero index $k<r$ gives a strictly positive scalar leading term of order $2k+1$, incompatible with $O(t^{2r+1})$. This proves \eqref{eq:temporal_depth_iff}, including the case in which all finite ancestral commutators vanish.
\end{proof}

For Pauli matrices $H_0=X$, $H_1=Z$, and $A=Z$, one has $[Z,Z]=0$ and $[[X,Z],Z]=4X$. Thus
$C_t(Z,Z)=\frac{16}{3}t^3I+O(t^4)$.
If the depth $r$ is known and $y_t$ approximates $v_R(t,A)$ within $\eta$, then
\begin{equation}\label{eq:onset_estimator}
 \widehat\kappa_r=(r!)^2(2r+1)t^{-(2r+1)}y_t,\qquad
 |\widehat\kappa_r-\kappa_r|
 \le C_{A,H,r}t+(r!)^2(2r+1)\eta t^{-(2r+1)}.
\end{equation}
The remainder and a lower bound on the leading coefficient determine the conditioning of depth detection. Odd short-time orders also occur in hypocoercivity \cite{AAC}; the observable here is the covariance of a specified operator target.

\begin{proposition}[One positive time and the noiseless algebra]\label{prop:noiseless_algebra}
For any $t_0>0$,
\begin{equation}\label{eq:noiseless_iff}
 RC_{t_0}(A,A)R=0\quad\Longleftrightarrow\quad A\in\mathcal N_\infty.
\end{equation}
On $\mathcal N_\infty$, $\Phi_tA=e^{itH_0}Ae^{-itH_0}$, and this action extends to a group of $*$-automorphisms.
\end{proposition}
\begin{proof}
If the anchored covariance is zero, its unanchored positive operator is zero because $\ran R$ is dense. Formula~\eqref{eq:timecov} is an integral of positive, norm-continuous operators. Testing against each vector shows that every value of the integrand is zero. The map $\Phi_{t_0-u}=e^{(t_0-u)\mathcal L}$ is injective as a bounded linear map; its inverse is not assumed positive. Therefore $\Gamma(\Phi_uA,\Phi_uA)=0$, and each $\delta_j\Phi_uA=0$ on $[0,t_0]$.

On this interval the evolution equation reduces to $\partial_u\Phi_uA=\delta_0\Phi_uA$, so $\Phi_uA=e^{u\delta_0}A$. The norm-analytic functions $\delta_je^{u\delta_0}A$ vanish on an interval and hence on all real $u$. The conjugation identity then gives every ancestral commutator zero, which means $A\in\mathcal N_\infty$.

Conversely, vanishing of all ancestral commutators makes the entire drift orbit commute with every noise operator. It solves the full bounded evolution equation with initial value $A$. Uniqueness yields $\Phi_tA=e^{t\delta_0}A$ and zero covariance. Moreover
\[
 \mathcal N_\infty
 =\{e^{-iuH_0}H_je^{iuH_0}:1\le j\le q,\ u\in\R\}',
\]
which is invariant under drift translation. Conjugation by $e^{itH_0}$ thus restricts to the asserted group of $*$-automorphisms.
\end{proof}
The ancestral-commutant description of the decoherence-free algebra is classical \cite[Proposition~3]{DFSU}. The finite-time anchored defect and the finite-level observation chain give the corresponding observation criteria.

\subsection{PBW transport}
In the free Lie algebra containing driving letters $h_j$ and target letters $a_i$, let $\partial_j=i\ad h_j$ and $\mathscr L=\partial_0+\frac12\sum_j\partial_j^2$. Write $\partial_j^\odot$ for the derivation acting on each symmetric Lie factor. PBW inversion $\beta=\sym^{-1}$ satisfies
\begin{equation}\label{eq:time_pbw}
 \beta\mathscr L=
 \left(\partial_0^\odot+\tfrac12\sum_j(\partial_j^\odot)^2\right)\beta
 =\mathscr L_\odot\beta.
\end{equation}
This follows by applying the Leibniz rule inside each symmetrized ordered product. The two orders of differentiation on distinct factors cancel the factor $1/2$, giving the replica cross terms.

The formal generator preserves the number of Lie primitive factors and raises total word degree by at least one and at most two on nonzero terms. Therefore each fixed output degree receives contributions from only finitely many powers, and
$\beta e^{t\mathscr L}=e^{t\mathscr L_\odot}\beta$
in the total-degree completion. Inner derivations preserve representation ideals. Descent of individual PBW projections additionally requires their kernel invariance, while convergence in an operator norm is supplied by the bounded semigroup above.

\subsection{Recovery with unknown drivers}
Observe the drivers and target jointly, together with their adjoints and the identity. For $\mathbf T=(H_0,\ldots,H_q,A)$, put
\[
 c=2\norm{H_0}+2\sum_j\norm{H_j}^2\ge\norm{\mathcal L},\qquad
 p_k(\mathbf T)=\mathcal L_H^kA.
\]
These are prescribed noncommutative polynomials of degree at most $1+2k$.
\begin{theorem}[Joint finite-time recovery]\label{thm:finite_time_recovery}
If $\norm A\le M$, then for $0\le t\le T$,
\begin{equation}\label{eq:time_taylor_budget}
 \left\|\Phi_t^HA-\sum_{k=0}^N\frac{t^k}{k!}p_k(\mathbf T)\right\|
 \le Me^{Tc}\frac{(Tc)^{N+1}}{(N+1)!}.
\end{equation}
For two true driver--target tuples with common budgets $M,c$, a spectral corner $P_b$ of the anchor satisfies
\begin{align}\label{eq:time_joint_budget}
 &\sup_{0\le t\le T}\norm{P_b(\Phi_t^HA-\Phi_t^{\widetilde H}\widetilde A)P_b}\notag\\
 &\quad\le r_{*,b}^{-2}\sum_{k=0}^N\frac{T^k}{k!}
       \norm{Rp_k(\mathbf T)R-Rp_k(\widetilde{\mathbf T})R}_1
       +2Me^{Tc}\frac{(Tc)^{N+1}}{(N+1)!}.
\end{align}
\end{theorem}
\begin{proof}
The exponential tail is bounded by $M\sum_{k>N}(Tc)^k/k!$; using $(N+1+j)!\ge(N+1)!j!$ proves \eqref{eq:time_taylor_budget}. On the finite corner, two-sided anchor inversion costs $r_{*,b}^{-2}$. Apply it termwise to the finite Taylor sum and add the two remainder bounds.
\end{proof}
The required polynomial observations have degree at most $1+2N$; their Gram diagonal observations have at most twice this degree. Known coordinate and dictionary weights must be inverted as part of the finite observation map.

\begin{corollary}[Uniform compact-time continuity]\label{cor:time_true_continuity}
Let true tuples $\mathbf T_n,\mathbf T$ be uniformly bounded, with their common generator--adjoint anchor observations converging in trace norm. Then, for each finite $T$ and $v\in H$,
\begin{equation}\label{eq:time_strong_star_path}
 \sup_{t\le T}\left(\norm{(\Phi_t^{H_n}A_n-\Phi_t^HA)v}
 +\norm{((\Phi_t^{H_n}A_n)^*-(\Phi_t^HA)^*)v}\right)\longrightarrow0.
\end{equation}
The output borders and every fixed averaged ordered-word block also converge uniformly on $[0,T]$ in anchored trace norm.
\end{corollary}
\begin{proof}
Theorem~\ref{thm:strongclosure} gives strong-$*$ convergence of the input tuples. Finite products are strongly-$*$ continuous on bounded sets, so every $p_k(\mathbf T_n)$ converges. A finite Taylor sum converges uniformly on compact time intervals, and the common $c$ makes its remainder uniform in $n$. Truncate first and then pass to the limit to prove \eqref{eq:time_strong_star_path}. Finite-rank approximation of $R$ converts uniform strong-$*$ convergence with a common bound into uniform Hilbert--Schmidt convergence of the response rows; their Grams converge in trace norm. Apply the same argument to any fixed initial word.
\end{proof}
For uniformly bounded random driver--target tuples, the resulting continuous mapping transfers weak convergence of their joint strong-$*$ laws to the laws of the averaged response paths. These paths average the auxiliary transport noise while retaining the original reconstructed randomness.


\section{Spectral observations and their completions}\label{sec:spectralgrowth}\label{sec:pfaffian}
Spectral observations retain a quotient of the operator structure. We first determine the exact growth of the spectral moments in the circular Gaussian model. Exterior--Pfaffian coordinates then recover nonzero skew spectra from finite data. Their inverse estimates and tail conditions identify the subcritical, critical, and determinant completions.

\subsection{A deterministic norm and the moment radius}
Let $m\ge1$, let the $m$ Gaussian families be independent and standard circular complex, and fix an integer $s_0\ge1$. For $K\in\Coeff_{m,2s_0}$, put
\[
 A_K=T_K^*T_K,\qquad Z_q(K)=\E\Tr A_K^q,\qquad q\ge s_0.
\]
All these moments are finite by the fixed-degree forward bounds and Schatten inclusions. On finite kernels define
\begin{equation}\label{eq:multilinear_tau}
 \tau(K)=\sup_{\norm{u_1}=\cdots=\norm{u_m}=1}
 \left\|\sum_{i_1,\ldots,i_m}(u_1)_{i_1}\cdots(u_m)_{i_m}
                 K_{i_1\ldots i_m}\right\|_{\op}.
\end{equation}
The all-input cut gives $\tau(K)\le\norm K_{\Coeff_{m,2s_0}}$. Consequently $\tau$ extends continuously to the original cut completion. It separates points, since its vanishing forces every source and physical matrix coefficient to vanish.

\begin{theorem}[Factorial radius and exponential threshold]\label{thm:factorial_radius}
In this model,
\begin{equation}\label{eq:factorial_root}
 \lim_{q\to\infty}\left(\frac{Z_q(K)}{(q!)^m}\right)^{1/(2q)}=\tau(K).
\end{equation}
The same limit holds with $Z_q(K)$ replaced by $\E\norm{T_K}_{\op}^{2q}$. Hence
\begin{equation}\label{eq:spectral_borel_radius}
 \mathcal B_{K,s_0}(z)=\sum_{q=s_0}^\infty\frac{Z_q(K)}{(q!)^m}z^q
 \quad\text{has radius}\quad R_K=\tau(K)^{-2},
\end{equation}
with $R_K=\infty$ when $K=0$. For $K\ne0$,
\begin{equation}\label{eq:stretched_threshold}
 \sup\{a\ge0:\E e^{a\norm{T_K}_{\op}^{2/m}}<\infty\}
 =\frac{m}{\tau(K)^{2/m}}.
\end{equation}
The exponential is integrable strictly below the threshold and nonintegrable strictly above it.
\end{theorem}
\begin{proof}
For finite source dimensions $d_j$ and physical rank at most $D$, write each circular Gaussian vector as $r_j\theta_j$, with independent radius and uniform direction. Since $r_j^2$ has Gamma distribution of shape $d_j$,
\[
 \E r_j^{2q}=q!\binom{q+d_j-1}{d_j-1}.
\]
Multilinearity separates the radial growth from the directional operator and gives
\begin{equation}\label{eq:finite_radial_identity}
 \frac{Z_q(K)}{(q!)^m}
 =\prod_{j=1}^m\binom{q+d_j-1}{d_j-1}
   \int\Tr[(T_K(\theta)^*T_K(\theta))^q]\,d\theta.
\end{equation}
The trace integrand is between $\norm{T_K(\theta)}^{2q}$ and $D\norm{T_K(\theta)}^{2q}$. The continuous function $\theta\mapsto\norm{T_K(\theta)}$ attains its maximum $\tau(K)$, and every neighborhood of a maximizer has positive spherical measure. Its $L^{2q}$ norm therefore tends to $\tau(K)$. The binomial factors and $D$ have $2q$th roots tending to one. This proves both normalized moment limits for finite kernels.

To pass to the original completion, we need an estimate uniform in the moment order. For a finite kernel $D_0$, the forward bound at probability and Schatten exponents $2q$, together with the inclusion $\Sch_{2s_0}\subset\Sch_{2q}$ and $q!\ge(q/e)^q$, gives
\begin{equation}\label{eq:uniform_radial_completion}
 \frac{\norm{T_{D_0}}_{L^{2q}(\Sch_{2q})}}{(q!)^{m/(2q)}}
 \le C_0^m(4e)^{m/2}\norm{D_0}_{\Coeff_{m,2s_0}},
 \qquad q\ge s_0.
\end{equation}
It also bounds the operator-norm moments. Thus either normalized moment root, denoted by $F_q$, satisfies
$|F_q(K)-F_q(L)|\le C_m\norm{K-L}_{\Coeff_{m,2s_0}}$
with $C_m$ independent of $q$. The same coefficient norm controls $|\tau(K)-\tau(L)|$. For finite kernels $K_n\to K$,
\[
 |F_q(K)-\tau(K)|
 \le C_m\norm{K-K_n}
 +|F_q(K_n)-\tau(K_n)|+|\tau(K_n)-\tau(K)|.
\]
Choose the finite approximation first, then let $q$ tend to infinity. This proves \eqref{eq:factorial_root} on the infinite-dimensional coefficient space, and Cauchy--Hadamard gives \eqref{eq:spectral_borel_radius}.

For the exponential threshold put $Y=\norm{T_K}_{\op}^2$ and suppose $K\ne0$. Stirling's formula turns the moment limit into
\[
 (\E Y^q)^{1/q}\sim\tau(K)^2(q/e)^m.
\]
Monotonicity of probability norms between neighboring integer exponents extends this asymptotic to real exponents. Taking the exponent $n/m$ and applying Stirling again gives
\[
 \lim_{n\to\infty}\left(\frac{\E Y^{n/m}}{n!}\right)^{1/n}
 =\frac{\tau(K)^{2/m}}m.
\]
Tonelli identifies $\E e^{aY^{1/m}}$ with its nonnegative moment series. The root test gives finiteness below and divergence above the threshold in \eqref{eq:stretched_threshold}. No assertion at the endpoint is needed.
\end{proof}
Radial asymptotics and large deviations for Gaussian chaoses have a substantial classical literature \cite{HKP,Ledoux}. The calculation above keeps the circular normalization and passes the exact constant through the directed Schatten-cut completion using the uniform estimate \eqref{eq:uniform_radial_completion}.

\begin{corollary}[Failure of a positive square exponential]\label{cor:positive_square_failure}
If $m\ge2$ and $K\ne0$, then for every $t>0$,
\[
 \E e^{t\norm{T_K}_{\op}^2}=\infty,\qquad
 \E\Tr(e^{tA_K}-I)=\infty,
\]
where the second trace is understood as an extended nonnegative quantity.
\end{corollary}
\begin{proof}
Put $Y=\norm{T_K}_{\op}^2$. By \eqref{eq:factorial_root}, the $q$th root of $\E Y^q/q!$ tends to infinity, so the positive exponential series diverges. The largest eigenvalue of $A_K$ is $Y$, and the trace difference is at least $e^{tY}-1$.
\end{proof}
For one scalar coefficient, $T_K$ is a product of $m$ standard circular Gaussians and $Z_q=(q!)^m$. Thus the factorial power already appears in the simplest nonzero model.

\subsection{Analytic transforms and averaged spectra}
The entire function
\[
 \Phi_{m,s_0}(w)=\sum_{q=s_0}^\infty\frac{w^q}{(q!)^m}
\]
satisfies, by spectral calculus and Tonelli,
\begin{equation}\label{eq:spectral_actual_function}
 \mathcal B_{K,s_0}(z)=\E\Tr\Phi_{m,s_0}(zA_K),\qquad z\ge0,
\end{equation}
with extended values permitted outside the convergence interval. The same identity holds by absolute convergence for $|z|<R_K$. Its order-$s_0$ zero at the origin makes it suitable for the stated infinite-dimensional trace assumption. Changing $s_0$ when both coefficient hypotheses hold changes the transform only by a polynomial and leaves the radius unchanged.

A transform finite on the whole positive half-axis is
\begin{equation}\label{eq:weighted_laplace}
 \mathcal L_{K,s_0}(t)=\E\Tr(A_K^{s_0}e^{-tA_K}),\qquad t\ge0.
\end{equation}
It is the Laplace transform of the finite positive measure
\begin{equation}\label{eq:weighted_mean_spectrum}
 \nu_{s_0}(D)=\E\sum_j\lambda_j(A_K)^{s_0}\mathbf1_D(\lambda_j(A_K)),
 \qquad \nu_{s_0}([0,\infty))=Z_{s_0}.
\end{equation}
The transform is analytic on $\Re t>0$ and its values for $t>0$ uniquely determine this weighted mean spectral measure. Domination by the already finite moments gives
\[
 (-1)^n\mathcal L_{K,s_0}^{(n)}(0+)=Z_{s_0+n}(K),\qquad n\ge0.
\]
These data specify an averaged spectral measure, rather than a law of full random spectral configurations.

Products of $m\ge3$ independent mean-one exponential variables have moments $(q!)^m$ and can be moment-indeterminate \cite{LS}. Consequently the factorial growth condition alone does not solve the moment problem among arbitrary measures. For $m\ge2$ and $K\ne0$, the ordinary Taylor series of \eqref{eq:weighted_laplace} at zero has radius zero, despite the existence of all right derivatives.

For a single positive trace-class operator $Q$, let $p_j=\Tr Q^j$, $e_j=\Tr(\wedge^jQ)$, and $h_j=\Tr(\Sym^jQ)$. Newton's identities read
\begin{equation}\label{eq:newton_spectral}
 je_j=\sum_{a=1}^j(-1)^{a-1}e_{j-a}p_a,\qquad
 jh_j=\sum_{a=1}^jh_{j-a}p_a,\qquad e_0=h_0=1.
\end{equation}
They follow by logarithmic differentiation of the Fredholm determinant and its reciprocal; see \cite{Macdonald,Bornemann}. These are invertible coordinates of the same unlabelled spectrum. The spectral radius $R_K$ in \eqref{eq:spectral_borel_radius} concerns a fixed chaos degree and growing spectral moments, whereas the radius $\rho$ in $\Nn\rho K$ weights the chaos degrees themselves.

\subsection{Exterior coordinates and the target quotient}
Let $E$ be a separable real Hilbert space. Exterior powers are normalized so that increasing wedges of orthonormal vectors are orthonormal. For a skew-adjoint $A\in\Sch_2(E)$, put
\begin{equation}\label{eq:pfaffian_form}
 \alpha_A=\sum_{i<j}\ip{Ae_j}{e_i}e_i\wedge e_j,
 \qquad \norm{\alpha_A}^2=\tfrac12\norm A_2^2.
\end{equation}
Write $a_1\ge a_2\ge\cdots\ge0$ for its positive skew block singular values, appending zeros after a finite spectrum. Each positive $a_j$ gives the two complex eigenvalues $\pm ia_j$ and occurs twice in the ordinary singular-value sequence. The target relation is
\[
 A\sim_{\rm sp}B\quad\Longleftrightarrow\quad a_j(A)=a_j(B)\quad(j\ge1).
\]
Thus the target records the positive skew singular values with multiplicity and forgets eigenplanes, zero-space dimension, and Gaussian-source coupling.

Define $F_0(\alpha)=1$ and $F_k(\alpha)=2^k\alpha^{\wedge k}/k!$. A tower $(\Xi_{2k})$ is \emph{Pfaffian-compatible} with $A$ when $\Xi_{2k}=F_k(\alpha_A)$ for every $k$. Pfaffian compatibility is imposed as an additional structural condition on the jointly realized tower.

\begin{lemma}[Exterior coordinates and their native completion]\label{lem:pfaffian_exterior}
For every $k\ge1$,
\begin{equation}\label{eq:pfaffian_elementary}
 \norm{\alpha_A^{\wedge k}}^2=(k!)^2e_k(a^2),\qquad
 c_k(A):=4^{-k}\norm{\Xi_{2k}}^2=e_k(a^2).
\end{equation}
For a compatible tower,
\begin{equation}\label{eq:pfaffian_determinant}
 \mathscr D_a(z)=\sum_{k\ge0}c_k(A)z^k
 =\prod_{j\ge1}(1+za_j^2),\qquad
 \det_F(I+zA^*A)=\mathscr D_a(z)^2.
\end{equation}
The product is entire, and its zeros $-a_j^{-2}$, with multiplicity, determine the nonzero target spectrum.

The map $\alpha\mapsto(F_k(\alpha))_{k\ge0}$ is continuous and has closed image in the tower space with seminorms
$\sum_{k\ge0}R^k\norm{u_k}$, $R>0$. Its inverse on the image is $u\mapsto u_1/2$.
\end{lemma}
\begin{proof}
In oriented invariant planes, $\alpha_A=\sum_ja_je_{2j-1}\wedge e_{2j}$. Two-forms commute under exterior multiplication. Finite expansion gives $k!$ times the product of the selected $k$ coefficients on each orthonormal wedge basis vector. Squaring and summing gives the first identity. Fixed-degree exterior multiplication is continuous; finite skew-block truncations and monotone convergence pass the identity to $\Sch_2$. In particular
\[
 e_k(a^2)\le\frac{(\sum_ja_j^2)^k}{k!},\qquad
 \norm{F_k(\alpha_A)}\le\frac{2^k\norm{\alpha_A}^k}{\sqrt{k!}}.
\]
These estimates prove local uniform convergence of the determinant series and the tower's finiteness at every radius. The Fredholm expansion and the doubled skew singular values give \eqref{eq:pfaffian_determinant}; see~\cite{Simon,Bornemann}. A nonzero factor of the convergent product vanishes exactly at its specified negative zero, so zeros and their orders give the asserted spectrum.

For convergence of $\alpha_l$ in $\wedge^2 E$, each fixed exterior power converges by continuity, and the displayed factorial majorant controls the tower tails uniformly on bounded sets of $\alpha_l$. Hence the convergence holds at every radius. Conversely convergence of generated towers implies convergence of their first component and therefore of $\alpha_l$. Continuity identifies every limiting component and every radius limit with the tower of that two-form. This proves closedness.
\end{proof}

One natural finite source of compatibility is chronological. For ordered vectors $d_1,\ldots,d_h$, let
$\Xi_k^{\rm disc}=\sum_{i_1<\cdots<i_k}d_{i_1}\wedge\cdots\wedge d_{i_k}$ and $\alpha^{\rm disc}=\Xi_2^{\rm disc}/2$. Then
\begin{equation}\label{eq:pfaffian_chronological}
 (\alpha^{\rm disc})^{\wedge k}=k!2^{-k}\Xi_{2k}^{\rm disc}.
\end{equation}
Terms with repeated labels vanish. On each ordered $2k$-label set the remaining coefficient is $k!2^{-k}$ times the Pfaffian of the skew sign matrix, whose upper triangular entries are one. Expansion along its first row gives an alternating sum of $2k-1$ unit terms, hence one, by induction. This proves \eqref{eq:pfaffian_chronological}. Fixed-degree limits of such towers retain the relation by continuity. Pfaffian identities in shuffle algebras provide a classical algebraic context~\cite{LuqueThibon}; a stochastic application must additionally establish convergence of the actual iterates on a common event.

For $0<\rho<2$, put $W_\rho(a)=\sup_jj^{1/\rho}a_j$. The following one-object gauges provide the spectral envelope used as the certificate class for the inverse theorem.
\begin{proposition}[Factorial and Fredholm size gauges]\label{prop:pfaffian_gauges}
For a compatible tower, define
\[
 \mathfrak F_\rho(A)=\sup_{k\ge1}
       \big[(k!)^{1/\rho}\norm{\Xi_{2k}}\big]^{1/k},\qquad
 M_\rho(A)=\sup_{x>0}x^{-\rho/2}\log\mathscr D_a(x).
\]
Then $\mathfrak F_\rho(A)\asymp_\rho W_\rho(a)$ and
$M_\rho(A)\asymp_\rho W_\rho(a)^\rho$, with infinite values allowed. In particular,
\begin{equation}\label{eq:pfaffian_class_certificate}
 a_k\le \frac{e^{1/\rho}\mathfrak F_\rho(A)}2\,k^{-1/\rho}.
\end{equation}
\end{proposition}
\begin{proof}
The factorial weight is inside the $k$th root. Since
$a_k^k\le\prod_{j\le k}a_j\le e_k(a^2)^{1/2}=2^{-k}\norm{\Xi_{2k}}$,
the inequality $k!\ge(k/e)^k$ proves \eqref{eq:pfaffian_class_certificate}.
For the reverse bound write $W=W_\rho(a)<\infty$ and $\beta=2/\rho>1$. Integral comparison gives
\[
 \log\mathscr D_a(x)\le\sum_j\log(1+xW^2j^{-\beta})
 \le B_\rho W^\rho x^{\rho/2},\qquad
 B_\rho=\int_0^\infty\log(1+t^{-2/\rho})\,dt<\infty.
\]
For $W>0$, choose $x^{\rho/2}=2k/(\rho B_\rho W^\rho)$. The coefficient bound becomes
\[
 e_k(a^2)\le x^{-k}\exp(B_\rho W^\rho x^{\rho/2})
 \le\frac{(C_\rho W^2)^k}{(k!)^{2/\rho}},
\]
after enlarging the constant. Equation~\eqref{eq:pfaffian_elementary} then gives $\mathfrak F_\rho\le C_\rho W$. Finally, at $x=a_k^{-2}$ the first $k$ determinant factors are at least two, so
$k\log2\le M_\rho a_k^{-\rho}$. This proves the remaining comparison. The zero operator is immediate.
\end{proof}

Each finite exterior coefficient contains the unresolved tail. For every $M$,
\[
 e_k(a^2)=\sum_{l=0}^k e_l(a_1^2,\ldots,a_M^2)
                       e_{k-l}(a_{M+1}^2,a_{M+2}^2,\ldots).
\]
Accordingly, the inverse below passes through the full power moments and a positive square-mass measure.

\subsection{Finite moment coordinates with their noise cost}
Fix $0<\rho<2$, $C>0$, and the sorted spectral class
\begin{equation}\label{eq:pfaffian_class}
 \mathcal K_\rho(C)=\{a_1\ge a_2\ge\cdots\ge0:\ a_j\le Cj^{-1/\rho}\}.
\end{equation}
For $a\in\mathcal K_\rho(C)$ set
\begin{equation}\label{eq:pfaffian_normalized}
 x_j=(a_j/C)^2,\quad E_k(a)=C^{-2k}c_k(a)=e_k(x),\quad
 p_l(x)=\sum_jx_j^l,\quad S_\rho=\frac2{2-\rho}.
\end{equation}
Then $0\le x_j\le1$ and $\sum_jx_j\le S_\rho$, by comparison with the integral of $t^{-2/\rho}$. We first allow any sorted $x\subset[0,1]$ with total mass at most $S$.
Define its square-mass measure and its Bernstein masses by
\begin{align}
 \mu_x&=\sum_jx_j\delta_{x_j},\label{eq:pfaffian_mass_measure}\\
 b_{k,n}(t)&=\binom nk t^k(1-t)^{n-k},&
 \beta_{k,n}(x)&=\int_0^1b_{k,n}(t)\,d\mu_x(t),\quad 0\le k\le n.
 \label{eq:pfaffian_bernstein_masses}
\end{align}
They are nonnegative and sum to $p_1(x)$. Write
$q_n(x,y)=\sum_{k=0}^n|\beta_{k,n}(x)-\beta_{k,n}(y)|$;
for spectra $a,b$, $q_n(a,b)$ always uses the normalization \eqref{eq:pfaffian_normalized}.

\begin{lemma}[Finite visibility and error amplification]\label{lem:pfaffian_newton}
The vector $\beta_n$ and $(e_1(x),\ldots,e_{n+1}(x))$ determine one another by finite polynomial transformations. Explicitly,
\begin{align}
 p_l&=\sum_{i=1}^{l-1}(-1)^{i-1}e_i p_{l-i}+(-1)^{l-1}le_l,
 \label{eq:pfaffian_newton}\\
 \beta_{k,n}&=\binom nk\sum_{j=0}^{n-k}(-1)^j\binom{n-k}j p_{k+j+1}.
 \label{eq:pfaffian_bernstein_moments}
\end{align}
If $|e_i(x)-e_i(y)|\le\eta_i$ for $1\le i\le n+1$ and both square-mass sums are at most $S$, define
\begin{align}
 D_1&=\eta_1,&
 D_l&=l\eta_l+\sum_{i=1}^{l-1}
       \left(\frac{S^i}{i!}D_{l-i}+S\eta_i\right)\quad(l\ge2),
 \label{eq:pfaffian_noise_recursion}\\
 Q_n(\eta;S)&=\sum_{r=0}^n\binom nr2^rD_{r+1}.
 \label{eq:pfaffian_noise_budget}
\end{align}
Then $|p_l(x)-p_l(y)|\le D_l$ and $q_n(x,y)\le Q_n(\eta;S)$.
\end{lemma}
\begin{proof}
Apply Newton's identities to the first $M$ coordinates. As $M\to\infty$, the elementary symmetric sums increase to $e_i(x)$ and the power sums increase to $p_l(x)$; the bounds $e_i(x)\le S^i/i!$ and $p_l(x)\le S$ justify passage to the limit in each finite identity. Expanding the Bernstein polynomials gives \eqref{eq:pfaffian_bernstein_moments}. Conversely they form a basis of polynomials of degree at most $n$, so their integrals determine the moments $p_1,\ldots,p_{n+1}$, and inverse Newton recursion determines $e_1,\ldots,e_{n+1}$.

Subtract \eqref{eq:pfaffian_newton} for two spectra, retaining $e_i(x)$ and $p_{l-i}(y)$ in the product difference. Their bounds give \eqref{eq:pfaffian_noise_recursion} inductively. Taking absolute values in \eqref{eq:pfaffian_bernstein_moments} and summing in $k$, the coefficient of $D_{r+1}$ is
\[
 \sum_{k=0}^r\binom nk\binom{n-k}{r-k}=\binom nr2^r.
\]
This proves the stated finite-order amplification bound.
\end{proof}

\subsection{A gap-free inverse from integer spectral counts}
\begin{theorem}[Finite data control every ordered spectral position]\label{thm:pfaffian_inverse}
For sorted summable $x,y\subset[0,1]$ of mass at most $S$ and every $n\ge1$,
\begin{equation}\label{eq:pfaffian_square_inverse}
 \norm{x-y}_\infty^2\le q_n(x,y)+\frac{4S}{\sqrt n}.
\end{equation}
Consequently, for $a,b\in\mathcal K_\rho(C)$, with
\begin{equation}\label{eq:pfaffian_Hn}
 H_n(a,b)=\min\left\{1,
     \left(q_n(a,b)+\frac{4S_\rho}{\sqrt n}\right)^{1/4}\right\},
\end{equation}
one has $\norm{a-b}_\infty\le CH_n(a,b)$. The bound is uniform in the rank and independent of spectral gaps and multiplicities.
\end{theorem}
\begin{proof}
We recover a spectral position by testing the fact that the counting function changes by at least one. Fix $j$, interchange $x,y$ if necessary, and suppose $d=x_j-y_j>0$. Set $t_0=y_j+d/2$ and take the ramp
\[
 f(u)=\begin{cases}
 0,&u\le t_0,\\
 2(u-t_0)/d,&t_0<u<x_j,\\
 1,&u\ge x_j.
 \end{cases}
\]
The first $j$ values of $x$ contribute one each, while at most the first $j-1$ values of $y$ contribute, each by at most one. Since $t_0>0$, summability makes both sums finite. Hence
\begin{equation}\label{eq:pfaffian_integer_gap}
 \sum_i f(x_i)-\sum_i f(y_i)\ge1.
\end{equation}
The observations concern the mass measure $\mu_x=\sum_ix_i\delta_{x_i}$, rather than the unweighted counting measure. Define $g(0)=0$ and $g(u)=f(u)/u$ for $u>0$. On the rising interval,
$g(u)=2d^{-1}(1-t_0/u)$ and $g'(u)=2t_0/(du^2)$; above $x_j$ it equals $u^{-1}$. The pieces are continuous, and $t_0\ge d/2$, $x_j\ge d$, so
\[
 \norm g_\infty\le d^{-1},\qquad \operatorname{Lip}(g)\le4d^{-2}.
\]

For any $L$-Lipschitz $g$, its Bernstein polynomial satisfies
\begin{equation}\label{eq:pfaffian_bernstein_error}
 \norm{B_ng-g}_\infty\le\frac{L}{2\sqrt n},\qquad
 \left|\int g\,d(\mu_x-\mu_y)\right|
 \le\norm g_\infty q_n(x,y)+\frac{LS}{\sqrt n}.
\end{equation}
Indeed $B_ng(t)=\E g(Z/n)$ for $Z$ binomial with parameters $(n,t)$, and
$\E|Z/n-t|\le\sqrt{t(1-t)/n}$. The integral of $B_ng$ is a linear combination of the observed Bernstein masses with coefficients $g(k/n)$. The two approximation errors each integrate against a measure of mass at most $S$, giving the second inequality.

Combining the ramp test with this finite observation estimate gives
\[
 1\le\int g\,d(\mu_x-\mu_y)
 \le\frac{q_n(x,y)}d+\frac{4S}{d^2\sqrt n}.
\]
Thus $d^2\le dq_n+4S/\sqrt n\le q_n+4S/\sqrt n$, since $d\le1$. Taking the supremum in $j$ proves \eqref{eq:pfaffian_square_inverse}. Finally
$|\sqrt u-\sqrt v|\le\sqrt{|u-v|}$
passes from squared values to skew values, and the bound $\norm{a-b}_\infty\le C$ gives \eqref{eq:pfaffian_Hn} with its stated cap. The argument uses no eigenvalue gap or rank bound.
\end{proof}

Retaining the factor $d$ in the last quadratic inequality gives the slightly sharper finite-data estimate
\begin{equation}\label{eq:pfaffian_quadratic_refinement}
 \norm{x-y}_\infty
 \le\frac{q_n(x,y)+\sqrt{q_n(x,y)^2+16S/\sqrt n}}2.
\end{equation}
Taking its square root gives a corresponding refinement for $\norm{a-b}_\infty/C$. The simpler quantity $H_n$ will be used below to keep the subsequent tail estimates in one power form.

This proof also proves separation by the full coefficient tower: equal elementary coefficients give $q_n=0$ for every $n$, and then $x=y$. General square-summable spectra are first rescaled into $[0,1]$. Polynomial moment inversion for spectral estimation has established precedents~\cite{KongValiant}. Here the square-mass encoding and the integer jump of the spectral counting function produce the gap-free estimate above. In the noiseless case the displayed argument yields the explicit position rate $O(n^{-1/8})$.

\subsection{Strictly subcritical topology and its completion}
For $r>0$ define
\begin{equation}\label{eq:pfaffian_weak_distance}
 d_r(a,b)=\sup_{k\ge1}k^{1/r}(|a-b|)^*_k,
\end{equation}
where the star denotes decreasing rearrangement of the absolute difference. This quasi-distance defines the usual separated weak-Lorentz uniform structure. For canonical real skew block representatives
$A(a)=\bigoplus_j a_jJ_0$, $J_0=\left(\begin{smallmatrix}0&1\\-1&0\end{smallmatrix}\right)$,
\begin{equation}\label{eq:pfaffian_canonical_distance}
 \norm{A(a)-A(b)}_{\Sch_{r,\infty}}=2^{1/r}d_r(a,b).
\end{equation}
The quasi-norm gauge in \eqref{eq:pfaffian_canonical_distance} is evaluated on these canonical block representatives.

\begin{theorem}[Quantitative spectral quotient completion]\label{thm:pfaffian_subcritical}
Fix $0<\rho<2$, $C>0$, and $r>\rho$. On $\mathcal K_\rho(C)$,
\begin{align}
 d_r(a,b)&\le C H_n(a,b)^{1-\rho/r},\label{eq:pfaffian_subcritical_inverse}\\
 q_n(a,b)&\le2(2n+1)S_\rho
       \left(\frac{\norm{a-b}_\infty}{C}\right)^{1-\rho/2}.
 \label{eq:pfaffian_forward_inverse}
\end{align}
The full elementary-coefficient map $J(a)=(E_k(a))_{k\ge1}$ is a uniform isomorphism onto its image, with the coordinate product uniformity on the data side. The class is compact and complete for the uniformity induced by $d_r$, and
\begin{equation}\label{eq:pfaffian_completed_iso}
 \widehat{\mathcal K_\rho^{\rm fin}(C)}^{\,d_r}
 \simeq\mathcal K_\rho(C)
 \xrightarrow[\ J\ ]{\simeq}
 \overline{J(\mathcal K_\rho^{\rm fin}(C))}^{\,\rm coord}.
\end{equation}
For every $p\ge1$ with $p>\rho$, the same coordinate topology is also the $\ell^p$ topology of the canonical spectral representatives. Quantitatively,
\begin{equation}\label{eq:pfaffian_strong_subcritical}
 \norm{a-b}_{\ell^p}
 \le\left(\frac{p}{p-\rho}\right)^{1/p}
       C H_n(a,b)^{1-\rho/p}.
\end{equation}
\end{theorem}
\begin{proof}
Put $\varepsilon=\norm{a-b}_\infty$. Positivity and the common envelope give
\begin{equation}\label{eq:pfaffian_interpolation_envelope}
 |a_j-b_j|\le\min\{\varepsilon,Cj^{-1/\rho}\}.
\end{equation}
The right side decreases with $j$, so it also bounds the decreasing rearrangement. Optimizing $j^{1/r}\min\{\varepsilon,Cj^{-1/\rho}\}$ gives
$d_r(a,b)\le C^{\rho/r}\varepsilon^{1-\rho/r}$.
For any $p>\rho$, comparison with the decreasing integral gives
\begin{equation}\label{eq:pfaffian_integral_interpolation}
 \sum_{j\ge1}|a_j-b_j|^p
 \le\int_0^\infty\min\{\varepsilon^p,C^pt^{-p/\rho}\}\,dt
 =\frac{p}{p-\rho}C^\rho\varepsilon^{p-\rho}.
\end{equation}
Theorem~\ref{thm:pfaffian_inverse} proves both inverse estimates.

For the forward bound, the vector $v_n(t)=(tb_{k,n}(t))_{k=0}^n$ has $\ell^1$ derivative at most $2n+1$, because
$b'_{k,n}=n(b_{k-1,n-1}-b_{k,n-1})$ and the Bernstein masses sum to one. Coupling entries in their specified order yields
\[
 q_n(a,b)\le(2n+1)\norm{x-y}_1.
\]
Cauchy--Schwarz and \eqref{eq:pfaffian_integral_interpolation} with $p=2$ give
\[
 \norm{x-y}_1
 \le C^{-2}(\norm a_2+\norm b_2)\norm{a-b}_2
 \le2S_\rho(\varepsilon/C)^{1-\rho/2}.
\]
This proves \eqref{eq:pfaffian_forward_inverse}. Finite Newton transformations and $\norm{a-b}_\infty\le d_r(a,b)$ give forward uniform continuity of the elementary coordinates. Conversely, for every target tolerance, choose a fixed large $n$ in \eqref{eq:pfaffian_subcritical_inverse} and then make its finite polynomial data discrepancy small. These two estimates give a uniform equivalence between the coordinate product uniformity and the weak-Lorentz uniformity generated by $d_r$. For general $r>0$, $d_r$ is a quasi-norm gauge and need not itself satisfy the metric triangle inequality. Compactness, Cauchy convergence, and completion below refer to its translation-invariant uniformity.

To prove compactness, start from any sequence in $\mathcal K_\rho(C)$ and extract a subsequence converging at every coordinate. The limit is nonnegative, decreasing, and satisfies the same envelope. Given $M$, the tail bound
\[
 \sup_l\sup_{j>M}|a_j^{(l)}-a_j|\le 2C(M+1)^{-1/\rho}
\]
controls all indices beyond $M$, while convergence on the finite head gives uniform convergence of the subsequence. Equations~\eqref{eq:pfaffian_interpolation_envelope}--\eqref{eq:pfaffian_integral_interpolation} then give convergence in $d_r$ and in every stated $\ell^p$ topology. Thus the class is compact. If a sequence is $d_r$-Cauchy, compactness gives a convergent subsequence and the Cauchy property forces the full sequence to have the same limit, proving completeness. Finally finite truncations satisfy
\[
 d_r(a,a^{[M]})\le C\sup_{k\ge1}k^{1/r}(M+k)^{-1/\rho}
 \le CM^{1/r-1/\rho}\longrightarrow0.
\]
Thus finite spectra are dense, the image is compact and closed, and \eqref{eq:pfaffian_completed_iso} follows. The same argument with $p>\rho$ proves the asserted equivalence of strong spectral topologies on this class.
\end{proof}

\subsection{The critical obstruction and relative-tail recovery}
\begin{proposition}[Critical data convergence does not give critical recovery]\label{prop:pfaffian_critical_failure}
For $0<\rho<2$, the class $\mathcal K_\rho(C)$ has no uniform inverse from coefficientwise or locally uniform determinant data to $d_\rho$. This remains true if the full compatible exterior tower is observed with all its exponential radius seminorms.
\end{proposition}
\begin{proof}
Use the decreasing flat blocks
\[
 a_j^{(L)}=CL^{-1/\rho}\mathbf1_{\{j\le L\}},\qquad A_L=A(a^{(L)}).
\]
Then $d_\rho(a^{(L)},0)=C$, whereas
$\sum_j(a_j^{(L)})^2=C^2L^{1-2/\rho}\to0$. In particular $A_L\to0$ in $\Sch_2$, so Lemma~\ref{lem:pfaffian_exterior} gives convergence of the whole generated tower in every radius. Also
\[
 \mathscr D_{a^{(L)}}(z)=(1+zC^2L^{-2/\rho})^L\longrightarrow1
\]
locally uniformly, since its difference from one is bounded by
$\exp(|z|C^2L^{1-2/\rho})-1$. The target distance nevertheless stays equal to $C$.
\end{proof}

The same issue prevents finite-rank density in the whole critical weak ideal. For compact operators and any $\rho>0$, its finite-rank closure is
\begin{equation}\label{eq:pfaffian_little_ideal}
 \Sch_{\rho,\infty}^0=\{A:n^{1/\rho}s_n(A)\longrightarrow0\}.
\end{equation}
Indeed singular-value truncation gives approximation for the displayed little-tail condition. Conversely, for a rank-$r$ operator $F$, the approximation-number inequality
$s_{n+r}(A)\le s_n(A-F)$ implies
$\limsup_n n^{1/\rho}s_n(A)\le\norm{A-F}_{\rho,\infty}$.
Arbitrarily accurate finite-rank approximation forces the limit to vanish. Thus a saturated tail $a_j=Cj^{-1/\rho}$ cannot be completed from finite ranks in its critical norm.

Let $a^\circ\in\mathcal K_\rho(C)$ be a known reference spectrum and let $b_M\downarrow0$. Set
\begin{equation}\label{eq:pfaffian_relative_class}
 \mathcal K(a^\circ,b)=\left\{a\in\mathcal K_\rho(C):
 \sup_{j>M}j^{1/\rho}|a_j-a_j^\circ|\le b_M\quad(M\ge1)\right\}.
\end{equation}
The reference may in particular have a saturated critical tail. On sorted spectra also define the auxiliary rank metric
$d_\rho^{\rm rk}(a,b)=\sup_jj^{1/\rho}|a_j-b_j|$.

\begin{theorem}[Critical completion on a relative-tail class]\label{thm:pfaffian_critical}
The class \eqref{eq:pfaffian_relative_class} is compact and complete in $d_\rho^{\rm rk}$. On that class the rank metric and the weak-ideal difference $d_\rho$ give the same uniformity, with
\begin{align}
 d_\rho(a,b)&\le d_\rho^{\rm rk}(a,b),\label{eq:pfaffian_rank_compare}\\
 d_\rho^{\rm rk}(a,b)&\le\inf_{M\ge1}
    \max\{M^{1/\rho}d_\rho(a,b),2b_M\}.
 \label{eq:pfaffian_rank_reverse}
\end{align}
The full invariant map has a uniformly continuous inverse on its closed image. An explicit finite-data version is
\begin{equation}\label{eq:pfaffian_critical_inverse}
 d_\rho(a,b)\le\inf_{M\ge1}
     \max\{M^{1/\rho}CH_n(a,b),2b_M\}.
\end{equation}
\end{theorem}
\begin{proof}
Suppose $d=d_\rho^{\rm rk}(a,b)$. The coordinate bound $|a_j-b_j|\le dj^{-1/\rho}$ also bounds the decreasing rearrangement, proving \eqref{eq:pfaffian_rank_compare}. Conversely, at each coordinate $|a_j-b_j|\le d_\rho(a,b)$. On the first $M$ coordinates this gives the weighted bound $M^{1/\rho}d_\rho(a,b)$; after $M$, the common relative-tail assumptions give the bound $2b_M$. Taking the maximum and then the infimum in $M$ proves \eqref{eq:pfaffian_rank_reverse}. Replacing the coordinate estimate by $CH_n$ proves \eqref{eq:pfaffian_critical_inverse}.

For compactness, take any sequence in the class and extract a subsequence converging at every spectral coordinate. Positivity, decreasing order, the envelope, and every relative-tail inequality pass to the limit $a$. For fixed $M$, the weighted difference on the first $M$ coordinates tends to zero. On the remaining coordinates it is at most $2b_M$, uniformly over the subsequence. Let the subsequence index tend to infinity first and then let $M$ increase. This gives convergence in $d_\rho^{\rm rk}$, proving compactness. A Cauchy sequence has a convergent subsequence by compactness, and then converges to the same limit, proving completeness.

The comparison inequalities give the same uniformity for $d_\rho$ and $d_\rho^{\rm rk}$. The elementary spectral coordinates are continuous, by their $\ell^2$ continuity on the envelope class. For the inverse, first choose $M$ to control $b_M$, then choose the moment order $n$, and finally choose the accuracy of its finite polynomial readout. Formula~\eqref{eq:pfaffian_critical_inverse} proves uniform continuity in precisely this order. The image is compact and therefore closed.
\end{proof}

Zero reference gives a uniform little-tail class. A known nonzero reference $a_j^\circ=cj^{-1/\rho}$ retains a saturated critical component. The relative-tail bound is supplied by the source or model class. On this class it gives the critical inverse and the corresponding completion, while the whole critical weak ideal retains the obstruction of Proposition~\ref{prop:pfaffian_critical_failure}.

\subsection{Unknown saturated tails as additional invariant coordinates}
The leading tail coefficient can also be treated as an invariant coordinate when the remainder is controlled uniformly. Consider spectra in \eqref{eq:pfaffian_class} for which some unknown $c\in[0,C]$ satisfies
\begin{equation}\label{eq:pfaffian_asymptotic_class}
 \sup_{j>M}|j^{1/\rho}a_j-c|\le b_M\downarrow0\qquad(M\ge1).
\end{equation}
We may take $b_M\le C$. The coordinate $c$ is unique. Recall the finite positive constant $B_\rho$ from Proposition~\ref{prop:pfaffian_gauges}; integration by parts gives
$B_\rho=\pi/\sin(\pi\rho/2)$.

\begin{proposition}[Recovering the critical tail coefficient]\label{prop:pfaffian_asymptotic}
For \eqref{eq:pfaffian_asymptotic_class},
\begin{equation}\label{eq:pfaffian_asymptotic_limit}
 \lim_{x\to\infty}x^{-\rho/2}\log\mathscr D_a(x)=B_\rho c^\rho.
\end{equation}
More precisely, for $M\ge1$ and $x>0$,
\begin{align}
 B_\rho(c-b_M)_+^\rho-E_{M,x}
 &\le x^{-\rho/2}\log\mathscr D_a(x)
 \le B_\rho(c+b_M)^\rho+E_{M,x},\label{eq:pfaffian_asymptotic_bounds}\\
 E_{M,x}&=x^{-\rho/2}\big[(M+1)\log(1+4C^2x)+2/\rho\big].
 \label{eq:pfaffian_asymptotic_error}
\end{align}
The class with common $C$ and $(b_M)$ is compact in $d_\rho^{\rm rk}$; its finite invariant data therefore have a uniform critical inverse.
\end{proposition}
\begin{proof}
For a pure power spectrum $a_j=dj^{-1/\rho}$, the decreasing function
$f(t)=\log(1+xd^2t^{-2/\rho})$ satisfies
\[
 \int_1^\infty f(t)\,dt\le\sum_jf(j)\le\int_0^\infty f(t)\,dt
 =B_\rho d^\rho x^{\rho/2}.
\]
The omitted integral over $(0,1)$ is at most $\log(1+xd^2)+2/\rho$, using
$1+xd^2t^{-2/\rho}\le(1+xd^2)t^{-2/\rho}$. This includes $d=0$.
For a general spectrum, after rank $M$ it lies between the pure powers with coefficients $(c-b_M)_+$ and $c+b_M$. Restoring or deleting their first $M$ terms and bounding the true head by $M\log(1+xC^2)$ gives \eqref{eq:pfaffian_asymptotic_bounds}. First let $x\to\infty$ and then $M\to\infty$ to obtain the limit, with both error choices uniform on the class.

For compactness extract $c_l\to c$ and then a coordinatewise convergent spectral subsequence. All class inequalities pass to the limit. Past rank $M$ the weighted difference is at most $2b_M+|c_l-c|$; the finite weighted head converges. This proves convergence in the rank metric. Continuous separating invariant coordinates and Proposition~\ref{prop:invariant_compact} give the final assertion.
\end{proof}

A finite approximation in this class keeps the estimated leading tail $\widehat c j^{-1/\rho}$ and truncates the weighted-$c_0$ remainder. The completed coordinate is therefore the finite-dimensional parameter $c$ together with that remainder, subject to positivity and spectral order.

\subsection{Finite exterior readings and certified spectral clusters}
Suppose the observed exterior levels obey
$\norm{\widehat\Xi_{2k}-\Xi_{2k}}\le u_k$. With
$\widehat c_k=4^{-k}\norm{\widehat\Xi_{2k}}^2$, the computable coefficient error is
\begin{equation}\label{eq:pfaffian_exterior_reading_error}
 |\widehat c_k-c_k|\le\epsilon_k
 :=4^{-k}u_k(2\norm{\widehat\Xi_{2k}}+u_k).
\end{equation}
This follows from the difference of squared norms. A direct scalar protocol can instead supply $\widehat c_k$ and $\epsilon_k$ themselves. If $\sum_ja_j^2\le S$, define
\begin{align}
 P_m(z)&=1+\sum_{k=1}^m\widehat c_kz^k,\label{eq:pfaffian_polynomial_data}\\
 \Delta_m(R)&=\sum_{k=1}^m\epsilon_kR^k
       +e^{RS}\frac{(RS)^{m+1}}{(m+1)!}.
 \label{eq:pfaffian_analytic_budget}
\end{align}
Then $\sup_{|z|\le R}|\mathscr D_a(z)-P_m(z)|\le\Delta_m(R)$, since
$c_k\le S^k/k!$ and $(m+1+h)!\ge(m+1)!h!$. The index $m$ is an exterior-order budget, not a bound on operator rank.

\begin{theorem}[Finite spectral and asymptotic certificates]\label{thm:pfaffian_contours}
For a simple closed contour $\Gamma\subset\{|z|\le R\}$, the condition
\begin{equation}\label{eq:pfaffian_rouche}
 \min_{z\in\Gamma}|P_m(z)|>\Delta_m(R)
\end{equation}
certifies equality of the zero counts of $P_m$ and $\mathscr D_a$ inside $\Gamma$, with multiplicities. For a circle of centre $-d<0$ and radius $s<d$, a count of $h$ certifies exactly $h$ positive skew values in
\begin{equation}\label{eq:pfaffian_spectral_interval}
 [(d+s)^{-1/2},(d-s)^{-1/2}].
\end{equation}
If disjoint certified inner circles lie inside a certified circle $|z|=R_0$, and their counts sum to the latter count $J\ge1$, they cover the largest $J$ skew block values. If, in addition, $a\in\mathcal K_\rho(C)$, the recovered head has coordinate errors at most $h_*$, and the unrecovered coordinates are set to zero,
\begin{equation}\label{eq:pfaffian_finite_head_error}
 \norm{a-\widehat a}_{\ell^p}
 \le\left(Jh_*^p+\frac{C^p\rho}{p-\rho}J^{1-p/\rho}\right)^{1/p},
 \qquad p\ge1,\quad p>\rho.
\end{equation}

Under \eqref{eq:pfaffian_asymptotic_class}, suppose also that
$|\widehat c-c|\le\epsilon_c$. Appending the blocks
$\widehat c j^{-1/\rho}J_0$ for $j>J$ gives an actual compact skew operator $\widehat A$ with
\begin{equation}\label{eq:pfaffian_critical_head_error}
 \norm{\widehat A-A(a)}_{\rho,\infty}
 \le2^{1/\rho}\max\{J^{1/\rho}h_*,b_J+\epsilon_c\}.
\end{equation}
The comparison uses the same canonical block planes, not unobserved eigenplanes of the original operator.
\end{theorem}
\begin{proof}
The homotopy $P_m+t(\mathscr D_a-P_m)$, $0\le t\le1$, has no zero on $\Gamma$ by \eqref{eq:pfaffian_rouche}. Its winding number is constant, so the argument principle gives the zero count. All true zeros are real and negative and equal to $-a_j^{-2}$, which proves \eqref{eq:pfaffian_spectral_interval}. The outer circle contains exactly those zeros with $a_j>R_0^{-1/2}$; summing the inner counts leaves none of these uncovered. When certified disks overlap in the spectral coordinate, they are treated as one cluster and only the total multiplicity is used.

The finite-head error sums to at most $Jh_*^p$. The envelope gives
$\sum_{j>J}a_j^p\le C^p\rho(p-\rho)^{-1}J^{1-p/\rho}$, proving \eqref{eq:pfaffian_finite_head_error}. In the critical case, the error sequence is bounded by
$\max\{J^{1/\rho}h_*,b_J+\epsilon_c\}j^{-1/\rho}$. Its decreasing rearrangement has the same bound. Each block contributes two equal singular values, giving \eqref{eq:pfaffian_critical_head_error}. The appended numerical sequence need not be decreasing at the joining point; it nevertheless defines the actual block operator whose difference is estimated.
\end{proof}

The asymptotic coefficient can itself be certified from finite data. Choose $x>0$ with $\Delta_m(x)<1/2$. Then $\mathscr D_a(x)\ge1$ and $P_m(x)>1/2$, so
$|\log\mathscr D_a(x)-\log P_m(x)|\le2\Delta_m(x)$.
For
\[
 \widehat h=x^{-\rho/2}\log P_m(x),\qquad
 u=E_{M,x}+2x^{-\rho/2}\Delta_m(x),
\]
Proposition~\ref{prop:pfaffian_asymptotic} gives the finite interval
\begin{equation}\label{eq:pfaffian_tail_interval}
 c\in\left[
 \left(\frac{(\widehat h-u)_+}{B_\rho}\right)^{1/\rho}-b_M,
 \left(\frac{(\widehat h+u)_+}{B_\rho}\right)^{1/\rho}+b_M
 \right]\cap[0,C].
\end{equation}
If the interval is empty, the data and the certificate class are incompatible. Otherwise its midpoint and half-width give $\widehat c,\epsilon_c$. First choose $M$ to reduce the model remainder, next $x$ to reduce $E_{M,x}$, and finally $m$ and the coefficient precision to reduce $\Delta_m(x)$. These are finite operations, although their required precision can be large.

The contour test requires separation only between the chosen contour and the approximating polynomial; near collisions it certifies the whole enclosed cluster. For example the squared spectra $(1+t,1-t)$ and $(1,1)$ have determinants differing by $t^2z^2$, while their positive skew values differ by order $|t|$. Thus no uniform coefficient-to-labelled-root Lipschitz estimate holds near this collision. The global inverse in Theorem~\ref{thm:pfaffian_inverse} and the local cluster certificates address different questions.

\begin{corollary}[A complete finite-observation budget]\label{cor:pfaffian_finite_budget}
Suppose the source or model certifies membership in $\mathcal K_\rho(C)$, and the normalized coefficients $E_1,\ldots,E_{n+1}$ are observed within errors $\varepsilon_1,\ldots,\varepsilon_{n+1}$. The nonempty fibre $\mathcal F_n$ of all fitting candidates in this class satisfies, for $r>\rho$,
\begin{equation}\label{eq:pfaffian_fibre_error}
 \operatorname{diam}_{d_r}\mathcal F_n
 \le C\min\left\{1,
 \left(Q_n(2\varepsilon;S_\rho)+\frac{4S_\rho}{\sqrt n}\right)^{(1-\rho/r)/4}\right\}.
\end{equation}
On a relative-tail class, Theorem~\ref{thm:pfaffian_critical} gives the corresponding critical bound. Equation~\eqref{eq:pfaffian_exterior_reading_error}, divided by $C^{2k}$, supplies these normalized errors from finite exterior readings.
\end{corollary}
\begin{proof}
Two fitting candidates differ by at most $2\varepsilon_k$ in their $k$th normalized coefficient. Apply Lemma~\ref{lem:pfaffian_newton}, then Theorem~\ref{thm:pfaffian_subcritical} or Theorem~\ref{thm:pfaffian_critical}. The fibre conclusions are precisely those of Theorem~\ref{thm:invariant_uniform}.
\end{proof}

The degree-dependent quantity $Q_n$ records the conditioning of the Newton--Bernstein transform. Quantitative recovery therefore fixes an observation order and then chooses its noise tolerance. General moment-based spectral approximation exhibits comparable high-precision costs~\cite{JinEtAl}.

\subsection{From finite Gaussian coefficients to random spectral certificates}
The deterministic spectral decoder can be connected to the actual source observations, without postulating noiseless invariant coordinates. For each $1\le k\le n+1$, suppose a compatible, jointly realized exterior level has a homogeneous representation
$\Xi_{2k}=I_{m_k}(L_k)$, where $m_k\ge0$ is its specified source degree. View an exterior vector as a column operator from $\C$ to the complexification of $\wedge^{2k}E$; its Schatten--Hilbert norm is exactly its exterior norm. A column can equivalently be embedded in the corresponding square block-operator space. Thus the finite Hermite decoder applies with $s=2$, retaining the physical column type.

Let $P_{a_k}L_k$ be its finite source and physical truncation, let $\delta_k$ be the Euclidean error of the retained Hermite marks, and let $\kappa_k$ be the inverse anchor cost. Suppose deterministic bounds $t_k$ and $B_k$ certify
\[
 \sqrt{m_k!}\norm{L_k-P_{a_k}L_k}_2\le t_k,
 \qquad \norm{\Xi_{2k}}_{L^2}\le B_k.
\]
The decoded finite column $\widehat L_k$ gives the actual finite random exterior vector $\widehat\Xi_{2k}=I_{m_k}(\widehat L_k)$.
The finite readings are fixed when applying the next proposition, so each decoded coefficient kernel is deterministic on the canonical source space. For readings obtained by independent repetitions, one first conditions on their certified measurement event and then evaluates these kernels on a fresh realization of that same source. The measurement-failure probability is added to the probability bounds below. This distinguishes population reconstruction from reuse of the training samples.

\begin{proposition}[End-to-end finite Gaussian spectral budget]\label{prop:pfaffian_source_budget}
Put $b_k=\kappa_k\delta_k+t_k$. Then
\begin{equation}\label{eq:pfaffian_source_error}
 \norm{\widehat\Xi_{2k}-\Xi_{2k}}_{L^2}\le b_k,\qquad
 \E|\widehat c_k-c_k|\le4^{-k}b_k(2B_k+b_k),
 \quad\widehat c_k=4^{-k}\norm{\widehat\Xi_{2k}}^2.
\end{equation}
Choose $\theta_k>0$ and a spectral budget $C>0$, and set
\begin{equation}\label{eq:pfaffian_source_tolerances}
 \varepsilon_k=\frac{4^{-k}C^{-2k}b_k(2B_k+b_k)}{\theta_k}.
\end{equation}
If $\Prob\{W_\rho(a)>C\}\le\delta_0$, then with probability at least
$1-\delta_0-\sum_{k=1}^{n+1}\theta_k$, the true spectrum belongs to the candidate fibre for the observations $C^{-2k}\widehat c_k$ with tolerances $\varepsilon_k$. On that event \eqref{eq:pfaffian_fibre_error} controls every fitting candidate. A critical conclusion additionally requires its relative or asymptotic tail certificate on the same event.
\end{proposition}
\begin{proof}
At $s=2$, the normalized coefficient map is an isometry and the finite rank loss in Theorem~\ref{thm:finitedecoder} is one. Its noise bound and the given truncation tail prove the first estimate in \eqref{eq:pfaffian_source_error}. Write the vector difference as $V_k$. The inequality
$|\norm{\Xi_{2k}+V_k}^2-\norm{\Xi_{2k}}^2|
 \le2\norm{\Xi_{2k}}\norm{V_k}+\norm{V_k}^2$
and Cauchy--Schwarz give the second estimate. When $b_k>0$, Markov's inequality at the tolerance in \eqref{eq:pfaffian_source_tolerances} gives failure probability at most $\theta_k$. When $b_k=0$, both reconstruction errors vanish almost surely, so the zero tolerance is valid directly. A union bound with the spectral-budget event gives the probability claim. Corollary~\ref{cor:pfaffian_finite_budget} applies pointwise on this event. No independence between the exterior levels is required.
\end{proof}

The source degree $m_k$ and the exterior order $2k$ are distinct indices; a model that links them must provide that relation explicitly. For a finite-chaos level, apply the same estimate to its components and sum their $L^2$ errors; an infinite expansion requires the all-radius tail certificate of Theorem~\ref{thm:main}. If the exterior levels are generated from recovered two-forms rather than observed separately, their finite-operation error follows from Proposition~\ref{prop:cascade}: the exterior product
$\wedge^pE\times\wedge^qE\to\wedge^{p+q}E$
has norm at most $\binom{p+q}p^{1/2}$, by antisymmetric tensor projection. Probability H\"older and the fixed-chaos bounds must also be charged at these operations. These estimates specify the source certificate that enters the spectral inverse.

The means of all $c_k$ alone do not determine the random spectral law. The models $A=J_0$ and $A=ZJ_0$, where $Z=\sqrt{1/2}$ or $\sqrt{3/2}$ equally likely, both satisfy $\E\mathscr D_A(z)=1+z$, but have different spectral laws. Same-sample bounded mixed coordinates use the separate law protocol of Appendix~\ref{sec:obsappendix}.

\subsection{The determinant closure under only a square-mass bound}
The stronger envelope $\rho<2$ excludes square mass escaping to zero. Without it, the observation completion is larger than the genuine spectral image. To describe it, use unnormalized squared values and put
\[
 \mathcal H_S=\{x_1\ge x_2\ge\cdots\ge0:\ \sum_jx_j\le S\},\qquad
 D_x(z)=\prod_j(1+zx_j),\quad S>0.
\]
The Hilbert--Schmidt bound supplies the mass budget, while the completed topology is the locally uniform determinant topology.

\begin{theorem}[Square mass at zero and the exact invariant closure]\label{thm:pfaffian_mass_defect}
The closure of $\{D_x:x\in\mathcal H_S\}$ in the topology of locally uniform convergence is
\begin{equation}\label{eq:pfaffian_defect_class}
 \mathfrak D_S=\left\{e^{\tau z}\prod_j(1+zx_j):
      x\in\mathcal H_S,\quad\tau\ge0,\quad\sum_jx_j+\tau\le S\right\}.
\end{equation}
The parameters $(x,\tau)$ are unique. This is also the coefficient completion of the finite spectra, and coefficientwise and locally uniform convergence agree on the bounded class.

If $D_{x^{(l)}}\to D_{x,\tau}$, then $x_j^{(l)}\to x_j$ for every $j$ and
\begin{equation}\label{eq:pfaffian_zero_mass_formula}
 \tau=\lim_l\sum_jx_j^{(l)}-\sum_jx_j
      =\lim_{M\to\infty}\lim_l\sum_{j>M}x_j^{(l)}.
\end{equation}
For such a convergent sequence the following are equivalent:
$\tau=0$; $\lim_M\sup_l\sum_{j>M}x_j^{(l)}=0$; and
$\norm{x^{(l)}-x}_{\ell^1}\to0$.
\end{theorem}
\begin{proof}
The mass bound controls every power moment except for the possible loss in the first one. From a sequence in $\mathcal H_S$, extract a subsequence with coordinate limits $x_j$ and total-mass limit $s\in[0,S]$. Fatou gives $\sum_jx_j\le s$; put $\tau=s-\sum_jx_j$. Decreasing order implies $x_j^{(l)}\le S/j$, and hence for $k\ge2$,
\begin{equation}\label{eq:pfaffian_higher_mass_tail}
 \sum_{j>M}(x_j^{(l)})^k
 \le S\left(\frac{S}{M+1}\right)^{k-1}.
\end{equation}
The tails of these higher moments are uniform. Thus $p_k(x^{(l)})\to p_k(x)$ for $k\ge2$, while $p_1(x^{(l)})\to\sum_jx_j+\tau$.

On $|z|<1/S$, the logarithmic determinant expansion is uniformly absolutely convergent on smaller disks. It follows that
\[
 \log D_{x^{(l)}}(z)
 \longrightarrow\tau z+\sum_j\log(1+zx_j).
\]
Equivalently, Newton recursion gives convergence of every determinant coefficient. The common bound $e_k(x^{(l)})\le S^k/k!$ upgrades this coefficient convergence to local uniform convergence on the entire plane. Therefore every possible limit belongs to \eqref{eq:pfaffian_defect_class}.

Conversely, given $(x,\tau)$ in that class, take the first $l$ values of $x$ together with $l$ copies of $\tau/l$, and sort the resulting finite spectrum. Its mass is at most $S$ and its determinant is
\[
 \prod_{j\le l}(1+zx_j)(1+\tau z/l)^l
 \longrightarrow e^{\tau z}\prod_j(1+zx_j)
\]
locally uniformly. This proves the exact closure and finite-spectrum density. The same coefficient majorant holds on the closure, so coefficientwise convergence there implies local uniform convergence. The reverse implication is Cauchy's coefficient formula. In particular this is the completion in either of the stated data topologies.

The nonzero zeros determine all positive $x_j$ with multiplicity. The derivative at zero then determines $\tau$ through $D'_{x,\tau}(0)=\sum_jx_j+\tau$. Thus the parameters are unique.

We also need continuity of a finite spectral head on the entire closure, not only along genuine spectra. Suppose $D_{x^{(l)},\tau_l}$ converges in $\mathfrak D_S$. Extract a subsequence with coordinate limits $x_j$ and total-mass limit $s=\lim_l(\tau_l+\sum_jx_j^{(l)})$. The uniform higher-moment bound \eqref{eq:pfaffian_higher_mass_tail} still applies. In the logarithmic expansion the first coefficient tends to $s$ and every higher coefficient tends to the power sum of $x$. The earlier coefficient-majorant argument therefore identifies the limit as $D_{x,s-\sum_jx_j}$. Uniqueness forces these same coordinate limits for every extracted subsequence. Hence each map $D_{x,\tau}\mapsto x_j$ is continuous on $\mathfrak D_S$, as is its total mass $D'_{x,\tau}(0)$. For a genuine approximating sequence, subtracting any finite head gives \eqref{eq:pfaffian_zero_mass_formula}.

Define the continuous residuals
\begin{equation}\label{eq:spectral_monotone_defect}
 r_M(D_{x,\tau})=D'_{x,\tau}(0)-\sum_{j\le M}x_j
               =\tau+\sum_{j>M}x_j.
\end{equation}
They lie in $[0,S]$ and decrease to $\tau$. Lemma~\ref{lem:monotone_defect} consequently identifies vanishing of the limit defect with uniform tail tightness of any convergent genuine sequence. Coordinate convergence and the common tail give $\ell^1$ convergence by a head--tail decomposition. Conversely $\ell^1$ convergence preserves the total mass and forces $\tau=0$. This proves all three equivalences. The same residuals show that the genuine spectral image is a dense $G_\delta$ subset of this determinant closure and that $\tau$ is upper semicontinuous there.
\end{proof}

For example, $x_j^{(L)}=L^{-1}\mathbf1_{j\le L}$ has determinant
$(1+z/L)^L\to e^z$, which corresponds to $(x,\tau)=(0,1)$. A genuine discrete spectrum with determinant $e^z$ would require $p_1=1$ and $p_2=0$, which is impossible. Thus $\tau$ records square mass concentrated at the spectral point zero.

The product form in \eqref{eq:pfaffian_defect_class} is classical Laguerre--P\'olya type~I structure~\cite{NguyenVishnyakova}. Under the fixed mass budget, this product representation identifies the observation completion and the criterion selecting genuine spectra. Equivalently, this topology is weak convergence of the measures
$\tau\delta_0+\sum_jx_j\delta_{x_j}$ on $[0,S]$: Newton's identities identify their polynomial moments, and polynomial approximation identifies weak convergence. This topology permits transfer of mass into the zero atom: the preceding finite spectra have $\tau_L=0$, whereas their limit has $\tau=1$. Hence the coordinates $(x,\tau)$ do not carry the product topology.
For $a\in\mathcal K_\rho(C)$,
\[
 \sum_{j>M}a_j^2\le\frac{C^2\rho}{2-\rho}M^{1-2/\rho}\longrightarrow0
\]
uniformly, so the defect is excluded. This supplies the true-image step in \eqref{eq:pfaffian_completed_iso}.

\subsection{Cross-invariants and operations on the quotient}\label{subsec:pfaffian_cross}
Scalar spectral towers discard relative eigenplanes. A common-frame protocol retains them sufficiently to generate the difference before taking its invariants. For real skew $A,B\in\Sch_2(E)$, define
\[
 U_{n,k}=F_k(\alpha_A)\wedge F_{n-k}(\alpha_B),\qquad
 G_n(A,B)_{kl}=\ip{U_{n,k}}{U_{n,l}},\qquad s_k=(-1)^{n-k}.
\]
\begin{proposition}[Polarized exterior differences]\label{prop:pfaffian_cross_difference}
If these common-frame cross-Gram matrices are observed, then
\begin{align}
 \sum_{k=0}^ns_kU_{n,k}&=F_n(\alpha_{A-B}),&
 \norm{F_n(\alpha_{A-B})}^2&=s^TG_n(A,B)s,\label{eq:pfaffian_difference_identity}\\
 \norm{A-B}_{\rho,\infty}
 &\asymp_\rho\sup_{n\ge1}(n!)^{1/(\rho n)}
                    (s^TG_n(A,B)s)^{1/(2n)},\quad 0<\rho<2.
 \label{eq:pfaffian_pair_gauge}
\end{align}
The two sides in the second statement may both be infinite.
\end{proposition}
\begin{proof}
The commuting two-form binomial formula, with $F_k(\alpha)=2^k\alpha^{\wedge k}/k!$, gives the first identity, including every binomial coefficient. Expanding its real squared norm gives the Gram quadratic form. This is a genuinely compatible tower for $A-B$, so Proposition~\ref{prop:pfaffian_gauges} applies to the difference. The doubled skew multiplicities convert $W_\rho$ to the operator weak-ideal norm.
\end{proof}

At a fixed level, an operator-norm error $\eta_n$ in $G_n$ gives squared-reading error at most $(n+1)\eta_n$, since $\norm s_2^2=n+1$. Passage to the supremum over all levels requires a joint noise and tail budget. The cross-Gram data retain relative eigenplane information that is absent from the two scalar determinants $\mathscr D_A$ and $\mathscr D_B$.

Only operations compatible with the spectral equivalence relation descend to the scalar spectral quotient. For example $J_0\sim_{\rm sp}-J_0$, but adding $J_0$ gives $2J_0$ and zero. Orthogonal direct sum does descend:
\[
 \mathscr D_{a\sqcup b}(z)=\mathscr D_a(z)\mathscr D_b(z).
\]
In the defect completion it sends $(x,\tau),(y,\sigma)$ to
$(\operatorname{sort}(x\cup y),\tau+\sigma)$, with added mass budgets. Local uniform multiplication is continuous, so this operation extends. Thus orthogonal direct sum is a continuous quotient operation, whereas \eqref{eq:pfaffian_pair_gauge} uses the stronger common-frame pair protocol.

The exterior observation mechanism therefore has three distinct completion regimes: a true strictly subcritical spectral completion, a critical completion on relative or asymptotic tail classes, and a determinant completion carrying square mass at zero.


\section{Geometric realizations and Wick operator systems}\label{sec:realizations}\label{sec:wickpipeline}\label{sec:completioncomparison}
The reconstruction theory has infinite-dimensional realizations with explicit source and physical tails. We begin with an all-degree noncommutative block model. Covariance geometry then constructs Wick operators on closed manifolds, and finite source-labelled observations recover their generated polynomial algebras and common-noise evolutions.

\subsection{An all-degree noncommutative model}
On $\ell^2(\N)\otimes\C^2$, let
\[
 X=\begin{pmatrix}0&1\\1&0\end{pmatrix},\qquad
 Z=\begin{pmatrix}1&0\\0&-1\end{pmatrix},
\]
and write $X_j,Z_j$ for their copies on the $j$th physical block. For independent real standard Gaussians $g_j$, define
\begin{equation}\label{eq:fullmodel}
 F_X=\bigoplus_{j\ge1}2^{-j}e^{g_j-1/2}X,\qquad
 F_Z=\bigoplus_{j\ge1}2^{-j}e^{g_j-1/2}Z.
\end{equation}
Their kernels are
\[
 K_m^X=\frac1{m!}\sum_{j\ge1}2^{-j}e_j^{\otimes m}\otimes X_j,
 \qquad K_m^Z=\frac1{m!}\sum_{j\ge1}2^{-j}e_j^{\otimes m}\otimes Z_j.
\]
Every cut has nonzero singular values $2^{-j}/m!$, each twice. For $s\ge2$ this immediately computes the original norm. For $s<2$, the single-cut upper bound and dual finite block-supported tests give the matching lower bound. Thus
\begin{equation}\label{eq:modelnorm}
 \norm{K_m^X}_{\Wc_{m,s}}=\norm{K_m^Z}_{\Wc_{m,s}}=\frac{C_s}{m!},
 \quad C_s=\left(\frac2{2^s-1}\right)^{1/s},\qquad
 \Nn\rho{K^X}=C_s\mathcal A(\rho),
\end{equation}
where $\mathcal A(u)=\sum_{m\ge0}u^m/\sqrt{m!}$ is finite for every $u$. These are elements of $\Alg_s$ for every $1\le s<\infty$.

At the Hilbert exponent, $\E\norm{F_X}_2^2=2e/3$. The operators are almost surely compact, bounded, and infinite-rank, while the first lognormal block excludes a deterministic samplewise operator bound. Their commutator is
\begin{equation}\label{eq:modelcommutator}
 [F_X,F_Z]=2\bigoplus_{j\ge1}4^{-j}e^{2g_j-1}XZ\ne0.
\end{equation}
For the normalized faithful anchor $R|_{E_j}=\sqrt{3/2}\,2^{-j}I_2$, retain the first $J$ source and physical blocks and degrees $m\le M$. With $\mathcal A_M(u)=\sum_{m=0}^Mu^m/\sqrt{m!}$,
\begin{equation}\label{eq:modeltail}
 \Nn\rho{K^X-P_{M,2J,2J,J}K^X}
 =C_s\big[2^{-J}\mathcal A_M(\rho)+\mathcal A(\rho)-\mathcal A_M(\rho)\big].
\end{equation}
The physical inverse costs $(2/3)4^J$. For Hilbert output radius $\sigma$, choose $\rho=\sqrt{1+2\sigma^2}$. If each input has coefficient error at most $\varepsilon$ at this radius, the commutator error is at most
$4\sqrt{2/3}\mathcal A(\rho)\varepsilon+2\varepsilon^2$.
Thus the model has explicit observation, truncation, and nonlinear-target budgets.

\subsection{Local-to-global assembly}
Let $A_{\lambda\mu}$ be operator blocks with orthogonal physical row and column supports. For $1\le s<\infty$,
\begin{equation}\label{eq:blockineq}
 \left\|\sum_{\lambda,\mu}A_{\lambda\mu}\right\|_s
 \le\left(\sum_{\lambda,\mu}\norm{A_{\lambda\mu}}_s^{t_s}\right)^{1/t_s},
 \qquad t_s=\min(s,2).
\end{equation}
To prove it, write $C_\mu=\sum_\lambda A_{\lambda\mu}$ and use
$AA^*=\sum_\mu C_\mu C_\mu^*$ and
$C_\mu^*C_\mu=\sum_\lambda A_{\lambda\mu}^*A_{\lambda\mu}$.
For $s\ge2$, apply the triangle inequality twice in $\Sch_{s/2}$. For $s<2$, apply trace subadditivity of the power $s/2$ to the positive sums. These are standard partitioned-operator estimates \cite{BK,BLP,Simon}.

A matching is a set of blocks with no repeated row or column. On a matching the $s$th powers of the Schatten norms add exactly. For a partition $I=\bigsqcup_\theta I_\theta$ into matchings, define
\begin{equation}\label{eq:matchinglattice}
 \norm b_{\mathsf M_s}=\sum_\theta\left(\sum_{i\in I_\theta}|b_i|^s\right)^{1/s}.
\end{equation}
This solid sequence norm controls block assembly and is order continuous: first retain finitely many outer summands, and then truncate within each $\ell^s$ summand. The universal exponent in \eqref{eq:blockineq} is optimal. Diagonal scalar blocks require $t\le s$, and the $N\times N$ all-ones matrix requires $t\le2$.

\begin{proposition}[Order-continuous assembly]\label{prop:assembly}
Suppose an order-continuous solid Banach sequence lattice $\mathsf E(I)$ controls finite block sums in every directed cut. If local kernels $K_i$ have majorant
$b_i=\max_S\norm{\Flat_S(K_i)}_s\in\mathsf E$,
then their sum converges unconditionally in the maximum-cut completion and hence in $\Coeff_{m,s}$. If $K_{n;i}$ are Cauchy in each block and share this majorant, their full kernels and same-source Gaussian evaluations are Cauchy. A finite-head bound $\delta_nG(R)$ and majorant tail $\eta(R)$ give
\begin{equation}\label{eq:localtailmodulus}
 C_{\mathsf E}\inf_R\{\delta_nG(R)+2\eta(R)\}
\end{equation}
as a coefficient bound, multiplied by the fixed-degree forward constant for evaluation.
\end{proposition}
\begin{proof}
For a finite set $F\subset I$, the assumed block estimate bounds every finite sub-sum over $I\setminus F$ by $C_{\mathsf E}\norm{b\mathbf1_{I\setminus F}}_{\mathsf E}$ in each cut. Order continuity makes this bound tend to zero as $F$ exhausts the index set. Thus the net of finite sub-sums is Cauchy simultaneously in all the finitely many cut norms. Their compatible completion gives one limiting kernel, independently of summation order. For $s<2$, the maximum-cut norm dominates the quotient-sum norm; for $s\ge2$ it is the original coefficient norm.

For a varying sequence, choose a common finite set making its majorant tail small before using convergence of the individual blocks. The finitely many retained blocks are Cauchy, and the two discarded tails contribute at most $2C_{\mathsf E}\eta(R)$. A finite-head bound $\delta_nG(R)$ therefore gives \eqref{eq:localtailmodulus}. Finally the fixed-degree Gaussian evaluation is continuous in the original coefficient norm, so the same Cauchy conclusion and its stated forward cost hold for the random responses.
\end{proof}

At $s<2$, one can instead use the physical corner rank $d_i$ and the weaker majorant
$b_i=d_i^{1/s-1/2}\norm{K_i}_2$.
On each evaluated block,
$\norm{I_mK_i}_s\le d_i^{1/s-1/2}\norm{I_mK_i}_2$.
For a matching and probability exponent $p\ge s$, Minkowski in $L^{p/s}$ gives the $\ell^s$ sum of the block $L^p(\Sch_s)$ norms. For smaller $p$, first increase it to $\max(p,s,2)$. Fixed-degree Gaussian moment comparison controls these block norms by $b_i$. Thus $\ell^s$ or $\mathsf M_s$ summability constructs the random response, and the fixed-degree inverse returns its kernel in the original quotient-sum space. Finitely overlapping spectral blocks can be separated into finitely many orthogonal colors.

\subsection{Covariance--Gram estimates}
Let $(X,\mu)$ be $\sigma$-finite, with strongly measurable vectors $\eta_x$ in a real Hilbert source, $a_x\in\cE$, and $b_x\in\cC$. Assume for $0\le j\le m$ that
\[
 \int\norm{\eta_x}^{2j}\norm{a_x}^2\,d\mu(x)<\infty,
 \qquad
 \int\norm{\eta_x}^{2j}\norm{b_x}^2\,d\mu(x)<\infty.
\]
Then $K=\int\eta_x^{\otimes m}\otimes\bar b_x\otimes a_x\,d\mu(x)$ is a Hilbert-tensor Bochner integral. Write
$\gamma(x,y)=\ip{\eta_x}{\eta_y}$,
$A(x,y)=\ip{a_x}{a_y}$, and $B(x,y)=\ip{b_x}{b_y}$.
A cut with $j$ input source legs factors as $\Flat_SK=UV^*$, with columns
$u_x=\eta_x^{\otimes(m-j)}\otimes a_x$ and
$d_x=\overline{\eta_x^{\otimes j}}\otimes b_x$.
Consequently
\begin{equation}\label{eq:covgram}
 \norm K_2^2=\int_{X^2}\gamma^m A\bar B\,d\mu^2,
 \qquad \ip{u_x}{u_y}=\gamma^{m-j}A,
 \qquad \ip{d_x}{d_y}=\bar\gamma^{\,j}B.
\end{equation}
Suppose $\int|\gamma|^m|A||B|\le H_0^2$ and the nonnegative kernels obey two-sided Schur bounds
\[
 \norm{|\gamma|^j|A|}_{\rm Schur}\le L_j^2,
 \qquad \norm{|\gamma|^j|B|}_{\rm Schur}\le R_j^2.
\]
The Schur test on the two column Grams gives $\norm{UV^*}\le L_{m-j}R_j$. Interpolation with the Hilbert--Schmidt bound yields
\begin{equation}\label{eq:covcut}
 \norm{\Flat_SK}_s\le H_0^{2/s}(L_{m-j}R_j)^{1-2/s},\qquad s\ge2.
\end{equation}

For two source kernels with the same $a_x,b_x$, let
\[
 \gamma_{n,n'}(x,y)=\ip{\eta_{n,x}}{\eta_{n',y}},\qquad
 D_{n,n'}=\gamma_{n,n}^m+\gamma_{n',n'}^m-\gamma_{n,n'}^m-\gamma_{n',n}^m.
\]
The four terms are exactly the paired expansion of the kernel difference. If $\varepsilon_{n,n'}=\int|D_{n,n'}||A||B|$, then
\begin{equation}\label{eq:covcutdiff}
 \norm{\Flat_S(K_n-K_{n'})}_s
 \le\varepsilon_{n,n'}^{1/s}
 (L_{n,m-j}R_{n,j}+L_{n',m-j}R_{n',j})^{1-2/s}.
\end{equation}
Mixed covariances retain the specified joint source realization. For each finite $s$, this estimate can give Schatten convergence, which also implies operator-norm convergence. The direct $s=\infty$ version loses the small factor $\varepsilon_{n,n'}^{1/s}$ and therefore supplies no convergence from the operator endpoint bound alone.

A change of source is governed by relative covariance. If $S:\mathcal A\to H$ is bounded, $Q=S^*S$, and $U$ is unitary on $\mathcal A$, then there is a bounded $V$ on $\overline{\ran S}$ satisfying $SU=VS$ if and only if
\begin{equation}\label{eq:relativecovariance}
 U^*QU\le c^2Q\quad\text{for some }c<\infty,
 \qquad \inf c=\norm V.
\end{equation}
Define $V(Sx)=SUx$ on $\ran S$; the quadratic inequality gives well-definedness, boundedness, and the exact norm. Applying this map on a source leg costs $\norm V$ in every directed cut; multiple legs multiply the costs. For real Gaussian evaluation the maps preserve the real source structure. Circular Gaussian contraction comparisons can also be obtained by adding independent complementary covariance and applying conditional Jensen.

\subsection{Uniform spectral localization on closed manifolds}
Let $(M,g)$ be a smooth compact boundaryless $d$-manifold. Set
\[
 \Lambda=(1-\Delta_g)^{1/2},\qquad \ell_R=\log(e+R),\qquad
 \mathscr L=\log(e+\Lambda).
\]
Fix a real eigenbasis and a smooth real dyadic partition of unity $\Delta_L=\rho_L(\Lambda)$, $L\in\{1,2,4,\ldots\}$, with fixed enlarged shell supports, finite overlap, and
$\sup_{L,\lambda}L^j|\partial_\lambda^j\rho_L(\lambda)|<\infty$ for each $j$.

\begin{lemma}[A uniform kernel estimate]\label{lem:spectralkernel}
Let $\beta_R$ be smooth and supported in $[cR,CR]$, $R\ge2$, with fixed $0<c<C$. For $J\ge0$ and differential operators $D,E$ of orders $q_x,q_y$, there is a finite integer $K$ such that the kernel of $\beta_R(\Lambda)$ satisfies
\begin{equation}\label{eq:spectralkernel}
 |D_xE_yK_R(x,y)|\le C_{J,D,E}
 \left(\max_{j\le K}\sup_\lambda R^j|\beta_R^{(j)}(\lambda)|\right)
 R^{d+q_x+q_y}(1+Rd_g(x,y))^{-J}.
\end{equation}
The corresponding two-sided Schur estimates hold. If the $y$-spectrum of $F$ lies in an enlarged $R$-shell, then
$\norm{E_yF}_{L^2(M^2)}\le CR^{q_y}\norm F_{L^2(M^2)}$.
\end{lemma}
\begin{proof}
We keep track of the multiplier dependence in the local semiclassical construction of \cite[Proposition~2.2(i), equation~(2.12), and Remark~2.3]{ORT}.

\emph{Fixed energy support.} Put $h=R^{-1}$, $P_h=-h^2\Delta_g$, and $b_h(t)=\beta_R(Rt)$. For sufficiently small $h$, choose a fixed $\omega\in C_c^\infty((0,\infty))$ equal to one on $[c^2/2,2C^2]$, and define
\[
 \theta_h(s)=\omega(s)b_h(\sqrt{h^2+s})\quad(s>0),\qquad \theta_h(s)=0\quad(s\le0).
\]
The spectral theorem gives $\theta_h(P_h)=\beta_R(\Lambda)$. The $\theta_h$ have a common compact positive energy support, and every fixed finite set of their derivatives is bounded by a finite rescaled seminorm of $\beta_R$. The remaining bounded interval of $R$ contains only finitely many possible eigenvalues in the multiplier support and is handled directly.

\emph{Symbols and uniform remainders.} In a fixed finite atlas with fixed left and right cutoffs, expand to order $A$. The coefficients have the form
\[
 a_\ell[\theta](z,\xi)=\sum_jc_{\ell,j}(z,\xi)\theta^{(j)}(p(z,\xi)),
 \qquad p(z,\xi)=g^{ab}(z)\xi_a\xi_b.
\]
The resolvent-symbol recursion starts with $q_0=(w-p)^{-1}$ and has the form
\[
 q_n=-(w-p)^{-1}\!
 \sum_{j+k+|\alpha|=n,\ k<n}
 c_\alpha\partial_\xi^\alpha d_j\,\partial_z^\alpha q_k.
\]
Each term is a finite sum of resolvent powers. The Helffer--Sj\"ostrand formula turns these powers into derivatives of $\theta$, yielding the displayed finite derivative dependence. The coefficients and any fixed set of their derivatives therefore obey a finite seminorm bound.

For fixed $\theta$, the remainder satisfies, for $A\ge2S$,
\[
 \sup_{0<h\le1}h^{-(A-2S)}
 \norm{\mathcal R_{A,h}[\theta]}_{H^{-S}\to H^S}<\infty.
\]
It is linear in $\theta$, and for each fixed $h$ is continuous on the Fr\'echet space of smooth functions with this common compact support: the spectral multiplier is finite-rank and the fixed finite symbol sum is continuous. Banach--Steinhaus applied to the normalized remainder operators makes the displayed supremum bounded by one finite $\theta$ seminorm. This supplies the uniform estimate for the moving family $\theta_h$.

\emph{Kernel bounds.} A local symbol term has kernel
\[
 (2\pi h)^{-d}h^\ell\int e^{i(z-y)\cdot\xi/h}A_{\ell,h}(z,y,\xi)\,d\xi.
\]
Differentiation costs at most $h^{-q_x-q_y}$. Repeated integration by parts in $\xi$ gives
$h^{-d-q_x-q_y}(1+|z-y|/h)^{-J}$,
with a finite seminorm of the amplitude. Distances in the fixed charts are uniformly comparable with geodesic distance on their supports.

Choose $S>d/2+\max(q_x,q_y)$ for the remainder. Differentiated point evaluations are uniformly bounded in $H^{-S}$, and Sobolev embedding on the output gives
\[
 \sup_{x,y}|D_xE_yK_{\mathcal R}(x,y)|
 \le C\norm{\mathcal R}_{H^{-S}\to H^S}.
\]
Taking $A\ge2S+J+d+q_x+q_y+2$ makes this remainder smaller than the required bound even away from the diagonal. A maximum over the finite atlas gives a single finite derivative order $K$, proving \eqref{eq:spectralkernel}.

Finally, integration over the $R^{-1}$ ball and its dyadic annuli gives, for $J>d$,
\[
 \sup_x\int R^d(1+Rd_g(x,y))^{-J}\,dy
 \le C\sum_{k\ge0}2^{(d-J)k}<\infty,
\]
and likewise with $x,y$ interchanged. This proves the Schur bounds. The Bernstein estimate follows from the bounded map $E:H^{q_y}\to L^2$ and the spectral theorem on the enlarged shell.
\end{proof}

Enlarged shells have rank $O(L^d)$, by applying the diagonal kernel bound to a nonnegative smooth multiplier equal to one on the shell. The kernel $\Gamma_{L,a}$ of $\Delta_L^2\Lambda^{-2a}$ therefore satisfies
\begin{equation}\label{eq:columnkernel}
 |\Gamma_{L,a}(x,y)|\lesssim L^{d-2a}(1+Ld_g(x,y))^{-J},\quad
 \norm{\Gamma_{L,a}}_{\rm Schur}\lesssim L^{-2a},\quad
 \norm{\Gamma_{L,a}}_{L^2(M^2)}\lesssim L^{d/2-2a}.
\end{equation}
The last estimate also follows from the rank bound and the spectral size on the shell.

\subsection{Covariance profiles and critical summability}
In a common real Gaussian source, let $\eta_{n,x}$ be smooth finite-dimensional source vectors, and set
$X_n(x)=W(\eta_{n,x})$ and $\Phi_{m,n}(x)=I_m(\eta_{n,x}^{\otimes m})$.
Assume that a nondecreasing $\mathfrak c(R)\ge1$ and an exponent $J>d$ satisfy, for every $0\le j\le m$ and all $n,n',R$,
\begin{equation}\label{eq:covprofileassumption}
 \max\left\{\sup_x\int_M,\ \sup_y\int_M\right\}
 |\gamma_{n,n'}(x,y)|^jR^d(1+Rd_g(x,y))^{-J}
 \le C_j\mathfrak c(R)^j.
\end{equation}
Each integral is in the other spatial variable. Assume also the common-source ultraviolet defect
\begin{equation}\label{eq:UVdefect}
 \varepsilon_{n,n'}^{(m)}=
 \norm{\gamma_{n,n}^m+\gamma_{n',n'}^m-\gamma_{n,n'}^m-\gamma_{n',n}^m}_{L^1(M^2)}
 \longrightarrow0
\end{equation}
as $n,n'\to\infty$.

Use the columns of $\Lambda^{-a}\Delta_L$ and $\Delta_Q\Lambda^{-b}$ to form
$K_{n;L,Q}=\int\eta_{n,x}^{\otimes m}\otimes\bar b_{Q,x}\otimes a_{L,x}\,dx$.
Its exact evaluation is
\[
 I_mK_{n;L,Q}=\Lambda^{-a}\Delta_LM_{\Phi_{m,n}}\Delta_Q\Lambda^{-b}.
\]
For $s\ge2$, \eqref{eq:covcut} and \eqref{eq:columnkernel} give
\begin{align}
 \max_S\norm{\Flat_SK_{n;L,Q}}_s
 &\lesssim\mathfrak c(L\vee Q)^{m/2}(L\wedge Q)^{d/s}L^{-a}Q^{-b},\label{eq:geomprofile}\\
 \max_S\norm{\Flat_S(K_{n;L,Q}-K_{n';L,Q})}_s
 &\lesssim(\varepsilon_{n,n'}^{(m)})^{1/s}
 \mathfrak c(L\vee Q)^{m(1/2-1/s)}L^{d/s-a}Q^{d/s-b}.\label{eq:geomdiff}
\end{align}
For example, when $L\ge Q$, the low-shell pointwise bound and the high-shell covariance profile give
\[
 \norm{K_{n;L,Q}}_2^2\lesssim\mathfrak c(L)^mQ^{d-2b}L^{-2a},\qquad
 \norm{K_{n;L,Q}-K_{n';L,Q}}_2^2
 \lesssim\varepsilon_{n,n'}^{(m)}L^{d-2a}Q^{d-2b}.
\]
The two Schur factors are bounded using $\mathfrak c(L\vee Q)$; interpolation proves the assertions.

Suppose $\mathfrak c(R)^{m/2}\lesssim R^\beta\ell_R^\kappa$ with $\beta,\kappa\ge0$. With outer logarithmic weights $\mathscr L^{-\alpha}$ and $\mathscr L^{-\zeta}$ the common majorant is
\begin{equation}\label{eq:geometricmajorant}
 \mathcal M_{L,Q}=(L\vee Q)^\beta\ell_{L\vee Q}^{\kappa}
 (L\wedge Q)^{d/s}L^{-a}Q^{-b}\ell_L^{-\alpha}\ell_Q^{-\zeta}.
\end{equation}
Each fixed dyadic gap is a matching, up to finitely many colors, so the appropriate assembly lattice is $\mathsf{Gap}_s=\ell^1_h(\ell^s_n)$.

\begin{theorem}[Wick operators in the original Schatten geometry]\label{thm:geometric}
Fix $1\le s<\infty$. Under \eqref{eq:covprofileassumption}--\eqref{eq:UVdefect} and the stated profile growth, either of the following conditions is sufficient:
\begin{align}
 a>\beta,\quad b>\beta,\quad a+b>\beta+d/s;
 &\hspace{8pt}\text{strict regime},\label{eq:strictface}\\
 a>\beta,\quad b>\beta,\quad a+b=\beta+d/s,\quad
 \alpha+\zeta>\kappa+1/s;
 &\hspace{8pt}\text{critical regime}.\label{eq:criticalface}
\end{align}
The local kernels assemble unconditionally in $\Wc_{m,s}$. Their evaluations are
\begin{equation}\label{eq:wickgeometry}
 T_n=\Lambda^{-a}\mathscr L^{-\alpha}
       M_{\Phi_{m,n}}\mathscr L^{-\zeta}\Lambda^{-b},
\end{equation}
and converge in every finite $L^p(\Sch_s)$ to a unique $T$. The limit is independent of the dyadic partition and of regularizations satisfying the joint covariance hypotheses.

If $\varepsilon_{N,N'}^{(m)}\lesssim N^{-\vartheta}$ for dyadic $N'\ge N$, the strict regime gives a positive power rate. In the critical regime,
\begin{equation}\label{eq:geomrate}
 \norm{T_N-T}_{L^p(\Sch_s)}\lesssim_{m,s}(p+s+2)^{m/2}\ell_N^{-\omega},
 \qquad \omega=\alpha+\zeta-\kappa-1/s>0,
\end{equation}
with the geometric parameters fixed. Under this quantitative ultraviolet hypothesis, both regimes also give almost-sure convergence along dyadic scales. Adjoints exchange $(a,\alpha)$ and $(b,\zeta)$.
\end{theorem}
\begin{proof}
For $s\ge2$, the majorant is \eqref{eq:geomprofile} with logarithmic weights. For $s<2$, each physical double block has rank at most $C(L\wedge Q)^d$. Multiply its Hilbert--Schmidt response estimate by the rank to the power $1/s-1/2$; this gives the same majorant. The low-exponent assembly argument above first constructs the random response, and the fixed-degree inverse then recovers its quotient-sum kernel.

To check summability write $L=2^{n+h}$, $Q=2^n$ on the upper half of the dyadic array. Up to fixed constants the majorant is
\begin{equation}\label{eq:gapsum}
 2^{-(a-\beta)h-(a+b-\beta-d/s)n}
 (1+n+h)^{\kappa-\alpha}(1+n)^{-\zeta},\qquad n,h\ge0.
\end{equation}
The lower half interchanges the two outer parameters. In the strict regime all exponential margins are positive and absorb every logarithmic factor. Thus the $\ell^1_h(\ell^s_n)$ norm is finite, with a high-scale tail of some order $R^{-\eta}$. In the critical regime the $n$ exponential disappears. The diagonal $h=0$ requires
\[
 \sum_{n\ge0}(1+n)^{s(\kappa-\alpha-\zeta)}<\infty,
\]
which is exactly the logarithmic condition. It suffices because
$(1+n+h)^t\le(1+h)^{t_+}(1+n)^t$
and the $h$ exponential remains summable. Splitting $L\vee Q>R$ into $h>\frac12\log_2R$ and its complement bounds the tail by $C\ell_R^{-\omega}$. The first part is exponentially small; in the second, $n>\frac12\log_2R-O(1)$. The diagonal gives the matching order for this majorant tail.

Order-continuous assembly now produces each regularized kernel. Fixed-block differences vanish by the common ultraviolet defect, and the same majorant controls all remaining blocks uniformly in $n,n'$. A finite-head--tail argument proves Cauchy convergence, in the maximum-cut norm at $s\ge2$ and by response inversion at $s<2$.

For a fixed $n$, the random multiplier $\Phi_{m,n}$ is smooth, and the two outer smoothing operators are bounded, so \eqref{eq:wickgeometry} is also a genuine bounded operator. Its finite spectral blocks coincide with those of the evaluated kernel. Almost-sure Schatten convergence of a subsequence of finite block sums, together with strong convergence of the spectral compressions of this bounded operator, identifies the two constructions. Injectivity of evaluation identifies their coefficients. Two dyadic partitions therefore produce the same regularized operator and kernel. For two regularization schemes, the joint mixed-covariance assumption applies to their interlaced family and identifies their limits as well.

For rates let $s_* =\max(s,2)$. Truncating at $L\vee Q\le R$ bounds the finite-head error by
$CN^{-\vartheta/s_*}R^A$
for some fixed $A>0$ absorbing shell counts and logarithms. In the strict regime balance this with $CR^{-\eta}$, for instance using $R=N^{\vartheta/(s_*(A+\eta))}$. In the critical regime choose $R=N^c$ with $0<c<\vartheta/(s_*A)$. The head decays as a power while the tail is $O(\ell_N^{-\omega})$. Fixed-degree hypercontractivity gives the displayed dependence on $p$.

At $N=2^j$, the critical error is $O(j^{-\omega})$. Choose a fixed $p>2/\omega$ and use Markov's inequality at threshold $j^{-\omega/2}$; the failure probabilities are summable. Borel--Cantelli proves the dyadic almost-sure convergence. The strict regime is easier. Adjoint identities hold on each finite cutoff and pass through the Schatten limit.
\end{proof}
The critical regime is nonempty only if $\beta<d/s$. Its logarithmic condition is necessary and sufficient for summability of the specified majorant, not a necessity assertion for every actual operator satisfying a possibly smaller covariance bound.

\subsection{The massive Gaussian free field}
Write $-\Delta_g\varphi_j=\lambda_j^2\varphi_j$ in the fixed real orthonormal basis. Choose nonnegative $\chi\in C_c^\infty([0,\infty))$ equal to one on $[0,1/4]$ and supported in $[0,1]$. In the same Gaussian realization define
\begin{equation}\label{eq:gff}
 X_N(x)=\sum_j\chi(N^{-2}\lambda_j^2)(1+\lambda_j^2)^{-1/2}g_j\varphi_j(x),
 \qquad \Phi_{m,N}={:X_N^m:}.
\end{equation}
Assume $m(d-2)<d$, and set
\[
 \beta_{d,m}=\tfrac m2(d-2)_+,\qquad
 \kappa_{d,m}=\tfrac m2\mathbf1_{\{d=2\}},\qquad
 \mathfrak w_d(R)=
 \begin{cases}1,&d=1,\\\ell_R,&d=2,\\R^{d-2},&d\ge3.\end{cases}
\]
\begin{proposition}[Joint free-field covariance bounds]\label{prop:gffcertificate}
Any fixed finite collection of cutoff profiles and all their dyadic scales, on the same source, satisfy \eqref{eq:covprofileassumption} with $\mathfrak c=\mathfrak w_d$. For every $0<\vartheta<\omega_{d,m}$,
\[
 \varepsilon_{N,N'}^{(m)}\lesssim\min(N,N')^{-\vartheta},\qquad
 \omega_{d,m}=
 \begin{cases}2,&d=1,2,\\d-m(d-2),&d\ge3.\end{cases}
\]
\end{proposition}
\begin{proof}
Decompose every self- and cross-covariance into dyadic spectral shells. The uniform kernel lemma gives
$|H_R(x,y)|\lesssim_A R^{d-2}(1+Rd_g(x,y))^{-A}$
for every fixed $A$, uniformly over the cutoff profiles. Summing shells gives a bounded envelope in dimension one, a logarithmic envelope in dimension two, and a distance-to-the-diagonal power $d_g(x,y)^{2-d}$ in higher dimension, with harmless bounded terms away from the diagonal. Integrating the $j$th powers against the localized kernel gives \eqref{eq:covprofileassumption}; near the diagonal the required integrability is $j(d-2)<d$.

The low-frequency plateau makes the difference from the full Green kernel consist only of shells $R\gtrsim N_*:=\min(N,N')$. Use
$|u^m-v^m|\le m|u-v|(|u|+|v|)^{m-1}$
and the same envelopes. The $L^1(M^2)$ contribution of a shell is bounded respectively by
$CR^{-2}$, $CR^{-2}\ell_R^{m-1}$, and $CR^{m(d-2)-d}$
in dimensions one, two, and at least three. Dyadic summation proves the asserted exponents, with a strict loss absorbing the two-dimensional logarithm. Mixed cutoff products have the same uniformly rescaled symbol bounds, so the estimates hold jointly for the entire finite profile family.
\end{proof}

Theorem~\ref{thm:geometric} now gives canonical fixed-degree Wick operators for every $1\le s<\infty$. Without logarithmic outer weights it suffices that
\begin{equation}\label{eq:gffcondition}
 a>\beta_{d,m},\qquad b>\beta_{d,m},\qquad a+b>\beta_{d,m}+d/s.
\end{equation}
On the critical face the logarithmic requirement is
$\alpha+\zeta>\kappa_{d,m}+1/s$.

There is also a common-source random distribution $\Phi_m$ such that, for every $u>\beta_{d,m}$, $\Phi_{m,N}\to\Phi_m$ in every finite $L^p(H^{-u}(M))$, with a positive power rate and dyadic almost-sure convergence. This follows directly from operator recovery. Choose $a_u>\beta_{d,m}$ and a sufficiently large finite $s$ with $a_u+u>\beta_{d,m}+d/s$. Define, for bounded $A$,
\[
 \langle\mathcal J_uA,f\rangle=\ip{A\Lambda^uf}{1},\qquad
 \norm{\mathcal J_uA}_{H^{-u}}\le\operatorname{Vol}(M)^{1/2}\norm A.
\]
At finite cutoff, $\mathcal J_u(T_N^{a_u,u})=\Phi_{m,N}$ because $\Lambda^{-a_u}1=1$. Operator convergence gives distributional convergence, and common cutoffs identify the limits for different $u$. For smooth $f,h$ the resulting operator satisfies
\begin{equation}\label{eq:distributionpair}
 \ip{T^{a,b}f}{h}
 =\left\langle\Phi_m,(\Lambda^{-b}f)\overline{\Lambda^{-a}h}\right\rangle.
\end{equation}
The remainder of this section uses these same-source operators as generators in the complete reconstruction theorem.

\subsection{Fixed-degree generators in the full algebra}
For $K\in\Wc_{m,s}$, regarded as a sequence supported only in degree $m$,
\begin{equation}\label{eq:fixeddegree_allradii}
 \Nn\rho K=\rho^m\sqrt{m!}\norm K_{\Wc_{m,s}}<\infty
 \qquad(\rho\ge1).
\end{equation}
Thus every fixed-degree geometric kernel already belongs to $\Alg_s$. Finitely many such generators and all their finite ordinary products lie in this same algebra. A further infinite sum over generator degrees would require a uniform weighted summability condition.

Fix $1\le s<\infty$, a closed manifold, and one real Gaussian realization. For $j=1,\ldots,d_0$, choose a degree $m_j\ge1$ and parameters $(a_j,b_j,\alpha_j,\zeta_j)$ satisfying Theorem~\ref{thm:geometric}. Write
\begin{equation}\label{eq:wick_family}
 T_{j,n}=\Lambda^{-a_j}\mathscr L^{-\alpha_j}
 M_{\Phi_{m_j,n}}\mathscr L^{-\zeta_j}\Lambda^{-b_j},
 \qquad T_j=\lim_nT_{j,n}=I_{m_j}K_j.
\end{equation}
Assume the geometric bounds provide known $B_j$ and $r_j(n)\downarrow0$ with
\[
 \norm{T_{j,n}-T_j}_{L^2(\Sch_s)}\le r_j(n),\qquad
 \sup_n\norm{T_{j,n}}_{L^2(\Sch_s)}\le B_j.
\]
The same $B_j$ bounds the limit. Each regularization has finite source support; let $L_n$ contain all source coordinates used by this finite family.

\subsection{From smooth shells to finite physical projections}
Define the nested finite-rank orthogonal spectral projections
\[
 \Pi_K=\mathbf1_{[1,K]}(\Lambda),\qquad d_K=\rank\Pi_K.
\]
Let $C_*$ be an enlarged dyadic support constant, so that the spectrum of $\Delta_L$ is at most $C_*L$. For the majorant $\mathcal M^j$ of the $j$th generator, absorb its fixed-degree evaluation and assembly constants in $C_j$ and set
\begin{equation}\label{eq:geometric_spatial_certificate}
 b_j(K)=2C_j\left\|\mathcal M^j_{L,Q}
             \mathbf1_{\{L\vee Q>K/C_*\}}\right\|_{\mathsf{Gap}_s}.
\end{equation}
\begin{lemma}[A geometric certificate for physical truncation]\label{lem:geometric_spatial_tail}
Uniformly over all regularizations and for the limit,
\begin{equation}\label{eq:physical_projection_tail}
 \norm{T_{j,n}-\Pi_KT_{j,n}\Pi_K}_{L^2(\Sch_s)}\le b_j(K),\qquad
 \norm{T_j-\Pi_KT_j\Pi_K}_{L^2(\Sch_s)}\le b_j(K).
\end{equation}
In the strict regime $b_j(K)\lesssim K^{-\eta_j}$ for some $\eta_j>0$; in the critical regime $b_j(K)\lesssim\ell_K^{-\omega_j}$, with $\omega_j=\alpha_j+\zeta_j-\kappa_j-1/s>0$.
\end{lemma}
\begin{proof}
Use the unconditional dyadic double-block expansion of $T_{j,n}$. Every block with $L\vee Q\le K/C_*$ has both outer spectral supports in $\Pi_K$, and hence is annihilated by $A\mapsto A-\Pi_KA\Pi_K$. The remaining block sum has $L^2(\Sch_s)$ norm at most $C_j$ times its $\mathsf{Gap}_s$ tail. The norm of the latter truncation-difference map is at most two, proving the uniform estimate. Pass to the limit in $n$ and use the tail estimates following \eqref{eq:gapsum}.
\end{proof}
Smooth shells supply analytic estimates; the sharp projections supply nested finite observation spaces. Lemma~\ref{lem:geometric_spatial_tail} connects these two roles without identifying the two types of cutoff.

\subsection{Input recovery and its three error terms}
Choose a faithful spectral anchor $Re_k=r_ke_k$, $\norm R_2=1$, in the same eigenbasis as $\Lambda$. Observe the finite same-sample Hermite correlations
\begin{equation}\label{eq:wick_pipeline_marks}
 z^j_{\alpha;k\ell}(n)=\E[\Psi_\alpha(g)(RT_{j,n}R)_{k\ell}],
 \quad |\alpha|=m_j,\quad\operatorname{supp}\alpha\subset[L_n],
 \quad e_k,e_\ell\in\Pi_KH.
\end{equation}
Let $\delta_j$ bound the Euclidean error of this entire finite vector. It may be supplied by the observation protocol or certified by independent sampling through Proposition~\ref{prop:empiricalmarks}. Put
\[
 \gamma_s=(1/s-1/2)_+,\qquad
 \kappa_K=\left(\min_{e_k\in\Pi_KH}r_k\right)^{-2},\qquad
 n_j(K)=d_K^{\gamma_s}\kappa_K\delta_j.
\]
Divide each retained coordinate by $r_kr_\ell\sqrt{m_j!}$ and set all unobserved coefficients to zero. This gives a finite kernel $\widehat K_j$ and response $\widehat T_j=I_{m_j}\widehat K_j$. Define
\begin{equation}\label{eq:wick_input_certificates}
 e_j=v_s^{m_j}\{r_j(n)+b_j(K)\}+n_j(K),\qquad
 \epsilon_j=r_j(n)+b_j(K)+u_s^{m_j}n_j(K).
\end{equation}
\begin{proposition}[All-radius input recovery]\label{prop:wick_pipeline_input}
The actual finite outputs satisfy
\begin{equation}\label{eq:wick_pipeline_input}
 \Nn\rho{\widehat K_j-K_j}\le\rho^{m_j}e_j\quad(\rho\ge1),\qquad
 \norm{\widehat T_j-T_j}_{L^2(\Sch_s)}\le\epsilon_j.
\end{equation}
If $r_j(n)+b_j(K)+d_K^{\gamma_s}\kappa_K\delta_j\to0$ for every $j$, all inputs converge jointly in the original all-radius topology.
\end{proposition}
\begin{proof}
Let $K_{j,n}$ be the kernel of $T_{j,n}$ and let $K_{j,n,K}$ be its sharp physical compression to $\Pi_K$. Source truncation is exact for this regularized family, by the definition of $L_n$. The exact finite inverse of the population marks therefore equals $K_{j,n,K}$.

The bias stays in degree $m_j$, and the two analytic approximation bounds give
\[
 \norm{I_{m_j}K_{j,n,K}-T_j}_{L^2(\Sch_s)}
 \le r_j(n)+b_j(K).
\]
The inverse Wick estimate bounds its normalized coefficient norm by $v_s^{m_j}(r_j(n)+b_j(K))$. Separately, Theorem~\ref{thm:finitedecoder} bounds the normalized coefficient norm of $\widehat K_j-K_{j,n,K}$ by $n_j(K)$. Forward evaluation bounds the latter noise response by $u_s^{m_j}n_j(K)$. Adding these errors proves both inequalities in \eqref{eq:wick_pipeline_input}. Since each difference has one fixed source degree, multiplication by $\rho^{m_j}$ gives every coefficient radius at once. The stated vanishing conditions consequently give joint convergence of the finite family in the original all-radius topology.
\end{proof}
The exact observations of the limit form an admissible array. The noisy regularized outputs themselves need not satisfy exact cross-resolution compatibility: the displayed bounds make them converge to that same completed object.

\subsection{Polynomial targets and their evolution}
For a nonempty $*$-word $w=i_1^{\varepsilon_1}\cdots i_k^{\varepsilon_k}$, where $\varepsilon_r\in\{1,*\}$, let
\[
 d(w)=\sum_{r=1}^km_{i_r},\qquad h(w)=(2k-1)^{d(w)/2}.
\]
Its computable error budget is
\begin{equation}\label{eq:wick_word_budget}
 E_w=h(w)\sum_{r=1}^k\epsilon_{i_r}
       \prod_{b<r}(B_{i_b}+\epsilon_{i_b})\prod_{b>r}B_{i_b}.
\end{equation}
For a fixed noncommutative $*$-polynomial $P=\sum_{w\ne\varnothing}a_ww$, set
$D_P=\max_{a_w\ne0}d(w)$ and $E_P=\sum_w|a_w|E_w$.
We state the Schatten conclusions for $P(0)=0$; for a nonzero constant term, subtract the same known scalar identity from both targets.

\begin{theorem}[Structural reconstruction of Wick operator systems]\label{thm:wick_pipeline}
There are unique kernels $K_P,\widehat K_P\in\Alg_s$, defined by the original ordered contraction product, such that
\begin{align}
 I(K_P)&=P(T_1,\ldots,T_{d_0}),&
 I(\widehat K_P)&=P(\widehat T_1,\ldots,\widehat T_{d_0}),\label{eq:wick_pipeline_identity}\\
 \norm{P(\widehat T)-P(T)}_{L^2(\Sch_s)}&\le E_P,&
 \Nn\rho{\widehat K_P-K_P}&\le C_{v_s\rho}\chi_{v_s\rho}^{D_P}E_P.
 \label{eq:wick_pipeline_output}
\end{align}
These estimates recover every fixed word, Lie expression, and PBW polynomial, including all lower-degree contraction terms. For any specified bounded deterministic drivers as in Section~\ref{sec:time},
\begin{align}
 \sup_{t\ge0}\norm{\Phi_t(P(\widehat T))-\Phi_t(P(T))}_{L^2(\Sch_s)}
 &\le E_P,\label{eq:wick_pipeline_time}\\
 \sup_{t\ge0}\Nn\rho{\widetilde\Phi_t(\widehat K_P-K_P)}
 &\le C_{v_s\rho}\chi_{v_s\rho}^{D_P}E_P.
 \label{eq:wick_pipeline_time_native}
\end{align}
As the input certificates tend to zero, all these targets converge to the operations and evolutions of the same regularization-independent objects.
\end{theorem}
\begin{proof}
Every fixed-degree input belongs to $\Alg_s$ by \eqref{eq:fixeddegree_allradii}. Theorem~\ref{thm:nativeproduct} therefore forms the actual kernels $K_P$ and $\widehat K_P$ by ordered contraction and identifies their evaluations with \eqref{eq:wick_pipeline_identity}. Their chaos degrees are at most $D_P$. This establishes existence of the output before estimating it; uniqueness follows from Gaussian coefficient injectivity.

For a word $w$, replace its true factors one at a time by the reconstructed factors, preserving order. Each difference term has one input error, reconstructed inputs to its left, and true inputs to its right. Schatten ideal inequalities and probability H\"older require $L^{2k}$ for the $k$ factors. A degree-$m_i$ input or difference contributes $(2k-1)^{m_i/2}$ by hypercontractivity, so the product of these costs is $h(w)$. The remaining size and error factors are exactly those in \eqref{eq:wick_word_budget}. Summing with the polynomial's absolute coefficients gives $E_P$. This argument uses the common source and requires no independence between generators.

To return the output to its original coefficient norm, write
$F=P(\widehat T)-P(T)=\sum_{j=0}^{D_P}F_j$. Degreewise inversion followed by Mehler extraction yields
\[
 \Nn\rho{\widehat K_P-K_P}
 \le\sum_{j=0}^{D_P}(v_s\rho)^j\norm{F_j}_{L^2(\Sch_s)}
 \le C_{v_s\rho}\chi_{v_s\rho}^{D_P}\norm F_{L^2(\Sch_s)}.
\]
This proves \eqref{eq:wick_pipeline_output}, retaining all lower-degree contractions as well as the highest degree. Contraction of $\Phi_t$ in $\Sch_s$ and of its coefficient lift at every radius proves the two time-uniform bounds. The common-noise replica identity identifies the transported target as $\Phi_t(P(T))$. Finally these continuous constructions act on the uniquely defined input limits, so they preserve regularization independence.
\end{proof}

Choose actual budgets $(n_l,K_l,J_l)$ with deterministic input error bounds tending to zero, and arrange sampling failure probability at most $2^{-l}$ at stage $l$ using Proposition~\ref{prop:empiricalmarks}. Borel--Cantelli then supplies one event of sampling probability one on which the certificates eventually hold. On this event the kernels converge at every radius and all fixed finite polynomial targets converge. Independence is required within each sample average; independence between different budget stages is unnecessary.

\subsection{Countable dictionaries}
We next retain a countable family of polynomial targets and its correlation kernel. Enumerate homogeneous $*$-polynomials of positive word degree as $p_l$, with degrees $k_l$ and absolute coefficient sums $a_l$. Put
\[
 m_* =\max_jm_j,\qquad H_k=(2k-1)^{m_*k/2},\qquad
 c_l=\frac{2^{-l}}{\max(1,a_l)k_l!H_{k_l}},
\]
and choose
$B\ge\max(1,\max_jB_j,\max_j(B_j+\epsilon_j))$,
$\Delta=\max_j\epsilon_j$.
A fixed $B$ suffices once all $\epsilon_j\le1$. Let $K_{p_l}$ be the corresponding coefficient kernel.
\begin{corollary}[A common certificate for the full dictionary]\label{cor:wick_dictionary}
Set $A_\rho=\chi_{v_s\rho}^{m_*}$. For $\rho\ge1$,
\begin{align}
 \sum_{l\ge1}c_l\Nn\rho{\widehat K_{p_l}-K_{p_l}}
 &\le C_{v_s\rho}A_\rho e^{A_\rho B}\Delta,\label{eq:wick_native_dictionary}\\
 \sum_{l>J}c_l\Nn\rho{K_{p_l}}
 &\le C_{v_s\rho}2^{-J}e^{A_\rho B}.\label{eq:wick_native_dictionary_tail}
\end{align}
With $p_0=I,c_0=1$, the actual averaged correlation kernel
\[
 \mathsf G_t(T)_{il}=c_ic_lR\Phi_t(p_i(T)^*p_l(T))R
\]
is a positive trace-class random operator and satisfies
\begin{equation}\label{eq:wick_actual_gram}
 \sup_{t\ge0}\norm{\mathsf G_t(\widehat T)-\mathsf G_t(T)}_{L^1(\Sch_1)}
 \le2(1+e^B)e^B\Delta.
\end{equation}
\end{corollary}
\begin{proof}
Proposition~\ref{prop:weightedwords} bounds the $l$th response difference by
$a_lk_lH_{k_l}B^{k_l-1}\Delta$.
Its chaos degree is at most $m_*k_l$. Applying the output extraction step and multiplying by $c_l$ bounds its coefficient contribution by
\[
 C_{v_s\rho}2^{-l}A_\rho
       \frac{(A_\rho B)^{k_l-1}}{(k_l-1)!}\Delta.
\]
Each factorial term is bounded by the full exponential series, so summation proves \eqref{eq:wick_native_dictionary}. Without taking a difference, the bound is
$C_{v_s\rho}2^{-l}(A_\rho B)^{k_l}/k_l!$,
which gives \eqref{eq:wick_native_dictionary_tail}.

Complete positivity makes every finite correlation block positive. For the trace estimate use $\norm{R\Phi_t(A)R}_1\le\norm A$. Expand each product difference into two terms and apply probability Cauchy--Schwarz. The weighted sums of $L^2(B(H))$ norms of the responses and their differences are bounded by $1+e^B$ and $e^B\Delta$, respectively. Absolute summation over the two block indices gives \eqref{eq:wick_actual_gram} and proves trace-norm convergence of the positive block kernel. The same calculation bounds the dictionary truncation tail by $2(1+e^B)e^B2^{-J}$.
\end{proof}
The Gram matrix of the averaged response row omits its covariance. The kernel above retains it through the product inside $\Phi_t$.

\subsection{A critical two-dimensional example}
On a closed surface take $s=2$, a common massive free field, and Wick degrees $m_1=1$, $m_2=2$. Set
\begin{equation}\label{eq:surface_pair}
 T_{j,n}=\Lambda^{-1/2}\mathscr L^{-1}
          M_{{:X_n^{m_j}:}}\mathscr L^{-1}\Lambda^{-1/2}.
\end{equation}
Both lie on the physical critical face $a_j+b_j=1=d/s$, with total logarithmic weight two. Their margins are
\[
 \omega_1=2-\tfrac12-\tfrac12=1,\qquad
 \omega_2=2-1-\tfrac12=\tfrac12.
\]
Thus $r_j(n)+b_j(K)\le C_j(\ell_n^{-\omega_j}+\ell_K^{-\omega_j})$. Choose the spectral anchor $R=c\Lambda^{-r}$, $r>1$, with $\norm R_2=1$. Spectral counting makes it Hilbert--Schmidt, and $\kappa_K\le c^{-2}K^{2r}$. Since $u_2=v_2=1$,
\begin{equation}\label{eq:surface_total_error}
 e_j=\epsilon_j\le C_j(\ell_n^{-\omega_j}+\ell_K^{-\omega_j})
                         +c^{-2}K^{2r}\delta_j.
\end{equation}
The regularized source dimension is $L_n=O(n^2)$, the physical dimension is $d_K=O(K^2)$, and the number of complex marked coordinates per response is
\[
 d_K^2\binom{L_n+m_j-1}{m_j}=O(K^4n^{2m_j}).
\]
Proposition~\ref{prop:empiricalmarks} supplies a sample budget for these coordinates, while \eqref{eq:surface_total_error} supplies the input error at every coefficient radius.

For the commutator, the mixed-degree hypercontractive factor is $3^{(1+2)/2}$. The resulting certificate is
\begin{equation}\label{eq:surface_commutator_budget}
 \norm{[\widehat T_1,\widehat T_2]-[T_1,T_2]}_{L^2(\Sch_2)}
 \le2\,3^{3/2}(B_2\epsilon_1+B_1\epsilon_2+\epsilon_1\epsilon_2).
\end{equation}
The same bound holds for its specified common-noise averaged evolution. The target is the actual ordered commutator of the model, irrespective of whether it vanishes for a particular realization or parameter choice.

This example connects covariance geometry, finite observations, statistical error, a concrete decoder, original completion, and ordinary noncommutative operations. Every step uses the same Gaussian realization, and the limit is independent of any regularization scheme satisfying the common covariance certificates.

\subsection{Different targets and independent obstructions}
The relevant spaces can be summarized as follows.
\begin{center}\small
\begin{tabular}{@{}>{\raggedright\arraybackslash}p{.20\linewidth}>{\raggedright\arraybackslash}p{.34\linewidth}>{\raggedright\arraybackslash}p{.40\linewidth}@{}}\toprule
Target & Specified topology or closure & Recovered information\\\midrule
$\Alg_s$ & All-radius original cut norms & Same-source coefficients and continuous ordinary operations\\[3pt]
Bounded operator tuples & Strong-$*$ topology in a fixed Hilbert space & Fixed-frame operators and their finite words\\[3pt]
Full positive word observations & Pointwise weak-operator UCP topology, encoded in trace class & Positive maps; zero Schwarz defect selects representations\\[3pt]
Weighted linear targets & Completion in $\|W^{1/2}P_Hx\|$ & The directly recoverable quotient selected by $W\lesssim S_N$\\[3pt]
Skew spectral configurations & Subcritical weak or strong spectral topology; critical relative-tail topology & Nonzero block spectra in the stated certificate class\\[3pt]
Square-mass determinant data & Locally uniform entire-function closure & A spectrum and possible mass at zero, selected by $\tau=0$ for the true image\\[3pt]
Averaged spectral responses & Weighted Laplace-transform topology & The weighted mean spectral measure, not its random configuration law\\\bottomrule
\end{tabular}
\end{center}
The boundary-table construction in Appendix~\ref{sec:physicalboundary} is another response quotient, determined by all permitted continuation experiments. Comparisons between these targets require a map that retains the relevant information.

\begin{proposition}[Independent failures of converse convergence]\label{prop:completionseparation}
Operator-norm convergence need not imply original Schatten convergence; random $L^2$ convergence need not imply all-radius coefficient convergence; positive anchored convergence need not have a genuine strong-$*$ limit; and full unlabelled spectra or marginal laws need not determine a fixed frame or fixed-source coefficients.
\end{proposition}
\begin{proof}
For the first assertion, on $\ell^2$ let $P_n$ have rank $n$ and set $T_n=n^{-1/s}P_n$. Then $\norm{T_n}_{\op}\to0$ but $\norm{T_n}_s=1$. As degree-zero kernels these retain norm one at every coefficient radius.

For the second, let $P$ be rank one and
$F_n=2^{-n}H_n(g)P/\sqrt{n!}$.
Then $\norm{F_n}_{L^2(\Sch_2)}=2^{-n}\to0$, whereas its coefficient norm is $(\rho/2)^n$, unbounded at $\rho=3$.

For the third, $T_n=e_ne_1^*$ has vanishing first-row anchor blocks and a surviving square block $r_1^2e_1e_1^*$. A true limit would have to be zero by its first row, contradicting its nonzero diagonal energy. Theorem~\ref{thm:strongclosure} identifies the defect as output escape.

Finally $e_1e_1^*$ and $e_2e_2^*$ have the same nonzero spectrum but differ in the fixed frame. The responses $gP$ and $-gP$ have the same marginal law but opposite source correlations $\E[g(gP)]$ and $\E[g(-gP)]$.
\end{proof}

The forward relations remain continuous. Convergence in $\Alg_s$ gives response convergence in every finite probability exponent; deterministic Schatten convergence gives operator-norm and strong-$*$ convergence; and strong-$*$ convergence on a bounded tuple ball gives convergence of finite words and uniformly weighted Gram observations. Reverse recovery is governed by the corresponding certificate: ideal mass, source correlation, degree tails, output tightness, or a critical spectral tail. Section~\ref{sec:pfaffian} adds two spectral regimes: a weak-$\rho$ envelope controls every strictly subcritical target, while critical recovery requires a tail class, and a square-mass budget admits determinant limits with defect parameter $\tau$. Theorem~\ref{thm:invariant_uniform} identifies completions from two-sided uniform control; Theorem~\ref{thm:target_short} gives the direct bounded factorization for linear targets.


\appendix
\section{Sharp observation costs on full Fock space}\label{sec:fockcosts}
The multiplier criterion of Theorem~\ref{thm:multiplier} has an exact realization on full Fock space. Branching controls the existence of finite-mass faithful anchors, while equal-length word orthogonality gives explicit target costs and the minimal observation depth.

\subsection{The exact branching threshold}
On the full Fock space
\[
 \mathcal F_d=\bigoplus_{n\ge0}(\C^d)^{\otimes n},\qquad
 S_ie_w=e_{iw},\qquad \Omega=e_\varnothing,\qquad d\ge2,
\]
the left creation operators are isometries with orthogonal ranges. For words of equal length, $S_u^*S_v=\delta_{uv}I$; see \cite{Popescu} for weighted versions of these shifts.

\begin{theorem}[Faithful finite-mass anchors]\label{thm:branch_threshold}
Given $c_i>0$, there exists an injective positive trace-one $Q$ with
$\vartheta_{Q^{1/2}}(S_i)\le c_i$ for all $i$ if and only if
\begin{equation}\label{eq:branch_threshold}
 s_c:=\sum_{i=1}^dc_i^{-2}<1.
\end{equation}
If in addition $\ip{Q\Omega}{\Omega}\ge a$, where $0<a<1$, the condition is $s_c\le1-a$. Both conclusions are realized by
\begin{equation}\label{eq:branch_anchor}
 Q_ce_w=(1-s_c)c_w^{-2}e_w,\qquad
 c_w=\prod_{j=1}^{|w|}c_{i_j},\qquad c_\varnothing=1.
\end{equation}
For $R=Q_c^{1/2}$,
\begin{equation}\label{eq:branch_factors}
 \Theta_R(S_i)=c_iS_i,\qquad
 \Theta_R(S_i^*)=c_i^{-1}S_i^*,\qquad \Lambda_R(S_i)=c_i.
\end{equation}
\end{theorem}
\begin{proof}
The necessity argument allows arbitrary, possibly nondiagonal, anchors. Let $q_w=\ip{Qe_w}{e_w}$. Faithfulness gives $q_\varnothing>0$. Test $S_iQS_i^*\le c_i^2Q$ on $e_{iv}$ to get $q_{iv}\ge c_i^{-2}q_v$. Iteration and summation over the orthonormal word basis yield
\[
 1=\Tr Q\ge q_\varnothing\sum_{n\ge0}s_c^n.
\]
Thus $s_c<1$ and $q_\varnothing\le1-s_c$, proving both necessary conditions.

Under these conditions \eqref{eq:branch_anchor} has strictly positive diagonal entries, trace $(1-s_c)\sum_ns_c^n=1$, and vacuum mass $1-s_c$. Direct evaluation on word vectors proves the two factorization formulas in \eqref{eq:branch_factors}. Their norms give the exact costs. Since $s_c<1$ implies every $c_i>1$, the formula for $\Lambda_R$ follows.
\end{proof}
For equal costs the condition is $c>\sqrt d$, with an unattained infimum among faithful trace-one anchors. If the vacuum mass is at least $a$, the minimum common cost is the attained value $\sqrt{d/(1-a)}$.

\subsection{Target costs in word and PBW coordinates}
\begin{theorem}[Exact homogeneous target cost]\label{thm:fock_target_cost}
For the anchor \eqref{eq:branch_anchor} and a homogeneous analytic free polynomial
$p(x)=\sum_{|w|=n}a_wx_w$, $n\ge1$,
\begin{equation}\label{eq:fock_target_cost}
 \norm{p(S)}^2=\sum_{|w|=n}|a_w|^2,\qquad
 \vartheta_R(p(S))^2=\sum_{|w|=n}|a_w|^2c_w^2.
\end{equation}
If a degree-$n$ PBW dictionary has word-coordinate matrix $C$ and $p=\sum_\beta z_\beta p_\beta$, then
\begin{equation}\label{eq:pbw_resource_form}
 \vartheta_R(p(S))^2=z^*C^*D_cCz,\qquad
 D_c=\diag(c_w^2)_{|w|=n}.
\end{equation}
\end{theorem}
\begin{proof}
Equal-length orthogonality gives
$p(S)^*p(S)=(\sum_w|a_w|^2)I$.
The multiplier homomorphism satisfies
$\Theta_R(p(S))=\sum_wa_wc_wS_w$.
Apply the same orthogonality to this operator, and then substitute $a=Cz$.
\end{proof}

The analytic free algebra is faithfully represented because $p(S)\Omega=\sum_wa_we_w$. Its $*$-extension has the relations $S_i^*S_j=\delta_{ij}I$. Cross terms in \eqref{eq:pbw_resource_form} remain present. In word order $(xxxy,xxyx,xyxx,yxxx)$, for example,
\[
 b_1=\sym(x,x,[x,y])=\tfrac13(1,0,0,-1),\qquad
 b_3=[x,[x,[x,y]]]=(1,-3,3,-1).
\]
They have factor deficits one and three, yet word inner product $2/3$ and weighted cross term $2c_1^6c_2^2/3$.

\begin{corollary}[Nested Lie targets and anchor design]\label{cor:lie_resource_design}
For $h_n=\ad_{S_1}^{n-1}S_2$, $n\ge2$,
\begin{align}
 h_n&=\sum_{j=0}^{n-1}(-1)^j\binom{n-1}jS_1^{n-1-j}S_2S_1^j,\\
 \norm{h_n}&=\binom{2n-2}{n-1}^{1/2},\qquad
 \vartheta_R(h_n)=c_1^{n-1}c_2\binom{2n-2}{n-1}^{1/2}.
 \label{eq:lie_exact_cost}
\end{align}
For $d=2$ and vacuum mass at least $a$, the minimum cost of $h_n/\norm{h_n}$ within the geometric anchor family \eqref{eq:branch_anchor} is
\begin{equation}\label{eq:lie_design_optimum}
 (1-a)^{-n/2}\left(\frac{n^n}{(n-1)^{n-1}}\right)^{1/2}.
\end{equation}
It is attained at $c_1^{-2}=(1-a)(n-1)/n$ and $c_2^{-2}=(1-a)/n$.
\end{corollary}
\begin{proof}
The bracket recurrence gives the binomial expansion. Its distinct words all have weight $c_1^{n-1}c_2$, and
$\sum_j\binom{n-1}j^2=\binom{2n-2}{n-1}$.
For design set $u_i=c_i^{-2}$. Minimizing the cost is equivalent to maximizing $u_1^{n-1}u_2$ subject to $u_1+u_2\le1-a$. The maximum uses the full budget, and differentiation of the logarithm gives the stated allocation.
\end{proof}
The design problem here is restricted to the geometric family; the branching existence criterion itself covered all faithful anchors.

\subsection{Exact observation depth and noise risk}
Let $P_N$ project onto words of length at most $N$, and let $D_N=1+d+\cdots+d^N$. Fix $M,N\ge0$ and a nonzero homogeneous analytic polynomial $p$ of degree $n\ge1$. The unknown is $X$ and the known right target is $p(S)$. Define
\begin{equation}\label{eq:fock_protocol}
 O_{M,\ell}(X)=P_MXRP_\ell,\qquad
 L_{M,N,p}(X)=P_MXp(S)RP_N.
\end{equation}
The observation contains $D_MD_\ell$ complex coordinates.

\begin{theorem}[Depth and minimax risk]\label{thm:fock_depth_risk}
Over all bounded $X$, the observation $O_{M,\ell}(X)$ determines $L_{M,N,p}(X)$ if and only if
\begin{equation}\label{eq:fock_depth}
 \ell\ge N+n.
\end{equation}
Under this condition, for $Y=O_{M,\ell}(X)+E$, $E=P_MEP_\ell$, $\norm E_2\le\delta$, use $YC_pP_N$, where $C_p=\Theta_R(p(S))$. For $\delta>0$,
\begin{equation}\label{eq:fock_minimax}
 \inf_D\sup_{X,E}\norm{D(O_{M,\ell}(X)+E)-L_{M,N,p}(X)}_2
 =\delta\left(\sum_{|w|=n}|a_w|^2c_w^2\right)^{1/2}.
\end{equation}
Both lower bounds already hold for finite-rank unknowns, with no additional size prior on $X$.
\end{theorem}
\begin{proof}
The operator $C_p$ raises degree by exactly $n$, so $C_pP_N=P_{N+n}C_pP_N$. If \eqref{eq:fock_depth} holds, factorization gives
\[
 L_{M,N,p}(X)=P_MXR C_pP_N=O_{M,\ell}(X)C_pP_N.
\]
Write $\beta_p=\norm{C_p}>0$. Orthogonality gives $C_p^*C_p=\beta_p^2I$, and the error is at most $\delta\beta_p$.

If $\ell<N+n$, choose a word vector $v$ of length $N$, put $z=p(S)v\ne0$, and take a unit vector $u\in P_M\mathcal F_d$. Then $X=uz^*$ has zero observation, because $z$ has degree $N+n$, but $L_{M,N,p}(X)v\ne0$. Thus the target is not determined.

For the noise lower bound choose a word vector $v\in P_N\mathcal F_d$ and put $z=C_pv/\beta_p$. This is a finite-support unit vector, so $R^{-1}z$ is well defined. Let $X_0=u(R^{-1}z)^*$. Then
\[
 O_{M,\ell}(X_0)=uz^*,\qquad
 L_{M,N,p}(X_0)=\beta_puv^*,
\]
where the second equality follows from $C_p^*z=\beta_pv$. The two unknowns $\pm\delta X_0$ admit the same zero observed datum with admissible opposite errors. Their targets are $2\delta\beta_p$ apart, proving the minimax lower bound.
\end{proof}

For finite-rank $X$ this experiment lies in the original Gaussian model. With a specified real standard Gaussian $g$ and $T_X=gX$,
\[
 \E[gP_MT_XRP_\ell]=P_MXRP_\ell,
 \qquad \norm{gA}_{L^2(\Sch_2)}=\norm A_2.
\]
The same depth and risk therefore apply to the random target $gL_{M,N,p}(X)$. Together, \eqref{eq:branch_threshold}, \eqref{eq:fock_target_cost}, and \eqref{eq:fock_depth} describe the observation mass, target norm, and finite depth required by this model.


\section{Bounded protocols and joint laws}\label{sec:obsappendix}
The main finite decoder uses source--response correlations. Bounded protocols can instead identify a fixed-frame Borel law or a scalar spectral quotient. This appendix gives their inverse coordinates and specifies the extra source-dependence and tail conditions needed to return to fixed-source coefficient recovery. A marginal law and a joint source law remain distinct targets.

\subsection{Bounded positive coordinates}
For a positive trace-class operator $Q$ on a fixed separable Hilbert space, define on its exterior Fock space
\begin{equation}\label{eq:extstate}
 \Gamma_\wedge(Q)=\bigoplus_{a\ge0}\wedge^aQ,\qquad
 \rho(Q)=\frac{\Gamma_\wedge(Q)}{\det_F(I+Q)}.
\end{equation}
The exterior-power expansion gives $\Tr\Gamma_\wedge(Q)=\det_F(I+Q)$, so $\rho(Q)$ is a positive trace-one operator \cite{Simon,Bornemann}. Let $P_0$ and $P_1$ denote the vacuum and one-particle sector projections.

\begin{proposition}[Bounded coordinates and vacuum inversion]\label{prop:boundedmarks}
One has
\begin{equation}\label{eq:vacuum-inverse}
 p_0(Q)=\Tr(P_0\rho(Q))=\det_F(I+Q)^{-1}>0,\qquad
 Q=\frac{P_1\rho(Q)P_1}{p_0(Q)}.
\end{equation}
For $0\le B\le I$ put
\begin{equation}\label{eq:detmarks}
 D_Q(B)=\Tr(\rho(Q)\Gamma_\wedge(B))
       =\frac{\det_F(I+QB)}{\det_F(I+Q)}\in(0,1].
\end{equation}
The values $D_Q(0)$ and $D_Q(P_v)$ on a countable dense set of unit vectors determine $Q$. For positive trace-class $Q,S$,
\begin{align}
 \norm{\rho(Q)-\rho(S)}_1&\le2\norm{Q-S}_1,\label{eq:boundedstateforward}\\
 \norm{Q-S}_1&\le e^L(1+L)\norm{\rho(Q)-\rho(S)}_1
 \quad\text{if }\Tr Q,\Tr S\le L.\label{eq:boundedstateinverse}
\end{align}
\end{proposition}
\begin{proof}
The zeroth and first exterior powers are $1$ and $Q$, proving \eqref{eq:vacuum-inverse}. Exterior powers preserve multiplication, and $\Gamma_\wedge(B)$ is a positive contraction. The Fredholm expansion then gives \eqref{eq:detmarks}. For a rank-one projection all higher exterior powers vanish, so
\[
 D_Q(P_v)=p_0(Q)(1+\ip{Qv}{v}).
\]
Density and complex polarization determine all matrix entries.

To prove the estimates, first work in finite dimension along $Q_t=(1-t)Q+tS$ and put $E=S-Q$. In a positive direction $E\ge0$, the derivative of each exterior sector is the antisymmetric compression of a sum of positive tensor products, so $D\Gamma_\wedge(Q_t)[E]\ge0$. For a general self-adjoint $E$, use its positive and negative parts:
\[
 \norm{D\Gamma_\wedge(Q_t)[E]}_1
 \le\Tr D\Gamma_\wedge(Q_t)[|E|]
 =\det(I+Q_t)\Tr((I+Q_t)^{-1}|E|)
 \le\det(I+Q_t)\norm E_1.
\]
Differentiating the normalization contributes at most another $\norm E_1$. Integration gives \eqref{eq:boundedstateforward}; common finite-rank compression and the Fredholm exterior expansion pass it to trace-class operators.

If both traces are at most $L$, then $p_0(Q),p_0(S)\ge e^{-L}$. Write the one-particle blocks as $A,C$ and the vacuum probabilities as $p,q$. The identity
$Q-S=(A-C)/p+(q-p)S/p$
and contractivity of the two block readouts prove \eqref{eq:boundedstateinverse}.
\end{proof}

The one-particle block here is a sector compression. A simpler bounded encoding, $(1\oplus Q)/(1+\Tr Q)$, also admits a vacuum inverse; exterior powers additionally organize spectral transforms.

For a bounded tuple $T$, form the row from the identity, its generators, and adjoints as in \eqref{eq:QH}. The first-row blocks of $Q_R(T)$ are $RT_jR$, giving
\[
 (T_j)_{k\ell}=\frac{(Q_R(T)_{0j})_{k\ell}}{r_kr_\ell}.
\]
The matrix-entry Borel structure on $B(H)^d$ is standard: on each bounded ball it is the strong-Borel structure, and these balls form a countable exhaustion. The observation is strongly-$*$ continuous on each ball and has the displayed Borel inverse on its image. Consequently the bounded coordinates above form a Borel injection for tuples which are bounded at each sample, without a common deterministic bound. Finite recovery from the normalized state retains both the vacuum-inversion cost and the finite anchor cost.

\subsection{Energy probes and output tails}\label{subsec:energyprobes}
Suppose the protocol permits applying the generators and adjoints to the same sample, taking finite coherent sums, and measuring output norms. For finite-support $h=(h_0,\ldots,h_b)$ write
\[
 Y_T(h)=Rh_0+\sum_{a=1}^bA_aRh_a.
\]
Let $z$ be an auxiliary standard circular scalar Gaussian, independent of the original sample. Since $|z|^2$ is exponential of mean one,
\begin{equation}\label{eq:bounded_energy_probe}
 U_h(Q(T)):=\E_z e^{-\norm{Y_T(zh)}^2}
 =\frac1{1+\ip{Q(T)h}h}\in(0,1].
\end{equation}
Basis vectors and their real and imaginary two-vector sums recover each matrix entry by inversion and polarization. On a fixed operator-norm ball, each fixed probe has a uniform positive lower bound, and the inverse derivative is bounded by the reciprocal square of that lower bound.

\begin{proposition}[Same-sample energy observations]\label{prop:energyprobes}
All finite readings
\begin{equation}\label{eq:energy_probe_replicas}
 \E_{T,z}\exp\left(-\sum_{j=1}^r\norm{Y_T(z_jh_j)}^2\right)
 =\E_T\prod_{j=1}^rU_{h_j}(Q(T))
\end{equation}
determine the fixed-frame Borel law of a samplewise bounded tuple. The auxiliary $z_j$ are independent, and all probes use the same original sample. On a uniform operator-norm ball, convergence of these readings gives a unique limiting law in the compact positive-border space; this law is supported on genuine tuples if and only if the mean escape score of \eqref{eq:lawescape} vanishes.
\end{proposition}
\begin{proof}
Conditioning on $T$ proves \eqref{eq:energy_probe_replicas}. Repeated probes give every mixed integer moment of a countable family of $[0,1]$ coordinates. Polynomial density on finite cubes determines their joint law. The inverse in \eqref{eq:bounded_energy_probe}, followed by faithful-anchor inversion, is Borel and identifies the tuple. These Borel statements can be checked on the countable exhaustion by bounded balls.

On a uniform ball, the probes are continuous on the compact positive-border space and separate its points. Their generated algebra contains constants, so Stone--Weierstrass determines its probability measures. Every subsequence has a weakly convergent further subsequence, and the limiting readings determine its unique limit; hence the original sequence converges. The genuine-support assertion is \eqref{eq:lawescape}.
\end{proof}

Output tails also admit finite-input bounded probes. For $\norm{A_a}\le M$, take independent cylindrical circular Gaussians $\xi_a$, used only through $RP_L$, and set
\[
 Z_{N,L}=\sum_{a=1}^b(I-P_N)A_aRP_L\xi_a,\qquad
 h_{N,L}=\E_\xi\norm{Z_{N,L}}^2.
\]
Then
$0\le h_N(T)-h_{N,L}\le bM^2\norm{R(I-P_L)}_2^2$.
Conditioned on $T$, the covariance $\Sigma$ satisfies
\[
 \E_\xi\norm Z^4=(\Tr\Sigma)^2+\Tr\Sigma^2\le2(bM^2)^2.
\]
Since $0\le w-(1-e^{-tw})/t\le tw^2/2$, for $t>0$,
\begin{equation}\label{eq:finite_energy_tail}
 0\le h_N(T)-\frac{1-\E_\xi e^{-t\norm{Z_{N,L}}^2}}t
 \le bM^2\norm{R(I-P_L)}_2^2+t(bM^2)^2.
\end{equation}
The errors correspond to finite input resolution and finite probe strength. With $J$ independent auxiliary repetitions, the bounded expectation has sample-mean error exceeding $\varepsilon$ with probability at most $1/(4J\varepsilon^2)$; its contribution to the tail reading is $\varepsilon/t$. This protocol assumes the stated output-norm access, separately from Hermite source observations.

\subsection{Joint laws and source dependence}
\begin{lemma}[Laws from countable bounded coordinates]\label{lem:boundedlaw}
Let $X$ be a standard Borel space and $F:X\to[-1,1]^\N$ a Borel injection with Borel inverse on its image. The expectations of all finite same-sample coordinate monomials determine the law of an $X$-valued random object. A candidate moment array is realizable if and only if its normalized functional is nonnegative on every cylinder polynomial nonnegative on the cube, and the resulting cube probability is concentrated on $F(X)$.
\end{lemma}
\begin{proof}
Polynomials are dense in the continuous functions on each finite cube, so the moments determine all finite-dimensional distributions and hence the countable-coordinate law. For existence, positivity and
$-\norm p_\infty\le p\le\norm p_\infty$
bound the normalized functional by the supremum norm. Cylinder polynomials are uniformly dense in the continuous functions on the compact countable product. Riesz representation gives a probability measure there; concentration on $F(X)$ permits the Borel inverse. Necessity follows by integration.
\end{proof}

For several positive observations, \eqref{eq:detmarks} gives such bounded coordinates. All their mixed moments specify the joint fixed-frame law. In particular
$\E[\rho(Q_1)\otimes\cdots\otimes\rho(Q_k)]$
is trace class of trace one and needs no unbounded response moments.

To recover a fixed Gaussian realization, also observe the bounded noise coordinates $b_j(g)=2\arctan(g_j)/\pi$. Let $f_k(T)$ be a countable complete bounded operator observation. The mixed expectations
\begin{equation}\label{eq:boundednoisemarks}
 \E\left[\prod_jb_j(g)^{\alpha_j}\prod_kf_k(T)^{\beta_k}\right]
\end{equation}
determine $\Law(g,T)$. If $T=t(g)$ and $\widetilde T=\widetilde t(g)$ have the same such law, their conditional laws given $g$ are identical Dirac measures, and hence $T=\widetilde T$ almost surely. Coefficient injectivity then recovers the same-source kernels.

For candidate data this deterministic source dependence is the condition
\begin{equation}\label{eq:graphcert}
 \E\left|f_k(T)-\E[f_k(T)\mid g]\right|^2=0\qquad(k\ge1).
\end{equation}
Together with the Gaussian source marginal and concentration on actual operator observations, it makes the full tuple a measurable function of the source, since the coordinates separate points and have a Borel inverse. Fixed-chaos recovery also requires the homogeneous finite-projection tests and integrability of Proposition~\ref{prop:chaosimage}; all-radius recovery requires the admissible Cauchy condition.

The difference between source and marginal information is already visible in $gP$ and $-gP$. For joint products, take Pauli matrices $X,Z$ and $A=gX$, $B_+=gZ$, $B_-=-gZ$. Every individual marginal law agrees between the two models, but $\E AB_+=XZ$ and $\E AB_-=-XZ$. The coupling is part of the operational information.

\subsection{Unlabelled spectra}
Taking only $B=bI$ in \eqref{eq:detmarks}, together with $b=0$, recovers $\det_F(I+bQ)$. Values on a set with an interior accumulation point determine this entire function and hence the nonzero eigenvalues with multiplicities. Kernel dimensions and relative eigenspaces are absent from this quotient.

For a rectangular operator $T:\cC\to\cE$, let
\[
 H_T=\begin{pmatrix}0&T^*\\T&0\end{pmatrix},\qquad
 J=\diag(I,-I),\qquad Q_T=(H_T)_+^{2a},
\]
where $2a\ge1$ and $T\in\Sch_{2a}$. The paired spectrum and $JH_TJ=-H_T$ give
\begin{equation}\label{eq:spectralpositive}
 H_T=Q_T^{1/(2a)}-JQ_T^{1/(2a)}J,\qquad
 \Tr Q_T=\norm T_{2a}^{2a},\qquad
 \det_F(I+zQ_T)=\det_F(I+z(T^*T)^a).
\end{equation}
Complete operator-valued coordinates of $Q_T$ recover $T$, whereas its scalar determinant coordinates recover only the nonzero singular values. The spectrum of a separate anchored Gram observation is generally a different spectrum.

Even a joint scalar source can erase order. For the Pauli triple let $G_j=2I+\sigma_j>0$ and $\widetilde G_j=G_j^T$. Then
\begin{align}
 \det\left(I-\sum_jz_jG_j\right)
 &=\det\left(I-\sum_jz_j\widetilde G_j\right)
 =\left(1-2\sum_jz_j\right)^2-\sum_jz_j^2,\label{eq:transposecharacter}\\
 \Tr(G_1G_2G_3)&=16+2i,\qquad
 \Tr(\widetilde G_1\widetilde G_2\widetilde G_3)=16-2i.\label{eq:transposeorientation}
\end{align}
The Pauli anticommutation relations and $\sigma_x\sigma_y\sigma_z=iI$ prove both identities. The circular Gaussian quadratic-form vectors $(g^*G_jg)_j$ and their transposed counterparts also have the same joint law, by $g\mapsto\bar g$. Matrix-valued cycle sources in Proposition~\ref{prop:matrix_word_source} retain the ordered traces that scalar sources remove.


\section{Continuation observability for Gaussian networks}\label{sec:physicalboundary}
A network module is observed through every permitted boundary continuation. Operational saturation therefore selects its response class, not a unique graph or fixed-source frame. Positive edge spectra generate separating continuation tests; finite elimination descends to the resulting orbit tables; and spectral-order tails complete those tables in their own norm. This is a concrete realization of operational observability with a coefficientwise product, distinct from the same-source Gaussian contraction product.

\subsection{Independent half-edges and directed cuts}
Let $G$ be a finite loopless graph, allowing parallel edges, with input anchor $c$, output anchor $e$, and $m\ge1$ Gaussian vertices. Each edge $b=\{u,v\}$ occupies independent half-edge Hilbert factors and carries a nonzero tensor
\[
 t_b=\sum_js_{b,j}x_{b,j}\otimes y_{b,j},\qquad
 B_b\in\Sch_2,\qquad h_b=\norm{B_b}_2,
\]
with orthonormal Schmidt families. At each vertex take the tensor product of its half-edge spaces; an isolated vertex has scalar space. Regroup $\bigotimes_bt_b$ by vertices to obtain $K_G$. Evaluate each Gaussian vertex by an independent standard circular Gaussian family in its full vertex space, using the Hilbert-kernel completion for infinite dimensions. Internal half-edges remain separate until this Gaussian evaluation.

\begin{proposition}[Edge spectra and all directed cuts]\label{prop:graphcuts}
Set $H_G=\prod_bh_b$. For $L_S=\{c\}\cup S$, $S\subset[m]$, the nonzero singular values of the associated cut are
\begin{equation}\label{eq:graphcuts}
 \left\{\left(\prod_{b\notin\partial L_S}h_b\right)
                  \prod_{b\in\partial L_S}s_{b,j_b}\right\}_{(j_b)}.
\end{equation}
For $2\le s\le\infty$, define $\kappa_b(s)=\log(h_b/\norm{B_b}_s)\ge0$. Then
\begin{equation}\label{eq:graphmincut}
 \Prof_{m,s}(K_G)=H_Ge^{-\lambda_s},\qquad
 \lambda_s=\min_{c\in L,\ e\notin L}\sum_{b\in\partial L}\kappa_b(s).
\end{equation}
\end{proposition}
\begin{proof}
A same-side edge is a vector or covector with one singular value $h_b$; a crossing edge is an operator with singular values $s_{b,j}$. Independence of the half-edge factors makes the full cut, after separate unitary rearrangements on its two sides, the tensor product of these local operators. Singular values multiply, proving \eqref{eq:graphcuts}. Their $\ell^s$ norm and maximization over all cuts give \eqref{eq:graphmincut}. Truncate every Schmidt expansion to pass from finite rank to infinite rank: the full tensor product converges in its Hilbert norm, and every $s\ge2$ cut norm is bounded by that norm.
\end{proof}
This places the network in the original cut geometry. At $s=\infty$, the formula concerns the deterministic coefficient norm; random operator-norm bounds still have the Gaussian dimensional behavior of Proposition~\ref{prop:sharpfinite}.

\subsection{Positive weights and boundary elimination}
Fix a spectral moment order $q\ge1$. Put $Q_b=B_b^*B_b$ and $\Omega_q=S_q$. The edge weight is the strictly positive class function
\begin{equation}\label{eq:physical_edge_weight}
 f_{Q,q}(\pi)=\Tr(Q^{\otimes q}R_q(\pi))
 =\prod_{a\in\operatorname{Cycles}(\pi)}\Tr(Q^{|a|}),
\end{equation}
where $R_q$ permutes tensor factors.

\begin{theorem}[Moments and boundary composition]\label{thm:physicalresponse}
A module with boundary vertices $B$ and Gaussian internal vertices $I$ has response
\begin{equation}\label{eq:physical_response}
 H_{G,q}(\pi_B)=\sum_{\pi_I\in S_q^I}
                \prod_{b=\{u,v\}}f_{Q_b,q}(\pi_u^{-1}\pi_v).
\end{equation}
For the two-anchor operator,
\[
 Z_q(K_G)=\E\Tr[(T_{K_G}^*T_{K_G})^q]=H_{G,q}(\tau_q,\Id),
 \qquad\tau_q=(1\,2\,\cdots\,q).
\]
Normalize $\widehat H_{G,q}=H_{G,q}/(q!)^{|I|}$. When modules have disjoint old internal vertices and a new shared vertex $v$ is eliminated once, their response is
\begin{equation}\label{eq:physical_eliminate}
 \widehat H_{{\rm new},q}(\pi_{B'})
 =\frac1{q!}\sum_{\sigma\in S_q}
       \prod_j\widehat H_{j,q}(\pi_{B_j\setminus\{v\}},\sigma).
\end{equation}
The final normalized scalar observation is $Z_q(K_G)/(q!)^m$.
\end{theorem}
\begin{proof}
First take finite edge ranks. Expanding the spectral trace gives $q$ copies of the tensor and $q$ conjugate copies. At a Gaussian vertex, a Wick pairing is a bijection between the unconjugated and conjugated coordinate copies, indexed by a permutation $\pi_v\in S_q$. Along an edge, the relative identification is $\pi_u^{-1}\pi_v$. A cycle of length $a$ identifies its Schmidt indices into
$\sum_js_{b,j}^{2a}=\Tr Q_b^a$.
All edge contributions multiply. The external Schatten trace supplies a full cycle at one anchor and the identity at the other. This proves \eqref{eq:physical_response}, the standard permutation-pairing mechanism for Gaussian tensor networks \cite{Cheng}, with the boundary labels retained.

Composition first multiplies the module responses and then sums the new internal label. Each internal Gaussian vertex contributes one factor $q!$ to the normalization; the new shared vertex contributes exactly one additional such factor. Regrouping the finite sums proves \eqref{eq:physical_eliminate}.

For infinite ranks, a Hilbert-valued circular Gaussian sum obeys
\[
 \E\left\|\sum_jg_jv_j\right\|^{2q}
 \le q!\left(\sum_j\norm{v_j}^2\right)^q.
\]
Diagonalizing its covariance expresses the left side as $q!$ times the complete homogeneous symmetric polynomial of its eigenvalues, bounded by the $q$th power of their sum. Applying the estimate to successive independent Gaussian legs, with Minkowski for the remaining probability norms, gives
\begin{equation}\label{eq:physical_hilbert_moment}
 \norm{T_K}_{L^{2q}(\Sch_2)}\le(q!)^{m/(2q)}\norm K_2.
\end{equation}
Schmidt truncations converge in the full Hilbert norm $\norm{K_G}_2=H_G$, hence in the required random $L^{2q}$ norm. Their spectral moments converge by H\"older. The edge power traces converge by trace-norm convergence of $Q_b$, and the permutation sum is finite for each $q$. All formulas pass to the limit.
\end{proof}

The index $q$ is a spectral moment order, independent of the source degree $m$. Parallel module multiplication takes place at the same $q$. On generating functions it is coefficientwise, or Hadamard, multiplication. For instance two scalar independent circular Gaussians each have normalized moment sequence one, as does their product after the appropriate source-degree normalization; Cauchy multiplication of $(1-z)^{-1}$ would instead produce coefficients $q+1$.

\subsection{All continuations determine the minimal state}
Let $\kappa(x,y)$ be the cycle type of $x^{-1}y$, and let
\[
 \Gamma_q=\{\gamma\in\operatorname{Sym}(\Omega_q):
            \kappa(\gamma x,\gamma y)=\kappa(x,y)\ \forall x,y\}.
\]
For $k$ ordered boundary points, define the real space
$V_{q,k}=\R^{\Omega_q^k/\Gamma_q}$
with the maximum norm. Let $\mathcal B_{q,k}$ be the real linear span of all finite physical-module responses, allowing arbitrary nonzero finite-rank positive edge spectra.

\begin{lemma}[Positive spectra span the color tests]\label{lem:positivecolors}
For fixed $q$, finitely many rank-$q$ positive operators with strictly positive rational eigenvalues have weights $f_{Q,q}$ spanning all real class functions on $S_q$.
\end{lemma}
\begin{proof}
For a cycle-type partition $\mu\vdash q$, the weight at spectrum $x=(x_1,\ldots,x_q)\in(0,\infty)^q$ is
$p_\mu(x)=\prod_{a\in\mu}\sum_jx_j^a$.
Newton's identities make $p_1,\ldots,p_q$ polynomial generators equivalent by an invertible triangular change to the elementary symmetric generators. They are algebraically independent, so the monomials $p_\mu$ for different partitions are linearly independent.

Their evaluation vectors on the positive open set therefore span the full partition-coordinate space. Otherwise a nonzero linear combination would vanish on an open set and hence identically. Choose as many evaluation points as there are partitions with invertible evaluation matrix. Its determinant remains nonzero in a neighborhood, allowing positive rational points by density. Inverting this finite matrix expresses every cycle-type indicator as a linear combination of the resulting positive edge weights.
\end{proof}
The linear combination may have signed coefficients. It is a postprocessing of positive physical readings, whereas composition of any single positive module continues to use positive weights.

\begin{theorem}[Continuation completeness and linear minimality]\label{thm:physicalcomplete}
For every $q,k\ge1$, $\mathcal B_{q,k}=V_{q,k}$. A linear summary of module responses supporting every finite physical continuation must be injective on this space. Its minimal dimension is
$\#(S_q^k/\Gamma_q)$.
\end{theorem}
\begin{proof}
Each edge weight depends only on the relative cycle color of its endpoints. A simultaneous color-preserving bijection of all vertex labels permutes the internal summation variables and leaves every edge weight unchanged. Hence every physical response lies in $V_{q,k}$.

To prove that these responses span all of $V_{q,k}$, it suffices to recover an orbit indicator. Fix a boundary configuration $a=(a_1,\ldots,a_k)$. Introduce an internal calibration vertex $y_s$ for every $s\in\Omega_q$, and impose all the pairwise colors of this labelled copy of $\Omega_q$. Connect boundary vertex $i$ to $y_{a_i}$ with the equality color. The resulting formal indicator response is
\begin{equation}\label{eq:physical_calibration}
 P_a(x)=\sum_{y\in\Omega_q^{\Omega_q}}
 \prod_{s<t}\mathbf1_{\{\kappa(y_s,y_t)=\kappa(s,t)\}}
 \prod_{i=1}^k\mathbf1_{\{x_i=y_{a_i}\}}.
\end{equation}
In a nonzero summand, distinct calibration vertices receive distinct labels: the color of $(s,t)$ is not equality when $s\ne t$. Their assignment is consequently a bijection, and the pairwise constraints make it an element of $\Gamma_q$. Therefore
\[
 P_a(x)=|\operatorname{Stab}_{\Gamma_q}(a)|\,
                  \mathbf1_{\{x\in\Gamma_qa\}}.
\]
Lemma~\ref{lem:positivecolors} expresses each indicator edge as a finite signed linear combination of positive physical edge weights. Expanding this finite product expresses the orbit indicator as a linear combination of actual module responses. Thus $\mathcal B_{q,k}=V_{q,k}$. This is the graph-algebra separation mechanism~\cite{Lovasz} realized by the specified spectral tests.

For minimality use the continuation pairing
\begin{equation}\label{eq:physical_pairing}
 \langle H,F\rangle_q
 =\sum_{x_2,\ldots,x_k}H(\Id,x_2,\ldots,x_k)F(\Id,x_2,\ldots,x_k).
\end{equation}
Common left multiplication belongs to $\Gamma_q$, so fixing the first label to $\Id$ meets every orbit. This pairing is therefore positive definite on the invariant response space. It is implementable by the permitted continuations: retain the first label as the identity anchor, Gaussianize each other shared boundary vertex once, and attach a full-cycle anchor through a rank-one edge of constant weight one. Normalization introduces only a known nonzero scalar.

If every physical continuation annihilates $H-\widetilde H$, the spanning result permits a linear combination of continuations equal to $H-\widetilde H$ itself. Its pairing square is zero, forcing equality of the responses. Any linear summary with a nonzero kernel would therefore lose a permitted continuation reading. The complete orbit table is injective on $V_{q,k}$ and supports the stated composition operations, attaining the minimal dimension.
\end{proof}
The theorem fixes the class of all finite continuations. Its calibration may use $q!$ internal vertices and an ill-conditioned signed inversion. Restrictions on ranks, graph classes, or experimental budgets lead to a different visible quotient.

\subsection{Finite inversion and operational errors}
Choose the orbit-indicator basis $\phi_1,\ldots,\phi_n$ of $V_{q,k}$. Completeness of the physical continuation tests supplies $n$ actual continuations $F_i$ with invertible matrix
\[
 A_q(i,a)=\langle\phi_a,F_i\rangle_q.
\]
Indeed, otherwise their rows would fail to span the dual and would annihilate a nonzero response. For $H=\sum_ah_a\phi_a$, its readings are $y=A_qh$. Therefore
\begin{equation}\label{eq:physical_finite_inverse}
 \widehat h=A_q^{-1}\widetilde y,\qquad
 \norm{\widehat H-H}_\infty\le\kappa_q\norm{\widetilde y-y}_{\ell^2},\qquad
 \kappa_q=\norm{A_q^{-1}}_{\ell^2\to\ell^\infty}.
\end{equation}
The finite output is in the real response space; noisy inversion need not preserve positivity.

Normalized elimination is contractive in the maximum norm, and pointwise multiplication has multilinear norm one. Both preserve the relevant color invariance. Thus the same error cascade of Section~\ref{sec:transport} applies to boundary composition: \eqref{eq:physical_finite_inverse} supplies input errors, local numerical errors supply residuals, and backward dual tests are the pulled-back continuations. Joint response tables are retained after elimination.

For perturbations of physical positive edges satisfying
$|\log(\widetilde f_b/f_b)|\le\epsilon_b$,
every configuration weight has relative factor between $e^{-\sum_b\epsilon_b}$ and $e^{\sum_b\epsilon_b}$. Positivity of the sum preserves this forward relative estimate. The signed finite inverse has the separate conditioning cost $\kappa_q$.

\subsection{A completion with spectral-order tails}
For a fixed boundary size $k$ and radius $R>0$ let
\begin{equation}\label{eq:physical_czero}
 \mathscr V^0_{R,k}=\{H=(H_q)_{q\ge1}:H_q\in V_{q,k},\
                       R^q\norm{H_q}_\infty\to0\},\qquad
 \norm H_R=\sup_{q\ge1}R^q\norm{H_q}_\infty.
\end{equation}
This weighted $c_0$ sum is the completion of finite-order response tables. The larger weighted $\ell^\infty$ space need not have dense finite-order truncations: for $H_q=R^{-q}\phi_q$ and $\norm{\phi_q}_\infty=1$, every truncation tail has norm one.

\begin{theorem}[Finite reconstruction and completed responses]\label{thm:physicalcompletion}
The finite inverses \eqref{eq:physical_finite_inverse} determine a unique completed object in \eqref{eq:physical_czero}. Suppose $R_+>R$ and
$\sup_qR_+^q\norm{H_q}_\infty\le B_+$.
Retaining only orders $q\le Q$ with reading errors at most $\delta_q$, and setting higher orders to zero, gives
\begin{equation}\label{eq:physical_full_error}
 \norm{\widehat H^{\le Q}-H}_R
 \le\max_{q\le Q}R^q\kappa_q\delta_q+(R/R_+)^{Q+1}B_+.
\end{equation}
Pointwise multiplication and normalized elimination extend continuously, with
\begin{equation}\label{eq:physical_radius_ops}
 \norm{HF}_{R_1R_2}\le\norm H_{R_1}\norm F_{R_2},\qquad
 \norm{\mathsf E_vH}_R\le\norm H_R.
\end{equation}
Boundary labels are pulled back to the common boundary set before multiplication. All finite-core composition identities persist under completion.
\end{theorem}
\begin{proof}
A Cauchy sequence in the weighted supremum norm converges in each finite-dimensional $V_{q,k}$. Its uniform small-tail property passes to the limit, so the limit lies in the weighted $c_0$ space. Conversely this tail condition is exactly convergence of finite-order truncations. The finite inverses separate the individual tables and thus the completed objects. The head error is \eqref{eq:physical_finite_inverse}; the tail uses
$R^q\norm{H_q}\le(R/R_+)^qB_+$.

The multiplication bound follows by multiplying the weights at each order. Normalized elimination averages $q!$ entries and is contractive. These operations preserve the vanishing weighted tails, so their finite-core identities extend by continuity. In particular finite elimination order is governed by rearrangement of finite sums and remains consistent in the completion.
\end{proof}

The physical modules supply the stronger-radius bound directly. Since $\Tr Q^a\le(\Tr Q)^a$, every edge weight is bounded by $(\Tr Q)^q$. Normalized internal summation is an average, hence
\begin{equation}\label{eq:physical_growth}
 \norm{\widehat H_{G,q}}_\infty\le\prod_b(\Tr Q_b)^q=H_G^{2q}.
\end{equation}
Choose $0<R<R_+<H_G^{-2}$. This gives an explicit stronger-radius budget (in particular $B_+=1$ is valid). For a final scalar table entry $h_q$ and $|z|<R_+$,
\[
 \left|\sum_{q>Q}h_qz^q\right|
 \le B_+\frac{(|z|/R_+)^{Q+1}}{1-|z|/R_+}.
\]
This sufficient joint-response radius preserves all future continuations. A final scalar spectral transform can have a larger exact radius.

\subsection{State dimensions and exact tree responses}
For $q\ge3$, the complete transposition graph automorphism theorem \cite[Theorem~1.1]{Ganesan} gives
\[
 \Gamma_q=\{x\mapsto axb,\ x\mapsto ax^{-1}b:a,b\in S_q\},
 \qquad |\Gamma_q|=2(q!)^2.
\]
These maps preserve every cycle color. Conversely a color-preserving map preserves transposition adjacency, so belongs to the group given by that theorem. Fixing the first boundary label leaves conjugation and conjugation followed by inversion. With
$c(g)=|C_{S_q}(g)|$ and $b(g)=\#\{y\in S_q:y^2=g^2\}$,
Burnside's formula yields
\begin{equation}\label{eq:physical_dimension}
 \dim V_{q,k}=\frac1{2q!}\sum_{g\in S_q}
                  \{c(g)^{k-1}+b(g)^{k-1}\}.
\end{equation}
Conjugation fixes $c(g)$ labels. The conjugation-inversion fixed-point equation is equivalent to $(xg)^2=g^2$, giving $b(g)$. For $q=1$ the dimension is one; for $q=2$ it is $2^{k-1}$. For example,
\begin{center}
\begin{tabular}{@{}c rrrr@{}}\toprule
Spectral trace degree $2q$ & $k=1$ & $k=2$ & $k=3$ & $k=4$\\\midrule
$6$ & $1$ & $3$ & $11$ & $46$\\
$8$ & $1$ & $5$ & $43$ & $550$\\\bottomrule
\end{tabular}
\end{center}

For an edge, the Schur--Weyl expansion \cite{Macdonald} is
\[
 f_{Q,q}(\pi)=\sum_{\lambda\vdash q}s_\lambda(Q)\chi^\lambda(\pi),\qquad
 (q!)^{-1}\sum_{\pi\in S_q}f_{Q,q}(\pi)=h_q(Q).
\]
\begin{theorem}[Tree formula and the exact scalar radius]\label{thm:physicaltree}
Suppose $G$ is a connected tree, let $P$ be its unique path between the two anchors, and let $\ell=|P|$. Then
\begin{equation}\label{eq:physical_tree_formula}
 \frac{Z_q(K_G)}{(q!)^m}
 =\left(\prod_{b\notin P}h_q(Q_b)\right)
 \sum_{j=0}^{q-1}(-1)^j\binom{q-1}j^{1-\ell}
                    \prod_{b\in P}s_{(q-j,1^j)}(Q_b).
\end{equation}
Moreover $\tau(K_G)^2=\prod_b\norm{Q_b}_{\op}$. The final normalized spectral generating function has exact radius $\tau(K_G)^{-2}$.
\end{theorem}
\begin{proof}
Eliminate leaves off the anchor path. Their relative permutation is averaged uniformly in $S_q$, giving one factor $h_q(Q_b)$ each. The remaining path gives normalized convolution of central functions. Character orthogonality multiplies their common $\lambda$-channel coefficients and contributes $d_\lambda^{1-\ell}$. At a full $q$-cycle, the only nonzero irreducible character values are those of hooks $\lambda=(q-j,1^j)$; their values are $(-1)^j$ and their dimensions are $\binom{q-1}j$. Substitution proves \eqref{eq:physical_tree_formula}. Infinite ranks follow from the Hilbert moment approximation \eqref{eq:physical_hilbert_moment}.

To compute $\tau$, contract the tensor with unit vectors at the Gaussian vertices and with unit input and output vectors at the anchors. Starting at any leaf, contraction across its adjacent edge costs at most $\norm{B_b}_{\op}$. Absorb the resulting vector into the adjacent vertex tensor and continue inductively. Partial contraction of a Hilbert tensor has norm at most the product of its norms, so this proves
$\tau(K_G)\le\prod_b\norm{B_b}_{\op}$.
For the reverse bound, choose top Schmidt vectors on every edge and their tensor products at every vertex. All these test vectors have norm one and give the product of top singular values. Each nonzero compact edge operator attains its largest singular value. Theorem~\ref{thm:factorial_radius} gives the exact generating radius.
\end{proof}

The hook restriction belongs to the final full-cycle readout. Intermediate response tables preserve the other channels needed by general continuations. Local determinants provide $e_j,h_j$, and the recursion
\[
 s_{(q)}=h_q,\qquad
 s_{(q-j,1^j)}=h_{q-j}e_j-s_{(q-j+1,1^{j-1})}
\]
computes the required hooks. Given the local spectra, all tree readings through order $Q$ require $O(|E(G)|Q^2)$ scalar arithmetic operations, excluding the acquisition of spectral data and bit complexity.

There is also an explicit final tail. Suppose every edge spectrum has finite rank and a simple largest eigenvalue $a_b$. Define
\[
 C_b=\prod_{\lambda<a_b}(1-\lambda/a_b)^{-1},\qquad
 A_G=\left(\prod_bC_b\right)
       \sum_{j=0}^{r_*-1}\prod_{b\in P}\frac{e_j(Q_b)}{a_b^j},
 \qquad r_* =\min_{b\in P}\rank Q_b.
\]
The bounds $h_q(Q_b)\le C_ba_b^q$ and $s_{(q-j,1^j)}\le h_{q-j}e_j$, together with $\ell\ge1$, give
$0\le Z_q/(q!)^m\le A_G\tau(K_G)^{2q}$.
Therefore
\begin{equation}\label{eq:physical_tree_tail}
 0\le\mathcal B_{K_G,1}(z)-\sum_{q=1}^Q\frac{Z_q(K_G)}{(q!)^m}z^q
 \le A_G\frac{(z\tau(K_G)^2)^{Q+1}}{1-z\tau(K_G)^2},
 \qquad0\le z<\tau(K_G)^{-2}.
\end{equation}
Repeated largest eigenvalues require an additional polynomial factor in the moment envelope.

The network application thus realizes each stage of structural reconstruction: edge spectra determine original norm costs; continuation experiments determine the minimal response space and a finite inverse; elimination supplies its operations; and stronger-radius bounds certify the completed response. On trees the resulting scalar target and its convergence radius are explicit.

\section*{Concluding perspective}
The recovered object is an algebra together with the norm scale in which its operations and approximations are controlled. Observable quotients specify its relations; Gaussian duality recovers its coefficient geometry; and ordered contractions transport ordinary multiplication. In the Hilbert geometry, symmetric site systems provide a microscopic realization of the same multiplication, including a quantitative finite-size correction and a compatible finite inverse.

The limiting topology remains part of this structure. Positive mass and source innovations select compact coefficient families, while anchored Schwarz defects and spectral mass defects distinguish genuine realizations inside larger observation closures. These mechanisms preserve different information, but use the same sequence of tasks: identify the target, establish its inverse estimate, transport its operations, and control its unresolved directions.

\paragraph{Acknowledgment.}
The author used ChatGPT (GPT-5.5, GPT-5.6, and GPT-6.0; OpenAI) interactively in developing proofs and refining the presentation, and takes full responsibility for the mathematical content.


\begin{thebibliography}{99}
\small\raggedright
\setlength{\itemsep}{2pt}

\bibitem{AAC} F. Achleitner, A. Arnold and E. A. Carlen,
The hypocoercivity index for the short time behavior of linear time-invariant ODE systems,
\emph{J. Differential Equations} \textbf{371} (2023), 83--115.
\doi{10.1016/j.jde.2023.06.027}; \arxiv{2109.10784}.

\bibitem{AT} W. N. Anderson, Jr.\ and G. E. Trapp,
Shorted operators. II, \emph{SIAM J. Appl. Math.} \textbf{28} (1975), 60--71. \doi{10.1137/0128007}.

\bibitem{Ando} T. Ando, Generalized Schur complements,
\emph{Linear Algebra Appl.} \textbf{27} (1979), 173--186. \doi{10.1016/0024-3795(79)90040-5}.

\bibitem{Aniello} P. Aniello, A. Kossakowski, G. Marmo and F. Ventriglia,
Brownian motion on Lie groups and open quantum systems,
\emph{J. Phys. A: Math. Theor.} \textbf{43} (2010), 265301.
\doi{10.1088/1751-8113/43/26/265301}; \arxiv{1002.3507}.

\bibitem{BBF} J. A. Ball, V. Bolotnikov and Q. Fang,
Multivariable backward-shift-invariant subspaces and observability operators,
\emph{Multidimens. Syst. Signal Process.} \textbf{18} (2007), 191--248.
\doi{10.1007/s11045-006-0011-y}; \arxiv{math/0610634}.

\bibitem{BCL} K. Ball, E. A. Carlen and E. H. Lieb,
Sharp uniform convexity and smoothness inequalities for trace norms,
\emph{Invent. Math.} \textbf{115} (1994), 463--482. \doi{10.1007/BF01231769}.

\bibitem{BL} J. Bergh and J. L\"ofstr\"om, \emph{Interpolation Spaces. An Introduction},
Grundlehren der mathematischen Wissenschaften 223, Springer, Berlin, 1976. \doi{10.1007/978-3-642-66451-9}.

\bibitem{BK} R. Bhatia and F. Kittaneh, Norm inequalities for partitioned operators and an application,
\emph{Math. Ann.} \textbf{287} (1990), 719--726. \doi{10.1007/BF01446925}.

\bibitem{Bornemann} F. Bornemann, On the numerical evaluation of Fredholm determinants,
\emph{Math. Comp.} \textbf{79} (2010), 871--915. \doi{10.1090/S0025-5718-09-02280-7}; \arxiv{0804.2543}.

\bibitem{BLP} J.-C. Bourin and E.-Y. Lee,
A Pythagorean theorem for partitioned matrices,
preprint, 2020. \arxiv{2011.13261}.

\bibitem{Buchholz} A. Buchholz, Operator Khintchine inequality in non-commutative probability,
\emph{Math. Ann.} \textbf{319} (2001), 1--16. \doi{10.1007/PL00004425}.

\bibitem{BCKP} S. Burgdorf, K. Cafuta, I. Klep and J. Povh,
The tracial moment problem and trace-optimization of polynomials,
\emph{Math. Program.} \textbf{137} (2013), 557--578.
\doi{10.1007/s10107-011-0505-8}.

\bibitem{CP} F. Chapoton and F. Patras,
Enveloping algebras of preLie algebras, Solomon idempotents and the Magnus formula,
\emph{Int. J. Algebra Comput.} \textbf{23} (2013), 853--861.
\doi{10.1142/S0218196713400134}; \arxiv{1201.2159}.

\bibitem{Cheng} N. Cheng, C. Lancien, G. Penington, M. Walter and F. Witteveen,
Random tensor networks with non-trivial links,
\emph{Ann. Henri Poincar\'e} \textbf{25} (2024), 2107--2212.
\doi{10.1007/s00023-023-01358-2}; \arxiv{2206.10482}.

\bibitem{Choi} M.-D. Choi, A Schwarz inequality for positive linear maps on $C^*$-algebras,
\emph{Illinois J. Math.} \textbf{18} (1974), 565--574. \doi{10.1215/ijm/1256051007}.

\bibitem{CJK} M.-D. Choi, N. Johnston and D. W. Kribs,
The multiplicative domain in quantum error correction,
\emph{J. Phys. A: Math. Theor.} \textbf{42} (2009), 245303.
\doi{10.1088/1751-8113/42/24/245303}; \arxiv{0811.0947}.

\bibitem{DPLS} P. Da Pelo, A. Lanconelli and A. I. Stan,
A sharp interpolation between the H\"older and Gaussian Young inequalities,
\emph{Infin. Dimens. Anal. Quantum Probab. Relat. Top.} \textbf{19} (2016), 1650001.
\doi{10.1142/S0219025716500016}; \arxiv{1504.06156}.

\bibitem{DFSU} J. Deschamps, F. Fagnola, E. Sasso and V. Umanit\`a,
Structure of uniformly continuous quantum Markov semigroups,
\emph{Rev. Math. Phys.} \textbf{28} (2016), 1650003.
\doi{10.1142/S0129055X16500033}; \arxiv{1412.3239}.

\bibitem{Douglas} R. G. Douglas,
On majorization, factorization, and range inclusion of operators on Hilbert space,
\emph{Proc. Amer. Math. Soc.} \textbf{17} (1966), 413--415. \doi{10.1090/S0002-9939-1966-0203464-1}.

\bibitem{EF} M. Ettehad and S. Foucart,
Instances of computational optimal recovery: dealing with observation errors,
\emph{SIAM/ASA J. Uncertain. Quantif.} \textbf{9} (2021), 1438--1456.
\doi{10.1137/20M1328476}.

\bibitem{FK} T. Formisano and E. Kissin,
Clarkson--McCarthy inequalities for $l_p$-spaces of operators in Schatten ideals,
\emph{Integral Equations Operator Theory} \textbf{79} (2014), 151--173.
\doi{10.1007/s00020-014-2145-x}.

\bibitem{FHS} V. Futorny, R. A. Horn and V. V. Sergeichuk,
Specht's criterion for systems of linear mappings,
\emph{Linear Algebra Appl.} \textbf{519} (2017), 278--295.
\doi{10.1016/j.laa.2017.01.006}; \arxiv{1701.08826}.

\bibitem{Ganesan} A. Ganesan, Automorphism group of the complete transposition graph,
\emph{J. Algebraic Combin.} \textbf{42} (2015), 793--801. \doi{10.1007/s10801-015-0602-5}; \arxiv{1404.7363}.

\bibitem{HP} U. Haagerup and G. Pisier, Bounded linear operators between $C^*$-algebras,
\emph{Duke Math. J.} \textbf{71} (1993), 889--925.
\doi{10.1215/S0012-7094-93-07134-7}; \arxiv{math/9302214}.

\bibitem{HKP} E. Hashorva, D. Korshunov and V. I. Piterbarg,
Asymptotic expansion of Gaussian chaos via probabilistic approach,
\emph{Extremes} \textbf{18} (2015), 315--347.
\doi{10.1007/s10687-015-0215-3}; \arxiv{1307.5857}.

\bibitem{Janson} S. Janson, \emph{Gaussian Hilbert Spaces}, Cambridge Tracts in Mathematics 129,
Cambridge University Press, Cambridge, 1997. \doi{10.1017/CBO9780511526169}.

\bibitem{JinEtAl} Y. Jin, C. Musco, A. Sidford and A. V. Singh,
Moments, random walks, and limits for spectrum approximation,
in \emph{Proceedings of the 36th Conference on Learning Theory},
Proceedings of Machine Learning Research \textbf{195} (2023), 5373--5394.
\arxiv{2307.00474}.

\bibitem{Jing} N. Jing, Unitary and orthogonal equivalence of sets of matrices,
\emph{Linear Algebra Appl.} \textbf{481} (2015), 235--242.
\doi{10.1016/j.laa.2015.04.036}; \arxiv{1504.06790}.

\bibitem{KS} I. Klep and M. Schweighofer,
Connes' embedding conjecture and sums of hermitian squares,
\emph{Adv. Math.} \textbf{217} (2008), 1816--1837.
\doi{10.1016/j.aim.2007.09.016}; \arxiv{math/0607615}.

\bibitem{KongValiant} W. Kong and G. Valiant,
Spectrum estimation from samples,
\emph{Ann. Statist.} \textbf{45} (2017), 2218--2247.
\doi{10.1214/16-AOS1525}; \arxiv{1602.00061}.

\bibitem{Ledoux} M. Ledoux, A note on large deviations for Wiener chaos,
in \emph{S\'eminaire de Probabilit\'es XXIV}, Lecture Notes in Mathematics 1426,
Springer, Berlin, 1990, 1--14.

\bibitem{LS} G. D. Lin and J. Stoyanov,
On the moment determinacy of products of non-identically distributed random variables,
\emph{Probab. Math. Statist.} \textbf{36} (2016), 21--33.
\arxiv{1406.1654}.

\bibitem{Lindblad} G. Lindblad, On the generators of quantum dynamical semigroups,
\emph{Commun. Math. Phys.} \textbf{48} (1976), 119--130. \doi{10.1007/BF01608499}.

\bibitem{Lovasz} L. Lov\'asz, The rank of connection matrices and the dimension of graph algebras,
\emph{European J. Combin.} \textbf{27} (2006), 962--970. \doi{10.1016/j.ejc.2005.04.012}; \arxiv{math/0408232}.

\bibitem{LuqueThibon} J.-G. Luque and J.-Y. Thibon,
Pfaffian and hafnian identities in shuffle algebras,
\emph{Adv.\ in Appl.\ Math.} \textbf{29} (2002), 620--646.
\arxiv{math/0204026}.

\bibitem{LP} F. Lust-Piquard and G. Pisier, Non-commutative Khintchine and Paley inequalities,
\emph{Ark. Mat.} \textbf{29} (1991), 241--260. \doi{10.1007/BF02384340}.

\bibitem{Maas} J. Maas, Malliavin calculus and decoupling inequalities in Banach spaces,
\emph{J. Math. Anal. Appl.} \textbf{363} (2010), 383--398.
\doi{10.1016/j.jmaa.2009.08.041}; \arxiv{0801.2899}.

\bibitem{Macdonald} I. G. Macdonald,
\emph{Symmetric Functions and Hall Polynomials}, 2nd ed., Oxford University Press,
Oxford, 1995. \doi{10.1093/oso/9780198534891.001.0001}.

\bibitem{MR} C. A. Micchelli and T. J. Rivlin, A survey of optimal recovery,
in \emph{Optimal Estimation in Approximation Theory}, Plenum Press, New York, 1977, 1--54. \doi{10.1007/978-1-4684-2388-4_1}.

\bibitem{Min} A. R. Minabutdinov,
A higher order asymptotic expansion of the Krawtchouk polynomials,
\emph{J. Math. Sci. (N.Y.)} \textbf{215} (2016), no. 6, 738--747.
\doi{10.1007/s10958-016-2879-x}; \arxiv{1508.07421}.

\bibitem{MS} B. Mourrain and K. Schm\"udgen, Flat extensions in $*$-algebras,
\emph{Proc. Amer. Math. Soc.} \textbf{144} (2016), 4873--4885.
\doi{10.1090/proc/13158}; \arxiv{1406.4975}.

\bibitem{NguyenVishnyakova} T. H. Nguyen and A. Vishnyakova,
On entire functions from the Laguerre--P\'olya I class with non-monotonic second quotients of Taylor coefficients,
\emph{Mat. Stud.} \textbf{56} (2021), no. 2, 149--161.
\doi{10.30970/ms.56.2.149-161}; \arxiv{2107.13061}.

\bibitem{Nourdin} I. Nourdin, Yet another proof of the Nualart--Peccati criterion,
\emph{Electron. Commun. Probab.} \textbf{16} (2011), 467--481.
\doi{10.1214/ECP.v16-1642}; \arxiv{1107.3252}.

\bibitem{NPR} I. Nourdin, G. Peccati and G. Reinert,
Invariance principles for homogeneous sums: universality of Gaussian Wiener chaos,
\emph{Ann. Probab.} \textbf{38} (2010), no. 5, 1947--1985.
\doi{10.1214/10-AOP531}; \arxiv{0904.1153}.

\bibitem{NPRR} I. Nourdin, G. Peccati and G. Reinert,
Stein's method and stochastic analysis of Rademacher functionals,
\emph{Electron. J. Probab.} \textbf{15} (2010), no. 55, 1703--1742.
\doi{10.1214/EJP.v15-823}; \arxiv{0810.2890}.

\bibitem{ORT} T. Oh, T. Robert and N. Tzvetkov,
Stochastic nonlinear wave dynamics on compact surfaces,
\emph{Ann. Henri Lebesgue} \textbf{6} (2023), 161--223.
\doi{10.5802/ahl.163}; \arxiv{1904.05277}.

\bibitem{Popescu} G. Popescu,
\emph{Noncommutative weighted shifts, joint similarity, and function theory in several variables},
preprint, 2024. \arxiv{2404.09084}.

\bibitem{PT} J. Potthoff and M. Timpel,
On a dual pair of spaces of smooth and generalized random variables,
\emph{Potential Anal.} \textbf{4} (1995), 637--654. \doi{10.1007/BF02345829}.

\bibitem{Ricard} \'E. Ricard, H\"older estimates for the noncommutative Mazur maps,
\emph{Arch. Math.} \textbf{104} (2015), 37--45.
\doi{10.1007/s00013-014-0710-9}; \arxiv{1407.8334}.

\bibitem{RX} \'E. Ricard and Q. Xu, A noncommutative martingale convexity inequality,
\emph{Ann. Probab.} \textbf{44} (2016), 867--882.
\doi{10.1214/14-AOP990}; \arxiv{1405.0431}.

\bibitem{Robinson} P. L. Robinson, \emph{Julia operators and Halmos dilations}, preprint, 2018.
\arxiv{1803.09329}.

\bibitem{Simon} B. Simon, \emph{Trace Ideals and Their Applications}, 2nd ed., Mathematical Surveys and
Monographs 120, American Mathematical Society, Providence, RI, 2005. \doi{10.1090/surv/120}.

\bibitem{Stinespring} W. F. Stinespring, Positive functions on $C^*$-algebras,
\emph{Proc. Amer. Math. Soc.} \textbf{6} (1955), 211--216.
\doi{10.1090/S0002-9939-1955-0069403-4}.
\end{thebibliography}

\end{document}
